% Generated mechanically from the frozen spaces-simplicial.tex target.
% Do not edit this generated file; edit the bound locale source instead.

% OK, start here.
%

\setcounter{chapter}{84}
\chapter{单纯空间}



\phantomsection
\label{spaces-simplicial-section-phantom}


\section{引言}
\label{spaces-simplicial-section-introduction}

\noindent
本章发展单纯拓扑空间、单纯赋环空间、单纯概形以及单纯代数空间的若干理论。
借助单纯空间理论，有时可以证明一些用其他方法看来难以证明的
局部到整体原理。若干应用实例见论文
\cite{faltings_finiteness}、\cite{Kiehl} 和 \cite{HodgeIII}。

\medskip\noindent
全文假定读者熟悉《单纯方法》一章的基本概念与结果；参见
《单纯对象》第 \ref{simplicial-section-introduction} 节。
特别地，对于单纯对象，我们仍记作 $X$ 而不记作 $X_\bullet$。








\section{单纯拓扑空间}
\label{spaces-simplicial-section-simplicial-top}

\noindent
所谓{\it 单纯空间}，是以拓扑空间为对象、以连续映射为态射的范畴中的
单纯对象。（我们用“单纯代数空间”指代数空间范畴中的单纯对象。）
可以把单纯空间 $X$ 描绘如下：
$$
\xymatrix{
X_2
\ar@<2ex>[r]
\ar@<0ex>[r]
\ar@<-2ex>[r]
&
X_1
\ar@<1ex>[r]
\ar@<-1ex>[r]
\ar@<1ex>[l]
\ar@<-1ex>[l]
&
X_0
\ar@<0ex>[l]
}
$$
这里有两个态射 $d^1_0, d^1_1 : X_1 \to X_0$，以及一个态射
$s^0_0 : X_0 \to X_1$，等等。务须牢记，应把
$d^n_i : X_n \to X_{n - 1}$ 看成“忘掉第 $i$ 个坐标的投影”，
并把 $s^n_j : X_n \to X_{n + 1}$ 看成重复第 $j$ 个坐标的对角映射。

\medskip\noindent
设 $X$ 为单纯空间。按如下方式为 $X$ 配置一个站点
$X_{Zar}$\footnote{此记号与《站点》例
\ref{sites-example-site-topological} 及《拓扑》第
\ref{topologies-definition-big-small-Zariski} 定义中的记号相似。}：
\begin{enumerate}
\item $X_{Zar}$ 的对象是某个 $X_n$ 的开集 $U$，其中 $n$ 为某个整数；
\item $X_{Zar}$ 的态射 $U \to V$ 由一个
$\varphi : [m] \to [n]$ 给出，其中 $n, m$ 满足
$U \subset X_n$、$V \subset X_m$，且 $\varphi$ 满足
$X(\varphi)(U) \subset V$；
\item $X_{Zar}$ 中的覆盖 $\{U_i \to U\}$ 意味着
$U, U_i \subset X_n$ 均为开集，映射 $U_i \to U$ 由
$\text{id} : [n] \to [n]$ 给出，并且 $U = \bigcup U_i$。
\end{enumerate}
特别注意，若 $X_{Zar}$ 的态射 $U \to V$ 由 $\varphi$ 给出，
则 $X(\varphi) : X_n \to X_m$ 确实诱导拓扑空间的连续映射 $U \to V$。

\noindent
显然，上述构造是为任意拓扑空间图式配置站点这一构造的特例。
我们陈述相应的例行引理。

\begin{lemma}
\label{spaces-simplicial-lemma-simplicial-site}
设 $X$ 为单纯空间。则上述定义的 $X_{Zar}$ 是一个站点。
\end{lemma}

\begin{proof}
略。
\end{proof}

\noindent
设 $X$ 为单纯空间，$\mathcal{F}$ 为 $X_{Zar}$ 上的层。
由覆盖的定义可知，把 $\mathcal{F}$ 限制到 $X_n$ 的开集上，
便得到拓扑空间 $X_n$ 上的层 $\mathcal{F}_n$。
对每个 $\varphi : [m] \to [n]$，当 $U \subset X_n$、
$V \subset X_m$ 且 $X(\varphi)(U) \subset V$ 时，
$\mathcal{F}$ 的限制映射定义一个 $X(\varphi)$-映射
$\mathcal{F}(\varphi) : \mathcal{F}_m \to \mathcal{F}_n$；参见
《层》第 \ref{sheaves-definition-f-map} 定义。
此外，给定 $\varphi : [m] \to [n]$ 和 $\psi : [l] \to [m]$，有
$$
\mathcal{F}(\varphi) \circ \mathcal{F}(\psi) =
\mathcal{F}(\varphi \circ \psi)
$$
(左边使用 $f$-映射的复合；参见《层》第
\ref{sheaves-definition-composition-f-maps} 定义）。
显然，逆命题也成立：若有如上的系统
$(\{\mathcal{F}_n\}_{n \geq 0},
\{\mathcal{F}(\varphi)\}_{\varphi \in \text{箭头}(\Delta)})$
满足所示等式，则可得到 $X_{Zar}$ 上的一个层。

\begin{lemma}
\label{spaces-simplicial-lemma-describe-sheaves-simplicial-site}
设 $X$ 为单纯空间。下列两个范畴等价：
\begin{enumerate}
\item $\Sh(X_{Zar})$；
\item 上述系统 $(\mathcal{F}_n, \mathcal{F}(\varphi))$ 的范畴。
\end{enumerate}
\end{lemma}

\begin{proof}
见上文讨论。
\end{proof}

\begin{lemma}
\label{spaces-simplicial-lemma-simplicial-space-site-functorial}
设 $f : Y \to X$ 为单纯空间的态射。函子 $u : X_{Zar} \to Y_{Zar}$
把开集 $U \subset X_n$ 映到开集 $f_n^{-1}(U) \subset Y_n$；
它定义站点态射 $f_{Zar} : Y_{Zar} \to X_{Zar}$。
\end{lemma}

\begin{proof}
显然 $u$ 是连续函子，故得到函子 $f_{Zar, *} = u^s$ 和
$f_{Zar}^{-1} = u_s$；参见《站点》第 \ref{sites-section-morphism-sites} 节。
为证明所得对象是站点态射，须证明 $u_s$ 正合。我们用《站点》引理
\ref{sites-lemma-directed-morphism} 来证明这一点。
设 $V \subset Y_n$ 为开子集。范畴 $\mathcal{I}_V^u$（参见《站点》第
\ref{sites-section-functoriality-PSh} 节）的对象是二元组 $(U, \varphi)$，
其中 $\varphi : [m] \to [n]$，$U \subset X_m$ 为开集，且
$Y(\varphi)(V) \subset f_m^{-1}(U)$。态射
$(U, \varphi) \to (U', \varphi')$ 则由满足
$X(\psi)(U) \subset U'$ 和 $\varphi \circ \psi = \varphi'$ 的
$\psi : [m'] \to [m]$ 给出。我们要证明 $\mathcal{I}_V^u$ 是余滤范畴。

\medskip\noindent
验证《范畴》第 \ref{categories-definition-codirected} 定义中的条件。
由于 $(X_n, \text{id}_{[n]})$ 是对象，条件 (1) 成立。
设 $(U, \varphi)$ 为对象。条件
$Y(\varphi)(V) \subset f_m^{-1}(U)$ 等价于
$V \subset f_n^{-1}(X(\varphi)^{-1}(U))$。因此，取
$\psi = \varphi$ 便得到态射
$(X(\varphi)^{-1}(U), \text{id}_{[n]}) \to (U, \varphi)$。
再者，对任意两个形如 $(U, \text{id}_{[n]})$ 和
$(U', \text{id}_{[n]})$ 的对象，对象
$(U \cap U', \text{id}_{[n]})$ 同时映到二者，故条件 (2) 成立。
为验证条件 (3)，设有两个由 $\psi, \psi' : [m'] \to [m]$ 给出的态射
$a, a': (U, \varphi) \to (U', \varphi')$。
把二者预复合以由 $\varphi$ 给出的态射
$(X(\varphi)^{-1}(U), \text{id}_{[n]}) \to (U, \varphi)$，
便使 $a,a'$ 相等，因为
$\varphi \circ \psi = \varphi' = \varphi \circ \psi'$。
证毕。
\end{proof}

\begin{lemma}
\label{spaces-simplicial-lemma-describe-functoriality}
设 $f : Y \to X$ 为单纯空间的态射。用引理
\ref{spaces-simplicial-lemma-describe-sheaves-simplicial-site} 对层的描述，
引理 \ref{spaces-simplicial-lemma-simplicial-space-site-functorial} 中的态射
$f_{Zar}$ 可描述如下：
\begin{enumerate}
\item 若 $\mathcal{G}$ 是 $Y$ 上的层，则
$(f_{Zar, *}\mathcal{G})_n = f_{n, *}\mathcal{G}_n$。
\item 若 $\mathcal{F}$ 是 $X$ 上的层，则
$(f_{Zar}^{-1}\mathcal{F})_n = f_n^{-1}\mathcal{F}_n$。
\end{enumerate}
\end{lemma}

\begin{proof}
第一部分由定义立即得到。对于第二部分，注意在引理
\ref{spaces-simplicial-lemma-simplicial-space-site-functorial} 的证明中已经证明：
对开集 $V \subset Y_n$，范畴 $(\mathcal{I}_V^u)^{opp}$ 包含一个余终子范畴，
其对象是满足 $f_n^{-1}(U) \supset V$ 的开集 $U \subset X_n$，
态射由包含映射给出。因此，把 $u_p\mathcal{F}$ 限制到 $Y_n$ 的开集上，
所得预层就是《层》引理 \ref{sheaves-lemma-pullback-presheaves}
中定义的 $f_{n, p}\mathcal{F}_n$。
由于 $f_{Zar}^{-1}\mathcal{F} = u_s\mathcal{F}$ 是
$u_p\mathcal{F}$ 的层化，层化只使用覆盖，而 $Y_{Zar}$ 中的覆盖只使用
同一个 $Y_n$ 上开集之间的包含映射，结论便由
$f_n^{-1}\mathcal{F}_n$ 相应地是 $f_{n, p}\mathcal{F}_n$ 的层化而得；
参见《层》第 \ref{sheaves-section-presheaves-functorial} 节。
\end{proof}

\noindent
设 $X$ 为拓扑空间。在《站点》例
\ref{sites-example-site-topological} 中，我们用 $X_{Zar}$ 表示如下站点：
对象为 $X$ 的开集，态射为包含映射，覆盖为开覆盖。
我们把拓扑斯 $\Sh(X_{Zar})$ 与 $X$ 上的层范畴等同。

\begin{lemma}
\label{spaces-simplicial-lemma-restriction-to-components}
设 $X$ 为单纯空间。函子
$X_{n, Zar} \to X_{Zar}$，$U \mapsto U$，既连续又余连续。
相应的拓扑斯态射 $g_n : \Sh(X_n) \to \Sh(X_{Zar})$ 满足：
\begin{enumerate}
\item $g_n^{-1}$ 把 $X$ 上的层 $\mathcal{F}$ 映到 $X_n$ 上的层
$\mathcal{F}_n$；
\item $g_n^{-1} : \Sh(X_{Zar}) \to \Sh(X_n)$ 有左伴随 $g^{Sh}_{n!}$；
\item $g^{Sh}_{n!}$ 与有限连通极限交换；
\item $g_n^{-1} : \textit{Ab}(X_{Zar}) \to \textit{Ab}(X_n)$
有左伴随 $g_{n!}$；
\item $g_{n!}$ 正合。
\end{enumerate}
\end{lemma}

\begin{proof}
除命题中所述性质外，范畴 $X_{n, Zar}$ 还有纤维积与等化子，
且该函子与它们交换（注意 $X_{Zar}$ 并非具有所有纤维积）。
因此，本引理由《站点》第 \ref{sites-section-cocontinuous-functors}、
\ref{sites-section-cocontinuous-morphism-topoi}
节以及《站点上的模》第
\ref{sites-modules-section-exactness-lower-shriek} 节的讨论得到。
更确切地说，使用《站点》引理
\ref{sites-lemma-cocontinuous-morphism-topoi}、
\ref{sites-lemma-when-shriek}、
\ref{sites-lemma-preserve-equalizers}，以及《站点上的模》引理
\ref{sites-modules-lemma-g-shriek-adjoint}、
\ref{sites-modules-lemma-exactness-lower-shriek}。
\end{proof}

\begin{lemma}
\label{spaces-simplicial-lemma-restriction-injective-to-component}
设 $X$ 为单纯空间。若 $\mathcal{I}$ 是 $X_{Zar}$ 上的内射阿贝尔层，
则 $\mathcal{I}_n$ 是 $X_n$ 上的内射阿贝尔层。
\end{lemma}

\begin{proof}
这由《同调代数》引理 \ref{homology-lemma-adjoint-preserve-injectives}
和引理 \ref{spaces-simplicial-lemma-restriction-to-components} 得到。
\end{proof}

\begin{lemma}
\label{spaces-simplicial-lemma-restriction-to-components-functorial}
设 $f : Y \to X$ 为单纯空间的态射。则
$$
\xymatrix{
\Sh(Y_n) \ar[d] \ar[r]_{f_n} & \Sh(X_n) \ar[d] \\
\Sh(Y_{Zar}) \ar[r]^{f_{Zar}} & \Sh(X_{Zar})
}
$$
是拓扑斯的交换图。
\end{lemma}

\begin{proof}
直接由引理 \ref{spaces-simplicial-lemma-describe-functoriality} 和
\ref{spaces-simplicial-lemma-restriction-to-components} 中对拉回函子的描述得到。
\end{proof}

\begin{lemma}
\label{spaces-simplicial-lemma-augmentation}
设 $Y$ 为单纯空间，$a : Y \to X$ 为一个增广
（《单纯对象》定义 \ref{simplicial-definition-augmentation}）。
设 $a_n : Y_n \to X$ 为相应的拓扑空间态射。存在典范拓扑斯态射
$$
a : \Sh(Y_{Zar}) \to \Sh(X)
$$
满足下列性质：
\begin{enumerate}
\item $a^{-1}\mathcal{F}$ 限制到 $Y_n$ 后是 $a_n^{-1}\mathcal{F}$；
\item 对所有 $\varphi : [m] \to [n]$，有
$a_m \circ Y(\varphi) = a_n$；
\item 作为拓扑斯态射有 $a \circ g_n = a_n$，其中 $g_n$ 如引理
\ref{spaces-simplicial-lemma-restriction-to-components}；
\item 对 $\mathcal{G} \in \Sh(Y_{Zar})$，$a_*\mathcal{G}$ 是两个映射
$a_{0, *}\mathcal{G}_0 \to a_{1, *}\mathcal{G}_1$ 的等化子。
\end{enumerate}
\end{lemma}

\begin{proof}
在任意范畴中，单纯对象的增广均满足 (2)。因此
$Y(\varphi)^{-1} a_m^{-1} \mathcal{F} = a_n^{-1}\mathcal{F}$，
这定义一个从 $a_m^{-1}\mathcal{F}$ 到 $a_n^{-1}\mathcal{F}$ 的
$Y(\varphi)$-映射。于是可用 (1) 定义 $a^{-1}\mathcal{F}$
（使用引理 \ref{spaces-simplicial-lemma-describe-sheaves-simplicial-site}），并用 (4)
定义 $a_*$。若这定义拓扑斯态射，则由构造有
$g_n^{-1} \circ a^{-1} = a_n^{-1}$，从而 (3) 成立。
为验证 $a$ 是拓扑斯态射，须证明 $a^{-1}$ 左伴随于 $a_*$，并证明
$a^{-1}$ 正合。后一事实由各函子 $a_n^{-1}$ 的正合性立即得到。

\medskip\noindent
设 $\mathcal{F}$ 是 $\Sh(X)$ 的对象，$\mathcal{G}$ 是
$\Sh(Y_{Zar})$ 的对象。给定
$\beta : a^{-1}\mathcal{F} \to \mathcal{G}$，考察其分量
$\beta_n : a_n^{-1}\mathcal{F} \to \mathcal{G}_n$。
这些映射伴随于映射
$\beta_n : \mathcal{F} \to a_{n, *}\mathcal{G}_n$。
与单纯结构的相容性说明 $\beta_0$ 的像落在 $a_*\mathcal{G}$ 中。
反之，设给定映射 $\alpha : \mathcal{F} \to a_*\mathcal{G}$。
对任意 $n$，选择 $\varphi : [0] \to [n]$，并考察复合
$$
\mathcal{F} \xrightarrow{\alpha} a_*\mathcal{G}
\to a_{0, *}\mathcal{G}_0 \xrightarrow{\mathcal{G}(\varphi)}
a_{n, *}\mathcal{G}_n
$$
这些复合伴随于映射 $a_n^{-1}\mathcal{F} \to \mathcal{G}_n$，
后者定义层态射 $a^{-1}\mathcal{F} \to \mathcal{G}$。
略去验证上述构造给出互逆双射
$$
\Mor_{\Sh(Y_{Zar})}(a^{-1}\mathcal{F}, \mathcal{G}) =
\Mor_{\Sh(X)}(\mathcal{F}, a_*\mathcal{G})
$$
由此证毕。值得注意的是，这个拓扑斯态射并不对应于定义这些拓扑斯的
站点之间任何显然的几何函子。
\end{proof}

\begin{lemma}
\label{spaces-simplicial-lemma-simplicial-resolution-Z}
设 $X$ 为单纯拓扑空间。$X_{Zar}$ 上的阿贝尔预层复形
$$
\ldots \to \mathbf{Z}_{X_2} \to \mathbf{Z}_{X_1} \to \mathbf{Z}_{X_0}
$$
以 $\sum (-1)^i d^n_i$ 为边界；它是常值预层 $\mathbf{Z}$ 的消解。
\end{lemma}

\begin{proof}
设 $U \subset X_m$ 为 $X_{Zar}$ 的对象。上述复形在 $U$ 上的值是
阿贝尔群复形
$$
\ldots \to
\mathbf{Z}[\Mor_\Delta([2], [m])] \to
\mathbf{Z}[\Mor_\Delta([1], [m])] \to
\mathbf{Z}[\Mor_\Delta([0], [m])]
$$
换言之，这是单纯集 $\Delta[m]$ 上自由阿贝尔群所伴随的复形；参见
《单纯对象》例 \ref{simplicial-example-simplex-simplicial-set}。
由于 $\Delta[m]$ 与 $\Delta[0]$ 同伦等价（参见《单纯对象》例
\ref{simplicial-example-simplex-contractible}），且“取自由阿贝尔群”是函子，
上述复形与 $\Delta[0]$ 上的自由阿贝尔群同伦等价
（《单纯对象》评注 \ref{simplicial-remark-homotopy-better} 及引理
\ref{simplicial-lemma-homotopy-equivalence-s-N}）。
该复形在正次数无环，在次数 $0$ 等于 $\mathbf{Z}$。
\end{proof}

\begin{lemma}
\label{spaces-simplicial-lemma-simplicial-sheaf-cohomology}
设 $X$ 为单纯拓扑空间，$\mathcal{F}$ 为 $X$ 上的阿贝尔层。
存在谱序列 $(E_r, d_r)_{r \geq 0}$，其中
$$
E_1^{p, q} = H^q(X_p, \mathcal{F}_p)
$$
它收敛到 $H^{p + q}(X_{Zar}, \mathcal{F})$，并且关于
$\mathcal{F}$ 具有函子性。
\end{lemma}

\begin{proof}
设 $\mathcal{F} \to \mathcal{I}^\bullet$ 为内射消解。考察项为
$$
A^{p, q} = \mathcal{I}^q(X_p)
$$
的双复形，其第一微分是沿映射 $d^{p + 1}_i$ 所给映射
$\mathcal{I}_p^q \to \mathcal{I}_{p + 1}^q$ 的交错和；参见引理
\ref{spaces-simplicial-lemma-describe-sheaves-simplicial-site}。注意
$$
A^{p, q} = \Gamma(X_p, \mathcal{I}_p^q) =
\Mor_{\textit{PSh}}(h_{X_p}, \mathcal{I}^q) =
\Mor_{\textit{PAb}}(\mathbf{Z}_{X_p}, \mathcal{I}^q)
$$
因此，由引理 \ref{spaces-simplicial-lemma-simplicial-resolution-Z} 和《站点上的上同调》引理
\ref{sites-cohomology-lemma-injective-abelian-sheaf-injective-presheaf}，
双复形的各行在正次数正合，并在次数 $0$ 给出
$\Gamma(X_{Zar}, \mathcal{I}^q)$。另一方面，由于限制函子正合
（引理 \ref{spaces-simplicial-lemma-restriction-to-components}），映射
$$
\mathcal{F}_p \to \mathcal{I}_p^\bullet
$$
是消解。各层 $\mathcal{I}_p^q$ 是 $X_p$ 上的内射阿贝尔层
（引理 \ref{spaces-simplicial-lemma-restriction-injective-to-component}）。
故各列的上同调计算群 $H^q(X_p, \mathcal{F}_p)$。
最后应用《同调代数》引理 \ref{homology-lemma-first-quadrant-ss} 和
\ref{homology-lemma-double-complex-gives-resolution} 即得结论。
\end{proof}

\begin{lemma}
\label{spaces-simplicial-lemma-augmentation-pushforward-higher-direct-image}
设 $X$ 为单纯空间，$a : X \to Y$ 为增广，并设 $\mathcal{F}$ 为
$X_{Zar}$ 上的阿贝尔层。则 $R^na_*\mathcal{F}$ 是预层
$$
V \longmapsto H^n((X \times_Y V)_{Zar}, \mathcal{F}|_{(X \times_Y V)_{Zar}})
$$
所伴随的层。
\end{lemma}

\begin{proof}
这是《上同调》引理 \ref{cohomology-lemma-describe-higher-direct-images}
或《站点上的上同调》引理
\ref{sites-cohomology-lemma-higher-direct-images} 的类似结果；
我们强烈建议读者跳过此证明。选取 $X_{Zar}$ 上
$\mathcal{F}$ 的内射消解并使用定义，只须证明：
(1) 把 $X_{Zar}$ 上的内射阿贝尔层限制到
$(X \times_Y V)_{Zar}$ 后仍是内射阿贝尔层；以及
(2) $a_*\mathcal{F}$ 等于规则
$$
V \longmapsto H^0((X \times_Y V)_{Zar}, \mathcal{F}|_{(X \times_Y V)_{Zar}})
$$
(2) 由下列事实得到：
\begin{enumerate}
\item[(2a)] 由引理 \ref{spaces-simplicial-lemma-augmentation}，$a_*\mathcal{F}$ 是两个映射
$a_{0, *}\mathcal{F}_0 \to a_{1, *}\mathcal{F}_1$
的等化子；
\item[(2b)]
$a_{0, *}\mathcal{F}_0(V) =
H^0(a_0^{-1}(V), \mathcal{F}_0)$ 且
$a_{1, *}\mathcal{F}_1(V) = H^0(a_1^{-1}(V), \mathcal{F}_1)$；
\item[(2c)] $X_0 \times_Y V = a_0^{-1}(V)$ 且
$X_1 \times_Y V = a_1^{-1}(V)$；
\item[(2d)]
$H^0((X \times_Y V)_{Zar}, \mathcal{F}|_{(X \times_Y V)_{Zar}})$
是两个映射
$H^0(X_0 \times_Y V, \mathcal{F}_0) \to H^0(X_1 \times_Y V, \mathcal{F}_1)$
的等化子，例如可由引理 \ref{spaces-simplicial-lemma-simplicial-sheaf-cohomology} 得到。
\end{enumerate}
定义限制函子
$\textit{Ab}(X_{Zar}) \to \textit{Ab}((X \times_Y V)_{Zar})$
的正合左伴随
$j_! : \textit{Ab}((X \times_Y V)_{Zar}) \to \textit{Ab}(X_{Zar})$
（零延拓），再应用《同调代数》引理
\ref{homology-lemma-adjoint-preserve-injectives}，即可得到 (1)。
\end{proof}

\noindent
设 $X$ 为拓扑空间。用 $X_\bullet$ 表示取常值 $X$ 的常值单纯拓扑空间。
由引理 \ref{spaces-simplicial-lemma-describe-sheaves-simplicial-site}，
$X_{\bullet, Zar}$ 上的层等同于 $X$ 上层范畴中的余单纯对象。

\begin{lemma}
\label{spaces-simplicial-lemma-constant-simplicial-space}
设 $X$ 为拓扑空间，$X_\bullet$ 为取常值 $X$ 的常值单纯拓扑空间。函子
$$
X_{\bullet, Zar} \longrightarrow X_{Zar},\quad
U \longmapsto U
$$
既连续又余连续，并定义拓扑斯态射
$g : \Sh(X_{\bullet, Zar}) \to \Sh(X)$，且 $g^{-1}$ 有左伴随 $g_!$。
有：
\begin{enumerate}
\item $g^{-1}$ 把 $X$ 上的层映到 $X$ 上的常值余单纯层；
\item $g_!$ 把 $X_{\bullet, Zar}$ 上的层 $\mathcal{F}$ 映到层
$\mathcal{F}_0$；
\item $g_*$ 把 $X_{\bullet, Zar}$ 上的层 $\mathcal{F}$ 映到两个映射
$\mathcal{F}_0 \to \mathcal{F}_1$ 的等化子。
\end{enumerate}
\end{lemma}

\begin{proof}
关于该函子的陈述可直接验证。$g$ 与 $g_!$ 的存在性由《站点》引理
\ref{sites-lemma-cocontinuous-morphism-topoi} 和
\ref{sites-lemma-when-shriek} 得到。$g^{-1}$ 的描述直接来自《站点》引理
\ref{sites-lemma-when-shriek}。所给函子分别是 $g^{-1}$ 的右、左伴随，
因而得到 $g_*$ 与 $g_!$ 的描述。
\end{proof}








\section{单纯站点与拓扑斯}
\label{spaces-simplicial-section-simplicial-sites}

\noindent
自然的做法似乎是把{\it 单纯站点}定义为如下（大）范畴中的单纯对象：
其对象是站点，态射是站点态射。参见《站点》定义
\ref{sites-definition-site}、\ref{sites-definition-morphism-sites}；
态射的复合如《站点》引理 \ref{sites-lemma-composition-morphisms-sites}。
不过，也可考虑下列变体：(a) 改用站点之间的余连续函子
（参见《站点》第 \ref{sites-section-cocontinuous-functors}、
\ref{sites-section-cocontinuous-morphism-topoi} 节）；
(b) 在适当的站点 $2$-范畴中工作，并引入站点态射之间的
$2$-态射概念；(c) 在由余连续函子构造的 $2$-范畴中工作。
我们不选取其中一种变体作为定义，而只按需要发展理论。

\medskip\noindent
当然，{\it 单纯拓扑斯}大概应定义为从 $\Delta^{opp}$ 到拓扑斯
$2$-范畴的伪函子。参见《范畴》定义
\ref{categories-definition-functor-into-2-category} 以及《站点》第
\ref{sites-section-topoi}、\ref{sites-section-2-category} 节。
我们将尽量避免直接处理这种庞然大物。

\medskip\noindent
{\bf 情形 A。}
设 $\mathcal{C}$ 是如下范畴中的单纯对象：其对象为站点，态射为站点态射。
这意味着，对 $\Delta$ 的每个态射 $\varphi : [m] \to [n]$，
都有站点态射 $f_\varphi : \mathcal{C}_n \to \mathcal{C}_m$。
该态射由反方向的连续函子给出，记为
$u_\varphi : \mathcal{C}_m \to \mathcal{C}_n$。

\begin{lemma}
\label{spaces-simplicial-lemma-simplicial-site-site}
设 $\mathcal{C}$ 为站点范畴中的单纯对象。沿用上述记号，按如下方式构造站点
$\mathcal{C}_{total}$：
\begin{enumerate}
\item $\mathcal{C}_{total}$ 的对象是某个 $\mathcal{C}_n$ 的对象 $U$，
其中 $n$ 为某个整数；
\item $\mathcal{C}_{total}$ 的态射 $(\varphi, f) : U \to V$ 由映射
$\varphi : [m] \to [n]$ 以及 $\mathcal{C}_n$ 中的态射
$f : U \to u_\varphi(V)$ 给出，其中
$U \in \Ob(\mathcal{C}_n)$、$V \in \Ob(\mathcal{C}_m)$；
\item $\mathcal{C}_{total}$ 中的覆盖
$\{(\text{id}, f_i) :  U_i \to U\}$ 由某个 $n$ 以及
$\mathcal{C}_n$ 中的覆盖 $\{f_i : U_i \to U\}$ 给出。
\end{enumerate}
\end{lemma}

\begin{proof}
$(\varphi, f) : U \to V$ 与 $(\psi, g) : V \to W$ 的复合为
$(\varphi \circ \psi, u_\varphi(g) \circ f)$。
这里使用了 $u_\varphi \circ u_\psi = u_{\varphi \circ \psi}$。

\medskip\noindent
设 $\{(\text{id}, f_i) :  U_i \to U\}$ 为 (3) 中的覆盖，
$(\varphi, g) : W \to U$ 为态射，其中
$W \in \Ob(\mathcal{C}_m)$。我们断言
$$
W \times_{(\varphi, g), U, (\text{id}, f_i)} U_i =
W \times_{g, u_\varphi(U), u_\varphi(f_i)} u_\varphi(U_i)
$$
在范畴 $\mathcal{C}_{total}$ 中成立。此式确有意义：按我们的站点态射定义，
$u_\varphi$ 把覆盖变成覆盖，故所需的纤维积在 $\mathcal{C}_m$ 中存在。
同样的推理证明该断言（略去细节）。因此覆盖族在基变换下稳定；
站点的其他公理立即成立。
\end{proof}

\noindent
{\bf 情形 B。}
设 $\mathcal{C}$ 是如下范畴中的单纯对象：其对象为站点，态射为余连续函子。
这意味着，对 $\Delta$ 的每个态射 $\varphi : [m] \to [n]$，
都有余连续函子 $u_\varphi : \mathcal{C}_n \to \mathcal{C}_m$。
相应的拓扑斯态射记为
$f_\varphi : \Sh(\mathcal{C}_n) \to \Sh(\mathcal{C}_m)$。

\begin{lemma}
\label{spaces-simplicial-lemma-simplicial-cocontinuous-site}
设 $\mathcal{C}$ 为以站点为对象、以余连续函子为态射的范畴中的单纯对象。
沿用上述记号，假设函子
$u_\varphi : \mathcal{C}_n \to \mathcal{C}_m$ 具有《站点》评注
\ref{sites-remark-cartesian-cocontinuous} 的性质 $P$。
则可按如下方式构造站点 $\mathcal{C}_{total}$：
\begin{enumerate}
\item $\mathcal{C}_{total}$ 的对象是某个 $\mathcal{C}_n$ 的对象 $U$，
其中 $n$ 为某个整数；
\item $\mathcal{C}_{total}$ 的态射 $(\varphi, f) : U \to V$ 由映射
$\varphi : [m] \to [n]$ 以及 $\mathcal{C}_m$ 中的态射
$f : u_\varphi(U) \to V$ 给出，其中
$U \in \Ob(\mathcal{C}_n)$、$V \in \Ob(\mathcal{C}_m)$；
\item $\mathcal{C}_{total}$ 中的覆盖
$\{(\text{id}, f_i) :  U_i \to U\}$ 由某个 $n$ 以及
$\mathcal{C}_n$ 中的覆盖 $\{f_i : U_i \to U\}$ 给出。
\end{enumerate}
\end{lemma}

\begin{proof}
$(\varphi, f) : U \to V$ 与 $(\psi, g) : V \to W$ 的复合为
$(\varphi \circ \psi, g \circ u_\psi(f))$。
这里使用了 $u_\psi \circ u_\varphi = u_{\varphi \circ \psi}$。

\medskip\noindent
设 $\{(\text{id}, f_i) :  U_i \to U\}$ 为 (3) 中的覆盖，
$(\varphi, g) : W \to U$ 为态射，其中
$W \in \Ob(\mathcal{C}_m)$。我们断言
$$
W \times_{(\varphi, g), U, (\text{id}, f_i)} U_i =
W \times_{g, U, f_i} U_i
$$
在范畴 $\mathcal{C}_{total}$ 中成立；右边是《站点》评注
\ref{sites-remark-cartesian-cocontinuous} 所定义的
$\mathcal{C}_m$ 的对象，并由性质 $P$ 保证存在。
这类纤维积与函子复合的相容性推出该断言（略去细节）。
由性质 $P$，族 $\{W \times_{g, U, f_i} U_i \to W\}$ 是
$\mathcal{C}_m$ 的覆盖，故覆盖族在基变换下稳定；
站点的其他公理立即成立。
\end{proof}

\begin{situation}
\label{spaces-simplicial-situation-simplicial-site}
此处考虑下列两种情形之一：
\begin{enumerate}
\item[(A)] $\mathcal{C}$ 是以站点为对象、以站点态射为态射的范畴中的
单纯对象。对 $\Delta$ 的每个态射 $\varphi : [m] \to [n]$，
都有由连续函子 $u_\varphi : \mathcal{C}_m \to \mathcal{C}_n$ 给出的
站点态射 $f_\varphi : \mathcal{C}_n \to \mathcal{C}_m$。
\item[(B)] $\mathcal{C}$ 是以站点为对象、以满足《站点》评注
\ref{sites-remark-cartesian-cocontinuous} 的性质 $P$ 的余连续函子为态射的
范畴中的单纯对象。对 $\Delta$ 的每个态射
$\varphi : [m] \to [n]$，都有余连续函子
$u_\varphi : \mathcal{C}_n \to \mathcal{C}_m$，它诱导拓扑斯态射
$f_\varphi : \Sh(\mathcal{C}_n) \to \Sh(\mathcal{C}_m)$。
\end{enumerate}
照例，用 $f_\varphi^{-1}$ 和 $f_{\varphi, *}$ 分别表示拉回与推前。
用 $\mathcal{C}_{total}$ 表示引理 \ref{spaces-simplicial-lemma-simplicial-site-site}
（情形 A）或引理 \ref{spaces-simplicial-lemma-simplicial-cocontinuous-site}
（情形 B）中定义的站点。
\end{situation}

\noindent
设 $\mathcal{C}$ 如情形 \ref{spaces-simplicial-situation-simplicial-site}，
$\mathcal{F}$ 为 $\mathcal{C}_{total}$ 上的层。
由覆盖的定义可知，把 $\mathcal{F}$ 限制到 $\mathcal{C}_n$ 的对象上，
便定义站点 $\mathcal{C}_n$ 上的层 $\mathcal{F}_n$。
对每个 $\varphi : [m] \to [n]$，$\mathcal{F}$ 沿态射
$(\varphi, f) : U \to V$ 的限制映射定义下列集合中的元素
$\mathcal{F}(\varphi)$，其中 $U \in \Ob(\mathcal{C}_n)$、
$V \in \Ob(\mathcal{C}_m)$：
$$
\Mor_{\Sh(\mathcal{C}_m)}(\mathcal{F}_m, f_{\varphi, *}\mathcal{F}_n) =
\Mor_{\Sh(\mathcal{C}_n)}(f_\varphi^{-1}\mathcal{F}_m, \mathcal{F}_n)
$$
此外，给定 $\varphi : [m] \to [n]$ 和 $\psi : [l] \to [m]$，下列图交换：
$$
\vcenter{
\xymatrix{
\mathcal{F}_l \ar[rr]_{\mathcal{F}(\varphi \circ \psi)}
\ar[rd]_{\mathcal{F}(\psi)}
& & f_{\varphi \circ \psi, *} \mathcal{F}_n \\
& f_{\psi, *}\mathcal{F}_m \ar[ur]_{f_{\psi, *}\mathcal{F}(\varphi)}
}
}
\quad\text{且}\quad
\vcenter{
\xymatrix{
f_{\varphi \circ \psi}^{-1}\mathcal{F}_l
\ar[rr]_{\mathcal{F}(\varphi \circ \psi)}
\ar[rd]_{f_\varphi^{-1}\mathcal{F}(\psi)}
& & \mathcal{F}_n \\
& f_\varphi^{-1}\mathcal{F}_m \ar[ur]_{\mathcal{F}(\varphi)}
}
}
$$
显然，逆命题也成立：若有系统
$(\{\mathcal{F}_n\}_{n \geq 0},
\{\mathcal{F}(\varphi)\}_{\varphi \in \text{箭头}(\Delta)})$
满足上述交换性约束，则可得到 $\mathcal{C}_{total}$ 上的一个层。

\begin{lemma}
\label{spaces-simplicial-lemma-describe-sheaves-simplicial-site-site}
在情形 \ref{spaces-simplicial-situation-simplicial-site} 中，下列两个范畴等价：
\begin{enumerate}
\item $\Sh(\mathcal{C}_{total})$；
\item 上述系统 $(\mathcal{F}_n, \mathcal{F}(\varphi))$ 的范畴。
\end{enumerate}
特别地，拓扑斯 $\Sh(\mathcal{C}_{total})$ 仅依赖于各拓扑斯
$\Sh(\mathcal{C}_n)$ 及拓扑斯态射 $f_\varphi$。
\end{lemma}

\begin{proof}
见上文讨论。
\end{proof}

\begin{lemma}
\label{spaces-simplicial-lemma-restriction-to-components-site}
在情形 \ref{spaces-simplicial-situation-simplicial-site} 中，函子
$\mathcal{C}_n \to \mathcal{C}_{total}$，$U \mapsto U$，
既连续又余连续。相应的拓扑斯态射
$g_n : \Sh(\mathcal{C}_n) \to \Sh(\mathcal{C}_{total})$ 满足：
\begin{enumerate}
\item $g_n^{-1}$ 把 $\mathcal{C}_{total}$ 上的层 $\mathcal{F}$ 映到
$\mathcal{C}_n$ 上的层 $\mathcal{F}_n$；
\item $g_n^{-1} : \Sh(\mathcal{C}_{total}) \to \Sh(\mathcal{C}_n)$
有左伴随 $g^{Sh}_{n!}$；
\item 对 $\Sh(\mathcal{C}_n)$ 中的 $\mathcal{G}$，
$g_{n!}^{Sh}\mathcal{G}$ 在 $\mathcal{C}_m$ 上的限制为
$\coprod\nolimits_{\varphi : [n] \to [m]} f_\varphi^{-1}\mathcal{G}$；
\item $g_{n!}^{Sh}$ 与有限连通极限交换；
\item $g_n^{-1} : \textit{Ab}(\mathcal{C}_{total}) \to
\textit{Ab}(\mathcal{C}_n)$ 有左伴随 $g_{n!}$；
\item 对 $\textit{Ab}(\mathcal{C}_n)$ 中的 $\mathcal{G}$，
$g_{n!}\mathcal{G}$ 在 $\mathcal{C}_m$ 上的限制为
$\bigoplus\nolimits_{\varphi : [n] \to [m]} f_\varphi^{-1}\mathcal{G}$；
\item $g_{n!}$ 正合。
\end{enumerate}
\end{lemma}

\begin{proof}
情形 A。若 $\{U_i \to U\}_{i \in I}$ 是 $\mathcal{C}_n$ 中的覆盖，
则按定义（引理 \ref{spaces-simplicial-lemma-simplicial-site-site}），其像
$\{U_i \to U\}_{i \in I}$ 是 $\mathcal{C}_{total}$ 中的覆盖。
对 $\mathcal{C}_n$ 的态射 $V \to U$，$\mathcal{C}_n$ 中的纤维积
$V \times_U U_i$ 也是 $\mathcal{C}_{total}$ 中的纤维积
（由引理 \ref{spaces-simplicial-lemma-simplicial-site-site} 证明中的断言）。
因此该函子连续。另一方面，该函子在 $\mathcal{C}_n$ 中 $U$ 的覆盖
与 $\mathcal{C}_{total}$ 中 $U$ 的覆盖之间给出双射，故它也余连续。

\medskip\noindent
情形 B。若 $\{U_i \to U\}_{i \in I}$ 是 $\mathcal{C}_n$ 中的覆盖，
则按定义（引理 \ref{spaces-simplicial-lemma-simplicial-cocontinuous-site}），其像
$\{U_i \to U\}_{i \in I}$ 是 $\mathcal{C}_{total}$ 中的覆盖。
对 $\mathcal{C}_n$ 的态射 $V \to U$，$\mathcal{C}_n$ 中的纤维积
$V \times_U U_i$ 也是 $\mathcal{C}_{total}$ 中的纤维积
（由引理 \ref{spaces-simplicial-lemma-simplicial-cocontinuous-site} 证明中的断言）。
因此该函子连续。另一方面，该函子在 $\mathcal{C}_n$ 中 $U$ 的覆盖
与 $\mathcal{C}_{total}$ 中 $U$ 的覆盖之间给出双射，故它也余连续。

\medskip\noindent
至此，在情形 A 与 B 中，(1) 以及 $g^{Sh}_{n!}$、$g_{n!}$ 的存在性
由《站点》引理 \ref{sites-lemma-cocontinuous-morphism-topoi}、
\ref{sites-lemma-when-shriek} 和《站点上的模》引理
\ref{sites-modules-lemma-g-shriek-adjoint} 得到。

\medskip\noindent
(3) 的证明。设 $\mathcal{G}$ 为 $\mathcal{C}_n$ 上的层。
考察 $\mathcal{C}_{total}$ 上的层 $\mathcal{H}$，其次数 $m$ 部分是层
$$
\mathcal{H}_m = \coprod\nolimits_{\varphi : [n] \to [m]}
f_\varphi^{-1}\mathcal{G}
$$
即引理陈述 (3) 中给出的层。给定映射 $\psi : [m] \to [m']$，
映射 $\mathcal{H}(\psi) : f_\psi^{-1}\mathcal{H}_m \to \mathcal{H}_{m'}$
在各分量上由下列等同给出：
$$
f_\psi^{-1} f_\varphi^{-1} \mathcal{G} \to
f_{\psi \circ \varphi}^{-1}\mathcal{G}
$$
注意，给定 $\mathcal{C}_{total}$ 上层的映射
$\alpha : \mathcal{H} \to \mathcal{F}$，可得到映射
$\mathcal{G} \to \mathcal{F}_n$；它对应于把 $\alpha_n$ 限制到
$\mathcal{H}_n$ 中的分量 $\mathcal{G}$。反之，给定
$\mathcal{C}_n$ 上层的映射 $\beta : \mathcal{G} \to \mathcal{F}_n$，
可定义 $\alpha : \mathcal{H} \to \mathcal{F}$：令 $\alpha_m$
在各分量
$$
f_\varphi^{-1}\mathcal{G} \to \mathcal{F}_m
$$
上使用伴随于 $\mathcal{F}(\varphi) \circ f_\varphi^{-1}\beta$ 的映射。
略去验证这两个构造给出互逆映射
$$
\Mor_{\Sh(\mathcal{C}_n)}(\mathcal{G}, \mathcal{F}_n) =
\Mor_{\Sh(\mathcal{C}_{total})}(\mathcal{H}, \mathcal{F})
$$
故 $\mathcal{H} = g^{Sh}_{n!}\mathcal{G}$，如所求。

\medskip\noindent
(4) 的证明。若 $\mathcal{G}$ 是 $\mathcal{C}_n$ 上的阿贝尔层，
则完全仿照上文，只是把 $\mathcal{H}$ 定义为
$\mathcal{C}_{total}$ 上的阿贝尔层，其次数 $m$ 部分为
$$
\bigoplus\nolimits_{\varphi : [n] \to [m]} f_\varphi^{-1}\mathcal{G}
$$
转移映射完全如上定义。双射
$$
\Mor_{\textit{Ab}(\mathcal{C}_n)}(\mathcal{G}, \mathcal{F}_n) =
\Mor_{\textit{Ab}(\mathcal{C}_{total})}(\mathcal{H}, \mathcal{F})
$$
的证明与上文完全相同。故 $\mathcal{H} = g_{n!}\mathcal{G}$，如所求。

\medskip\noindent
$g^{Sh}_{n!}$ 与 $g_{n!}$ 的正合性由这些函子的所给公式得到。
\end{proof}

\begin{lemma}
\label{spaces-simplicial-lemma-restriction-injective-to-component-site}
\begin{slogan}
单纯站点上的内射阿贝尔层在每个分量上均内射
\end{slogan}
在情形 \ref{spaces-simplicial-situation-simplicial-site} 中，若 $\mathcal{I}$ 在
$\textit{Ab}(\mathcal{C}_{total})$ 中内射，则 $\mathcal{I}_n$ 在
$\textit{Ab}(\mathcal{C}_n)$ 中内射。若 $\mathcal{I}^\bullet$ 是
$\textit{Ab}(\mathcal{C}_{total})$ 中的 K-内射复形，则
$\mathcal{I}_n^\bullet$ 是 $\textit{Ab}(\mathcal{C}_n)$ 中的 K-内射复形。
\end{lemma}

\begin{proof}
第一条陈述由《同调代数》引理
\ref{homology-lemma-adjoint-preserve-injectives} 和引理
\ref{spaces-simplicial-lemma-restriction-to-components-site} 得到。
第二条陈述由《导出范畴》引理
\ref{derived-lemma-adjoint-preserve-K-injectives} 和引理
\ref{spaces-simplicial-lemma-restriction-to-components-site} 得到。
\end{proof}







\section{单纯站点的增广}
\label{spaces-simplicial-section-augmentation-simplicial-sites}

\noindent
我们继续采用第 \ref{spaces-simplicial-section-simplicial-sites} 节所述方法，
阐明该节情形 A 与 B 中增广的含义。

\begin{remark}
\label{spaces-simplicial-remark-augmentation-site}
在情形 \ref{spaces-simplicial-situation-simplicial-site} 中，所谓{\it 指向站点
$\mathcal{D}$ 的增广 $a_0$}，意指：
\begin{enumerate}
\item[(A)] $a_0 : \mathcal{C}_0 \to \mathcal{D}$ 是由连续函子
$u_0 : \mathcal{D} \to \mathcal{C}_0$ 给出的站点态射，且对所有
$\varphi, \psi : [0] \to [n]$，有
$u_\varphi \circ u_0 = u_\psi \circ u_0$；
\item[(B)] $a_0 : \Sh(\mathcal{C}_0) \to \Sh(\mathcal{D})$ 是由余连续函子
$u_0 : \mathcal{C}_0 \to \mathcal{D}$ 给出的拓扑斯态射，且对所有
$\varphi, \psi : [0] \to [n]$，有
$u_0 \circ u_\varphi = u_0 \circ u_\psi$。
\end{enumerate}
\end{remark}

\begin{lemma}
\label{spaces-simplicial-lemma-augmentation-site}
在情形 \ref{spaces-simplicial-situation-simplicial-site} 中，设 $a_0$ 为评注
\ref{spaces-simplicial-remark-augmentation-site} 所述指向站点 $\mathcal{D}$ 的增广。
则 $a_0$ 诱导：
\begin{enumerate}
\item 对所有 $n \geq 0$，拓扑斯态射
$a_n : \Sh(\mathcal{C}_n) \to \Sh(\mathcal{D})$；
\item 拓扑斯态射 $a : \Sh(\mathcal{C}_{total}) \to \Sh(\mathcal{D})$。
\end{enumerate}
它们满足：
\begin{enumerate}
\item 对所有 $\varphi : [m] \to [n]$，有
$a_m \circ f_\varphi = a_n$；
\item 若 $g_n : \Sh(\mathcal{C}_n) \to \Sh(\mathcal{C}_{total})$
如引理 \ref{spaces-simplicial-lemma-restriction-to-components-site}，则
$a \circ g_n = a_n$；
\item 对 $\mathcal{F} \in \Sh(\mathcal{C}_{total})$，
$a_*\mathcal{F}$ 是两个映射
$a_{0, *}\mathcal{F}_0 \to a_{1, *}\mathcal{F}_1$ 的等化子。
\end{enumerate}
\end{lemma}

\begin{proof}
情形 A。令 $u_n : \mathcal{D} \to \mathcal{C}_n$ 为诸函子
$u_\varphi \circ u_0$（$\varphi : [0] \to [n]$）的共同值。
则 $u_n$ 对应于站点态射 $a_n : \mathcal{C}_n \to \mathcal{D}$；
参见《站点》引理 \ref{sites-lemma-composition-morphisms-sites}。
同一引理说明，对所有 $\varphi : [m] \to [n]$，有
$a_m \circ f_\varphi = a_n$。

\medskip\noindent
情形 B。令 $u_n : \mathcal{C}_n \to \mathcal{D}$ 为诸函子
$u_0 \circ u_\varphi$（$\varphi : [0] \to [n]$）的共同值。
则 $u_n$ 余连续，因而定义拓扑斯态射
% P10-E0172
$a_n : \Sh(\mathcal{C}_n) \to \Sh(\mathcal{D)}$；参见
《站点》引理 \ref{sites-lemma-composition-cocontinuous}。
同一引理说明，对所有 $\varphi : [m] \to [n]$，有
$a_m \circ f_\varphi = a_n$。

\medskip\noindent
考察函子 $a^{-1} : \Sh(\mathcal{D}) \to \Sh(\mathcal{C}_{total})$。
它把集合层 $\mathcal{G}$ 映到层 $\mathcal{F} = a^{-1}\mathcal{G}$；
后者的分量为 $a_n^{-1}\mathcal{G}$，转移映射
$\mathcal{F}(\varphi)$ 为下列等同：
$$
f_\varphi^{-1}\mathcal{F}_m =
f_\varphi^{-1} a_m^{-1}\mathcal{G} =
a_n^{-1}\mathcal{G} =
\mathcal{F}_n
$$
其中 $\varphi : [m] \to [n]$；参见引理
\ref{spaces-simplicial-lemma-describe-sheaves-simplicial-site-site} 对
$\Sh(\mathcal{C}_{total})$ 的描述。由于各函子 $a_n^{-1}$ 正合，
$a^{-1}$ 也是正合函子。最后，把
$a_* : \Sh(\mathcal{C}_{total}) \to \Sh(\mathcal{D})$ 取为如下函子：
% P10-E0173
它把 $\Sh(\mathcal{D})$ 上的层 $\mathcal{F}$ 映到
$$
\xymatrix{
a_*\mathcal{F} \ar@{=}[r] &
\text{等化子}(a_{0, *}\mathcal{F}_0
\ar@<1ex>[r] \ar@<-1ex>[r] &
a_{1, *}\mathcal{F}_1)
}
$$
这里两个映射由两个映射 $\varphi : [0] \to [1]$ 经下式给出：
$$
a_{0, *}\mathcal{F}_0 \to
a_{0, *}f_{\varphi, *} f_\varphi^{-1}\mathcal{F}_0
\xrightarrow{\mathcal{F}(\varphi)}
a_{0, *}f_{\varphi, *} \mathcal{F}_1 = a_{1, *}\mathcal{F}_1
$$
其中第一个箭头来自 $1 \to f_{\varphi, *} f_\varphi^{-1}$。
% P10-E0174
用 $\mathcal{G}_\bullet$ 表示取常值 $\mathcal{G}$ 的常值余单纯层，
并用 $a_{\bullet, *}\mathcal{F}$ 表示次数 $n$ 分量为
$a_{n, *}\mathcal{F}_n$ 的余单纯层。
% P10-E0175
由拓扑斯态射 $a_n$ 的通常伴随关系可知，映射
$a^{-1}\mathcal{G} \to \mathcal{F}$ 等同于余单纯层的映射
$$
\mathcal{G}_\bullet \longrightarrow a_{\bullet, *}\mathcal{F}
$$
由《单纯对象》引理 \ref{simplicial-lemma-augmentation-howto} 的对偶，
这又等同于映射 $\mathcal{G} \to a_*\mathcal{F}$。
故 $a^{-1}$ 与 $a_*$ 互为伴随函子，从而得到拓扑斯态射
$a$\footnote{在情形 B 中，态射 $a$ 对应于余连续函子
$\mathcal{C}_{total} \to \mathcal{D}$；它把
$\mathcal{C}_n$ 中的 $U$ 映到 $u_n(U)$。}。
% P10-E0176
等式 $a \circ g_n = f_n$ 由定义立即得到。
\end{proof}



\section{单纯站点的态射}
\label{spaces-simplicial-section-morphism-simplicial-sites}

\noindent
我们继续采用第 \ref{spaces-simplicial-section-simplicial-sites} 节所述方法，
阐明该节情形 A 与 B 中单纯站点态射的含义。

\begin{remark}
\label{spaces-simplicial-remark-morphism-simplicial-sites}
设 $\mathcal{C}_n, f_\varphi, u_\varphi$ 和
$\mathcal{C}'_n, f'_\varphi, u'_\varphi$ 如情形
\ref{spaces-simplicial-situation-simplicial-site}。所谓{\it 单纯站点之间的态射 $h$}，意指：
\begin{enumerate}
\item[(A)] 站点态射 $h_n : \mathcal{C}_n \to \mathcal{C}'_n$，
使得对所有 $\varphi : [m] \to [n]$，作为站点态射有
$f'_\varphi \circ h_n = h_m \circ f_\varphi$；
\item[(B)] 余连续函子 $v_n : \mathcal{C}_n \to \mathcal{C}'_n$，
它诱导拓扑斯态射
$h_n : \Sh(\mathcal{C}_n) \to \Sh(\mathcal{C}'_n)$，
使得对所有 $\varphi : [m] \to [n]$，作为函子有
$u'_\varphi \circ v_n = v_m \circ u_\varphi$。
\end{enumerate}
在两种情形中，作为拓扑斯态射均有
$f'_\varphi \circ h_n = h_m \circ f_\varphi$；
情形 B 参见《站点》引理
\ref{sites-lemma-composition-cocontinuous}，情形 A 参见《站点》定义
\ref{sites-definition-composition-morphisms-sites}。
\end{remark}

\begin{lemma}
\label{spaces-simplicial-lemma-morphism-simplicial-sites}
设 $\mathcal{C}_n, f_\varphi, u_\varphi$ 和
$\mathcal{C}'_n, f'_\varphi, u'_\varphi$ 如情形
\ref{spaces-simplicial-situation-simplicial-site}，并设 $h$ 为评注
\ref{spaces-simplicial-remark-morphism-simplicial-sites} 所述的单纯站点之间的态射。
则得到拓扑斯态射
$$
h_{total} : \Sh(\mathcal{C}_{total}) \to \Sh(\mathcal{C}'_{total})
$$
以及交换图
$$
\xymatrix{
\Sh(\mathcal{C}_n) \ar[d]_{g_n} \ar[r]_{h_n} &
\Sh(\mathcal{C}'_n) \ar[d]^{g'_n} \\
\Sh(\mathcal{C}_{total}) \ar[r]^{h_{total}} &
\Sh(\mathcal{C}'_{total})
}
$$
此外，有 $(g'_n)^{-1} \circ h_{total, *} = h_{n, *} \circ g_n^{-1}$。
\end{lemma}

\begin{proof}
情形 A。设 $h_n$ 对应于连续函子
$v_n : \mathcal{C}'_n \to \mathcal{C}_n$。在次数 $n$ 使用 $v_n$，
即可定义函子 $v_{total} : \mathcal{C}'_{total} \to \mathcal{C}_{total}$。
它显然连续（参见引理 \ref{spaces-simplicial-lemma-simplicial-site-site} 中覆盖的定义）。令
$h_{total}^{-1} = v_{total, s} :
\Sh(\mathcal{C}'_{total}) \to \Sh(\mathcal{C}_{total})$，以及
$h_{total, *} = v_{total}^s = v_{total}^p :
\Sh(\mathcal{C}_{total}) \to \Sh(\mathcal{C}'_{total})$
为《站点》第 \ref{sites-section-continuous-functors}、
\ref{sites-section-morphism-sites} 节所构造并研究的一对伴随函子。
为证明 $h_{total}$ 是拓扑斯态射，还须验证 $h_{total}^{-1}$ 正合。
先注意
$(g'_n)^{-1} \circ h_{total, *} = h_{n, *} \circ g_n^{-1}$；
这可通过计算 $\mathcal{C}'_n$ 的对象 $U$ 上的截面立即看出。
因此，若按引理 \ref{spaces-simplicial-lemma-describe-sheaves-simplicial-site-site}
把 $\mathcal{C}_{total}$ 上的层 $\mathcal{F}$ 看成系统
$(\mathcal{F}_n, \mathcal{F}(\varphi))$，则
$h_{total, *}\mathcal{F}$ 对应于系统
$(h_{n, *}\mathcal{F}_n, h_{n, *}\mathcal{F}(\varphi))$。
显然，函子
$(\mathcal{F}'_n, \mathcal{F}'(\varphi)) \to
(h_n^{-1}\mathcal{F}'_n, h_n^{-1}\mathcal{F}'(\varphi))$
是它的左伴随。由伴随的唯一性，$h_{total}^{-1}$ 在系统上由此规则给出。
特别地，由该引理对 $\mathcal{C}_{total}$ 上层的描述以及各函子
$h_n^{-1}$ 的正合性，$h_{total}^{-1}$ 正合，从而得到所需拓扑斯态射。
最后还有 $g_n^{-1} \circ h_{total}^{-1} =
h_n^{-1} \circ (g'_n)^{-1}$，这证明引理所示图交换。

\medskip\noindent
情形 B。在次数 $n$ 使用 $v_n$，得到函子
$v_{total} : \mathcal{C}_{total} \to \mathcal{C}'_{total}$。
它显然余连续（参见引理 \ref{spaces-simplicial-lemma-simplicial-cocontinuous-site}
中覆盖的定义）。令 $h_{total}$ 为 $v_{total}$ 所伴随的拓扑斯态射。
引理所示图的交换性直接来自《站点》引理
\ref{sites-lemma-composition-cocontinuous}。
取左伴随后，引理最后的等式化为
$$
h_{total}^{-1} \circ (g'_n)^{Sh}_! = g^{Sh}_{n!} \circ h_n^{-1}
$$
这直接由引理 \ref{spaces-simplicial-lemma-restriction-to-components-site} 中对函子
$(g'_n)^{Sh}_!$ 与 $g^{Sh}_{n!}$ 的显式描述、对
$\varphi : [m] \to [n]$ 成立的等式
$h_n^{-1} \circ (f'_\varphi)^{-1} =
f_\varphi^{-1} \circ h_m^{-1}$，以及已知 $h_{total}^{-1}$
与限制到单纯站点次数 $n$ 部分交换这一事实得到。
\end{proof}

\begin{lemma}
\label{spaces-simplicial-lemma-direct-image-morphism-simplicial-sites}
采用引理 \ref{spaces-simplicial-lemma-morphism-simplicial-sites} 的记号与假设。
对 $K \in D(\mathcal{C}_{total})$，有
$(g'_n)^{-1}Rh_{total, *}K = Rh_{n, *}g_n^{-1}K$。
\end{lemma}

\begin{proof}
设 $\mathcal{I}^\bullet$ 是 $\mathcal{C}_{total}$ 上表示 $K$ 的
K-内射复形。则 $g_n^{-1}K$ 由
$g_n^{-1}\mathcal{I}^\bullet = \mathcal{I}_n^\bullet$ 表示；
由引理 \ref{spaces-simplicial-lemma-restriction-injective-to-component-site}，后者是 K-内射的。
% P10-E0177
我们有 $(g'_n)^{-1}h_{total, *}\mathcal{I}^\bullet =
h_{n, *}g_n^{-1}\mathcal{I}_n^\bullet$；
由引理 \ref{spaces-simplicial-lemma-morphism-simplicial-sites} 即得所需等式。
\end{proof}

\begin{remark}
\label{spaces-simplicial-remark-morphism-augmentation-simplicial-sites}
设 $\mathcal{C}_n, f_\varphi, u_\varphi$ 和
$\mathcal{C}'_n, f'_\varphi, u'_\varphi$ 如情形
\ref{spaces-simplicial-situation-simplicial-site}。设 $a_0$ 和 $a'_0$ 分别为评注
\ref{spaces-simplicial-remark-augmentation-site} 所述指向站点 $\mathcal{D}$ 和
$\mathcal{D}'$ 的增广。设 $h$ 为评注
\ref{spaces-simplicial-remark-morphism-simplicial-sites} 所述的单纯站点之间的态射。
若拓扑斯态射 $h_{-1} : \Sh(\mathcal{D}) \to \Sh(\mathcal{D}')$
满足下列条件，则称它{\it 与 $h$、$a_0$、$a'_0$ 相容}：
\begin{enumerate}
\item[(A)] $h_{-1}$ 来自站点态射
$h_{-1} : \mathcal{D} \to \mathcal{D}'$，且作为站点态射有
$a'_0 \circ h_0 = h_{-1} \circ a_0$；
\item[(B)] $h_{-1}$ 来自余连续函子
$v_{-1} : \mathcal{D} \to \mathcal{D}'$，且作为函子有
$u'_0 \circ v_0 = v_{-1} \circ u_0$。
\end{enumerate}
在两种情形中，作为拓扑斯态射均有
$a'_0 \circ h_0 = h_{-1} \circ a_0$；情形 B 参见《站点》引理
\ref{sites-lemma-composition-cocontinuous}，情形 A 参见《站点》定义
\ref{sites-definition-composition-morphisms-sites}。
\end{remark}

\begin{lemma}
\label{spaces-simplicial-lemma-morphism-augmentation-simplicial-sites}
设 $\mathcal{C}_n, f_\varphi, u_\varphi, \mathcal{D}, a_0$，
$\mathcal{C}'_n, f'_\varphi, u'_\varphi, \mathcal{D}', a'_0$，以及
$h_n$（$n \geq -1$）如评注
\ref{spaces-simplicial-remark-morphism-augmentation-simplicial-sites}。
则得到交换图
$$
\xymatrix{
\Sh(\mathcal{C}_{total}) \ar[d]_a \ar[r]_{h_{total}} &
\Sh(\mathcal{C}'_{total}) \ar[d]^{a'} \\
\Sh(\mathcal{D}) \ar[r]^{h_{-1}} &
\Sh(\mathcal{D}')
}
$$
\end{lemma}

\begin{proof}
态射 $h$ 由引理 \ref{spaces-simplicial-lemma-morphism-simplicial-sites} 定义，
态射 $a$ 与 $a'$ 由引理 \ref{spaces-simplicial-lemma-augmentation-site} 定义。
故只须证明该图交换。为此，证明
$a^{-1} \circ h_{-1}^{-1} = h_{total}^{-1} \circ (a')^{-1}$。
由引理 \ref{spaces-simplicial-lemma-morphism-simplicial-sites} 中的交换图，
以及引理 \ref{spaces-simplicial-lemma-describe-sheaves-simplicial-site-site}
借助分量对 $\Sh(\mathcal{C}_{total})$ 和
$\Sh(\mathcal{C}'_{total})$ 的描述，只须证明
$$
\xymatrix{
\Sh(\mathcal{C}_n) \ar[d]_{a_n} \ar[r]_{h_n} &
\Sh(\mathcal{C}'_n) \ar[d]^{a'_n} \\
\Sh(\mathcal{D}) \ar[r]^{h_{-1}} &
\Sh(\mathcal{D}')
}
$$
对所有 $n$ 均交换。$n=0$ 的情形是评注
\ref{spaces-simplicial-remark-morphism-augmentation-simplicial-sites} 中的假设。
当 $n>0$ 时，选取 $\varphi : [0] \to [n]$；所需交换性由
$n=0$ 的情形、关系 $a_n = a_0 \circ f_\varphi$ 与
$a'_n = a'_0 \circ f'_\varphi$，以及交换关系
$f'_\varphi \circ h_n = h_0 \circ f_\varphi$ 得到。
\end{proof}




\section{赋环单纯站点}
\label{spaces-simplicial-section-simplicial-sites-modules}

\noindent
现在为单纯拓扑斯赋予一个环层。

\begin{lemma}
\label{spaces-simplicial-lemma-restriction-module-to-components-site}
在情形 \ref{spaces-simplicial-situation-simplicial-site} 中，设 $\mathcal{O}$ 为
$\mathcal{C}_{total}$ 上的环层。存在典范赋环拓扑斯态射
$g_n : (\Sh(\mathcal{C}_n), \mathcal{O}_n) \to
(\Sh(\mathcal{C}_{total}), \mathcal{O})$
其底层拓扑斯态射与引理 \ref{spaces-simplicial-lemma-restriction-to-components-site}
中的 $g_n$ 一致。函子
$g_n^* : \textit{Mod}(\mathcal{O}) \to \textit{Mod}(\mathcal{O}_n)$
有左伴随 $g_{n!}$。对 $\textit{Mod}(\mathcal{O}_n)$ 中的
$\mathcal{G}$，$g_{n!}\mathcal{G}$ 在 $\mathcal{C}_m$ 上的限制为
$$
\bigoplus\nolimits_{\varphi : [n] \to [m]} f_\varphi^*\mathcal{G}
$$
其中 $f_\varphi : (\Sh(\mathcal{C}_m), \mathcal{O}_m) \to
(\Sh(\mathcal{C}_n), \mathcal{O}_n)$ 是赋环拓扑斯态射；
其底层拓扑斯态射与先前定义的 $f_\varphi$ 一致，并在环层上使用映射
$\mathcal{O}(\varphi) : f_\varphi^{-1}\mathcal{O}_n \to \mathcal{O}_m$
。
\end{lemma}

\begin{proof}
由引理 \ref{spaces-simplicial-lemma-restriction-to-components-site}，有
$g_n^{-1}\mathcal{O} = \mathcal{O}_n$，从而得到上述赋环拓扑斯态射。
由《站点上的模》引理 \ref{sites-modules-lemma-lower-shriek-modules}，
得到伴随 $g_{n!}$。为证明 $g_{n!}$ 的公式，先定义
$\mathcal{C}_{total}$ 上的 $\mathcal{O}$-模层 $\mathcal{H}$，
其次数 $m$ 分量为 $\mathcal{O}_m$-模
$$
\mathcal{H}_m =
\bigoplus\nolimits_{\varphi : [n] \to [m]} f_\varphi^*\mathcal{G}
$$
给定映射 $\psi : [m] \to [m']$，映射
$\mathcal{H}(\psi) : f_\psi^{-1}\mathcal{H}_m \to \mathcal{H}_{m'}$
在各分量上由下式给出：
$$
f_\psi^{-1} f_\varphi^*\mathcal{G} \to
f_\psi^* f_\varphi^*\mathcal{G} \to
f_{\psi \circ \varphi}^*\mathcal{G}
$$
由于映射 $f_\psi^{-1}\mathcal{H}_m \to \mathcal{H}_{m'}$
关于 $\mathcal{O}(\psi) : f_\psi^{-1}\mathcal{O}_m \to \mathcal{O}_{m'}$
是半线性的，这确实定义一个 $\mathcal{O}$-模
（使用引理 \ref{spaces-simplicial-lemma-describe-sheaves-simplicial-site-site}）。
随后直接证明
$$
\Mor_{\mathcal{O}_n}(\mathcal{G}, \mathcal{F}_n) =
\Mor_{\mathcal{O}}(\mathcal{H}, \mathcal{F})
$$
其论证同引理 \ref{spaces-simplicial-lemma-restriction-to-components-site} 的证明。
故 $\mathcal{H} = g_{n!}\mathcal{G}$，如所求。
\end{proof}

\begin{lemma}
\label{spaces-simplicial-lemma-restriction-injective-to-component-limp}
在情形 \ref{spaces-simplicial-situation-simplicial-site} 中，设 $\mathcal{O}$ 为
$\mathcal{C}_{total}$ 上的环层。若 $\mathcal{I}$ 在
$\textit{Mod}(\mathcal{O})$ 中内射，则 $\mathcal{I}_n$ 是
$\mathcal{C}_n$ 上的完全无环层。
\end{lemma}

\begin{proof}
把《站点上的上同调》引理
\ref{sites-cohomology-lemma-pullback-injective-limp} 应用于包含函子
$\mathcal{C}_n \to \mathcal{C}_{total}$，并使用引理
\ref{spaces-simplicial-lemma-restriction-to-components-site} 所证明的该函子性质，即得结论。
\end{proof}

\begin{lemma}
\label{spaces-simplicial-lemma-exactness-g-shriek-modules}
在引理 \ref{spaces-simplicial-lemma-restriction-module-to-components-site} 的假设下，
若对所有 $\varphi : [n] \to [m]$，映射
$f_\varphi^{-1}\mathcal{O}_n \to \mathcal{O}_m$ 均平坦，则函子
$g_{n!} : \textit{Mod}(\mathcal{O}_n) \to \textit{Mod}(\mathcal{O})$
正合。
\end{lemma}

\begin{proof}
回忆 $g_{n!}\mathcal{G}$ 是如下 $\mathcal{O}$-模：其次数 $m$ 部分是
$\mathcal{O}_m$-模
$$
\bigoplus\nolimits_{\varphi : [n] \to [m]} f_\varphi^*\mathcal{G}
$$
这里赋环拓扑斯态射
$f_\varphi : (\Sh(\mathcal{C}_m), \mathcal{O}_m) \to
(\Sh(\mathcal{C}_n), \mathcal{O}_n)$ 使用引理陈述中的映射
$f_\varphi^{-1}\mathcal{O}_n \to \mathcal{O}_m$。
若这些映射平坦，则 $f_\varphi^*$ 正合
（《站点上的模》引理 \ref{sites-modules-lemma-flat-pullback-exact}）。
由站点 $\mathcal{C}_{total}$ 的定义可知，这些函子具有所需正合性，故得结论。
\end{proof}

\begin{lemma}
\label{spaces-simplicial-lemma-restriction-injective-to-component-site-module}
在情形 \ref{spaces-simplicial-situation-simplicial-site} 中，设 $\mathcal{O}$ 为
$\mathcal{C}_{total}$ 上的环层，并假设对所有
$\varphi : [n] \to [m]$，映射
$f_\varphi^{-1}\mathcal{O}_n \to \mathcal{O}_m$ 均平坦。
若 $\mathcal{I}$ 在 $\textit{Mod}(\mathcal{O})$ 中内射，
则 $\mathcal{I}_n$ 在 $\textit{Mod}(\mathcal{O}_n)$ 中内射。
\end{lemma}

\begin{proof}
这由《同调代数》引理 \ref{homology-lemma-adjoint-preserve-injectives}
和引理 \ref{spaces-simplicial-lemma-exactness-g-shriek-modules} 得到。
\end{proof}







\section{赋环单纯站点的态射}
\label{spaces-simplicial-section-morphism-simplicial-sites-modules}

\noindent
我们继续第 \ref{spaces-simplicial-section-morphism-simplicial-sites} 节的讨论。

\begin{remark}
\label{spaces-simplicial-remark-morphism-simplicial-sites-modules}
设 $\mathcal{C}_n, f_\varphi, u_\varphi$ 与
$\mathcal{C}'_n, f'_\varphi, u'_\varphi$ 如
情形 \ref{spaces-simplicial-situation-simplicial-site} 所述。
设 $\mathcal{O}$ 与 $\mathcal{O}'$
分别为 $\mathcal{C}_{total}$ 与 $\mathcal{C}'_{total}$ 上的环层。
若 $h$ 是注 \ref{spaces-simplicial-remark-morphism-simplicial-sites}
所述的单纯站点间的态射，且
$h^\sharp : h_{total}^{-1}\mathcal{O}' \to \mathcal{O}$，
等价地，$h^\sharp : \mathcal{O}' \to h_{total, *}\mathcal{O}$
是环层的同态，则称 $(h, h^\sharp)$ 是
{\it 赋环单纯站点间的态射}。
\end{remark}

\begin{lemma}
\label{spaces-simplicial-lemma-morphism-simplicial-sites-modules}
设 $\mathcal{C}_n, f_\varphi, u_\varphi$ 与
$\mathcal{C}'_n, f'_\varphi, u'_\varphi$ 如
情形 \ref{spaces-simplicial-situation-simplicial-site} 所述。
设 $\mathcal{O}$ 与 $\mathcal{O}'$
分别为 $\mathcal{C}_{total}$ 与 $\mathcal{C}'_{total}$ 上的环层。
设 $(h, h^\sharp)$ 是注
\ref{spaces-simplicial-remark-morphism-simplicial-sites-modules} 所述的单纯站点间的态射。
则我们得到赋环拓扑斯的态射
$$
h_{total} :
(\Sh(\mathcal{C}_{total}), \mathcal{O})
\to
(\Sh(\mathcal{C}'_{total}), \mathcal{O}')
$$
以及交换图
$$
\xymatrix{
(\Sh(\mathcal{C}_n), \mathcal{O}_n) \ar[d]_{g_n} \ar[r]_{h_n} &
(\Sh(\mathcal{C}'_n), \mathcal{O}'_n) \ar[d]^{g'_n} \\
(\Sh(\mathcal{C}_{total}), \mathcal{O}) \ar[r]^{h_{total}} &
(\Sh(\mathcal{C}'_{total}), \mathcal{O}')
}
$$
；这里它们是赋环拓扑斯的图，而 $g_n$ 与 $g'_n$ 如
引理 \ref{spaces-simplicial-lemma-restriction-module-to-components-site} 所述。
此外，作为函子 $\textit{Mod}(\mathcal{O}) \to
\textit{Mod}(\mathcal{O}'_n)$，我们有
$(g'_n)^* \circ h_{total, *} = h_{n, *} \circ g_n^*$
。
\end{lemma}

\begin{proof}
由引理 \ref{spaces-simplicial-lemma-morphism-simplicial-sites} 与
\ref{spaces-simplicial-lemma-restriction-module-to-components-site}，并追踪环层即可得出。
有一个小细节：为了将 $h_n$ 定义为赋环拓扑斯的态射，我们令
$h_n^\sharp = g_n^{-1}h^\sharp :
g_n^{-1}h_{total}^{-1}\mathcal{O}' \to g_n^{-1}\mathcal{O}$
；这是有意义的，因为
$g_n^{-1}h_{total}^{-1}\mathcal{O}' = h_n^{-1}(g'_n)^{-1}\mathcal{O}' =
h_n^{-1}\mathcal{O}'_n$ 且 $g_n^{-1}\mathcal{O} = \mathcal{O}_n$。
注意，对任意 $\mathcal{O}$-模层 $\mathcal{F}$，均有
$g_n^*\mathcal{F} = g_n^{-1}\mathcal{F}$；
对 $g'_n$ 亦然。这也有助于说明，为什么
$(g'_n)^* \circ h_{total, *} = h_{n, *} \circ g_n^*$
可由引理 \ref{spaces-simplicial-lemma-morphism-simplicial-sites} 中的相应陈述推出。
\end{proof}

\begin{lemma}
\label{spaces-simplicial-lemma-direct-image-morphism-simplicial-sites-modules}
采用引理 \ref{spaces-simplicial-lemma-morphism-simplicial-sites-modules} 中的记号与假设。
对 $K \in D(\mathcal{O})$，有
$(g'_n)^*Rh_{total, *}K = Rh_{n, *}g_n^*K$。
\end{lemma}

\begin{proof}
回忆一下，由引理 \ref{spaces-simplicial-lemma-restriction-module-to-components-site}
中的构造，$g_n^{-1}\mathcal{O} = \mathcal{O}_n$，故
$g_n^* = g_n^{-1}$。特别地，$g_n^*$ 是正合的，而 $Lg_n^*$
可通过将 $g_n^*$ 作用于任意一个模复形代表来给出。对 $g'_n$ 亦然。
存在典范基变换映射
$(g'_n)^*Rh_{total, *}K \to Rh_{n, *}g_n^*K$，见
《站点上的上同调》，注 \ref{sites-cohomology-remark-base-change}。
由《站点上的上同调》，引理
\ref{sites-cohomology-lemma-modules-abelian-unbounded}
可知，它在 $D(\mathcal{C}'_n)$ 中的像为映射
$(g'_n)^{-1}Rh_{total, *}K_{ab} \to Rh_{n, *}g_n^{-1}K_{ab}$
，其中 $K_{ab}$ 是 $K$ 在 $D(\mathcal{C}_{total})$ 中的像。
引理 \ref{spaces-simplicial-lemma-direct-image-morphism-simplicial-sites}
已证明此映射是同构，故结论成立。
\end{proof}








\section{单纯站点上的上同调}
\label{spaces-simplicial-section-cohomology-simplicial-sites}

\noindent
设 $\mathcal{C}$ 如情形 \ref{spaces-simplicial-situation-simplicial-site} 所述。
在以下各引理的陈述中，我们以
$g_n : \Sh(\mathcal{C}_n) \to \Sh(\mathcal{C}_{total})$ 表示
引理 \ref{spaces-simplicial-lemma-restriction-to-components-site} 中的拓扑斯态射。
若 $\varphi : [m] \to [n]$ 是 $\Delta$ 的一个态射，则拓扑斯图
$$
\xymatrix{
\Sh(\mathcal{C}_n) \ar[rd]_{g_n} \ar[rr]_{f_\varphi} & &
\Sh(\mathcal{C}_m) \ar[ld]^{g_m} \\
& \Sh(\mathcal{C}_{total})
}
$$
并不交换，但存在来自映射
$\mathcal{F}(\varphi) : f_\varphi^{-1}\mathcal{F}_m \to \mathcal{F}_n$
的 $2$-态射 $g_n \to g_m \circ f_\varphi$。
参见《站点》第 \ref{sites-section-2-category} 节。

\begin{lemma}
\label{spaces-simplicial-lemma-simplicial-resolution-Z-site}
在情形 \ref{spaces-simplicial-situation-simplicial-site} 中，并采用上述记号，存在复形
$$
\ldots \to g_{2!}\mathbf{Z} \to g_{1!}\mathbf{Z} \to g_{0!}\mathbf{Z}
$$
；它由 $\mathcal{C}_{total}$ 上的阿贝尔层组成，并构成
$\mathcal{C}_{total}$ 上取值为 $\mathbf{Z}$ 的常值层的一个消解。
\end{lemma}

\begin{proof}
以下不再说明地使用引理 \ref{spaces-simplicial-lemma-restriction-to-components-site}
中对函子 $g_{n!}$ 的描述。
取 $\sum (-1)^i d^n_i$ 为该复形的映射，其中
$d^n_i : g_{n!}\mathbf{Z} \to g_{n - 1!}\mathbf{Z}$ 是映射 $\mathbf{Z} \to
\bigoplus_{[n - 1] \to [n]} \mathbf{Z} = g_n^{-1}g_{n - 1!}\mathbf{Z}$
的伴随；这里该映射对应于标记为
$\delta^n_i : [n - 1] \to [n]$ 的因子。
于是将 $g_m^{-1}$ 作用于该复形，得到 $\mathcal{C}_m$ 上的复形
$$
\ldots \to
\bigoplus\nolimits_{\alpha \in \Mor_\Delta([2], [m])]} \mathbf{Z} \to
\bigoplus\nolimits_{\alpha \in \Mor_\Delta([1], [m])]} \mathbf{Z} \to
\bigoplus\nolimits_{\alpha \in \Mor_\Delta([0], [m])]} \mathbf{Z}
$$
。
换言之，这是与单纯集合 $\Delta[m]$ 上的自由阿贝尔层相关联的复形；
参见《单纯对象》例 \ref{simplicial-example-simplex-simplicial-set}。
由于 $\Delta[m]$ 与 $\Delta[0]$ 同伦等价（参见《单纯对象》例
\ref{simplicial-example-simplex-contractible}），且“取其上的自由阿贝尔层”
是一个函子，故上述复形与 $\Delta[0]$ 上的自由阿贝尔层同伦等价
（《单纯对象》注 \ref{simplicial-remark-homotopy-better} 及
引理 \ref{simplicial-lemma-homotopy-equivalence-s-N}）。
此复形在正次数上无环，并在次数 $0$ 上等于 $\mathbf{Z}$。
\end{proof}

\begin{lemma}
\label{spaces-simplicial-lemma-cech-complex}
在情形 \ref{spaces-simplicial-situation-simplicial-site} 中，设 $\mathcal{F}$ 是
$\mathcal{C}_{total}$ 上的阿贝尔层。存在典范复形
$$
0 \to \Gamma(\mathcal{C}_{total}, \mathcal{F}) \to
\Gamma(\mathcal{C}_0, \mathcal{F}_0) \to
\Gamma(\mathcal{C}_1, \mathcal{F}_1) \to
\Gamma(\mathcal{C}_2, \mathcal{F}_2) \to \ldots
$$
；它在次数 $-1, 0$ 上正合；若 $\mathcal{F}$ 是内射的，
则它处处正合。
\end{lemma}

\begin{proof}
注意
$\Hom(\mathbf{Z}, \mathcal{F}) = \Gamma(\mathcal{C}_{total}, \mathcal{F})$
且
$\Hom(g_{n!}\mathbf{Z}, \mathcal{F}) = \Gamma(\mathcal{C}_n, \mathcal{F}_n)$.
因此，由引理 \ref{spaces-simplicial-lemma-simplicial-resolution-Z-site}
以及当 $\mathcal{F}$ 内射时 $\Hom(-, \mathcal{F})$ 正合这一事实，
立即得到本引理。
\end{proof}

\begin{lemma}
\label{spaces-simplicial-lemma-simplicial-sheaf-cohomology-site}
在情形 \ref{spaces-simplicial-situation-simplicial-site} 中，对属于
$D^+(\mathcal{C}_{total})$ 的 $K$，存在谱序列
$(E_r, d_r)_{r \geq 0}$，其中
$$
E_1^{p, q} = H^q(\mathcal{C}_p, K_p),\quad
d_1^{p, q} : E_1^{p, q} \to E_1^{p + 1, q}
$$
，它收敛到 $H^{p + q}(\mathcal{C}_{total}, K)$。
此谱序列关于 $K$ 是函子性的。
\end{lemma}

\begin{proof}
设 $\mathcal{I}^\bullet$ 是表示 $K$ 的下有界内射复形。
考虑各项为
$$
A^{p, q} = \Gamma(\mathcal{C}_p, \mathcal{I}^q_p)
$$
的双复形；其水平箭头来自引理 \ref{spaces-simplicial-lemma-cech-complex}，
竖直箭头来自复形 $\mathcal{I}^\bullet$ 的微分。
该双复形的各行在正次数上正合，并在次数 $0$ 上给出
$\Gamma(\mathcal{C}_{total}, \mathcal{I}^q)$。
另一方面，由于到 $\mathcal{C}_p$ 的限制是正合的
（引理 \ref{spaces-simplicial-lemma-restriction-to-components-site}），复形
$\mathcal{I}_p^\bullet$ 在 $D(\mathcal{C}_p)$ 中表示 $K_p$。
诸层 $\mathcal{I}_p^q$ 是 $\mathcal{C}_p$ 上的内射阿贝尔层
（引理 \ref{spaces-simplicial-lemma-restriction-injective-to-component-site}）。
因此，各列的上同调计算群 $H^q(\mathcal{C}_p, K_p)$。
应用《同调代数》引理 \ref{homology-lemma-first-quadrant-ss} 与
\ref{homology-lemma-double-complex-gives-resolution} 即得结论。
\end{proof}

\begin{remark}
\label{spaces-simplicial-remark-simplicial-cohomology-site}
采用引理 \ref{spaces-simplicial-lemma-simplicial-sheaf-cohomology-site}
中的假设与记号，但不要求 $D(\mathcal{C}_{total})$ 中的 $K$ 下有界。
我们断言，此时也存在自然谱序列。具体而言，设
$\mathcal{I}^\bullet$ 是 $\mathcal{C}_{total}$ 上表示 $K$
且各项内射的层 K-内射复形。则
\begin{align*}
R\Gamma(\mathcal{C}_{total}, K)
& =
R\Hom(\mathbf{Z}, K) \\
& =
R\Hom(
\ldots \to g_{2!}\mathbf{Z} \to g_{1!}\mathbf{Z} \to g_{0!}\mathbf{Z}, K) \\
& =
\Gamma(\mathcal{C}_{total},
\SheafHom^\bullet(
\ldots \to g_{2!}\mathbf{Z} \to g_{1!}\mathbf{Z} \to g_{0!}\mathbf{Z},
\mathcal{I}^\bullet)) \\
& =
\text{Tot}_\pi(A^{\bullet, \bullet})
\end{align*}
其中 $A^{\bullet, \bullet}$ 是各项为
$A^{p, q} = \Gamma(\mathcal{C}_p, \mathcal{I}^q_p)$ 的双复形，
$\text{Tot}_\pi$ 表示该双复形的乘积全化。
这里，第一个等式在任意站点上成立。第二个等式由引理
\ref{spaces-simplicial-lemma-simplicial-resolution-Z-site} 成立。
第三个等式成立是因为 $\mathcal{I}^\bullet$ 是 K-内射的；参见
《站点上的上同调》第 \ref{sites-cohomology-section-hom-complexes} 节
与第 \ref{sites-cohomology-section-internal-hom} 节。
最后一个等式由 $\SheafHom^\bullet$ 的构造以及等式
$\Hom(g_{p!}\mathbf{Z}, \mathcal{I}^q) =
\Gamma(\mathcal{C}_p, \mathcal{I}^q_p)$ 成立。
现在将 $\text{Tot}_\pi(A^{\bullet, \bullet})$ 看成滤过复形，其中
$F^i\text{Tot}^n_\pi(A^{\bullet, \bullet}) =
\prod_{p + q = n,\ p \geq i} A^{p, q}$，由此得到所需谱序列。
除正则性、有界性、收敛性与极限项等问题外，我们得到的谱序列
与《同调代数》第 \ref{homology-section-double-complex} 节中的谱序列
$({}'E_r, {}'d_r)_{r \geq 0}$ 具有同样的表现
（该节讨论的是直和全化的情形）。特别地，正如引理
\ref{spaces-simplicial-lemma-simplicial-sheaf-cohomology-site} 中一样，我们得到
$E_1^{p, q} = H^q(\mathcal{C}_p, K_p)$。
\end{remark}

\begin{lemma}
\label{spaces-simplicial-lemma-simplicial-sheaf-cohomology-site-zero}
在情形 \ref{spaces-simplicial-situation-simplicial-site} 中，设 $K$ 是
$D(\mathcal{C}_{total})$ 的一个对象。
\begin{enumerate}
\item 若对所有 $p \geq 0$ 都有 $H^{-p}(\mathcal{C}_p, K_p) = 0$，则
$H^0(\mathcal{C}_{total}, K) = 0$。
\item 若对所有 $p \geq 0$ 都有 $R\Gamma(\mathcal{C}_p, K_p) = 0$，
则 $R\Gamma(\mathcal{C}_{total}, K) = 0$。
\end{enumerate}
\end{lemma}

\begin{proof}
采用注 \ref{spaces-simplicial-remark-simplicial-cohomology-site} 中的记号，
可见 $R\Gamma(\mathcal{C}_{total}, K)$ 由
$\text{Tot}_\pi(A^{\bullet, \bullet})$ 表示。
(1) 中的假设说明 $H^{-p}(A^{p, \bullet}) = 0$。
于是 (1) 中的消失由《代数续论》引理
\ref{more-algebra-lemma-vanishing-coh-prod-totalization} 得出。
(2) 由 (1) 及平移得到。
\end{proof}

\begin{lemma}
\label{spaces-simplicial-lemma-sanity-check}
设 $\mathcal{C}$ 如情形 \ref{spaces-simplicial-situation-simplicial-site} 所述，
$U \in \Ob(\mathcal{C}_n)$，且
$\mathcal{F} \in \textit{Ab}(\mathcal{C}_{total})$。
则 $H^p(U, \mathcal{F}) = H^p(U, g_n^{-1}\mathcal{F})$，
其中在左端，$U$ 被视为 $\mathcal{C}_{total}$ 的一个对象。
\end{lemma}

\begin{proof}
注意，“将 $U$ 视为 $\mathcal{C}_{total}$ 的对象”这一说法，
在情形 (A) 下由引理 \ref{spaces-simplicial-lemma-simplicial-site-site}
中的 $\mathcal{C}_{total}$ 构造得到解释，在情形 (B) 下则由引理
\ref{spaces-simplicial-lemma-simplicial-cocontinuous-site} 中的构造得到解释。
于是，该等式由引理 \ref{spaces-simplicial-lemma-restriction-injective-to-component-site}
及上同调的定义得出。
\end{proof}






\section{单纯站点的上同调与增广}
\label{spaces-simplicial-section-cohomology-augmentation-simplicial-sites}

\noindent
考虑情形 \ref{spaces-simplicial-situation-simplicial-site} 所述的单纯站点 $\mathcal{C}$。
设 $a_0$ 是注 \ref{spaces-simplicial-remark-augmentation-site}
所述的朝向站点 $\mathcal{D}$ 的增广。
由引理 \ref{spaces-simplicial-lemma-augmentation-site}，我们得到拓扑斯的态射
$$
a : \Sh(\mathcal{C}_{total}) \longrightarrow \Sh(\mathcal{D})
$$
以及拓扑斯的态射
$g_n : \Sh(\mathcal{C}_n) \to \Sh(\mathcal{C}_{total})$
，后者如引理 \ref{spaces-simplicial-lemma-restriction-to-components-site} 所述。
将复合 $a \circ g_n$ 记为
$a_n : \Sh(\mathcal{C}_n) \to \Sh(\mathcal{D})$。
此外，单纯结构给出拓扑斯的态射
$f_\varphi : \Sh(\mathcal{C}_n) \to \Sh(\mathcal{C}_m)$
，使得对所有 $\varphi : [m] \to [n]$，均有
% P10-E0180
$a_n \circ f_\varphi = a_m$。

\begin{lemma}
\label{spaces-simplicial-lemma-simplicial-resolution-augmentation}
在情形 \ref{spaces-simplicial-situation-simplicial-site} 中，设
$a_0$ 是注 \ref{spaces-simplicial-remark-augmentation-site}
所述的朝向站点 $\mathcal{D}$ 的增广。
对 $\mathcal{D}$ 上任意阿贝尔层 $\mathcal{G}$，存在正合复形
$$
\ldots \to
g_{2!}(a_2^{-1}\mathcal{G}) \to
g_{1!}(a_1^{-1}\mathcal{G}) \to
g_{0!}(a_0^{-1}\mathcal{G}) \to
a^{-1}\mathcal{G} \to 0
$$
；它由 $\mathcal{C}_{total}$ 上的阿贝尔层组成。
\end{lemma}

\begin{proof}
建议读者先阅读引理 \ref{spaces-simplicial-lemma-simplicial-resolution-Z-site} 的证明。
以下不再说明地使用引理 \ref{spaces-simplicial-lemma-augmentation-site}，
以及引理 \ref{spaces-simplicial-lemma-restriction-to-components-site}
中对函子 $g_{n!}$ 的描述。
特别地，$g_{n!}(a_n^{-1}\mathcal{G})$ 是
$\mathcal{C}_{total}$ 上的层，其到 $\mathcal{C}_m$ 的限制为层
$$
\bigoplus\nolimits_{\varphi : [n] \to [m]} f_\varphi^{-1}a_n^{-1}\mathcal{G} =
\bigoplus\nolimits_{\varphi : [n] \to [m]} a_m^{-1}\mathcal{G}
$$
。取 $\sum (-1)^i d^n_i$ 为该复形的映射，其中
$d^n_i : g_{n!}(a_n^{-1}\mathcal{G}) \to g_{n - 1!}(a_{n - 1}^{-1}\mathcal{G})$
是映射
$a_n^{-1}\mathcal{G} \to \bigoplus_{[n - 1] \to [n]} a_n^{-1}\mathcal{G} =
g_n^{-1}g_{n - 1!}(a_{n - 1}^{-1}\mathcal{G})$
的伴随；这里该映射对应于标记为
$\delta^n_i : [n - 1] \to [n]$ 的因子。
映射 $g_{0!}(a_0^{-1}\mathcal{G}) \to a^{-1}\mathcal{G}$
与 $a_0^{-1}\mathcal{G}$ 的恒等映射伴随。
于是，将 $g_m^{-1}$ 作用于次数为 $\ldots, 2, 1, 0$ 的链复形，
得到 $\mathcal{C}_m$ 上的复形
$$
\ldots \to
\bigoplus\nolimits_{\alpha \in \Mor_\Delta([2], [m])]} a_m^{-1}\mathcal{G} \to
\bigoplus\nolimits_{\alpha \in \Mor_\Delta([1], [m])]} a_m^{-1}\mathcal{G} \to
\bigoplus\nolimits_{\alpha \in \Mor_\Delta([0], [m])]} a_m^{-1}\mathcal{G}
$$
。这等于 $a_m^{-1}\mathcal{G}$ 在常值层 $\mathbf{Z}$ 上与复形
$$
\ldots \to
\bigoplus\nolimits_{\alpha \in \Mor_\Delta([2], [m])]} \mathbf{Z} \to
\bigoplus\nolimits_{\alpha \in \Mor_\Delta([1], [m])]} \mathbf{Z} \to
\bigoplus\nolimits_{\alpha \in \Mor_\Delta([0], [m])]} \mathbf{Z}
$$
的张量积；该复形已在引理 \ref{spaces-simplicial-lemma-simplicial-resolution-Z-site}
的证明中讨论过。那里已经看到，此复形同伦等价于置于次数 $0$
的 $\mathbf{Z}$，证明完毕。
\end{proof}

\begin{lemma}
\label{spaces-simplicial-lemma-augmentation-cech-complex}
在情形 \ref{spaces-simplicial-situation-simplicial-site} 中，设
$a_0$ 是注 \ref{spaces-simplicial-remark-augmentation-site}
所述的朝向站点 $\mathcal{D}$ 的增广。
对 $\mathcal{C}_{total}$ 上的阿贝尔层 $\mathcal{F}$，
存在 $\mathcal{D}$ 上的典范复形
$$
0 \to a_*\mathcal{F} \to a_{0, *}\mathcal{F}_0 \to a_{1, *}\mathcal{F}_1 \to
a_{2, *}\mathcal{F}_2 \to \ldots
$$
；它在次数 $-1, 0$ 上正合；若 $\mathcal{F}$ 是内射的，
则它处处正合。
\end{lemma}

\begin{proof}
为构造该复形，由 Yoneda 引理，只需对 $\mathcal{D}$ 上任意
阿贝尔层 $\mathcal{G}$ 构造复形
$$
0 \to \Hom(\mathcal{G}, a_*\mathcal{F}) \to
\Hom(\mathcal{G}, a_{0, *}\mathcal{F}_0) \to
\Hom(\mathcal{G}, a_{1, *}\mathcal{F}_1) \to \ldots
$$
，并使其关于 $\mathcal{G}$ 是函子性的。为此，将
$\Hom(-, \mathcal{F})$ 作用于引理
\ref{spaces-simplicial-lemma-simplicial-resolution-augmentation} 的正合复形，
并使用拉回与推出的伴随性。
由于 $\Hom(-, \mathcal{F})$ 左正合，次数 $-1, 0$
上的正合性由该构造得出。若 $\mathcal{F}$ 是内射阿贝尔层，
则 $\Hom(-, \mathcal{F})$ 正合，故此复形正合。
\end{proof}

\begin{lemma}
\label{spaces-simplicial-lemma-augmentation-spectral-sequence}
在情形 \ref{spaces-simplicial-situation-simplicial-site} 中，设
$a_0$ 是注 \ref{spaces-simplicial-remark-augmentation-site}
所述的朝向站点 $\mathcal{D}$ 的增广。
对任意属于 $D^+(\mathcal{C}_{total})$ 的 $K$，存在谱序列
$(E_r, d_r)_{r \geq 0}$，其中
$$
E_1^{p, q} = R^qa_{p, *} K_p,\quad
d_1^{p, q} : E_1^{p, q} \to E_1^{p + 1, q}
$$
；它收敛到 $R^{p + q}a_*K$。此谱序列关于 $K$ 是函子性的。
\end{lemma}

\begin{proof}
设 $\mathcal{I}^\bullet$ 是表示 $K$ 的下有界内射复形。
考虑各项为
$$
A^{p, q} = a_{p, *}\mathcal{I}^q_p
$$
的双复形；其水平箭头来自引理
\ref{spaces-simplicial-lemma-augmentation-cech-complex}，
竖直箭头来自复形 $\mathcal{I}^\bullet$ 的微分。
该双复形的各行在正次数上正合，并在次数 $0$ 上给出
$a_*\mathcal{I}^q$。另一方面，由于到 $\mathcal{C}_p$ 的限制正合
（引理 \ref{spaces-simplicial-lemma-restriction-to-components-site}），复形
$\mathcal{I}_p^\bullet$ 在 $D(\mathcal{C}_p)$ 中表示 $K_p$。
诸层 $\mathcal{I}_p^q$ 是 $\mathcal{C}_p$ 上的内射阿贝尔层
（引理 \ref{spaces-simplicial-lemma-restriction-injective-to-component-site}）。
因此，各列的上同调计算 $R^qa_{p, *}K_p$。
应用《同调代数》引理 \ref{homology-lemma-first-quadrant-ss} 与
\ref{homology-lemma-double-complex-gives-resolution} 即得结论。
\end{proof}





\section{赋环单纯站点上的上同调}
\label{spaces-simplicial-section-cohomology-simplicial-sites-modules}

\noindent
本节是第 \ref{spaces-simplicial-section-cohomology-simplicial-sites} 节
对于模层的类似版本。

\medskip\noindent
在情形 \ref{spaces-simplicial-situation-simplicial-site} 中，设 $\mathcal{O}$ 是
$\mathcal{C}_{total}$ 上的环层。
在以下各引理的陈述中，我们以
$g_n : (\Sh(\mathcal{C}_n), \mathcal{O}_n) \to
(\Sh(\mathcal{C}_{total}), \mathcal{O})$
表示引理 \ref{spaces-simplicial-lemma-restriction-module-to-components-site}
中的赋环拓扑斯态射。
若 $\varphi : [m] \to [n]$ 是 $\Delta$ 的一个态射，则赋环拓扑斯图
$$
\xymatrix{
(\Sh(\mathcal{C}_n), \mathcal{O}_n) \ar[rd]_{g_n} \ar[rr]_{f_\varphi} & &
(\Sh(\mathcal{C}_m), \mathcal{O}_m) \ar[ld]^{g_m} \\
& (\Sh(\mathcal{C}_{total}), \mathcal{O})
}
$$
并不交换，但存在来自映射
$\mathcal{F}(\varphi) : f_\varphi^{-1}\mathcal{F}_m \to \mathcal{F}_n$
的 $2$-态射 $g_n \to g_m \circ f_\varphi$。
参见《站点》第 \ref{sites-section-2-category} 节。

\begin{lemma}
\label{spaces-simplicial-lemma-simplicial-resolution-ringed}
在情形 \ref{spaces-simplicial-situation-simplicial-site} 中，设 $\mathcal{O}$ 是
$\mathcal{C}_{total}$ 上的环层。存在 $\mathcal{O}$-模复形
$$
\ldots \to g_{2!}\mathcal{O}_2 \to g_{1!}\mathcal{O}_1 \to g_{0!}\mathcal{O}_0
$$
；它构成 $\mathcal{O}$ 的一个消解。
这里 $g_{n!}$ 如引理
\ref{spaces-simplicial-lemma-restriction-module-to-components-site} 所述。
\end{lemma}

\begin{proof}
我们将使用引理 \ref{spaces-simplicial-lemma-restriction-to-components-site}
中对 $g_{n!}$ 的描述。取 $\sum (-1)^i d^n_i$ 为该复形的映射，
其中 $d^n_i : g_{n!}\mathcal{O}_n \to g_{n - 1!}\mathcal{O}_{n - 1}$
是映射
$\mathcal{O}_n \to \bigoplus_{[n - 1] \to [n]} \mathcal{O}_n =
g_n^*g_{n - 1!}\mathcal{O}_{n - 1}$
的伴随；这里该映射对应于标记为
$\delta^n_i : [n - 1] \to [n]$ 的因子。
于是将 $g_m^{-1}$ 作用于该复形，得到 $\mathcal{C}_m$ 上的复形
$$
\ldots \to
\bigoplus\nolimits_{\alpha \in \Mor_\Delta([2], [m])]} \mathcal{O}_m \to
\bigoplus\nolimits_{\alpha \in \Mor_\Delta([1], [m])]} \mathcal{O}_m \to
\bigoplus\nolimits_{\alpha \in \Mor_\Delta([0], [m])]} \mathcal{O}_m
$$
。
换言之，这是与单纯集合 $\Delta[m]$ 上的自由
$\mathcal{O}_m$-模相关联的复形；参见《单纯对象》例
\ref{simplicial-example-simplex-simplicial-set}。
由于 $\Delta[m]$ 与 $\Delta[0]$ 同伦等价（参见《单纯对象》例
\ref{simplicial-example-simplex-contractible}），
% P10-E0178
且“取其上的自由阿贝尔层”是一个函子，故上述复形同伦等价于
$\Delta[0]$ 上的自由阿贝尔层
（《单纯对象》注 \ref{simplicial-remark-homotopy-better} 及
引理 \ref{simplicial-lemma-homotopy-equivalence-s-N}）。
此复形在正次数上无环，并在次数 $0$ 上等于 $\mathcal{O}_m$。
\end{proof}

\begin{lemma}
\label{spaces-simplicial-lemma-cech-complex-modules}
在情形 \ref{spaces-simplicial-situation-simplicial-site} 中，设 $\mathcal{O}$ 是环层，
$\mathcal{F}$ 是 $\mathcal{O}$-模层。存在典范复形
$$
0 \to \Gamma(\mathcal{C}_{total}, \mathcal{F}) \to
\Gamma(\mathcal{C}_0, \mathcal{F}_0) \to
\Gamma(\mathcal{C}_1, \mathcal{F}_1) \to
\Gamma(\mathcal{C}_2, \mathcal{F}_2) \to \ldots
$$
；它在次数 $-1, 0$ 上正合；若 $\mathcal{F}$ 是内射
$\mathcal{O}$-模，则它处处正合。
\end{lemma}

\begin{proof}
注意
$\Hom(\mathcal{O}, \mathcal{F}) = \Gamma(\mathcal{C}_{total}, \mathcal{F})$
且
$\Hom(g_{n!}\mathcal{O}_n, \mathcal{F}) = \Gamma(\mathcal{C}_n, \mathcal{F}_n)$.
因此，由引理 \ref{spaces-simplicial-lemma-simplicial-resolution-ringed}
以及当 $\mathcal{F}$ 内射时 $\Hom(-, \mathcal{F})$ 正合这一事实，
立即得到本引理。
\end{proof}

\begin{lemma}
\label{spaces-simplicial-lemma-simplicial-module-cohomology-site}
在情形 \ref{spaces-simplicial-situation-simplicial-site} 中，设 $\mathcal{O}$ 是环层。
对属于 $D^+(\mathcal{O})$ 的 $K$，存在谱序列
$(E_r, d_r)_{r \geq 0}$，其中
$$
E_1^{p, q} = H^q(\mathcal{C}_p, K_p),\quad
d_1^{p, q} : E_1^{p, q} \to E_1^{p + 1, q}
$$
；它收敛到 $H^{p + q}(\mathcal{C}_{total}, K)$。
此谱序列关于 $K$ 是函子性的。
\end{lemma}

\begin{proof}
设 $\mathcal{I}^\bullet$ 是表示 $K$ 的下有界内射
$\mathcal{O}$-模复形。考虑各项为
$$
A^{p, q} = \Gamma(\mathcal{C}_p, \mathcal{I}^q_p)
$$
的双复形；其水平箭头来自引理
\ref{spaces-simplicial-lemma-cech-complex-modules}，竖直箭头来自复形
$\mathcal{I}^\bullet$ 的微分。注意
% P10-E0179
$\Gamma(\mathcal{D}, -) =
\Hom_{\mathcal{O}_\mathcal{D}}(\mathcal{O}_\mathcal{D}, -)$
在 $\textit{Mod}(\mathcal{O}_\mathcal{D})$ 上成立。因此，该引理说明
双复形的各行在正次数上正合，并在次数 $0$ 上给出
$\Gamma(\mathcal{C}_{total}, \mathcal{I}^q)$。
于是，由《同调代数》引理
\ref{homology-lemma-double-complex-gives-resolution}，
与该双复形相关联的总复形计算
$R\Gamma(\mathcal{C}_{total}, K)$。
另一方面，由于到 $\mathcal{C}_p$ 的限制正合
（引理 \ref{spaces-simplicial-lemma-restriction-to-components-site}），复形
$\mathcal{I}_p^\bullet$ 在 $D(\mathcal{C}_p)$ 中表示 $K_p$。
诸层 $\mathcal{I}_p^q$ 在 $\mathcal{C}_p$ 上完全无环
（引理 \ref{spaces-simplicial-lemma-restriction-injective-to-component-limp}）。
因此，由 Leray 无环性引理（《导出范畴》引理
\ref{derived-lemma-leray-acyclicity}）以及《站点上的上同调》引理
\ref{sites-cohomology-lemma-limp-acyclic}，各列的上同调计算群
$H^q(\mathcal{C}_p, K_p)$。
应用《同调代数》引理 \ref{homology-lemma-first-quadrant-ss}
即得结论。
\end{proof}

\begin{lemma}
\label{spaces-simplicial-lemma-sanity-check-modules}
在情形 \ref{spaces-simplicial-situation-simplicial-site} 中，设 $\mathcal{O}$ 是环层，
$U \in \Ob(\mathcal{C}_n)$，且
$\mathcal{F} \in \textit{Mod}(\mathcal{O})$。
则 $H^p(U, \mathcal{F}) = H^p(U, g_n^*\mathcal{F})$，
其中在左端，$U$ 被视为 $\mathcal{C}_{total}$ 的一个对象。
\end{lemma}

\begin{proof}
注意，“将 $U$ 视为 $\mathcal{C}_{total}$ 的对象”这一说法，
在情形 (A) 下由引理 \ref{spaces-simplicial-lemma-simplicial-site-site}
中的 $\mathcal{C}_{total}$ 构造得到解释，在情形 (B) 下则由引理
\ref{spaces-simplicial-lemma-simplicial-cocontinuous-site} 中的构造得到解释。
在两种情形下，函子 $\mathcal{C}_n \to \mathcal{C}$
均连续且余连续；参见引理 \ref{spaces-simplicial-lemma-restriction-to-components-site}。
此外，按定义有 $g_n^{-1}\mathcal{O} = \mathcal{O}_n$。
因此，该结论是《站点上的上同调》引理
\ref{sites-cohomology-lemma-pullback-same-cohomology} 的特例。
\end{proof}






\section{赋环单纯站点的上同调与增广}
\label{spaces-simplicial-section-cohomology-augmentation-ringed-simplicial-sites}

\noindent
本节是第 \ref{spaces-simplicial-section-cohomology-augmentation-simplicial-sites} 节
对于模层的类似版本。

\medskip\noindent
考虑情形 \ref{spaces-simplicial-situation-simplicial-site} 所述的单纯站点 $\mathcal{C}$。
设 $a_0$ 是注 \ref{spaces-simplicial-remark-augmentation-site}
所述的朝向站点 $\mathcal{D}$ 的增广。
设 $\mathcal{O}$ 是 $\mathcal{C}_{total}$ 上的环层，
$\mathcal{O}_\mathcal{D}$ 是 $\mathcal{D}$ 上的环层。
假设给定态射
$$
a^\sharp : \mathcal{O}_\mathcal{D} \longrightarrow a_*\mathcal{O}
$$
，其中 $a$ 如引理 \ref{spaces-simplicial-lemma-augmentation-site} 所述。
从而，我们得到赋环拓扑斯的态射
$$
a :
(\Sh(\mathcal{C}_{total}), \mathcal{O})
\longrightarrow
(\Sh(\mathcal{D}), \mathcal{O}_\mathcal{D})
$$
我们按照引理 \ref{spaces-simplicial-lemma-restriction-module-to-components-site}，
将 $g_n : (\Sh(\mathcal{C}_n), \mathcal{O}_n) \to
(\Sh(\mathcal{C}_{total}), \mathcal{O})$ 视为赋环拓扑斯的态射。
取赋环拓扑斯态射的复合 $a_n = a \circ g_n$
（引理 \ref{spaces-simplicial-lemma-augmentation-site}），便得到
$$
a_n :
(\Sh(\mathcal{C}_n), \mathcal{O}_n)
\longrightarrow
(\Sh(\mathcal{D}), \mathcal{O}_\mathcal{D})
$$
利用转移映射 $f_\varphi^{-1}\mathcal{O}_m \to \mathcal{O}_n$，
我们得到赋环拓扑斯的态射
$$
f_\varphi : (\Sh(\mathcal{C}_n), \mathcal{O}_n) \to
(\Sh(\mathcal{C}_m), \mathcal{O}_m)
$$
，使得对所有 $\varphi : [m] \to [n]$，作为赋环拓扑斯的态射均有
$a_m \circ f_\varphi = a_n$。

\begin{lemma}
\label{spaces-simplicial-lemma-flat-augmentation-modules}
采用上述记号。态射
$a : (\Sh(\mathcal{C}_{total}), \mathcal{O}) \to
(\Sh(\mathcal{D}), \mathcal{O}_\mathcal{D})$
平坦，当且仅当对 $n \geq 0$，态射
$a_n : (\Sh(\mathcal{C}_n), \mathcal{O}_n) \to
(\Sh(\mathcal{D}), \mathcal{O}_\mathcal{D})$
平坦。
\end{lemma}

\begin{proof}
由于 $g_n : (\Sh(\mathcal{C}_n), \mathcal{O}_n) \to
(\Sh(\mathcal{C}_{total}), \mathcal{O})$ 平坦，可见若 $a$ 平坦，
则作为复合的 $a_n = a \circ g_n$ 也平坦。
反之，假设对所有 $n$，$a_n$ 都平坦。
我们须验证 $\mathcal{O}$ 作为 $a^{-1}\mathcal{O}_\mathcal{D}$-模层是平坦的。
设 $\mathcal{F} \to \mathcal{G}$ 是
$a^{-1}\mathcal{O}_\mathcal{D}$-模的单射。我们须证明
$$
\mathcal{F} \otimes_{a^{-1}\mathcal{O}_\mathcal{D}} \mathcal{O}
\to
\mathcal{G} \otimes_{a^{-1}\mathcal{O}_\mathcal{D}} \mathcal{O}
$$
是单射。这可以在 $\mathcal{C}_n$ 上检验，即在施用
$g_n^{-1}$ 之后检验。由于 $g_n^{-1}\mathcal{O} = \mathcal{O}_n$，
故 $g_n^* = g_n^{-1}$，于是得到
$$
g_n^{-1}\mathcal{F} \otimes_{g_n^{-1}a^{-1}\mathcal{O}_\mathcal{D}}
\mathcal{O}_n
\to
g_n^{-1}\mathcal{G} \otimes_{g_n^{-1}a^{-1}\mathcal{O}_\mathcal{D}}
\mathcal{O}_n
$$
；它是单射，因为
$g_n^{-1}a^{-1}\mathcal{O}_\mathcal{D} = a_n^{-1}\mathcal{O}_\mathcal{D}$
，且我们假设了 $a_n$ 平坦。
\end{proof}

\begin{lemma}
\label{spaces-simplicial-lemma-simplicial-resolution-augmentation-modules}
采用上述记号。对 $\mathcal{O}_\mathcal{D}$-模 $\mathcal{G}$，
存在正合复形
$$
\ldots \to
g_{2!}(a_2^*\mathcal{G}) \to
g_{1!}(a_1^*\mathcal{G}) \to
g_{0!}(a_0^*\mathcal{G}) \to
a^*\mathcal{G} \to 0
$$
；它由 $\mathcal{C}_{total}$ 上的 $\mathcal{O}$-模层组成。
这里 $g_{n!}$ 如引理
\ref{spaces-simplicial-lemma-restriction-module-to-components-site} 所述。
\end{lemma}

\begin{proof}
注意，$a^*\mathcal{G}$ 是 $\mathcal{C}_{total}$ 上的
$\mathcal{O}$-模，其到 $\mathcal{C}_m$ 的限制为
$\mathcal{O}_m$-模 $a_m^*\mathcal{G}$。
引理 \ref{spaces-simplicial-lemma-restriction-module-to-components-site}
中对模上的函子 $g_{n!}$ 的描述表明，
$g_{n!}(a_n^*\mathcal{G})$ 是 $\mathcal{C}_{total}$ 上的
$\mathcal{O}$-模，其到 $\mathcal{C}_m$ 的限制为 $\mathcal{O}_m$-模
$$
\bigoplus\nolimits_{\varphi : [n] \to [m]} f_\varphi^*a_n^*\mathcal{G} =
\bigoplus\nolimits_{\varphi : [n] \to [m]} a_m^*\mathcal{G}
$$
。其余证明与引理 \ref{spaces-simplicial-lemma-simplicial-resolution-augmentation}
的证明完全相同，只需将 $a_m^{-1}\mathcal{G}$ 替换为
$a_m^*\mathcal{G}$。
\end{proof}

\begin{lemma}
\label{spaces-simplicial-lemma-augmentation-cech-complex-modules}
采用上述记号。对 $\mathcal{C}_{total}$ 上的
$\mathcal{O}$-模 $\mathcal{F}$，存在典范复形
$$
0 \to a_*\mathcal{F} \to a_{0, *}\mathcal{F}_0 \to a_{1, *}\mathcal{F}_1 \to
a_{2, *}\mathcal{F}_2 \to \ldots
$$
；它由 $\mathcal{O}_\mathcal{D}$-模组成，并在次数 $-1, 0$ 上正合。
若 $\mathcal{F}$ 是内射 $\mathcal{O}$-模，则该复形在所有次数上正合，
且对任意 $\mathcal{O}_\mathcal{D}$-模 $\mathcal{G}$，
施用函子 $\Hom_{\mathcal{O}_\mathcal{D}}(\mathcal{G}, -)$ 后仍正合。
\end{lemma}

\begin{proof}
为构造该复形，由 Yoneda 引理，只需对 $\mathcal{D}$ 上任意
$\mathcal{O}_\mathcal{D}$-模 $\mathcal{G}$ 构造复形
$$
0 \to \Hom_{\mathcal{O}_\mathcal{D}}(\mathcal{G}, a_*\mathcal{F}) \to
\Hom_{\mathcal{O}_\mathcal{D}}(\mathcal{G}, a_{0, *}\mathcal{F}_0) \to
\Hom_{\mathcal{O}_\mathcal{D}}(\mathcal{G}, a_{1, *}\mathcal{F}_1) \to \ldots
$$
，并使其关于 $\mathcal{G}$ 是函子性的。为此，将
$\Hom_\mathcal{O}(-, \mathcal{F})$ 作用于引理
\ref{spaces-simplicial-lemma-simplicial-resolution-augmentation-modules} 的正合复形，
并使用拉回与推出的伴随性。
由于 $\Hom_\mathcal{O}(-, \mathcal{F})$ 左正合，次数 $-1, 0$
上的正合性由该构造得出。若 $\mathcal{F}$ 是内射
$\mathcal{O}$-模，则 $\Hom_\mathcal{O}(-, \mathcal{F})$ 正合，
故此复形正合。
\end{proof}

\begin{lemma}
\label{spaces-simplicial-lemma-augmentation-spectral-sequence-modules}
采用上述记号。对任意属于 $D^+(\mathcal{O})$ 的 $K$，
在 $\textit{Mod}(\mathcal{O}_\mathcal{D})$ 中存在谱序列
$(E_r, d_r)_{r \geq 0}$，其中
$$
E_1^{p, q} = R^qa_{p, *} K_p
$$
；它收敛到 $R^{p + q}a_*K$。此谱序列关于 $K$ 是函子性的。
\end{lemma}

\begin{proof}
设 $\mathcal{I}^\bullet$ 是表示 $K$ 的下有界内射
$\mathcal{O}$-模复形。考虑各项为
$$
A^{p, q} = a_{p, *}\mathcal{I}^q_p
$$
的双复形；其水平箭头来自引理
\ref{spaces-simplicial-lemma-augmentation-cech-complex-modules}，
竖直箭头来自复形 $\mathcal{I}^\bullet$ 的微分。
该引理说明，双复形的各行在正次数上正合，并在次数 $0$ 上给出
$a_*\mathcal{I}^q$。于是，由《同调代数》引理
\ref{homology-lemma-double-complex-gives-resolution}，
与该双复形相关联的总复形计算 $Ra_*K$。
另一方面，由于到 $\mathcal{C}_p$ 的限制正合
（引理 \ref{spaces-simplicial-lemma-restriction-to-components-site}），复形
$\mathcal{I}_p^\bullet$ 在 $D(\mathcal{C}_p)$ 中表示 $K_p$。
诸层 $\mathcal{I}_p^q$ 在 $\mathcal{C}_p$ 上完全无环
（引理 \ref{spaces-simplicial-lemma-restriction-injective-to-component-limp}）。
因此，由 Leray 无环性引理（《导出范畴》引理
\ref{derived-lemma-leray-acyclicity}）以及《站点上的上同调》引理
\ref{sites-cohomology-lemma-limp-acyclic}，各列的上同调为层
$R^qa_{p, *}K_p$。应用《同调代数》引理
\ref{homology-lemma-first-quadrant-ss} 即得结论。
\end{proof}





\section{笛卡尔层与模}
\label{spaces-simplicial-section-cartesian}

\noindent
定义如下。

\begin{definition}
\label{spaces-simplicial-definition-cartesian-sheaf}
在情形 \ref{spaces-simplicial-situation-simplicial-site} 中。
\begin{enumerate}
\item 若对所有 $\varphi : [m] \to [n]$，映射
$\mathcal{F}(\varphi) : f_\varphi^{-1}\mathcal{F}_m \to \mathcal{F}_n$
都是同构，则称 $\mathcal{C}_{total}$ 上的集合层或阿贝尔群层
$\mathcal{F}$ 是{\it 笛卡尔的}。
\item 若 $\mathcal{O}$ 是 $\mathcal{C}_{total}$ 上的环层，
且对所有 $\varphi : [m] \to [n]$，映射
$f_\varphi^*\mathcal{F}_m \to \mathcal{F}_n$ 都是同构，
则称 $\mathcal{O}$-模层 $\mathcal{F}$ 是{\it 笛卡尔的}。
\item 若对所有 $\varphi : [m] \to [n]$，映射
$f_\varphi^{-1}K_m \to K_n$
都是同构，则称 $D(\mathcal{C}_{total})$ 的对象 $K$
是{\it 笛卡尔的}。
\item 若 $\mathcal{O}$ 是 $\mathcal{C}_{total}$ 上的环层，
且对所有 $\varphi : [m] \to [n]$，映射
$Lf_\varphi^*K_m \to K_n$
都是同构，则称 $D(\mathcal{O})$ 的对象 $K$ 是{\it 笛卡尔的}。
\end{enumerate}
\end{definition}

\noindent
当然，纤维化范畴的笛卡尔截面有一个一般概念，而以上仅是其若干例子。
% P10-E0181
关于拉回的这一性质只需对退化映射检验。

\begin{lemma}
\label{spaces-simplicial-lemma-check-cartesian-module}
在情形 \ref{spaces-simplicial-situation-simplicial-site} 中。
\begin{enumerate}
\item 集合层或阿贝尔群层 $\mathcal{F}$ 是笛卡尔的，当且仅当映射
$(f_{\delta^n_j})^{-1}\mathcal{F}_{n - 1} \to \mathcal{F}_n$
是同构。
\item $D(\mathcal{C}_{total})$ 的对象 $K$ 是笛卡尔的，
当且仅当映射
$(f_{\delta^n_j})^{-1}K_{n - 1} \to K_n$
是同构。
\item 若 $\mathcal{O}$ 是 $\mathcal{C}_{total}$ 上的环层，
则 $\mathcal{O}$-模层 $\mathcal{F}$ 是笛卡尔的，当且仅当映射
$(f_{\delta^n_j})^*\mathcal{F}_{n - 1} \to \mathcal{F}_n$
是同构。
\item 若 $\mathcal{O}$ 是 $\mathcal{C}_{total}$ 上的环层，
则 $D(\mathcal{O})$ 的对象 $K$ 是笛卡尔的，当且仅当映射
$L(f_{\delta^n_j})^*K_{n - 1} \to K_n$
是同构。
% P10-E0182
\item 在此补充更多内容。
\end{enumerate}
\end{lemma}

\begin{proof}
在每一种情形下，关键都是拉回函子的复合仍为拉回函子；
关于 (4)，参见《站点上的上同调》引理
\ref{sites-cohomology-lemma-derived-pullback-composition}。
我们说明 (1) 中的论证，并略去其余情形的证明。
% P10-E0183
范畴 $\Delta$ 由态射 $\delta^n_j$ 与 $\sigma^n_j$ 生成；
参见《单纯对象》引理 \ref{simplicial-lemma-face-degeneracy}。
因此，只需检验映射
$(f_{\delta^n_j})^{-1}\mathcal{F}_{n - 1} \to \mathcal{F}_n$
与 $(f_{\sigma^n_j})^{-1}\mathcal{F}_{n + 1} \to \mathcal{F}_n$
是同构；关于记号，参见《单纯对象》引理
\ref{simplicial-lemma-characterize-simplicial-object}。
由于 $\sigma^n_j \circ \delta_j^{n + 1} = \text{id}_{[n]}$，复合
$$
\mathcal{F}_n =
(f_{\sigma^n_j})^{-1}
(f_{\delta_j^{n + 1}})^{-1}
\mathcal{F}_n \to
(f_{\sigma^n_j})^{-1}
\mathcal{F}_{n + 1} \to
\mathcal{F}_n
$$
是恒等映射。因此，对 $\delta^{n + 1}_j$ 的结论蕴含
对 $\sigma^n_j$ 的结论。
\end{proof}

\begin{lemma}
\label{spaces-simplicial-lemma-augmentation-cartesian-module}
在情形 \ref{spaces-simplicial-situation-simplicial-site} 中，设 $a_0$ 是注
\ref{spaces-simplicial-remark-augmentation-site} 所述的朝向站点 $\mathcal{D}$ 的增广。
\begin{enumerate}
\item $\mathcal{D}$ 上集合层或阿贝尔群层的拉回
$a^{-1}\mathcal{G}$ 是笛卡尔的。
\item $D(\mathcal{D})$ 的对象 $K$ 的拉回 $a^{-1}K$ 是笛卡尔的。
\end{enumerate}
设 $\mathcal{O}$ 是 $\mathcal{C}_{total}$ 上的环层，
$\mathcal{O}_\mathcal{D}$ 是 $\mathcal{D}$ 上的环层，且
$a^\sharp : \mathcal{O}_\mathcal{D} \to a_*\mathcal{O}$ 是第
\ref{spaces-simplicial-section-cohomology-augmentation-ringed-simplicial-sites} 节所述的态射。
\begin{enumerate}
\item[(3)] $\mathcal{O}_\mathcal{D}$-模层的拉回
$a^*\mathcal{F}$ 是笛卡尔的。
\item[(4)] $D(\mathcal{O}_\mathcal{D})$ 的对象 $K$
的导出拉回 $La^*K$ 是笛卡尔的。
\end{enumerate}
\end{lemma}

\begin{proof}
这立即由对所有 $\varphi : [m] \to [n]$ 成立的等式
$a_m \circ f_\varphi = a_n$ 得出。参见引理
\ref{spaces-simplicial-lemma-augmentation-site} 以及第
\ref{spaces-simplicial-section-cohomology-augmentation-ringed-simplicial-sites} 节中的讨论。
\end{proof}

\begin{lemma}
\label{spaces-simplicial-lemma-characterize-cartesian}
在情形 \ref{spaces-simplicial-situation-simplicial-site} 中，
笛卡尔集合层（相应地，笛卡尔阿贝尔群层）范畴等价于
二元组 $(\mathcal{F}, \alpha)$ 的范畴，其中 $\mathcal{F}$ 是
$\mathcal{C}_0$ 上的集合层（相应地，阿贝尔群层），且
$$
\alpha :
(f_{\delta_1^1})^{-1}\mathcal{F}
\longrightarrow (f_{\delta_0^1})^{-1}\mathcal{F}
$$
是 $\mathcal{C}_1$ 上集合层（相应地，阿贝尔群层）的同构，满足
$(f_{\delta^2_1})^{-1}\alpha =
(f_{\delta^2_0})^{-1}\alpha \circ (f_{\delta^2_2})^{-1}\alpha$
；这里等式两端都是 $\mathcal{C}_2$ 上层的映射。
\end{lemma}

\begin{proof}
简记
$d^n_j = f_{\delta^n_j} : \Sh(\mathcal{C}_n) \to \Sh(\mathcal{C}_{n - 1})$.
引理陈述中关于 $\alpha$ 的条件是有意义的，因为
$$
d^1_1 \circ d^2_2 = d^1_1 \circ d^2_1, \quad
d^1_1 \circ d^2_0 = d^1_0 \circ d^2_2, \quad
d^1_0 \circ d^2_0 = d^1_0 \circ d^2_1
$$
是拓扑斯态射 $\Sh(\mathcal{C}_2) \to \Sh(\mathcal{C}_0)$ 的等式；
参见《单纯对象》注 \ref{simplicial-remark-relations}。
因此，可将这些映射表示为
$$
\xymatrix{
& (d^2_0)^{-1}(d^1_1)^{-1}\mathcal{F} \ar[r]_-{(d^2_0)^{-1}\alpha} &
(d^2_0)^{-1}(d^1_0)^{-1}\mathcal{F} \ar@{=}[rd] & \\
(d^2_2)^{-1}(d^1_0)^{-1}\mathcal{F} \ar@{=}[ru] & & &
(d^2_1)^{-1}(d^1_0)^{-1}\mathcal{F} \\
& (d^2_2)^{-1}(d^1_1)^{-1}\mathcal{F} \ar[lu]^{(d^2_2)^{-1}\alpha} \ar@{=}[r] &
(d^2_1)^{-1}(d^1_1)^{-1}\mathcal{F} \ar[ru]_{(d^2_1)^{-1}\alpha}
}
$$
，而该条件表示此图交换。显然，给定 $\mathcal{C}_{total}$ 上的
笛卡尔集合层（相应地，阿贝尔群层）$\mathcal{G}$，
可以令 $\mathcal{F} = \mathcal{G}_0$，并令 $\alpha$ 等于复合
% P10-E0184
$$
(d_1^1)^{-1}\mathcal{G}_0 \to \mathcal{G}_1
\leftarrow (d_1^0)^{-1}\mathcal{G}_0
$$
；由于 $\mathcal{G}$ 是笛卡尔的，图中箭头均可逆。
为证明此函子是等价，我们构造一个拟逆。
拟逆的构造类似于《下降》第 \ref{descent-section-descent-modules} 节
讨论的构造；我们借用其中的记号
$\tau^n_i : [0] \to [n]$, $0 \mapsto i$ 以及
$\tau^n_{ij} : [1] \to [n]$, $0 \mapsto i$, $1 \mapsto j$。
具体而言，给定引理所述的二元组 $(\mathcal{F}, \alpha)$，令
$\mathcal{G}_n = (f_{\tau^n_n})^{-1}\mathcal{F}$。
给定 $\varphi : [n] \to [m]$，用下图定义
$\mathcal{G}(\varphi) : (f_\varphi)^{-1}\mathcal{G}_n \to \mathcal{G}_m$：
$$
\xymatrix{
(f_\varphi)^{-1}\mathcal{G}_n \ar@{=}[r] &
(f_\varphi)^{-1}(f_{\tau^n_n})^{-1}\mathcal{F} \ar@{=}[r] &
(f_{\tau^m_{\varphi(n)}})^{-1}\mathcal{F} \ar@{=}[r] &
(f_{\tau^m_{\varphi(n)m}})^{-1}(d^1_1)^{-1}\mathcal{F}
\ar[d]^{(f_{\tau^m_{\varphi(n)m}})^{-1}\alpha} \\
&
\mathcal{G}_m \ar@{=}[r] &
(f_{\tau^m_m})^{-1}\mathcal{F} \ar@{=}[r] &
(f_{\tau^m_{\varphi(n)m}})^{-1}(d^1_0)^{-1}\mathcal{F}
}
$$
我们略去如下验证：上图的交换性蕴含这些映射能够正确复合，
因而给出 $\mathcal{C}_{total}$ 上的层；参见引理
\ref{spaces-simplicial-lemma-describe-sheaves-simplicial-site-site}。
我们也略去这两个函子互为拟逆的验证。
\end{proof}

\begin{lemma}
\label{spaces-simplicial-lemma-characterize-cartesian-modules}
在情形 \ref{spaces-simplicial-situation-simplicial-site} 中，设 $\mathcal{O}$ 是
$\mathcal{C}_{total}$ 上的环层。
笛卡尔 $\mathcal{O}$-模范畴等价于二元组
$(\mathcal{F}, \alpha)$ 的范畴，其中 $\mathcal{F}$ 是
$\mathcal{O}_0$-模，且
$$
\alpha :
(f_{\delta_1^1})^*\mathcal{F}
\longrightarrow (f_{\delta_0^1})^*\mathcal{F}
$$
是 $\mathcal{O}_1$-模的同构，满足
$(f_{\delta^2_1})^*\alpha =
(f_{\delta^2_0})^*\alpha \circ (f_{\delta^2_2})^*\alpha$
；这里等式两端都是 $\mathcal{O}_2$-模映射。
\end{lemma}

\begin{proof}
证明与引理 \ref{spaces-simplicial-lemma-characterize-cartesian} 的证明相同，
只需将阿贝尔群层的拉回替换为模的拉回。
\end{proof}

\begin{lemma}
\label{spaces-simplicial-lemma-Serre-subcat-cartesian-modules}
在情形 \ref{spaces-simplicial-situation-simplicial-site} 中。
\begin{enumerate}
\item 笛卡尔阿贝尔层的满子范畴构成
$\textit{Ab}(\mathcal{C}_{total})$ 的弱 Serre 子范畴。
笛卡尔阿贝尔层系统的余极限仍是笛卡尔的。
\item 设 $\mathcal{O}$ 是 $\mathcal{C}_{total}$ 上的环层，
并且态射
$$
f_{\delta^n_j} : (\Sh(\mathcal{C}_n), \mathcal{O}_n)
\to (\Sh(\mathcal{C}_{n - 1}), \mathcal{O}_{n - 1})
$$
平坦。则笛卡尔 $\mathcal{O}$-模的满子范畴构成
$\textit{Mod}(\mathcal{O})$ 的弱 Serre 子范畴。
笛卡尔 $\mathcal{O}$-模系统的余极限仍是笛卡尔的。
\end{enumerate}
\end{lemma}

\begin{proof}
为说明 (1) 中得到弱 Serre 子范畴，我们检验《同调代数》引理
\ref{homology-lemma-characterize-weak-serre-subcategory} 所列的条件。
首先，若 $\varphi : \mathcal{F} \to \mathcal{G}$ 是笛卡尔阿贝尔层间的映射，
则 $\Ker(\varphi)$ 与 $\Coker(\varphi)$ 也是笛卡尔的，
因为限制函子
$\Sh(\mathcal{C}_{total}) \to \Sh(\mathcal{C}_n)$
以及函子 $f_\varphi^{-1}$ 都正合。类似地，若
$$
0 \to \mathcal{F} \to \mathcal{H} \to \mathcal{G} \to 0
$$
是 $\mathcal{C}_{total}$ 上阿贝尔层的短正合列，且
$\mathcal{F}$ 与 $\mathcal{G}$ 都是笛卡尔的，则由五引理可知
$\mathcal{H}$ 也是笛卡尔的。关于余极限的性质，使用如下事实：
拉回是左伴随，故余极限与拉回交换。
在模的情形下，采用同样的论证，并使用平坦拉回的正合性
（《站点上的模》引理 \ref{sites-modules-lemma-flat-pullback-exact}），
以及只需对 $f_{\delta^n_j}$ 检验该条件这一事实；
参见引理 \ref{spaces-simplicial-lemma-check-cartesian-module}。
\end{proof}

\begin{remark}[警告]
\label{spaces-simplicial-remark-warning-cartesian-modules}
尽管有引理 \ref{spaces-simplicial-lemma-Serre-subcat-cartesian-modules}，
笛卡尔 $\mathcal{O}$-模范畴仍可能是阿贝尔范畴，却并非
$\textit{Mod}(\mathcal{O})$ 的 Serre 子范畴。
具体而言，假设只知道 $f_{\delta_1^1}$ 与 $f_{\delta_0^1}$ 平坦。
由引理 \ref{spaces-simplicial-lemma-characterize-cartesian-modules} 很容易推出，
笛卡尔 $\mathcal{O}$-模范畴是阿贝尔范畴。
但若（例如）$f_{\delta_0^2}$ 不平坦，则没有理由认为从
笛卡尔 $\mathcal{O}$-模范畴到全体 $\mathcal{O}$-模范畴的包含函子正合。
\end{remark}

\begin{lemma}
\label{spaces-simplicial-lemma-derived-cartesian-modules}
在情形 \ref{spaces-simplicial-situation-simplicial-site} 中。
\begin{enumerate}
\item $D(\mathcal{C}_{total})$ 的对象 $K$ 是笛卡尔的，
当且仅当对所有 $q$，$H^q(K)$ 都是笛卡尔阿贝尔层。
\item 设 $\mathcal{O}$ 是 $\mathcal{C}_{total}$ 上的环层，
并且态射
$f_{\delta^n_j} : (\Sh(\mathcal{C}_n), \mathcal{O}_n)
\to (\Sh(\mathcal{C}_{n - 1}), \mathcal{O}_{n - 1})$ 平坦。
则 $D(\mathcal{O})$ 的对象 $K$ 是笛卡尔的，
当且仅当对所有 $q$，$H^q(K)$ 都是笛卡尔 $\mathcal{O}$-模。
\end{enumerate}
\end{lemma}

\begin{proof}
(1) 成立是因为拉回函子 $(f_\varphi)^{-1}$ 正合。
(2) 由引理 \ref{spaces-simplicial-lemma-check-cartesian-module} 中的刻画，
以及平坦性所给出的
$L(f_{\delta^n_j})^* = (f_{\delta^n_j})^*$ 得出。
\end{proof}

\begin{lemma}
\label{spaces-simplicial-lemma-derived-cartesian-shriek}
在情形 \ref{spaces-simplicial-situation-simplicial-site} 中。
\begin{enumerate}
\item $D(\mathcal{C}_{total})$ 的对象 $K$ 是笛卡尔的，
当且仅当对所有 $n$，典范映射
$$
g_{n!}K_n \longrightarrow
g_{n!}\mathbf{Z} \otimes^\mathbf{L}_\mathbf{Z} K
$$
是同构。
\item 设 $\mathcal{O}$ 是 $\mathcal{C}_{total}$ 上的环层，
并且对所有 $\varphi : [n] \to [m]$，态射
$f_\varphi^{-1}\mathcal{O}_n \to \mathcal{O}_m$ 都平坦。
则 $D(\mathcal{O})$ 的对象 $K$ 是笛卡尔的，
当且仅当对所有 $n$，典范映射
$$
g_{n!}K_n \longrightarrow
g_{n!}\mathcal{O}_n \otimes^\mathbf{L}_\mathcal{O} K
$$
是同构。
\end{enumerate}
\end{lemma}

\begin{proof}
证明 (1)。由于 $g_{n!}$ 正合，它在导出范畴上诱导一个与
$g_n^{-1}$ 伴随的函子。该映射是映射
$K_n \to (g_n^{-1}g_{n!}\mathbf{Z}) \otimes^\mathbf{L}_\mathbf{Z} K_n$
的伴随；后者对应于 $\mathbf{Z} \to g_n^{-1}g_{n!}\mathbf{Z}$，
而这一映射又与
$\text{id} : g_{n!}\mathbf{Z} \to g_{n!}\mathbf{Z}$ 伴随。
利用引理 \ref{spaces-simplicial-lemma-restriction-to-components-site}
中对 $g_{n!}$ 的描述，可见该映射到 $\mathcal{C}_m$ 的限制为
$$
\bigoplus\nolimits_{\varphi : [n] \to [m]} f_\varphi^{-1}K_n
\longrightarrow
\bigoplus\nolimits_{\varphi : [n] \to [m]} K_m
$$
因此，该陈述是清楚的。

\medskip\noindent
证明 (2)。由于 $g_{n!}$ 正合
（引理 \ref{spaces-simplicial-lemma-exactness-g-shriek-modules}），
它在导出范畴上诱导一个与同样正合的 $g_n^*$ 伴随的函子。
该映射是映射
$K_n \to (g_n^*g_{n!}\mathcal{O}_n) \otimes^\mathbf{L}_{\mathcal{O}_n} K_n$
的伴随；后者对应于 $\mathcal{O}_n \to g_n^*g_{n!}\mathcal{O}_n$，
而这一映射又与
$\text{id} : g_{n!}\mathcal{O}_n \to g_{n!}\mathcal{O}_n$ 伴随。
利用引理 \ref{spaces-simplicial-lemma-restriction-module-to-components-site}
中对 $g_{n!}$ 的描述，可见该映射到 $\mathcal{C}_m$ 的限制为
$$
\bigoplus\nolimits_{\varphi : [n] \to [m]} f_\varphi^*K_n
\longrightarrow
\bigoplus\nolimits_{\varphi : [n] \to [m]}
f_\varphi^*\mathcal{O}_n \otimes_{\mathcal{O}_m} K_m =
\bigoplus\nolimits_{\varphi : [n] \to [m]} K_m
$$
因此，该陈述是清楚的。
\end{proof}

\begin{lemma}
\label{spaces-simplicial-lemma-quasi-coherent-sheaf}
在情形 \ref{spaces-simplicial-situation-simplicial-site} 中，设 $\mathcal{O}$ 是
$\mathcal{C}_{total}$ 上的环层，$\mathcal{F}$ 是 $\mathcal{O}$-模层。
则 $\mathcal{F}$ 在《站点上的模》定义
\ref{sites-modules-definition-site-local} 的意义下拟凝聚，
当且仅当 $\mathcal{F}$ 是笛卡尔的，且对所有 $n$，
$\mathcal{F}_n$ 都是拟凝聚 $\mathcal{O}_n$-模。
\end{lemma}

\begin{proof}
假设 $\mathcal{F}$ 拟凝聚。由于拟凝聚模的拉回仍拟凝聚
（《站点上的模》引理 \ref{sites-modules-lemma-local-pullback}），
可见对所有 $n$，$\mathcal{F}_n$ 都是拟凝聚
$\mathcal{O}_n$-模。为证明 $\mathcal{F}$ 是笛卡尔的，
设 $U$ 是 $\mathcal{C}_n$ 的对象，其中 $n$ 为某个整数，并将 $U$
视为 $\mathcal{C}_{total}$ 的对象。由于 $\mathcal{F}$ 拟凝聚，
存在覆盖 $\{U_i \to U\}$，并且对每个 $i$，存在表现
$$
\bigoplus\nolimits_{j \in J_i} \mathcal{O}_{\mathcal{C}_{total}/U_i} \to
\bigoplus\nolimits_{k \in K_i} \mathcal{O}_{\mathcal{C}_{total}/U_i} \to
\mathcal{F}|_{\mathcal{C}_{total}/U_i} \to 0
$$
。由站点 $\mathcal{C}_{total}$ 的构造可知，
$\{U_i \to U\}$ 是 $\mathcal{C}_n$ 中的覆盖。
接着，设 $V$ 是 $\mathcal{C}_m$ 的对象，其中 $m$ 为某个整数，并设
$V \to U$ 是 $\mathcal{C}_{total}$ 中位于
$\varphi : [n] \to [m]$ 之上的态射。纤维积
$V_i = V \times_U U_i$ 存在，并在 $\mathcal{C}_m$ 中得到诱导覆盖
$\{V_i \to V\}$。将上述表现限制到站点
$\mathcal{C}_n/U_i$ 与 $\mathcal{C}_m/V_i$，得到表现
% P10-E0185
$$
\bigoplus\nolimits_{j \in J_i} \mathcal{O}_{\mathcal{C}_m/U_i} \to
\bigoplus\nolimits_{k \in K_i} \mathcal{O}_{\mathcal{C}_m/U_i} \to
\mathcal{F}_n|_{\mathcal{C}_n/U_i} \to 0
$$
以及
$$
\bigoplus\nolimits_{j \in J_i} \mathcal{O}_{\mathcal{C}_m/V_i} \to
\bigoplus\nolimits_{k \in K_i} \mathcal{O}_{\mathcal{C}_m/V_i} \to
\mathcal{F}_m|_{\mathcal{C}_m/V_i} \to 0
$$
。这些表现与映射
$\mathcal{F}(\varphi) : f_\varphi^*\mathcal{F}_n \to \mathcal{F}_m$
相容（因为该映射利用 $\mathcal{F}$ 沿
$\mathcal{C}_{total}$ 中位于 $\varphi$ 之上的态射的限制映射来定义）。
由此可知 $\mathcal{F}(\varphi)|_{\mathcal{C}_m/V_i}$ 是同构。
由于 $\{V_i \to V\}$ 是覆盖，可知
$\mathcal{F}(\varphi)|_{\mathcal{C}_m/V}$ 是同构。
因 $V$ 与 $U$ 任意，这证明了 $\mathcal{F}$ 是笛卡尔的。
（在情形 A 下使用《站点》引理
\ref{sites-lemma-morphism-of-sites-covering}。）

\medskip\noindent
反之，假设对所有 $n$，$\mathcal{F}_n$ 都拟凝聚，
且 $\mathcal{F}$ 是笛卡尔的。于是，对任意 $n$ 及
$\mathcal{C}_n$ 的对象 $U$，可以选取 $\mathcal{C}_n$ 中的覆盖
$\{U_i \to U\}$，并对每个 $i$ 选取表现
% P10-E0186
$$
\bigoplus\nolimits_{j \in J_i} \mathcal{O}_{\mathcal{C}_m/U_i} \to
\bigoplus\nolimits_{k \in K_i} \mathcal{O}_{\mathcal{C}_m/U_i} \to
\mathcal{F}_n|_{\mathcal{C}_n/U_i} \to 0
$$
。将其拉回到 $\mathcal{C}_{total}/U_i$，得到模复形
$$
\bigoplus\nolimits_{j \in J_i} \mathcal{O}_{\mathcal{C}_{total}/U_i} \to
\bigoplus\nolimits_{k \in K_i} \mathcal{O}_{\mathcal{C}_{total}/U_i} \to
\mathcal{F}|_{\mathcal{C}_{total}/U_i} \to 0
$$
。这是 $\mathcal{C}_{total}/U_i$ 上的模复形。
$\mathcal{F}$ 的笛卡尔性蕴含此复形正合。细节从略。
\end{proof}











\section{导出范畴的单纯系统}
\label{spaces-simplicial-section-glueing}

\noindent
本节将在阿贝尔层导出范畴的框架下证明
\cite[Proposition 3.2.9]{BBD} 的一个特例。
模的情形在第 \ref{spaces-simplicial-section-glueing-modules} 节中讨论。

\begin{definition}
\label{spaces-simplicial-definition-cartesian-derived}
在情形 \ref{spaces-simplicial-situation-simplicial-site} 中，
一个{\it 导出范畴的单纯系统}由下列数据组成：
\begin{enumerate}
\item 对每个 $n$，给定 $D(\mathcal{C}_n)$ 的对象 $K_n$；
\item 对每个 $\varphi : [m] \to [n]$，给定
$D(\mathcal{C}_n)$ 中的映射
$K_\varphi : f_\varphi^{-1}K_m \to K_n$。
\end{enumerate}
这些数据须满足：
$$
K_{\varphi \circ \psi} = K_\varphi \circ f_\varphi^{-1}K_\psi :
f_{\varphi \circ \psi}^{-1}K_l = f_\varphi^{-1} f_\psi^{-1}K_l
\longrightarrow
K_n
$$
对 $\Delta$ 的任意态射 $\varphi : [m] \to [n]$ 与
$\psi : [l] \to [m]$，上述等式均成立。
若对所有 $\varphi$，映射 $K_\varphi$ 都是同构，
则称该单纯系统是{\it 笛卡尔的}。
对两个导出范畴的单纯系统，显然可以定义
{\it 导出范畴的单纯系统的态射}。
\end{definition}

\noindent
我们有意给这个概念取了一个长得荒唐的名字。
目标是证明：在某些假设下，导出范畴的单纯系统来自
$D(\mathcal{C}_{total})$ 的一个对象。

\begin{lemma}
\label{spaces-simplicial-lemma-cartesian-objects-derived}
在情形 \ref{spaces-simplicial-situation-simplicial-site} 中，若
$K \in D(\mathcal{C}_{total})$，则
$(K_n, K(\varphi))$ 是导出范畴的单纯系统。
若 $K$ 是笛卡尔的，则该系统也是笛卡尔的。
\end{lemma}

\begin{proof}
这是显然的。
\end{proof}

\begin{lemma}
\label{spaces-simplicial-lemma-construct-cartesian-objects-derived}
在情形 \ref{spaces-simplicial-situation-simplicial-site} 中，假设给定
$K_0 \in D(\mathcal{C}_0)$ 以及同构
$$
\alpha :
f_{\delta_1^1}^{-1}K_0
\longrightarrow
f_{\delta_0^1}^{-1}K_0
$$
，它满足余圈条件。令
$\tau^n_i : [0] \to [n]$, $0 \mapsto i$，并令
$K_n = f_{\tau^n_n}^{-1}K_0$。则诸 $K_n$
构成导出范畴的一个笛卡尔单纯系统。
\end{lemma}

\begin{proof}
请与引理 \ref{spaces-simplicial-lemma-characterize-cartesian} 及其证明比较
（其中也明确写出了余圈条件）。
该构造类似于《下降》第 \ref{descent-section-descent-modules} 节
讨论的构造；我们借用其中的记号
$\tau^n_i : [0] \to [n]$, $0 \mapsto i$ 以及
$\tau^n_{ij} : [1] \to [n]$, $0 \mapsto i$, $1 \mapsto j$。
给定 $\varphi : [n] \to [m]$，用下图定义
$K_\varphi : f_\varphi^{-1}K_n \to K_m$：
$$
\xymatrix{
f_\varphi^{-1}K_n \ar@{=}[r] &
f_\varphi^{-1} f_{\tau^n_n}^{-1}K_0 \ar@{=}[r] &
f_{\tau^m_{\varphi(n)}}^{-1}K_0 \ar@{=}[r] &
f_{\tau^m_{\varphi(n)m}}^{-1}f_{\delta^1_1}^{-1}K_0
\ar[d]_{f_{\tau^m_{\varphi(n)m}}^{-1}\alpha} \\
&
K_m \ar@{=}[r] &
f_{\tau^m_m}^{-1}K_0 \ar@{=}[r] &
f_{\tau^m_{\varphi(n)m}}^{-1}f_{\delta^1_0}^{-1}K_0
}
$$
我们略去如下验证：余圈条件蕴含这些映射
（在各自的导出范畴中）能够正确复合，因而给出
导出范畴的一个单纯系统。
\end{proof}

\begin{lemma}
\label{spaces-simplicial-lemma-abelian-postnikov}
在情形 \ref{spaces-simplicial-situation-simplicial-site} 中，设 $K$ 是
$D(\mathcal{C}_{total})$ 的对象。令
$$
X_n = (g_{n!}\mathbf{Z})
\otimes^\mathbf{L}_\mathbf{Z} K
\quad\text{且}\quad
Y_n =
(g_{n!}\mathbf{Z} \to \ldots \to g_{0!}\mathbf{Z})[-n]
\otimes^\mathbf{L}_\mathbf{Z} K
$$
为 $D(\mathcal{C}_{total})$ 的对象，其中的映射如引理
\ref{spaces-simplicial-lemma-simplicial-resolution-Z-site} 所述。
配备显然的典范映射 $Y_n \to X_n$ 以及
$Y_0 \to Y_1[1] \to Y_2[2] \to \ldots$ 后，有：
\begin{enumerate}
\item 区别三角 $Y_n \to X_n \to Y_{n - 1} \to Y_n[1]$
为 $\ldots \to X_2 \to X_1 \to X_0$ 定义一个 Postnikov 系统
（《导出范畴》定义 \ref{derived-definition-postnikov-system}）；
\item 在 $D(\mathcal{C}_{total})$ 中，
$K = \text{hocolim} Y_n[n]$。
\end{enumerate}
\end{lemma}

\begin{proof}
首先，若 $K = \mathbf{Z}$，则这正是将《导出范畴》例
\ref{derived-example-key-postnikov} 的构造应用于
$\textit{Ab}(\mathcal{C}_{total})$ 中的复形
$$
\ldots \to
g_{2!}\mathbf{Z} \to
g_{1!}\mathbf{Z} \to
g_{0!}\mathbf{Z}
$$
，再结合如下事实：由引理 \ref{spaces-simplicial-lemma-simplicial-resolution-Z-site}，
此复形在 $D(\mathcal{C}_{total})$ 中表示 $K = \mathbf{Z}$。
一般情形由此以及下列两个事实得出：正合函子
$- \otimes^\mathbf{L}_\mathbf{Z} K$ 将 Postnikov 系统映为
Postnikov 系统；函子
$- \otimes^\mathbf{L}_\mathbf{Z} K$ 与同伦余极限交换。
\end{proof}

\begin{lemma}
\label{spaces-simplicial-lemma-nullity-cartesian-objects-derived}
在情形 \ref{spaces-simplicial-situation-simplicial-site} 中。
% P10-E0187
设 $K, K' \in D(\mathcal{C}_{total})$。假设
\begin{enumerate}
\item $K$ 是笛卡尔的；
\item 当 $i > 0$ 时，$\Hom(K_i[i], K'_i) = 0$；
\item 当 $i \geq 0$ 时，$\Hom(K_i[i + 1], K'_i) = 0$。
\end{enumerate}
则任何诱导零映射 $K_0 \to K'_0$ 的映射 $K \to K'$ 都是零映射。
\end{lemma}

\begin{proof}
考虑引理 \ref{spaces-simplicial-lemma-abelian-postnikov} 中与 $K$ 相关联的对象
$X_n$ 及 Postnikov 系统 $Y_n$。
由于 $K = \text{hocolim} Y_n[n]$，映射 $K \to K'$
诱导相容的态射族 $Y_n[n] \to K'$。
由 (1) 与引理 \ref{spaces-simplicial-lemma-derived-cartesian-shriek}，有
$X_n = g_{n!}K_n$。由于 $Y_0 = X_0$，可见
$K_0 \to K'_0$ 为零蕴含 $Y_0 \to K'$ 为零。
假设已经证明对某个 $n \geq 0$，映射 $Y_n[n] \to K'$ 为零。
由区别三角
$$
Y_n[n] \to Y_{n + 1}[n + 1] \to X_{n + 1}[n + 1] \to Y_n[n + 1]
$$
得到正合列
$$
\Hom(X_{n + 1}[n + 1], K') \to
\Hom(Y_{n + 1}[n + 1], K') \to
\Hom(Y_n[n], K')
$$
由于 $X_{n + 1}[n + 1] = g_{n + 1!}K_{n + 1}[n + 1]$，
第一个群等于
$$
\Hom(K_{n + 1}[n + 1], K'_{n + 1})
$$
，它由假设 (2) 为零。归纳可知所有映射
$Y_n[n] \to K'$ 都为零。考虑同伦余极限的定义区别三角
$$
\bigoplus Y_n[n] \to
\bigoplus Y_n[n] \to
K \to
(\bigoplus Y_n[n])[1]
$$
。如上论证，可见只需证明
$$
\Hom((\bigoplus Y_n[n])[1], K') = \prod \Hom(Y_n[n + 1], K')
$$
对所有 $n \geq 0$ 均为零。为此，再次如上论证，只需证明
$$
\Hom(K_n[n + 1], K'_n)  = 0
$$
对所有 $n \geq 0$ 成立；这由条件 (3) 得出。
\end{proof}

\begin{lemma}
\label{spaces-simplicial-lemma-hom-cartesian-objects-derived}
在情形 \ref{spaces-simplicial-situation-simplicial-site} 中。
% P10-E0188
设 $K, K' \in D(\mathcal{C}_{total})$。假设
\begin{enumerate}
\item $K$ 是笛卡尔的；
\item 当 $i > 1$ 时，$\Hom(K_i[i - 1], K'_i) = 0$。
\end{enumerate}
则 $K$ 与 $K'$ 所关联的单纯系统之间的任何映射
$\{K_n \to K'_n\}$，均来自 $D(\mathcal{C}_{total})$
中的一个映射 $K \to K'$。
\end{lemma}

\begin{proof}
设 $\{K_n \to K'_n\}_{n \geq 0}$ 是导出范畴单纯系统的一个态射。
考虑引理 \ref{spaces-simplicial-lemma-abelian-postnikov} 中与 $K$ 相关联的对象
$X_n$ 及 Postnikov 系统 $Y_n$。
由 (1) 与引理 \ref{spaces-simplicial-lemma-derived-cartesian-shriek}，有
$X_n = g_{n!}K_n$。特别地，映射 $K_0 \to K'_0$
诱导态射 $X_0 \to K'$。由于 $\{K_n \to K'_n\}$
是系统的态射，一个从略的计算表明复合
$$
X_1 \to X_0 \to K'
$$
为零。由于 $Y_0 = X_0$，且 $Y_1$ 位于区别三角
$$
Y_1 \to X_1 \to Y_0 \to Y_1[1]
$$
中，可知存在态射 $Y_1[1] \to K'$，它与
$X_0 = Y_0 \to Y_1[1]$ 的复合正是上述态射 $X_0 \to K'$。
假设已对 $n \geq 1$ 给定映射 $Y_n[n] \to K'$。
由区别三角
$$
X_{n + 1}[n] \to Y_n[n] \to Y_{n + 1}[n + 1] \to X_{n + 1}[n + 1]
$$
得到正合列
$$
\Hom(Y_{n + 1}[n + 1], K') \to
\Hom(Y_n[n], K') \to
\Hom(X_{n + 1}[n], K')
$$
由于 $X_{n + 1}[n] = g_{n + 1!}K_{n + 1}[n]$，最后一个群等于
$$
\Hom(K_{n + 1}[n], K'_{n + 1})
$$
，它由假设 (2) 为零。归纳得到与转移映射相容的映射系统
$Y_n[n] \to K'$，且它在 $Y_0$ 上约化为给定映射。
这产生映射
$$
\gamma :
K = \text{hocolim} Y_n[n]
\longrightarrow
K'
$$
无论如何，此映射使图
$$
\xymatrix{
X_0 \ar[rd] \ar[r] &
K \ar[d]^\gamma \\
& K'
}
$$
交换。限制到 $\mathcal{C}_0$，可知映射
$\gamma_0 : K_0 \to K'_0$ 与该单纯系统态射的第一个映射
$K_0 \to K'_0$ 相同。由于 $K$ 是笛卡尔的，由此容易看出
$\{\gamma_n\}$ 正是我们最初给定的单纯系统态射。
\end{proof}

\begin{lemma}
\label{spaces-simplicial-lemma-cartesian-object-derived-from-simplicial}
在情形 \ref{spaces-simplicial-situation-simplicial-site} 中，设
$(K_n, K_\varphi)$ 是导出范畴的一个单纯系统。假设
\begin{enumerate}
\item $(K_n, K_\varphi)$ 是笛卡尔的；
\item 当 $i \geq 0$ 且 $t > 0$ 时，
$\Hom(K_i[t], K_i) = 0$。
\end{enumerate}
则存在 $D(\mathcal{C}_{total})$ 的一个笛卡尔对象 $K$，
其所关联的单纯系统同构于 $(K_n, K_\varphi)$。
\end{lemma}

\begin{proof}
在 $D(\mathcal{C}_{total})$ 中令 $X_n = g_{n!}K_n$。
对每个 $n \geq 1$，有
$$
\Hom(X_n, X_{n - 1}) =
\Hom(K_n, g_n^{-1}g_{n - 1!}K_{n - 1}) =
\bigoplus\nolimits_{\varphi : [n - 1] \to [n]}
\Hom(K_n, f_\varphi^{-1}K_{n - 1})
$$
因此，我们得到映射 $X_n \to X_{n - 1}$，它对应于映射
$K_\varphi^{-1} : K_n \to f_\varphi^{-1}K_{n - 1}$
在 $\varphi$ 遍历 $\delta^n_0, \ldots, \delta^n_n$ 时的交错和。
这可以进行，因为由假设 (1)，$K_\varphi$ 可逆。
请注意，这与引理 \ref{spaces-simplicial-lemma-simplicial-resolution-Z-site}
证明中映射的定义相似。由此得到
$D(\mathcal{C}_{total})$ 中的复形
$$
\ldots \to X_2 \to X_1 \to X_0
$$
。我们略去说明各复合为零的计算。由《导出范畴》引理
\ref{derived-lemma-existence-postnikov-system}，若有
$$
\Hom(X_i[i - j - 2], X_j) = 0\text{，当 }i > j + 2
$$
，则可将此复形延拓为一个 Postnikov 系统。该群等于
$$
\Hom(K_i[i - j - 2], g_i^{-1}g_{j!}K_j)
$$
再次利用 $(K_n, K_\varphi)$ 是笛卡尔的，可见
$g_i^{-1}g_{j!}K_j$ 同构于有限个 $K_i$ 副本的直和。
因此，该群由假设 (2) 消失。
令 Postnikov 系统由 $Y_0 = X_0$ 以及对 $n \geq 1$
的区别三角 $Y_n \to X_n \to Y_{n - 1} \to Y_n[1]$ 给出。
令
$$
K = \text{hocolim} Y_n[n]
$$
为完成证明，须证明对所有 $m$，$g_m^{-1}K$ 同构于 $K_m$，
且这些同构与映射 $K_\varphi$ 相容。注意
$$
g_m^{-1} K = \text{hocolim} g_m^{-1}Y_n[n]
$$
，且 $g_m^{-1}Y_n[n]$ 是 $g_m^{-1}X_n$ 的一个 Postnikov 系统。
考虑同构
$$
g_m^{-1}X_n =
\bigoplus\nolimits_{\varphi : [n] \to [m]} f_\varphi^{-1}K_n
\xrightarrow{\bigoplus K_\varphi}
\bigoplus\nolimits_{\varphi : [n] \to [m]} K_m
$$
这些映射定义复形的同构
$$
\xymatrix{
\ldots \ar[r] &
g_m^{-1}X_2 \ar[r] \ar[d] &
g_m^{-1}X_1 \ar[r] \ar[d] &
g_m^{-1}X_0 \ar[d] \\
\ldots \ar[r] &
\bigoplus\nolimits_{\varphi : [2] \to [m]} K_m \ar[r] &
\bigoplus\nolimits_{\varphi : [1] \to [m]} K_m \ar[r] &
\bigoplus\nolimits_{\varphi : [0] \to [m]} K_m
}
$$
；它位于 $D(\mathcal{C}_m)$ 中，而下行的箭头如引理
\ref{spaces-simplicial-lemma-simplicial-resolution-Z-site} 的证明所述。
由于我们对复形 $\ldots \to X_2 \to X_1 \to X_0$
中箭头的选取，各方块交换；计算从略。
下行复形有由下式给出的 Postnikov 塔：
$$
Y'_{m, n} =
\left(\bigoplus\nolimits_{\varphi : [n] \to [m]} \mathbf{Z} \to
\ldots \to
\bigoplus\nolimits_{\varphi : [0] \to [m]} \mathbf{Z}\right)[-n]
\otimes^\mathbf{L}_\mathbf{Z} K_m
$$
，并且 $\text{hocolim} Y'_{m, n} = K_m$
（请与引理 \ref{spaces-simplicial-lemma-abelian-postnikov} 的证明以及《导出范畴》例
\ref{derived-example-key-postnikov} 比较）。
应用《导出范畴》引理
\ref{derived-lemma-existence-postnikov-system} 的第二部分，
只要有
$$
\Hom(g_m^{-1}X_i[i - j - 1], \bigoplus\nolimits_{\varphi : [j] \to [m]} K_m)
= 0\text{，当 }i > j + 1
$$
，大图中的竖直映射便可延拓为 Postnikov 系统的同构。
% P10-E0189
若当 $i > j + 1$ 时
$\Hom(K_m[i - j - 1], K_m) = 0$，上述消失成立；
而这由假设 (2) 得到。在 $D(\mathcal{C}_m)$ 中，选取由
$\gamma_{m, n} : g_m^{-1}Y_n \to Y'_{m, n}$ 给出的
Postnikov 系统同构。由同伦余极限的唯一性，可以找到同构
$$
g_m^{-1} K = \text{hocolim} g_m^{-1}Y_n[n]
\xrightarrow{\gamma_m}
\text{hocolim} Y'_{m, n} = K_m
$$
，它与 $\gamma_{m, n}$ 相容。

\medskip\noindent
还须证明，对每个 $\varphi : [m] \to [n]$，映射 $\gamma_m$
构成交换图
$$
\xymatrix{
f_\varphi^{-1}g_m^{-1}K \ar[d]_{f_\varphi^{-1}\gamma_m} \ar[r]_{K(\varphi)} &
g_n^{-1}K \ar[d]^{\gamma_n} \\
f_\varphi^{-1}K_m \ar[r]^{K_\varphi} &
K_n
}
$$
。考虑图
$$
\xymatrix{
f_\varphi^{-1}(\bigoplus_{\psi : [0] \to [m]} f_\psi^{-1}K_0)
\ar@{=}[r] \ar[d]_{f_\varphi^{-1}(\bigoplus K_\psi)} &
f_\varphi^{-1}g_m^{-1}X_0 \ar[d] \ar[r]_{X_0(\varphi)} &
g_n^{-1}X_0 \ar[d] &
\bigoplus_{\chi : [0] \to [n]} f_\chi^{-1}K_0
\ar@{=}[l] \ar[d]^{\bigoplus K_\chi} \\
f_\varphi^{-1}(\bigoplus_{\psi : [0] \to [m]} K_m) \ar@{=}[d] &
f_\varphi^{-1}g_m^{-1}K \ar[d]_{f_\varphi^{-1}\gamma_m} \ar[r]_{K(\varphi)} &
g_n^{-1}K \ar[d]^{\gamma_n} &
\bigoplus_{\chi : [0] \to [n]} K_n \ar@{=}[d] \\
f_\varphi^{-1}Y'_{0, m} \ar[r] &
f_\varphi^{-1}K_m \ar[r]^{K_\varphi} &
K_n &
Y'_{0, n} \ar[l]
}
$$
由于 $X_0 \to K$ 是单纯对象的态射，上方中间的方块交换。
由于 $\gamma_m$（相应地，$\gamma_n$）与
$\gamma_{0, m}$（相应地，$\gamma_{0, n}$）相容，
左矩形（相应地，右矩形）交换；后两者正是图中的箭头
$\bigoplus K_\psi$ 与 $\bigoplus K_\chi$。
% P10-E0190
由于 $(K_n, K_\varphi)$ 是单纯系统，且映射 $X_0(\varphi)$
由显然的恒等式
$f_\varphi^{-1}f_\psi^{-1}K_0 = f_{\varphi \circ \psi}^{-1}K_0$
给出，故沿该图的外矩形绕行是交换的。
注意，箭头 $\bigoplus_\psi K_m \to Y'_{0, m} \to K_m$
在任意一个直和分量上诱导同构
（这是因为上面对 Postnikov 系统 $Y'_{i, j}$ 的显式构造）。
% P10-E0191
因此，若取左上角的一个直和分量，则由于 $\gamma_m$ 是同构，
该分量同构地映到 $f_\varphi^{-1}g_m^{-1}K$。
展开上述陈述，但只考察这个直和分量，即可推出下方中间的方块也交换。
证明完毕。
\end{proof}









\section{导出范畴的单纯系统：模}
\label{spaces-simplicial-section-glueing-modules}

\noindent
本节将在 $\mathcal{O}$-模导出范畴的框架下证明
\cite[Proposition 3.2.9]{BBD} 的一个特例。
（稍微）容易一些的阿贝尔层情形在第
\ref{spaces-simplicial-section-glueing} 节中讨论。

\begin{definition}
\label{spaces-simplicial-definition-cartesian-derived-modules}
在情形 \ref{spaces-simplicial-situation-simplicial-site} 中，设 $\mathcal{O}$ 是
$\mathcal{C}_{total}$ 上的环层。一个
{\it 模导出范畴的单纯系统}由下列数据组成：
\begin{enumerate}
\item 对每个 $n$，给定 $D(\mathcal{O}_n)$ 的对象 $K_n$；
\item 对每个 $\varphi : [m] \to [n]$，给定
$D(\mathcal{O}_n)$ 中的映射
$K_\varphi : Lf_\varphi^*K_m \to K_n$。
\end{enumerate}
这些数据须满足：
$$
K_{\varphi \circ \psi} = K_\varphi \circ Lf_\varphi^*K_\psi :
Lf_{\varphi \circ \psi}^*K_l = Lf_\varphi^* Lf_\psi^*K_l
\longrightarrow
K_n
$$
对 $\Delta$ 的任意态射 $\varphi : [m] \to [n]$ 与
$\psi : [l] \to [m]$，上述等式均成立。
若对所有 $\varphi$，映射 $K_\varphi$ 都是同构，
则称该单纯系统是{\it 笛卡尔的}。
对两个导出范畴的单纯系统，显然可以定义
{\it 模导出范畴的单纯系统的态射}。
\end{definition}

\noindent
我们有意给这个概念取了一个长得荒唐的名字。
目标是证明：在某些假设下，模导出范畴的单纯系统来自
$D(\mathcal{O})$ 的一个对象。

\begin{lemma}
\label{spaces-simplicial-lemma-cartesian-objects-derived-modules}
在情形 \ref{spaces-simplicial-situation-simplicial-site} 中，设 $\mathcal{O}$ 是
$\mathcal{C}_{total}$ 上的环层。
若 $K \in D(\mathcal{O})$，则 $(K_n, K(\varphi))$
是模导出范畴的一个单纯系统。
若 $K$ 是笛卡尔的，则该系统也是笛卡尔的。
\end{lemma}

\begin{proof}
这立即由定义得出。
\end{proof}

\begin{lemma}
\label{spaces-simplicial-lemma-construct-cartesian-objects-derived-modules}
在情形 \ref{spaces-simplicial-situation-simplicial-site} 中，设 $\mathcal{O}$ 是
$\mathcal{C}_{total}$ 上的环层。假设给定
$K_0 \in D(\mathcal{O}_0)$ 以及同构
$$
\alpha :
L(f_{\delta_1^1})^*K_0
\longrightarrow
L(f_{\delta_0^1})^*K_0
$$
，它满足余圈条件。令 $\tau^n_i : [0] \to [n]$, $0 \mapsto i$，
并令 $K_n = Lf_{\tau^n_n}^*K_0$。诸对象 $K_n$
构成模导出范畴的一个笛卡尔单纯系统的成员。
\end{lemma}

\begin{proof}
请与引理 \ref{spaces-simplicial-lemma-construct-cartesian-objects-derived}、
引理 \ref{spaces-simplicial-lemma-characterize-cartesian} 及其证明比较
（其中也明确写出了余圈条件）。
该构造类似于《下降》第 \ref{descent-section-descent-modules} 节
讨论的构造；我们借用其中的记号
$\tau^n_i : [0] \to [n]$, $0 \mapsto i$ 以及
$\tau^n_{ij} : [1] \to [n]$, $0 \mapsto i$, $1 \mapsto j$。
给定 $\varphi : [n] \to [m]$，用下图定义
$K_\varphi : L(f_\varphi)^*K_n \to K_m$：
$$
\xymatrix{
L(f_\varphi)^*K_n \ar@{=}[r] &
L(f_\varphi)^* L(f_{\tau^n_n})^*K_0 \ar@{=}[r] &
L(f_{\tau^m_{\varphi(n)}})^*K_0 \ar@{=}[r] &
L(f_{\tau^m_{\varphi(n)m}})^* L(f_{\delta^1_1})^*K_0
\ar[d]_{L(f_{\tau^m_{\varphi(n)m}})^*\alpha} \\
&
K_m \ar@{=}[r] &
L(f_{\tau^m_m})^*K_0 \ar@{=}[r] &
L(f_{\tau^m_{\varphi(n)m}})^* L(f_{\delta^1_0})^*K_0
}
$$
% P10-E0192
我们略去如下验证：余圈条件蕴含这些映射
（在各自的导出范畴中）能够正确复合，因而给出
模导出范畴的一个单纯系统。
\end{proof}

\begin{lemma}
\label{spaces-simplicial-lemma-modules-postnikov}
在情形 \ref{spaces-simplicial-situation-simplicial-site} 中，设 $\mathcal{O}$ 是
$\mathcal{C}_{total}$ 上的环层。
% P10-E0193
设 $K$ 是 $D(\mathcal{C}_{total})$ 的对象。令
$$
X_n = (g_{n!}\mathcal{O}_n)
\otimes^\mathbf{L}_\mathcal{O} K
\quad\text{且}\quad
Y_n =
(g_{n!}\mathcal{O}_n \to \ldots \to g_{0!}\mathcal{O}_0)[-n]
\otimes^\mathbf{L}_\mathcal{O} K
$$
为 $D(\mathcal{O})$ 的对象，其中的映射如
% P10-E0194
引理 \ref{spaces-simplicial-lemma-simplicial-resolution-Z-site} 所述。
配备显然的典范映射 $Y_n \to X_n$ 以及
$Y_0 \to Y_1[1] \to Y_2[2] \to \ldots$ 后，有：
\begin{enumerate}
\item 区别三角 $Y_n \to X_n \to Y_{n - 1} \to Y_n[1]$
为 $\ldots \to X_2 \to X_1 \to X_0$ 定义一个 Postnikov 系统
（《导出范畴》定义 \ref{derived-definition-postnikov-system}）；
\item 在 $D(\mathcal{O})$ 中，$K = \text{hocolim} Y_n[n]$。
\end{enumerate}
\end{lemma}

\begin{proof}
首先，若 $K = \mathcal{O}$，则这正是将《导出范畴》例
\ref{derived-example-key-postnikov} 的构造应用于
$\textit{Ab}(\mathcal{C}_{total})$ 中的复形
$$
\ldots \to
g_{2!}\mathcal{O}_2 \to
g_{1!}\mathcal{O}_1 \to
g_{0!}\mathcal{O}_0
$$
，再结合如下事实：由引理 \ref{spaces-simplicial-lemma-simplicial-resolution-ringed}，
此复形在 $D(\mathcal{C}_{total})$ 中表示 $K = \mathcal{O}$。
一般情形由此以及下列两个事实得出：正合函子
$- \otimes^\mathbf{L}_\mathcal{O} K$ 将 Postnikov 系统映为
Postnikov 系统；函子
$- \otimes^\mathbf{L}_\mathcal{O} K$ 与同伦余极限交换。
\end{proof}

\begin{lemma}
\label{spaces-simplicial-lemma-nullity-cartesian-modules-derived}
在情形 \ref{spaces-simplicial-situation-simplicial-site} 中，设 $\mathcal{O}$ 是
$\mathcal{C}_{total}$ 上的环层。
% P10-E0195
设 $K, K' \in D(\mathcal{O})$。假设
\begin{enumerate}
% P10-E0196
\item 对 $\varphi : [m] \to [n]$，
$f_\varphi^{-1}\mathcal{O}_n \to \mathcal{O}_m$ 平坦；
\item $K$ 是笛卡尔的；
\item 当 $i > 0$ 时，$\Hom(K_i[i], K'_i) = 0$；
\item 当 $i \geq 0$ 时，$\Hom(K_i[i + 1], K'_i) = 0$。
\end{enumerate}
则任何诱导零映射 $K_0 \to K'_0$ 的映射 $K \to K'$ 都是零映射。
\end{lemma}

\begin{proof}
证明与引理 \ref{spaces-simplicial-lemma-nullity-cartesian-objects-derived} 的证明完全相同，
只需用引理 \ref{spaces-simplicial-lemma-modules-postnikov} 代替
引理 \ref{spaces-simplicial-lemma-abelian-postnikov}。
\end{proof}

\begin{lemma}
\label{spaces-simplicial-lemma-hom-cartesian-modules-derived}
在情形 \ref{spaces-simplicial-situation-simplicial-site} 中，设 $\mathcal{O}$ 是
$\mathcal{C}_{total}$ 上的环层。
% P10-E0197
设 $K, K' \in D(\mathcal{O})$。假设
\begin{enumerate}
% P10-E0196
\item 对 $\varphi : [m] \to [n]$，
$f_\varphi^{-1}\mathcal{O}_n \to \mathcal{O}_m$ 平坦；
\item $K$ 是笛卡尔的；
\item 当 $i > 1$ 时，$\Hom(K_i[i - 1], K'_i) = 0$。
\end{enumerate}
则 $K$ 与 $K'$ 所关联的单纯系统之间的任何映射
$\{K_n \to K'_n\}$，均来自 $D(\mathcal{O})$ 中的一个映射
$K \to K'$。
\end{lemma}

\begin{proof}
证明与引理 \ref{spaces-simplicial-lemma-hom-cartesian-objects-derived} 的证明完全相同，
只需用引理 \ref{spaces-simplicial-lemma-modules-postnikov} 代替
引理 \ref{spaces-simplicial-lemma-abelian-postnikov}。
\end{proof}

\begin{lemma}
\label{spaces-simplicial-lemma-cartesian-module-derived-from-simplicial}
在情形 \ref{spaces-simplicial-situation-simplicial-site} 中，设 $\mathcal{O}$ 是
$\mathcal{C}_{total}$ 上的环层，且
$(K_n, K_\varphi)$ 是模导出范畴的一个单纯系统。假设
\begin{enumerate}
% P10-E0196
\item 对 $\varphi : [m] \to [n]$，
$f_\varphi^{-1}\mathcal{O}_n \to \mathcal{O}_m$ 平坦；
\item $(K_n, K_\varphi)$ 是笛卡尔的；
\item 当 $i \geq 0$ 且 $t > 0$ 时，
$\Hom(K_i[t], K_i) = 0$。
\end{enumerate}
则存在 $D(\mathcal{O})$ 的一个笛卡尔对象 $K$，
其所关联的单纯系统同构于 $(K_n, K_\varphi)$。
\end{lemma}

\begin{proof}
证明与引理 \ref{spaces-simplicial-lemma-cartesian-object-derived-from-simplicial}
的证明完全相同，但作如下改动：
\begin{enumerate}
\item 处处使用 $g_n^* = Lg_n^*$ 代替 $g_n^{-1}$；
\item 处处使用 $f_\varphi^* = Lf_\varphi^*$ 代替 $f_\varphi^{-1}$；
\item 引用引理 \ref{spaces-simplicial-lemma-simplicial-resolution-ringed}，
而非引理 \ref{spaces-simplicial-lemma-simplicial-resolution-Z-site}；
\item 在 $Y'_{m, n}$ 的构造中使用 $\mathcal{O}_m$，
而非 $\mathbf{Z}$；
\item 与引理 \ref{spaces-simplicial-lemma-modules-postnikov} 的证明比较，
而非与引理 \ref{spaces-simplicial-lemma-abelian-postnikov} 的证明比较。
\end{enumerate}
证明完毕。
\end{proof}







\section{与半可表示对象关联的站点}
\label{spaces-simplicial-section-semi-representable}

\noindent
设 $\mathcal{C}$ 是一个站点。回忆，$\mathcal{C}$ 的一个
{\it 半可表示对象}，就是 $\mathcal{C}$ 的对象族
$\{U_i\}_{i \in I}$。一个
{\it 半可表示对象的态射
$\{U_i\}_{i \in I} \to \{V_j\}_{j \in J}$}
由映射 $\alpha : I \to J$ 以及对每个 $i \in I$
给定的 $\mathcal{C}$ 中的态射
$f_i : U_i \to V_{\alpha(i)}$ 组成。
$\mathcal{C}$ 的半可表示对象范畴记为
$\text{SR}(\mathcal{C})$。更多信息参见《超覆盖》定义
\ref{hypercovering-definition-SR} 及其所在一节。

\medskip\noindent
对 $\mathcal{C}$ 的半可表示对象 $K = \{U_i\}_{i \in I}$，令
$$
\mathcal{C}/K = \coprod\nolimits_{i \in I} \mathcal{C}/U_i
$$
为 $\mathcal{C}$ 在诸 $U_i$ 处的局部化之不交并。
此范畴上有自然的站点结构，其覆盖继承自各局部化
$\mathcal{C}/U_i$。站点 $\mathcal{C}/K$ 称为
{\it $\mathcal{C}$ 在 $K$ 处的局部化}。
注意，$\mathcal{C}/K$ 上的层等同于各
$\mathcal{C}/U_i$ 上的层 $\mathcal{F}_i$ 所成的族，即
$$
\Sh(\mathcal{C}/K) = \prod\nolimits_{i \in I} \Sh(\mathcal{C}/U_i)
$$
这一描述有时有助于理解所发生的事情。

\medskip\noindent
设 $\mathcal{C}$ 是一个站点，且
$K = \{U_i\}_{i \in I}$ 是 $\text{SR}(\mathcal{C})$ 的对象。
存在连续且余连续的局部化函子
$j : \mathcal{C}/K \to \mathcal{C}$；它是诸局部化函子的积
% P10-E0198
$j_i : \mathcal{C}/V_i \to \mathcal{C}$.
正如《站点》第 \ref{sites-section-localize} 节中一样，
我们得到函子 $j_!$、$j^{-1}$、$j_*$。
用积分解
$\Sh(\mathcal{C}/K) = \prod\nolimits_{i \in I} \Sh(\mathcal{C}/U_i)$
表示时，有
$$
\begin{matrix}
j_! & : &
(\mathcal{F}_i)_{i \in I} &
\longmapsto &
\coprod j_{i, !}\mathcal{F}_i \\
j^{-1} & : &
\mathcal{G} &
\longmapsto &
(j_i^{-1}\mathcal{G})_{i \in I} \\
j_* & : &
(\mathcal{F}_i)_{i \in I} &
\longmapsto &
\prod j_{i, *}\mathcal{F}_i
\end{matrix}
$$
，读者容易验证这些公式。

\medskip\noindent
设 $f : K \to L$ 是 $\text{SR}(\mathcal{C})$ 的一个态射。
则逐分量应用《站点》引理 \ref{sites-lemma-relocalize}
的构造，得到连续且余连续的函子
$$
v : \mathcal{C}/K \longrightarrow \mathcal{C}/L
$$
。更准确地，假设 $f = (\alpha, f_i)$，其中
$K = \{U_i\}_{i \in I}$、$L = \{V_j\}_{j \in J}$、
$\alpha : I \to J$，且 $f_i : U_i \to V_{\alpha(i)}$。
则函子 $v$ 通过上述引理的构造，将分量
$\mathcal{C}/U_i$ 映入分量 $\mathcal{C}/V_{\alpha(i)}$。
特别地，得到拓扑斯的态射
$$
f : \Sh(\mathcal{C}/K) \to \Sh(\mathcal{C}/L)
$$
。用积分解
$\Sh(\mathcal{C}/K) = \prod\nolimits_{i \in I} \Sh(\mathcal{C}/U_i)$ and
$\Sh(\mathcal{C}/L) = \prod\nolimits_{j \in J} \Sh(\mathcal{C}/V_j)$
表示时，读者可验证
$$
\begin{matrix}
f_! & : &
(\mathcal{F}_i)_{i \in I} &
\longmapsto &
(\coprod\nolimits_{i \in I, \alpha(i) = j} f_{i, !}\mathcal{F}_i)_{j \in J} \\
f^{-1} & : &
(\mathcal{G}_j)_{j \in J} &
\longmapsto &
(f_i^{-1}\mathcal{G}_{\alpha(i)})_{i \in I} \\
f_* & : &
(\mathcal{F}_i)_{i \in I} &
\longmapsto &
(\prod\nolimits_{i \in I, \alpha(i) = j} f_{i, *}\mathcal{F}_i)_{j \in J}
\end{matrix}
$$
，其中 $f_i : \Sh(\mathcal{C}/U_i) \to
\Sh(\mathcal{C}/V_{\alpha(i)})$ 是与
$f_i : U_i \to V_{\alpha(i)}$ 所对应的局部化函子
$\mathcal{C}/U_i \to \mathcal{C}/V_{\alpha(i)}$ 相关联的态射。

\begin{lemma}
\label{spaces-simplicial-lemma-has-P}
设 $\mathcal{C}$ 是一个站点。
\begin{enumerate}
\item 对 $\text{SR}(\mathcal{C})$ 中的 $K$，函子
$j : \mathcal{C}/K \to \mathcal{C}$ 连续、余连续，并具有
《站点》注 \ref{sites-remark-cartesian-cocontinuous} 的性质 P。
\item 对 $\text{SR}(\mathcal{C})$ 中的 $f : K \to L$，
上述函子 $v : \mathcal{C}/K \to \mathcal{C}/L$
连续、余连续，并具有《站点》注
\ref{sites-remark-cartesian-cocontinuous} 的性质 P。
\end{enumerate}
\end{lemma}

\begin{proof}
证明 (2)。采用引理前讨论中的记号。由《站点》第
\ref{sites-section-localize} 节，局部化函子
$\mathcal{C}/U_i \to \mathcal{C}/V_{\alpha(i)}$
连续且余连续；由《站点》注
\ref{sites-remark-localization-cartesian-cocontinuous}，它们满足 $P$。
由此形式地推出 $v$ 连续、余连续并具有 $P$。细节从略。
(1) 的证明也从略。
\end{proof}

\begin{lemma}
\label{spaces-simplicial-lemma-push-pull-localization}
设 $\mathcal{C}$ 是一个站点，且
$K$ 属于 $\text{SR}(\mathcal{C})$。对
$\Sh(\mathcal{C})$ 中的 $\mathcal{F}$，有
$$
j_*j^{-1}\mathcal{F} = \SheafHom(F(K)^\#, \mathcal{F})
$$
，其中 $F$ 如《超覆盖》定义
\ref{hypercovering-definition-SR-F} 所述。
\end{lemma}

\begin{proof}
设 $K = \{U_i\}_{i \in I}$。
利用上面对函子 $j^{-1}$ 与 $j_*$ 的描述，可见
$$
j_*j^{-1}\mathcal{F} =
\prod\nolimits_{i \in I} j_{i, *}(\mathcal{F}|_{\mathcal{C}/U_i}) =
\prod\nolimits_{i \in I} \SheafHom(h_{U_i}^\#, \mathcal{F})
$$
% P10-E0199
第二个等式由《站点》引理 \ref{sites-lemma-hom-sheaf-hU} 得出。
% P10-E0200
由于在 $\textit{PSh}(\mathcal{C}$ 中
$F(K) = \coprod h_{U_i}$，故在 $\Sh(\mathcal{C})$ 中
$F(K)^\# = \coprod h_{U_i}^\#$。又因
$\SheafHom(-, \mathcal{F})$ 将余积变为积
（这立即由《站点》第 \ref{sites-section-glueing-sheaves} 节
的构造得出），结论成立。
\end{proof}

\begin{lemma}
\label{spaces-simplicial-lemma-localize-compare}
设 $\mathcal{C}$ 是一个站点。
\begin{enumerate}
\item 对 $\text{SR}(\mathcal{C})$ 中的 $K$，函子 $j_!$ 给出等价
$\Sh(\mathcal{C}/K) \to \Sh(\mathcal{C})/F(K)^\#$，
其中 $F$ 如《超覆盖》定义
\ref{hypercovering-definition-SR-F} 所述。
\item 经由 (1) 中的等同，函子
$j^{-1} : \Sh(\mathcal{C}) \to \Sh(\mathcal{C}/K)$ 对应于
$\mathcal{F} \mapsto (\mathcal{F} \times F(K)^\# \to F(K)^\#)$.
\item 对 $\text{SR}(\mathcal{C})$ 中的 $f : K \to L$，
经由 (1) 中的等同，函子 $f^{-1}$ 对应于函子
$\Sh(\mathcal{C})/F(L)^\# \to \Sh(\mathcal{C})/F(K)^\#$,
$(\mathcal{G} \to F(L)^\#) \mapsto
(\mathcal{G} \times_{F(L)^\#} F(K)^\# \to F(K)^\#)$.
\end{enumerate}
\end{lemma}

\begin{proof}
注意，若 $K = \{U_i\}_{i \in I}$，则范畴
$\Sh(\mathcal{C}/K)$ 分解为诸范畴
$\Sh(\mathcal{C}/U_i)$ 的积。又注意
$F(K)^\# = \coprod_{i \in I} h_{U_i}^\#$（层范畴中的余积）。
因此，$\Sh(\mathcal{C})/F(K)^\#$ 是诸范畴
$\Sh(\mathcal{C})/h_{U_i}^\#$ 的积。
于是，(1) 与 (2) 由每个 $i$ 对应的陈述得出；参见《站点》引理
\ref{sites-lemma-essential-image-j-shriek} 与
\ref{sites-lemma-compute-j-shriek-restrict}。
类似地，若 $L = \{V_j\}_{j \in J}$，且 $f$ 由
$\alpha : I \to J$ 及 $f_i : U_i \to V_{\alpha(i)}$ 给出，
则对每个再局部化态射
$\mathcal{C}/U_i \to \mathcal{C}/V_{\alpha(i)}$
应用《站点》引理 \ref{sites-lemma-relocalize-explicit}，即得 (3)。
\end{proof}

\begin{lemma}
\label{spaces-simplicial-lemma-localize-injective}
设 $\mathcal{C}$ 是一个站点。对
$K$ 属于 $\text{SR}(\mathcal{C})$，函子 $j^{-1}$ 将内射阿贝尔层
映为内射阿贝尔层。类似地，函子 $j^{-1}$ 将阿贝尔层的
K-内射复形映为阿贝尔层的 K-内射复形。
\end{lemma}

\begin{proof}
第一个陈述是《站点上的上同调》引理
\ref{sites-cohomology-lemma-cohomology-of-open}
向半可表示对象的自然推广。事实上，利用
$\Sh(\mathcal{C}/K)$ 的积分解以及上面对函子 $j^{-1}$ 的描述，
它由该引理得出。第二个陈述是《站点上的上同调》引理
\ref{sites-cohomology-lemma-restrict-K-injective-to-open}
的自然推广，并由拓扑斯的积分解从该引理得出。

\medskip\noindent
另一种证明：由引理 \ref{spaces-simplicial-lemma-localize-compare} 的 (1)，
$j$ 诱导拓扑斯的一个局部化，故通过扩大站点，结论也立即由
《站点上的上同调》引理
\ref{sites-cohomology-lemma-cohomology-of-open} 与
\ref{sites-cohomology-lemma-restrict-K-injective-to-open} 得出；
在内射层的情形下，可与《站点上的上同调》引理
\ref{sites-cohomology-lemma-cohomology-on-sheaf-sets} 的证明比较。
\end{proof}

\begin{remark}[对象之上的变体]
\label{spaces-simplicial-remark-semi-representable-over-object}
设 $\mathcal{C}$ 是一个站点，且 $X \in \Ob(\mathcal{C})$。
{\it $X$ 之上的半可表示对象}范畴
$\text{SR}(\mathcal{C}, X)$ 由公式
$\text{SR}(\mathcal{C}, X) = \text{SR}(\mathcal{C}/X)$ 定义。
参见《超覆盖》定义 \ref{hypercovering-definition-SR}。
因此，可以将上述讨论应用于站点 $\mathcal{C}/X$。
简言之，上述构造给出：
\begin{enumerate}
\item 对 $\text{SR}(\mathcal{C}, X)$ 中的 $K$，给出站点
$\mathcal{C}/K$；
\item 若 $K = \{U_i/X\}$，给出分解
$\Sh(\mathcal{C}/K) = \prod \Sh(\mathcal{C}/U_i)$；
\item 给出局部化函子 $j : \mathcal{C}/K \to \mathcal{C}/X$；
\item 对 $\text{SR}(\mathcal{C}, X)$ 中的 $f : K \to L$，
给出态射 $f : \Sh(\mathcal{C}/K) \to \Sh(\mathcal{C}/L)$。
\end{enumerate}
将 $\mathcal{C}$ 处处替换为 $\mathcal{C}/X$，本节所有结果
在此情形下均成立。
\end{remark}

\begin{remark}[赋环变体]
\label{spaces-simplicial-remark-semi-representable-ringed}
设 $\mathcal{C}$ 是一个站点，且 $\mathcal{O}_\mathcal{C}$
是 $\mathcal{C}$ 上的环层。在此情形下，对 $\mathcal{C}$
的任意半可表示对象 $K$，站点 $\mathcal{C}/K$ 是赋环站点，
其环层为 $\mathcal{O}_K = j^{-1}\mathcal{O}_\mathcal{C}$。
上述构造给出：
\begin{enumerate}
\item 对 $\text{SR}(\mathcal{C})$ 中的 $K$，给出赋环站点
$(\mathcal{C}/K, \mathcal{O}_K)$；
\item 若 $K = \{U_i\}$，给出分解
$\textit{Mod}(\mathcal{O}_K) =
\prod \textit{Mod}(\mathcal{O}_{U_i})$；
\item 给出赋环拓扑斯的局部化态射
$j : (\Sh(\mathcal{C}/K), \mathcal{O}_K) \to
(\Sh(\mathcal{C}), \mathcal{O}_\mathcal{C})$
；
\item 对 $\text{SR}(\mathcal{C})$ 中的 $f : K \to L$，
给出赋环拓扑斯的态射
$f : (\Sh(\mathcal{C}/K), \mathcal{O}_K) \to
(\Sh(\mathcal{C}/L), \mathcal{O}_L)$。
\end{enumerate}
上述许多结果在此框架下成立。例如，函子 $j^*$ 有正合左伴随
$$
j_! : \textit{Mod}(\mathcal{O}_K) \to \textit{Mod}(\mathcal{O}_\mathcal{C}),
$$
；用 (2) 的积分解表示时，它将
$(\mathcal{F}_i)_{i \in I}$ 映到
$\bigoplus j_{i, !}\mathcal{F}_i$。
类似地，给定上述 $f : K \to L$，函子 $f^*$ 有正合左伴随
$f_! : \textit{Mod}(\mathcal{O}_K) \to \textit{Mod}(\mathcal{O}_L)$。
因此，函子 $j^*$ 与 $f^*$ 正合，即 $j$ 与 $f$
是赋环拓扑斯的平坦态射（这也由等式
$\mathcal{O}_K = j^{-1}\mathcal{O}_\mathcal{C}$ 与
$\mathcal{O}_K = f^{-1}\mathcal{O}_L$ 得出）。
\end{remark}

\begin{remark}[对象之上的赋环变体]
\label{spaces-simplicial-remark-semi-representable-ringed-over-object}
设 $\mathcal{C}$ 是一个站点，$\mathcal{O}_\mathcal{C}$
是 $\mathcal{C}$ 上的环层，且 $X \in \Ob(\mathcal{C})$。
% P10-E0201
记 $\mathcal{O}_X = \mathcal{O}_\mathcal{C}|_{\mathcal{C}/U}$。
结合注 \ref{spaces-simplicial-remark-semi-representable-over-object} 与
\ref{spaces-simplicial-remark-semi-representable-ringed} 中的构造，得到：
\begin{enumerate}
\item 对 $\text{SR}(\mathcal{C}, X)$ 中的 $K$，给出赋环站点
$(\mathcal{C}/K, \mathcal{O}_K)$；
\item 若 $K = \{U_i\}$，给出分解
$\textit{Mod}(\mathcal{O}_K) =
\prod \textit{Mod}(\mathcal{O}_{U_i})$；
\item 给出赋环拓扑斯的局部化态射
$j : (\Sh(\mathcal{C}/K), \mathcal{O}_K) \to
(\Sh(\mathcal{C}/X), \mathcal{O}_X)$
；
\item 对 $\text{SR}(\mathcal{C}, X)$ 中的 $f : K \to L$，
给出赋环拓扑斯的态射
$f : (\Sh(\mathcal{C}/K), \mathcal{O}_K) \to
(\Sh(\mathcal{C}/L), \mathcal{O}_L)$。
\end{enumerate}
当然，注 \ref{spaces-simplicial-remark-semi-representable-ringed}
提到的所有结果在此框架下也成立。
\end{remark}





% P10-E0202
\section{与单纯半可表示对象关联的站点}
\label{spaces-simplicial-section-simplicial-semi-representable}

\noindent
设 $\mathcal{C}$ 是一个站点，且 $K$ 是
$\text{SR}(\mathcal{C})$ 的单纯对象。照例，令
$K_n = K([n])$，并以 $K(\varphi) : K_n \to K_m$
表示与 $\varphi : [m] \to [n]$ 相关联的态射。
由第 \ref{spaces-simplicial-section-semi-representable} 节中的构造，
在以站点为对象、以余连续函子为态射的范畴中，
我们得到单纯对象 $n \mapsto \mathcal{C}/K_n$。
换言之，得到第 \ref{spaces-simplicial-section-simplicial-sites} 节情形 B 所述的资料。
由引理 \ref{spaces-simplicial-lemma-has-P}，这些函子满足性质 P。
因此，可应用引理 \ref{spaces-simplicial-lemma-simplicial-cocontinuous-site}，
得到站点 $(\mathcal{C}/K)_{total}$。

\medskip\noindent
站点 $(\mathcal{C}/K)_{total}$ 可显式描述如下。
设 $K_n = \{U_{n, i}\}_{i \in I_n}$。对
$\varphi : [m] \to [n]$，态射
$K(\varphi) : K_n \to K_m$ 由映射
$\alpha(\varphi) : I_n \to I_m$ 以及对每个 $i \in I_n$
给定的态射
$f_{\varphi, i} : U_{n, i} \to U_{m, \alpha(\varphi)(i)}$ 所组成。
于是：
\begin{enumerate}
\item $(\mathcal{C}/K)_{total}$ 的对象对应于某个 $n$
及某个 $i \in I_n$ 下 $\mathcal{C}/U_{n, i}$ 的对象
$(U/U_{n, i})$；
\item $U/U_{n, i}$ 与 $V/U_{m, j}$ 之间的态射是二元组
$(\varphi, f)$，其中 $\varphi : [m] \to [n]$、
$j = \alpha(\varphi)(i)$，且 $f : U \to V$ 是
$\mathcal{C}$ 的态射，使得图
$$
\vcenter{
\xymatrix{
U \ar[r]_f \ar[d] & V \ar[d] \\
U_{n, i} \ar[r]^-{f_{\varphi, i}} &
U_{m, j}
}
}
$$
交换；
\item 对象 $U/U_{n, i}$ 的覆盖按如下方式构造：
从 $\mathcal{C}$ 中的覆盖 $\{f_j : U_j \to U\}$ 出发，
并令 $\{(\text{id}, f_j) : U_j/U_{n, i} \to U/U_{n, i}\}$
为 $(\mathcal{C}/K)_{total}$ 中的覆盖。
\end{enumerate}
我们为单纯站点建立的全部一般理论均适用于
$(\mathcal{C}/K)_{total}$。注意，显然的遗忘函子
$$
j_{total} : (\mathcal{C}/K)_{total} \longrightarrow \mathcal{C}
$$
连续且余连续。事实表明，与之关联的拓扑斯态射来自一个
（显然的）增广。

\begin{lemma}
\label{spaces-simplicial-lemma-augmentation-simplicial-semi-representable}
设 $\mathcal{C}$ 是一个站点，且 $K$ 是
$\text{SR}(\mathcal{C})$ 的单纯对象。局部化函子
$j_0 : \mathcal{C}/K_0 \to \mathcal{C}$ 定义增广
$a_0 : \Sh(\mathcal{C}/K_0) \to \Sh(\mathcal{C})$，
如注 \ref{spaces-simplicial-remark-augmentation-site} 的情形 (B) 所述。
引理 \ref{spaces-simplicial-lemma-augmentation-site} 中相应的拓扑斯态射
$$
a_n : \Sh(\mathcal{C}/K_n) \longrightarrow \Sh(\mathcal{C}),\quad
a : \Sh((\mathcal{C}/K)_{total}) \longrightarrow \Sh(\mathcal{C})
$$
等于与下列连续且余连续的局部化函子相关联的拓扑斯态射：
$j_n : \mathcal{C}/K_n \to \mathcal{C}$ 以及
$j_{total} : (\mathcal{C}/K)_{total} \to \mathcal{C}$。
\end{lemma}

\begin{proof}
逐一展开定义即可立即得出。关于 $a$ 与 $j_{total}$ 的关系，
特别参见引理 \ref{spaces-simplicial-lemma-augmentation-site} 证明中的脚注。
\end{proof}

\begin{lemma}
\label{spaces-simplicial-lemma-comparison}
采用引理 \ref{spaces-simplicial-lemma-augmentation-simplicial-semi-representable}
中的假设与记号，具有下列性质：
\begin{enumerate}
\item 存在函子
$a^{Sh}_! : \Sh((\mathcal{C}/K)_{total}) \to \Sh(\mathcal{C})$
，它左伴随于
$a^{-1} : \Sh(\mathcal{C}) \to \Sh((\mathcal{C}/K)_{total})$；
\item 存在函子
$a_! : \textit{Ab}((\mathcal{C}/K)_{total}) \to \textit{Ab}(\mathcal{C})$
，它左伴随于
$a^{-1} : \textit{Ab}(\mathcal{C}) \to
\textit{Ab}((\mathcal{C}/K)_{total})$；
\item 函子 $a^{-1}$ 将 $\Sh(\mathcal{C})$ 中的 $\mathcal{F}$
映到 $(\mathcal{C}/K)_{total}$ 上的层，该层在次数 $n$
上等于 $a_n^{-1}\mathcal{F}$；
\item 函子 $a_*$ 将
$\textit{Ab}((\mathcal{C}/K)_{total})$ 中的 $\mathcal{G}$
映到两个映射
$j_{0, *}\mathcal{G}_0 \to j_{1, *}\mathcal{G}_1$ 的等化子。
\end{enumerate}
\end{lemma}

\begin{proof}
(3) 与 (4) 对单纯站点的任意增广均成立；参见引理
\ref{spaces-simplicial-lemma-augmentation-site}。
由于 $j_{total}$ 连续且余连续，(1) 与 (2) 成立。
函子 $a^{Sh}_!$ 构造于《站点》引理
\ref{sites-lemma-when-shriek}，函子 $a_!$ 构造于
《站点上的模》引理
\ref{sites-modules-lemma-g-shriek-adjoint}。
\end{proof}

\begin{lemma}
\label{spaces-simplicial-lemma-sanity-check-simplicial-semi-representable}
设 $\mathcal{C}$ 是一个站点，$K$ 是
$\text{SR}(\mathcal{C})$ 的单纯对象，$U/U_{n, i}$ 是
$\mathcal{C}/K_n$ 的对象，且
$\mathcal{F} \in \textit{Ab}((\mathcal{C}/K)_{total})$。则
$$
H^p(U, \mathcal{F}) = H^p(U, \mathcal{F}_{n, i})
$$
，其中
\begin{enumerate}
% P10-E0203
\item 左端的 $U$ 被视为 $\mathcal{C}_{total}$ 的对象；
\item 右端的 $\mathcal{F}_{n, i}$ 是层 $\mathcal{F}_n$
在第 \ref{spaces-simplicial-section-semi-representable} 节的分解
$\Sh(\mathcal{C}/K_n) = \prod \Sh(\mathcal{C}/U_{n, i})$
中的第 $i$ 个分量；这里该层是
$\mathcal{C}/K_n$ 上的层。
\end{enumerate}
\end{lemma}

\begin{proof}
这立即由引理 \ref{spaces-simplicial-lemma-sanity-check} 以及第
\ref{spaces-simplicial-section-semi-representable} 节的积分解得出。
\end{proof}

\begin{remark}[对象之上的变体]
\label{spaces-simplicial-remark-augmentation-over-object}
设 $\mathcal{C}$ 是一个站点，且 $X \in \Ob(\mathcal{C})$。
回忆，我们有 $X$ 之上的半可表示对象范畴
$\text{SR}(\mathcal{C}, X) = \text{SR}(\mathcal{C}/X)$
；参见注 \ref{spaces-simplicial-remark-semi-representable-over-object}。
可以将上述讨论应用于站点 $\mathcal{C}/X$。
简言之，上述构造给出：
\begin{enumerate}
\item 对 $\text{SR}(\mathcal{C}, X)$ 的单纯对象 $K$，
给出站点 $(\mathcal{C}/K)_{total}$；
\item 给出局部化函子
$j_{total} : (\mathcal{C}/K)_{total} \to \mathcal{C}/X$；
\item 给出局部化函子 $j_n : \mathcal{C}/K_n \to \mathcal{C}/X$；
\item 给出拓扑斯态射
$a : \Sh((\mathcal{C}/K)_{total}) \to \Sh(\mathcal{C}/X)$；
\item 给出拓扑斯态射
$a_n : \Sh(\mathcal{C}/K_n) \to \Sh(\mathcal{C}/X)$；
\item 给出函子
$a^{Sh}_! : \Sh((\mathcal{C}/K)_{total}) \to \Sh(\mathcal{C}/X)$
，它左伴随于 $a^{-1}$；
\item 给出函子
$a_! : \textit{Ab}((\mathcal{C}/K)_{total}) \to \textit{Ab}(\mathcal{C}/X)$
，它左伴随于 $a^{-1}$。
\end{enumerate}
本节所有结果在此框架下均成立。为证明这一点，
只需将站点 $\mathcal{C}$ 处处替换为 $\mathcal{C}/X$。
\end{remark}

\begin{remark}[赋环变体]
\label{spaces-simplicial-remark-augmentation-ringed}
设 $\mathcal{C}$ 是一个站点，且 $\mathcal{O}_\mathcal{C}$ 是环层。
给定 $\mathcal{C}$ 的单纯半可表示对象 $K$，令
$\mathcal{O} = a^{-1}\mathcal{O}_\mathcal{C}$，其中 $a$ 如引理
\ref{spaces-simplicial-lemma-augmentation-simplicial-semi-representable} 与
\ref{spaces-simplicial-lemma-comparison} 所述。
如注 \ref{spaces-simplicial-remark-semi-representable-ringed} 中一样追踪环层，
上述构造给出：
\begin{enumerate}
\item 对 $\text{SR}(\mathcal{C})$ 的单纯对象 $K$，
给出赋环站点 $((\mathcal{C}/K)_{total}, \mathcal{O})$；
\item 给出赋环拓扑斯的态射
$a : (\Sh((\mathcal{C}/K)_{total}), \mathcal{O}) \to
(\Sh(\mathcal{C}), \mathcal{O}_\mathcal{C})$；
\item 给出赋环拓扑斯的态射
$a_n : (\Sh(\mathcal{C}/K_n), \mathcal{O}_n) \to
(\Sh(\mathcal{C}), \mathcal{O}_\mathcal{C})$；
\item 给出函子
$a_! : \textit{Mod}(\mathcal{O}) \to \textit{Mod}(\mathcal{O}_\mathcal{C})$
，它左伴随于 $a^*$。
\end{enumerate}
函子 $a_!$ 存在（但一般并不正合），因为
$a^{-1}\mathcal{O}_\mathcal{C} = \mathcal{O}$，且可以把
《站点上的模》引理 \ref{sites-modules-lemma-g-shriek-adjoint}
在引理 \ref{spaces-simplicial-lemma-comparison} 证明中的使用，替换为
《站点上的模》引理
\ref{sites-modules-lemma-lower-shriek-modules}。
如注 \ref{spaces-simplicial-remark-semi-representable-ringed} 所讨论的，
存在正合函子
$a_{n!} : \textit{Mod}(\mathcal{O}_n) \to
\textit{Mod}(\mathcal{O}_\mathcal{C})$
，它左伴随于 $a_n^*$。从而，态射 $a$ 与 $a_n$ 平坦。
注 \ref{spaces-simplicial-remark-semi-representable-ringed} 蕴含：对
$\varphi : [m] \to [n]$，赋环拓扑斯的态射
$f_\varphi : (\Sh(\mathcal{C}/K_n), \mathcal{O}_n) \to
(\Sh(\mathcal{C}/K_m), \mathcal{O}_m)$
平坦，且存在正合函子
$f_{\varphi !} : \textit{Mod}(\mathcal{O}_n) \to \textit{Mod}(\mathcal{O}_m)$
，它左伴随于 $f_\varphi^*$。这又蕴含：对赋环拓扑斯的平坦态射
$g_n : (\Sh(\mathcal{C}/K_n), \mathcal{O}_n) \to
(\Sh((\mathcal{C}/K)_{total}), \mathcal{O})$
，左伴随于 $g_n^*$ 的函子
$g_{n!} : \textit{Mod}(\mathcal{O}_n) \to
\textit{Mod}(\mathcal{O})$ 正合；参见引理
\ref{spaces-simplicial-lemma-exactness-g-shriek-modules}。
\end{remark}

\begin{remark}[对象之上的赋环变体]
\label{spaces-simplicial-remark-augmentation-ringed-over-object}
设 $\mathcal{C}$ 是一个站点，$\mathcal{O}_\mathcal{C}$ 是环层，
且 $X \in \Ob(\mathcal{C})$。记
$\mathcal{O}_X = \mathcal{O}_\mathcal{C}|_{\mathcal{C}/X}$。
结合注 \ref{spaces-simplicial-remark-augmentation-over-object} 与
\ref{spaces-simplicial-remark-augmentation-ringed} 中的构造，得到：
\begin{enumerate}
\item 对 $\text{SR}(\mathcal{C}, X)$ 的单纯对象 $K$，
给出赋环站点 $((\mathcal{C}/K)_{total}, \mathcal{O})$；
\item 给出赋环拓扑斯的态射
$a : (\Sh((\mathcal{C}/K)_{total}), \mathcal{O}) \to
(\Sh(\mathcal{C}/X), \mathcal{O}_X)$；
\item 给出赋环拓扑斯的态射
$a_n : (\Sh(\mathcal{C}/K_n), \mathcal{O}_n) \to
(\Sh(\mathcal{C}/X), \mathcal{O}_X)$；
\item 给出函子
$a_! : \textit{Mod}(\mathcal{O}) \to \textit{Mod}(\mathcal{O}_X)$
，它左伴随于 $a^*$。
\end{enumerate}
当然，注 \ref{spaces-simplicial-remark-augmentation-ringed}
提到的所有结果在此框架下也成立。
\end{remark}






\section{超覆盖的上同调下降}
\label{spaces-simplicial-section-cohomological-descent-hypercoverings}

\noindent
设 $\mathcal{C}$ 是一个站点。本节假设 $\mathcal{C}$
有等化子与纤维积。设 $K$ 是《超覆盖》定义
\ref{hypercovering-definition-hypercovering-variant}
所定义的超覆盖。我们将研究第
\ref{spaces-simplicial-section-simplicial-semi-representable} 节中的增广
$$
a : \Sh((\mathcal{C}/K)_{total}) \longrightarrow \Sh(\mathcal{C})
$$
。

\begin{lemma}
\label{spaces-simplicial-lemma-hypercovering-descent-sheaves}
设 $\mathcal{C}$ 是具有等化子与纤维积的站点，且 $K$ 是超覆盖。则
\begin{enumerate}
\item $a^{-1} : \Sh(\mathcal{C}) \to \Sh((\mathcal{C}/K)_{total})$
满忠实，其本质像为笛卡尔集合层；
\item $a^{-1} : \textit{Ab}(\mathcal{C}) \to
\textit{Ab}((\mathcal{C}/K)_{total})$
满忠实，其本质像为笛卡尔阿贝尔群层。
\end{enumerate}
在两种情形下，$a_*$ 都给出拟逆函子。
\end{lemma}

\begin{proof}
由于函子 $a^{-1}$ 与积交换，阿贝尔层的情形立即由集合层的情形得出。
在其余证明中，我们只处理集合层。
由引理 \ref{spaces-simplicial-lemma-augmentation-cartesian-module}，对
$\Sh(\mathcal{C})$ 中的 $\mathcal{F}$，$a^{-1}\mathcal{F}$ 是笛卡尔的。
只需证明：伴随映射 $\mathcal{F} \to a_*a^{-1}\mathcal{F}$
% P10-E0204
对 $\Sh(\mathcal{C})$ 中的 $\mathcal{F}$ 是同构；并且对
$(\mathcal{C}/K)_{total}$ 上的笛卡尔层 $\mathcal{G}$，
伴随映射 $a^{-1}a_*\mathcal{G} \to \mathcal{G}$ 是同构。

\medskip\noindent
设 $\mathcal{F}$ 是 $\mathcal{C}$ 上的层。
回忆，$a_*a^{-1}\mathcal{F}$ 是两个映射
$a_{0, *}a_0^{-1}\mathcal{F} \to a_{1, *}a_1^{-1}\mathcal{F}$
的等化子；参见引理 \ref{spaces-simplicial-lemma-comparison}。
由引理 \ref{spaces-simplicial-lemma-push-pull-localization}，
$$
a_{0, *}a_0^{-1}\mathcal{F} = \SheafHom(F(K_0)^\#, \mathcal{F})
\quad\text{且}\quad
a_{1, *}a_1^{-1}\mathcal{F} = \SheafHom(F(K_1)^\#, \mathcal{F})
$$
另一方面，我们知道
$$
\xymatrix{
F(K_1)^\# \ar@<1ex>[r] \ar@<-1ex>[r] &
F(K_0)^\# \ar[r] & \Sh(\mathcal{C})\text{ 的终对象 }*
}
$$
按超覆盖的定义，这是集合层范畴中的余等化子图。
因此，只需证明 $\SheafHom(-, \mathcal{F})$ 将余等化子变为等化子；
这立即由《站点》第 \ref{sites-section-glueing-sheaves} 节
的构造得出。

\medskip\noindent
设 $\mathcal{G}$ 是 $(\mathcal{C}/K)_{total}$ 上的笛卡尔层。
我们将证明，存在 $\mathcal{C}$ 上的层 $\mathcal{F}$，使得
$\mathcal{G} = a^{-1}\mathcal{F}$。这将完成证明，因为由上一段的结果，
此时 $a^{-1}a_*\mathcal{G} = a^{-1}a_*a^{-1}\mathcal{F} =
a^{-1}\mathcal{F} = \mathcal{G}$。
对 $n \geq 0$，置 $\mathcal{K}_n = F(K_n)^\#$。于是有层的映射
$$
\xymatrix{
\mathcal{K}_2
\ar@<1ex>[r]
\ar@<0ex>[r]
\ar@<-1ex>[r]
&
\mathcal{K}_1
\ar@<0.5ex>[r]
\ar@<-0.5ex>[r]
&
\mathcal{K}_0
}
$$
它们来自 $K$ 是单纯半可表示对象这一事实。
$K$ 是超覆盖这一条件意味着
$$
\mathcal{K}_1 \to \mathcal{K}_0 \times \mathcal{K}_0
\quad\text{且}\quad
\mathcal{K}_2 \to
\left(\text{cosk}_1(
\xymatrix{
\mathcal{K}_1
\ar@<0.5ex>[r]
\ar@<-0.5ex>[r]
&
\mathcal{K}_0 \ar[l]
})\right)_2
$$
是层的满射。利用引理 \ref{spaces-simplicial-lemma-characterize-cartesian}
对 $(\mathcal{C}/K)_{total}$ 上笛卡尔层的描述，以及引理
\ref{spaces-simplicial-lemma-localize-compare} 对 $\Sh(\mathcal{C}/K_n)$ 的描述，
可见我们的问题完全可以用如下两项来表述\footnote{虽然具体表述并不重要，
我们仍将其写出：问题是证明，给定
$\Sh(\mathcal{C})/\mathcal{K}_0$ 的对象
$\mathcal{G}_0/\mathcal{K}_0$ 以及同构
$$
\alpha :
\mathcal{G}_0 \times_{\mathcal{K}_0, \mathcal{K}(\delta^1_1)} \mathcal{K}_1 \to
\mathcal{G}_0 \times_{\mathcal{K}_0, \mathcal{K}(\delta^1_0)} \mathcal{K}_1
$$
它位于 $\mathcal{K}_1$ 上方并在
$\Sh(\mathcal{C})/\mathcal{K}_2$ 中满足余圈条件；则存在
$\Sh(\mathcal{C})$ 中的 $\mathcal{F}$ 和 $\mathcal{K}_0$ 上方与
$\alpha$ 相容的同构
$\mathcal{F} \times \mathcal{K}_0 \to \mathcal{G}_0$。}：
\begin{enumerate}
\item 拓扑斯 $\Sh(\mathcal{C})$；
\item $\Sh(\mathcal{C})$ 中以 $\mathcal{K}_n$ 为各项的单纯对象
$\mathcal{K}$。
\end{enumerate}
因此，按《站点》引理 \ref{sites-lemma-topos-good-site}，以另一个站点
$\mathcal{C}'$ 替换 $\mathcal{C}$ 后，可以假设 $\mathcal{C}$
具有所有有限极限，$\mathcal{C}$ 上的拓扑是亚典范的；并且
$\mathcal{C}$ 的一族态射 $\{V_j \to V\}$ 是覆盖，当且仅当
$\coprod h_{V_j} \to V$ 是满射；此外，$\mathcal{C}$ 中存在单纯对象
$U$，使得作为单纯层有 $\mathcal{K}_n = h_{U_n}$。
沿上述等价反向追溯，可以假设对所有 $n$ 都有 $K_n = \{U_n\}$。

\medskip\noindent
设 $X$ 是 $\mathcal{C}$ 的终对象。于是
$\{U_0 \to X\}$ 是覆盖，$\{U_1 \to U_0 \times U_0\}$ 是覆盖，
且 $\{U_2 \to (\text{cosk}_1 \text{sk}_1 U)_2\}$ 是覆盖。
沿用《单纯对象》定义 \ref{simplicial-definition-face-degeneracy}
的记号，以 $d^n_i : U_n \to U_{n - 1}$ 和
$s^n_j : U_n \to U_{n + 1}$ 分别表示对应于
$\delta^n_i$ 和 $\sigma^n_j$ 的态射。
稍微滥用记号，对 $\mathcal{C}$ 的态射 $c : V \to W$，
仍以同一字母表示拓扑斯态射
$c : \Sh(\mathcal{C}/V) \to \Sh(\mathcal{C}/W)$。
现在，$\mathcal{G}$ 由 $\mathcal{C}/U_0$ 上的层 $\mathcal{G}_0$
和同构
$\alpha : (d^1_1)^{-1}\mathcal{G}_0 \to (d^1_0)^{-1}\mathcal{G}_0$
给出；该同构满足引理 \ref{spaces-simplicial-lemma-characterize-cartesian}
在 $\mathcal{C}/U_2$ 上表述的余圈条件。
由于 $\{U_2 \to (\text{cosk}_1 \text{sk}_1 U)_2\}$ 是覆盖，
层上的相应拉回函子是忠实的（略去一个小细节）。
因此，可以用 $\text{cosk}_1 \text{sk}_1 U$ 替换 $U$，因为这会以
$(\text{cosk}_1 \text{sk}_1 U)_2$ 替换 $U_2$，而保持
$U_1$ 和 $U_0$ 不变。于是
$$
(d^2_0, d^2_1, d^2_2) : U_2 \to U_1 \times U_1 \times U_1
$$
% P10-E0205
是单态射，其在 $T$-值点上的像由《单纯对象》引理
\ref{simplicial-lemma-work-out} 描述。特别地，存在态射 $c$
嵌入如下交换图：
$$
\xymatrix{
U_1 \times_{(d^1_1, d^1_0), U_0 \times U_0, (d^1_1, d^1_0)} U_1
\ar[d] \ar[rr]_c & & U_2 \ar[d] \\
U_1 \times U_1
\ar[rr]^{(\text{pr}_1, \text{pr}_2, s^0_0 \circ d^1_1 \circ \text{pr}_1)} & &
U_1 \times U_1 \times U_1
}
$$
因为沿图的另一方向复合会给出 $U_2$ 的一个点。
% P10-E0206
将 $U_2$ 上 $\alpha$ 的余圈条件拉回，得到如下条件：
经两个投影拉回到
$U_1 \times_{(d^1_1, d^1_0), U_0 \times U_0, (d^1_1, d^1_0)} U_1$
的 $\alpha$ 的拉回，与经
$s^0_0 \circ d^1_1 \circ \text{pr}_1$ 拉回的 $\alpha$ 相同，
而后者为恒同映射（即 $\alpha$ 经 $s^0_0$ 的拉回是恒同映射）。
由《站点》引理 \ref{sites-lemma-glue-maps}，这意味着
$\alpha$ 来自 $\mathcal{C}/U_0 \times U_0$ 上层的同构
$$
\alpha' : \text{pr}_1^{-1}\mathcal{G}_0 \to \text{pr}_2^{-1}\mathcal{G}_0
$$
。
最后，利用 $U_1 \to U_0 \times U_0$ 的满射性和上面对 $U_2$
的描述，容易看出态射 $U_2 \to U_0 \times U_0 \times U_0$
在伴随层上是满射。因此，$\alpha'$ 在
$U_0 \times U_0 \times U_0$ 上满足余圈条件。
将《站点》引理 \ref{sites-lemma-mapping-property-glue}
应用于覆盖 $\{U_0 \to X\}$，即可完成证明。
\end{proof}

\begin{lemma}
\label{spaces-simplicial-lemma-hypercovering-cech-complex}
设 $\mathcal{C}$ 是具有等化子与纤维积的站点，且 $K$ 是超覆盖。
引理 \ref{spaces-simplicial-lemma-augmentation-cech-complex} 中与
$a^{-1}\mathcal{F}$ 关联的 {\v C}ech 复形
$$
a_{0, *}a_0^{-1}\mathcal{F} \to a_{1, *}a_1^{-1}\mathcal{F} \to
a_{2, *}a_2^{-1}\mathcal{F} \to \ldots
$$
等于复形 $\SheafHom(s(\mathbf{Z}_{F(K)}^\#), \mathcal{F})$。
这里的 $s(\mathbf{Z}_{F(K)}^\#)$ 如《超覆盖》定义
\ref{hypercovering-definition-homology} 所述。
\end{lemma}

\begin{proof}
由引理 \ref{spaces-simplicial-lemma-push-pull-localization}，有
$$
a_{n, *}a_n^{-1}\mathcal{F} = \SheafHom'(F(K_n)^\#, \mathcal{F})
$$
其中 $\SheafHom'$ 如《站点》第 \ref{sites-section-glueing-sheaves} 节所述。
引理 \ref{spaces-simplicial-lemma-augmentation-cech-complex} 的复形中的边界映射
来自单纯结构。因此，复形的相等性来自对
$\Sh(\mathcal{C})$ 中 $
\mathcal{G}$ 的典范恒等式
$\SheafHom'(\mathcal{G}, \mathcal{F}) =
\SheafHom(\mathbf{Z}_\mathcal{G}, \mathcal{F})$。
\end{proof}

\begin{lemma}
\label{spaces-simplicial-lemma-hypercovering-descent-bounded-abelian}
设 $\mathcal{C}$ 是具有等化子与纤维积的站点，且 $K$ 是超覆盖。
对 $E \in D(\mathcal{C})$，映射
$$
E \longrightarrow Ra_*a^{-1}E
$$
是同构。
\end{lemma}

\begin{proof}
首先，设 $\mathcal{I}$ 是 $\mathcal{C}$ 上的内射阿贝尔层。
引理 \ref{spaces-simplicial-lemma-augmentation-spectral-sequence}
对层 $a^{-1}\mathcal{I}$ 给出的谱序列退化，因为由引理
\ref{spaces-simplicial-lemma-localize-injective}，
$(a^{-1}\mathcal{I})_p = a_p^{-1}\mathcal{I}$ 是内射的。因此，复形
$$
a_{0, *}a_0^{-1}\mathcal{I} \to
a_{1, *}a_1^{-1}\mathcal{I} \to
a_{2, *}a_2^{-1}\mathcal{I} \to\ldots
$$
计算 $Ra_*a^{-1}\mathcal{I}$。由引理
\ref{spaces-simplicial-lemma-hypercovering-cech-complex}，它等于复形
$\SheafHom(s(\mathbf{Z}_{F(K)}^\#), \mathcal{I})$。
由于 $K$ 是超覆盖，将《超覆盖》引理
\ref{hypercovering-lemma-acyclic-hypercover-sheaves}
应用于单纯预层 $F(K)$，可知 $s(\mathbf{Z}_{F(K)}^\#)$
在次数 $> 0$ 处正合。由于 $\mathcal{I}$ 是内射的，函子
$\SheafHom(-, \mathcal{I})$ 正合，从而
$\SheafHom(s(\mathbf{Z}_{F(K)}^\#), \mathcal{I})$
在正次数处正合。因此，对 $p > 0$，有
$R^pa_*a^{-1}\mathcal{I} = 0$。另一方面，由引理
\ref{spaces-simplicial-lemma-hypercovering-descent-sheaves}，有
$\mathcal{I} = a_*a^{-1}\mathcal{I}$。

\medskip\noindent
下有界情形。设 $E \in D^+(\mathcal{C})$。
选取表示 $E$ 的下有界内射复形 $\mathcal{I}^\bullet$。
由第一段的结果和 Leray 无环性引理
（《导出范畴》引理 \ref{derived-lemma-leray-acyclicity}），
$Ra_*a^{-1}\mathcal{I}^\bullet$ 由复形
$a_*a^{-1}\mathcal{I}^\bullet = \mathcal{I}^\bullet$
计算。因此，本情形下引理成立。

\medskip\noindent
无界情形。我们建议读者略过这一部分，因为论证与上文相同，
只是为绕开收敛性问题，我们使用双复形的显式表示。
设 $E \in D(\mathcal{C})$。为了证明映射
$E \to Ra_*a^{-1}E$ 是同构，只需证明对
$\mathcal{C}$ 的每个对象 $U$ 都有
$$
R\Gamma(U, E) = R\Gamma(U, Ra_*a^{-1}E)
$$
我们将计算两边，并证明映射 $E \to Ra_*a^{-1}E$ 诱导同构。
选取表示 $E$ 的 K-内射复形 $\mathcal{I}^\bullet$。
再选取拟同构
$a^{-1}\mathcal{I}^\bullet \to \mathcal{J}^\bullet$，
其中 $\mathcal{J}^\bullet$ 是 $(\mathcal{C}/K)_{total}$
上的某个 K-内射复形。我们有
$$
R\Gamma(U, E) = R\Hom(\mathbf{Z}_U^\#, E)
$$
以及
$$
R\Gamma(U, Ra_*a^{-1}E) = R\Hom(\mathbf{Z}_U^\#, Ra_*a^{-1}E) =
R\Hom(a^{-1}\mathbf{Z}_U^\#, a^{-1}E)
$$
由引理 \ref{spaces-simplicial-lemma-simplicial-resolution-augmentation}，有拟同构
$$
\Big(\ldots \to
g_{2!}(a_2^{-1}\mathbf{Z}_U^\#) \to
g_{1!}(a_1^{-1}\mathbf{Z}_U^\#) \to
g_{0!}(a_0^{-1}\mathbf{Z}_U^\#)\Big)
\longrightarrow
a^{-1}\mathbf{Z}_U^\#
$$
因此，$R\Hom(a^{-1}\mathbf{Z}_U^\#, a^{-1}E)$ 等于
$$
R\Gamma((\mathcal{C}/K)_{total},
R\SheafHom(
\ldots \to
g_{2!}(a_2^{-1}\mathbf{Z}_U^\#) \to
g_{1!}(a_1^{-1}\mathbf{Z}_U^\#) \to
g_{0!}(a_0^{-1}\mathbf{Z}_U^\#),
\mathcal{J}^\bullet))
$$
根据《站点上的上同调》第
\ref{sites-cohomology-section-internal-hom} 节中的构造，并利用
$\mathcal{J}^\bullet$ 是 K-内射的，可知它由具有如下各项的
阿贝尔群复形表示：
$$
\prod\nolimits_{p + q = n}
\Hom(g_{p!}(a_p^{-1}\mathbf{Z}_U^\#), \mathcal{J}^q) =
\prod\nolimits_{p + q = n}
\Hom(a_p^{-1}\mathbf{Z}_U^\#, g_p^{-1}\mathcal{J}^q)
$$
更多细节见《站点上的上同调》引理
\ref{sites-cohomology-lemma-RHom-into-K-injective} 和
\ref{sites-cohomology-lemma-section-RHom-over-U}。
因此，$R\Gamma(U, Ra_*a^{-1}E)$ 由乘积全复形
$\text{Tot}_\pi(B^{\bullet, \bullet})$ 计算，其中
$B^{p, q} = \Hom(a_p^{-1}\mathbf{Z}_U^\#, g_p^{-1}\mathcal{J}^q)$。
对另一边作类似论证。首先注意到
$$
s(\mathbf{Z}_{F(K)}^\#) \longrightarrow \mathbf{Z}
$$
由《超覆盖》引理 \ref{hypercovering-lemma-acyclic-hypercover-sheaves}，
这是 $\mathcal{C}$ 上复形的拟同构。由于 $\mathbf{Z}_U^\#$
是平坦的 $\mathbf{Z}$-模层，可知
$$
s(\mathbf{Z}_{F(K)}^\#) \otimes_\mathbf{Z} \mathbf{Z}_U^\#
\longrightarrow
\mathbf{Z}_U^\#
$$
是拟同构。因此，$R\Hom(\mathbf{Z}_U^\#, E)$ 等于
$$
R\Gamma(\mathcal{C}, R\SheafHom(
s(\mathbf{Z}_{F(K)}^\#) \otimes_\mathbf{Z} \mathbf{Z}_U^\#,
\mathcal{I}^\bullet))
$$
根据 $R\SheafHom$ 的构造，并利用 $\mathcal{I}^\bullet$
是 K-内射的，它由具有如下各项的阿贝尔群复形表示：
$$
\prod\nolimits_{p + q = n}
\Hom(\mathbf{Z}^\#_{K_p} \otimes_\mathbf{Z} \mathbf{Z}_U^\#, \mathcal{I}^q)
=
\prod\nolimits_{p + q = n}
\Hom(a_p^{-1}\mathbf{Z}_U^\#, a_p^{-1}\mathcal{I}^q)
$$
各项相等是因为由《站点上的模》评注
\ref{sites-modules-remark-j-shriek-tensor}，有
$\mathbf{Z}^\#_{K_p} \otimes_\mathbf{Z} \mathbf{Z}_U^\# =
a_{p!}a_p^{-1}\mathbf{Z}_U^\#$。
因此，$R\Gamma(U, E)$ 由乘积全复形
$\text{Tot}_\pi(A^{\bullet, \bullet})$ 计算，其中
$A^{p, q} = \Hom(a_p^{-1}\mathbf{Z}_U^\#, a_p^{-1}\mathcal{I}^q)$。

\medskip\noindent
由于 $\mathcal{I}^\bullet$ 是 K-内射的，由引理
\ref{spaces-simplicial-lemma-localize-injective}，$a_p^{-1}\mathcal{I}^\bullet$
是 K-内射的。由于 $\mathcal{J}^\bullet$ 是 K-内射的，由引理
\ref{spaces-simplicial-lemma-restriction-injective-to-component-site}，
$g_p^{-1}\mathcal{J}^\bullet$ 是 K-内射的。二者都表示对象
$a_p^{-1}E$。因此，对每个 $p \geq 0$，由给定映射
$a^{-1}\mathcal{I}^\bullet \to \mathcal{J}^\bullet$
在 $g_p^{-1}$ 下的像诱导的复形映射
$$
A^{p, \bullet} = \Hom(a_p^{-1}\mathbf{Z}_U^\#, a_p^{-1}\mathcal{I}^\bullet)
\longrightarrow
\Hom(a_p^{-1}\mathbf{Z}_U^\#, g_p^{-1}\mathcal{J}^\bullet) = B^{p, \bullet}
$$
% P10-E0207
是拟同构，因为这两个复形都计算
$$
R\Hom(a_p^{-1}\mathbf{Z}_U^\#, a_p^{-1}E)
$$
由《代数进阶》引理 \ref{more-algebra-lemma-prod-qis-gives-qis}，
可知如下交换图中的右侧竖直箭头是拟同构：
$$
\xymatrix{
R\Gamma(U, E) \ar[r] \ar[d] &
\text{Tot}_\pi(A^{\bullet, \bullet}) \ar[d] \\
R\Gamma(U, Ra_*a^{-1}E) \ar[r] &
\text{Tot}_\pi(B^{\bullet, \bullet})
}
$$
由于上面已经看到水平箭头是拟同构，左侧竖直箭头也是拟同构。
\end{proof}

\begin{lemma}
\label{spaces-simplicial-lemma-compare-cohomology-hypercovering}
设 $\mathcal{C}$ 是具有等化子与纤维积的站点，且 $K$ 是超覆盖。
则有典范同构
$$
R\Gamma(\mathcal{C}, E) =
R\Gamma((\mathcal{C}/K)_{total}, a^{-1}E)
$$
其中 $E \in D(\mathcal{C})$。
\end{lemma}

\begin{proof}
这由引理 \ref{spaces-simplicial-lemma-hypercovering-descent-bounded-abelian} 得出，
因为按《站点上的上同调》评注
\ref{sites-cohomology-remark-before-Leray}，有
$R\Gamma((\mathcal{C}/K)_{total}, -) =
R\Gamma(\mathcal{C}, -) \circ Ra_*$。
\end{proof}

\begin{lemma}
\label{spaces-simplicial-lemma-hypercovering-equivalence-bounded}
设 $\mathcal{C}$ 是具有等化子与纤维积的站点，且 $K$ 是超覆盖。
以 $\mathcal{A} \subset \textit{Ab}((\mathcal{C}/K)_{total})$
表示笛卡尔阿贝尔层组成的弱 Serre 子范畴。则函子 $a^{-1}$ 给出等价
$$
D^+(\mathcal{C}) \longrightarrow D_\mathcal{A}^+((\mathcal{C}/K)_{total})
$$
，其拟逆为 $Ra_*$。
\end{lemma}

\begin{proof}
由引理 \ref{spaces-simplicial-lemma-Serre-subcat-cartesian-modules}，
$\mathcal{A}$ 是弱 Serre 子范畴。该等价是迄今所得结果的形式推论。
应用引理 \ref{spaces-simplicial-lemma-hypercovering-descent-sheaves}、
\ref{spaces-simplicial-lemma-hypercovering-descent-bounded-abelian} 以及
《站点上的上同调》引理
\ref{sites-cohomology-lemma-equivalence-bounded}。
\end{proof}

\noindent
我们建议读者略过下面的评注。

\begin{remark}
\label{spaces-simplicial-remark-compare-cohomology-hypercovering-presheaf}
设 $\mathcal{C}$ 是站点，$\mathcal{G}$ 是 $\mathcal{C}$
上的集合预层。如果 $\mathcal{C}$ 具有等化子与纤维积，
那么我们已在《超覆盖》定义
\ref{hypercovering-definition-hypercovering-variant}
中定义了 $\mathcal{G}$ 的超覆盖。我们断言，本节的所有结果
在这一框架下都有有效的对应版本。为此，完全仿照《站点》引理
\ref{sites-lemma-localize-topos-site}，定义 $\mathcal{C}$
在 $\mathcal{G}$ 处的局部化 $\mathcal{C}/\mathcal{G}$
（该引理只对层表述；拓扑斯 $\Sh(\mathcal{C}/\mathcal{G})$
等于拓扑斯 $\Sh(\mathcal{C})$ 在层 $\mathcal{G}^\#$ 处的局部化）。
读者容易证明，站点 $\mathcal{C}/\mathcal{G}$ 具有纤维积与等化子，
并且 $\mathcal{C}$ 中 $\mathcal{G}$ 的超覆盖与站点
$\mathcal{C}/\mathcal{G}$ 的超覆盖是同一回事。因此，在上面关于
超覆盖的引理中用 $\mathcal{C}/\mathcal{G}$ 替换站点 $\mathcal{C}$，
便得到关于 $\mathcal{G}$ 的超覆盖的相应结果的证明。
例如，对 $\mathcal{G}$ 的超覆盖 $K$，有
$$
R\Gamma(\mathcal{C}/\mathcal{G}, E) =
R\Gamma((\mathcal{C}/K)_{total}, a^{-1}E)
$$
其中 $E \in D^+(\mathcal{C}/\mathcal{G})$，而
$a : \Sh((\mathcal{C}/K)_{total}) \to \Sh(\mathcal{C}/\mathcal{G})$
是典范增广。这就是引理 \ref{spaces-simplicial-lemma-compare-cohomology-hypercovering}。
把 $R\Gamma(\mathcal{G}, -) : D(\mathcal{C}) \to D(\textit{Ab})$
定义为函子 $H^0(\mathcal{G}, -) = H^0(\mathcal{G}^\#, -)$
的导出函子；该函子见《超覆盖》第
\ref{hypercovering-section-hypercoverings-verdier} 节和
《站点上的上同调》第 \ref{sites-cohomology-section-limp} 节。我们有
$$
R\Gamma(\mathcal{G}, E) = R\Gamma(\mathcal{C}/\mathcal{G}, j^{-1}E)
$$
% P10-E0208
这是《站点上的上同调》引理
\ref{sites-cohomology-lemma-cohomology-of-open}
对局部化函子 $j : \mathcal{C}/\mathcal{G} \to \mathcal{C}$
的类似版本。综合以上各点，得到
$$
R\Gamma(\mathcal{G}, E) =
R\Gamma((\mathcal{C}/K)_{total}, a^{-1}j^{-1}E) =
R\Gamma((\mathcal{C}/K)_{total}, g^{-1}E)
$$
其中 $E \in D^+(\mathcal{C})$，而
$g : \Sh((\mathcal{C}/K)_{total}) \to \Sh(\mathcal{C})$
是 $a$ 与 $j$ 的复合。
\end{remark}







\section{超覆盖的上同调下降：模}
\label{spaces-simplicial-section-cohomological-descent-hypercoverings-modules}

\noindent
设 $\mathcal{C}$ 是站点，$\mathcal{O}_\mathcal{C}$ 是环层。
假设 $\mathcal{C}$ 具有等化子与纤维积，并设 $K$ 是《超覆盖》定义
\ref{hypercovering-definition-hypercovering-variant}
所定义的超覆盖。我们将研究评注 \ref{spaces-simplicial-remark-augmentation-ringed}
中的增广
$$
a :
(\Sh((\mathcal{C}/K)_{total}), \mathcal{O})
\longrightarrow
(\Sh(\mathcal{C}), \mathcal{O}_\mathcal{C})
$$
的上同调下降。

\begin{lemma}
\label{spaces-simplicial-lemma-hypercovering-descent-modules}
设 $\mathcal{C}$ 是具有等化子与纤维积的站点，
$\mathcal{O}_\mathcal{C}$ 是环层，且 $K$ 是超覆盖。
采用上述记号，
$$
a^* : \textit{Mod}(\mathcal{O}_\mathcal{C}) \to \textit{Mod}(\mathcal{O})
$$
满忠实，其本质像为笛卡尔 $\mathcal{O}$-模。
函子 $a_*$ 给出拟逆。
\end{lemma}

\begin{proof}
由于 $a^{-1}\mathcal{O}_\mathcal{C} = \mathcal{O}$，有
$a^* = a^{-1}$。因此，本引理立即由引理
\ref{spaces-simplicial-lemma-hypercovering-descent-sheaves} 得出。
\end{proof}

\begin{lemma}
\label{spaces-simplicial-lemma-hypercovering-descent-bounded-modules}
设 $\mathcal{C}$ 是具有等化子与纤维积的站点，
$\mathcal{O}_\mathcal{C}$ 是环层，且 $K$ 是超覆盖。
对 $E \in D(\mathcal{O}_\mathcal{C})$，映射
$$
E \longrightarrow Ra_*La^*E
$$
是同构。
\end{lemma}

\begin{proof}
由于 $a^{-1}\mathcal{O}_\mathcal{C} = \mathcal{O}$，有
$La^* = a^* = a^{-1}$。此外，模上的 $Ra_*$ 与阿贝尔层上的
% P10-E0209
$Ra_*$ 一致；参见《站点上的上同调》引理
\ref{sites-cohomology-lemma-modules-abelian-unbounded}。
因此，本引理立即由引理
\ref{spaces-simplicial-lemma-hypercovering-descent-bounded-abelian} 得出。
\end{proof}

\begin{lemma}
\label{spaces-simplicial-lemma-compare-cohomology-hypercovering-modules}
设 $\mathcal{C}$ 是具有等化子与纤维积的站点，
$\mathcal{O}_\mathcal{C}$ 是环层，且 $K$ 是超覆盖。
则有典范同构
$$
R\Gamma(\mathcal{C}, E) =
R\Gamma((\mathcal{C}/K)_{total}, La^*E)
$$
其中 $E \in D(\mathcal{O}_\mathcal{C})$。
\end{lemma}

\begin{proof}
这由引理 \ref{spaces-simplicial-lemma-hypercovering-descent-bounded-modules} 得出，
因为按《站点上的上同调》评注
\ref{sites-cohomology-remark-before-Leray}，或按《站点上的上同调》引理
\ref{sites-cohomology-lemma-Leray-unbounded}，有
$R\Gamma((\mathcal{C}/K)_{total}, -) =
R\Gamma(\mathcal{C}, -) \circ Ra_*$。
\end{proof}

\begin{lemma}
\label{spaces-simplicial-lemma-hypercovering-equivalence-bounded-modules}
设 $\mathcal{C}$ 是具有等化子与纤维积的站点，
$\mathcal{O}_\mathcal{C}$ 是环层，且 $K$ 是超覆盖。
以 $\mathcal{A} \subset \textit{Mod}(\mathcal{O})$
表示笛卡尔 $\mathcal{O}$-模组成的弱 Serre 子范畴。
则函子 $La^*$ 给出等价
$$
D^+(\mathcal{O}_\mathcal{C}) \longrightarrow D_\mathcal{A}^+(\mathcal{O})
$$
，其拟逆为 $Ra_*$。
\end{lemma}

\begin{proof}
由引理 \ref{spaces-simplicial-lemma-Serre-subcat-cartesian-modules}，
$\mathcal{A}$ 是弱 Serre 子范畴（评注
\ref{spaces-simplicial-remark-augmentation-ringed} 中的讨论表明所需假设成立）。
该等价是迄今所得结果的形式推论。应用引理
\ref{spaces-simplicial-lemma-hypercovering-descent-modules}、
\ref{spaces-simplicial-lemma-hypercovering-descent-bounded-modules} 以及
《站点上的上同调》引理 \ref{sites-cohomology-lemma-equivalence-bounded}。
\end{proof}






\section{对象的超覆盖之上同调下降}
\label{spaces-simplicial-section-cohomological-descent-hypercoverings-X}

\noindent
本节假设 $\mathcal{C}$ 具有纤维积，并设 $X \in \Ob(\mathcal{C})$。
设 $K$ 是《超覆盖》定义
\ref{hypercovering-definition-hypercovering} 所定义的 $X$ 的超覆盖。
我们将研究评注 \ref{spaces-simplicial-remark-augmentation-over-object} 中的增广
$$
a : \Sh((\mathcal{C}/K)_{total}) \longrightarrow \Sh(\mathcal{C}/X)
$$
。注意，$\mathcal{C}/X$ 是具有等化子与纤维积的站点；并且
$K$ 是站点 $\mathcal{C}/X$ 的超覆盖\footnote{逆命题未必成立。
也就是说，若 $K$ 是
$\text{SR}(\mathcal{C}, X) = \text{SR}(\mathcal{C}/X)$
的单纯对象，并按《超覆盖》定义
\ref{hypercovering-definition-hypercovering-variant}
定义了站点 $\mathcal{C}/X$ 的超覆盖，则 $K$ 未必定义 $X$
的超覆盖。例如，若 $K_0 = \{U_{0, i}\}_{i \in I_0}$，
则后一个条件保证 $\{U_{0, i} \to X\}$ 是 $\mathcal{C}$ 的覆盖，
而前一个条件仅要求
$\coprod h_{U_{0, i}}^\# \to h_X^\#$ 是层的满射。}；
这由《超覆盖》引理 \ref{hypercovering-lemma-hypercovering-F} 得出。
这意味着，第 \ref{spaces-simplicial-section-cohomological-descent-hypercoverings}
节对超覆盖证明的每个结果，在本节的情形下都有直接的对应版本。

\begin{lemma}
\label{spaces-simplicial-lemma-hypercovering-X-descent-sheaves}
设 $\mathcal{C}$ 是具有纤维积的站点，$X \in \Ob(\mathcal{C})$，
且 $K$ 是 $X$ 的超覆盖。则
\begin{enumerate}
\item $a^{-1} : \Sh(\mathcal{C}/X) \to \Sh((\mathcal{C}/K)_{total})$
满忠实，其本质像为笛卡尔集合层；
\item $a^{-1} : \textit{Ab}(\mathcal{C}/X) \to
\textit{Ab}((\mathcal{C}/K)_{total})$
满忠实，其本质像为笛卡尔阿贝尔群层。
\end{enumerate}
在两种情形下，$a_*$ 都给出拟逆函子。
\end{lemma}

\begin{proof}
由评注 \ref{spaces-simplicial-remark-semi-representable-over-object}、
\ref{spaces-simplicial-remark-augmentation-over-object} 以及本节引言中的讨论，
这由引理 \ref{spaces-simplicial-lemma-hypercovering-descent-sheaves} 得出。
\end{proof}

\begin{lemma}
\label{spaces-simplicial-lemma-hypercovering-X-descent-bounded-abelian}
% P10-E0210
设 $\mathcal{C}$ 是具有纤维积的站点，$X \in \Ob(\mathcal{C})$，
且 $K$ 是 $X$ 的超覆盖。对 $E \in D(\mathcal{C}/X)$，映射
$$
E \longrightarrow Ra_*a^{-1}E
$$
是同构。
\end{lemma}

\begin{proof}
由评注 \ref{spaces-simplicial-remark-semi-representable-over-object}、
\ref{spaces-simplicial-remark-augmentation-over-object} 以及本节引言中的讨论，
这由引理 \ref{spaces-simplicial-lemma-hypercovering-descent-bounded-abelian} 得出。
\end{proof}

\begin{lemma}
\label{spaces-simplicial-lemma-compare-cohomology-hypercovering-X}
设 $\mathcal{C}$ 是具有纤维积的站点，$X \in \Ob(\mathcal{C})$，
且 $K$ 是 $X$ 的超覆盖。则有典范同构
$$
R\Gamma(X, E) = R\Gamma((\mathcal{C}/K)_{total}, a^{-1}E)
$$
其中 $E \in D(\mathcal{C}/X)$。
\end{lemma}

\begin{proof}
由评注 \ref{spaces-simplicial-remark-semi-representable-over-object} 和
\ref{spaces-simplicial-remark-augmentation-over-object}，这由引理
\ref{spaces-simplicial-lemma-compare-cohomology-hypercovering} 得出。
\end{proof}

\begin{lemma}
\label{spaces-simplicial-lemma-hypercovering-X-equivalence-bounded}
设 $\mathcal{C}$ 是具有纤维积的站点，$X \in \Ob(\mathcal{C})$，
且 $K$ 是 $X$ 的超覆盖。以
$\mathcal{A} \subset \textit{Ab}((\mathcal{C}/K)_{total})$
表示笛卡尔阿贝尔层组成的弱 Serre 子范畴。
则函子 $a^{-1}$ 给出等价
$$
D^+(\mathcal{C}/X) \longrightarrow D_\mathcal{A}^+((\mathcal{C}/K)_{total})
$$
，其拟逆为 $Ra_*$。
\end{lemma}

\begin{proof}
由评注 \ref{spaces-simplicial-remark-semi-representable-over-object} 和
\ref{spaces-simplicial-remark-augmentation-over-object}，这由引理
\ref{spaces-simplicial-lemma-hypercovering-equivalence-bounded} 得出。
\end{proof}








\section{对象的超覆盖之上同调下降：模}
\label{spaces-simplicial-section-cohomological-descent-hypercoverings-X-modules}

\noindent
本节假设 $\mathcal{C}$ 具有纤维积，并设 $X \in \Ob(\mathcal{C})$。
设 $K$ 是《超覆盖》定义
\ref{hypercovering-definition-hypercovering} 所定义的 $X$ 的超覆盖。
设 $\mathcal{O}_\mathcal{C}$ 是 $\mathcal{C}$ 上的环层，并置
$\mathcal{O}_X = \mathcal{O}_\mathcal{C}|_{\mathcal{C}/X}$。
我们将研究评注 \ref{spaces-simplicial-remark-augmentation-ringed-over-object} 中的增广
$$
a :
(\Sh((\mathcal{C}/K)_{total}), \mathcal{O})
\longrightarrow
(\Sh(\mathcal{C}/X), \mathcal{O}_X)
$$
。注意，$\mathcal{C}/X$ 是具有等化子与纤维积的站点，
而 $K$ 是站点 $\mathcal{C}/X$ 的超覆盖。因此，本节的结果
立即由第 \ref{spaces-simplicial-section-cohomological-descent-hypercoverings-modules}
节中的相应结果得出。

\begin{lemma}
\label{spaces-simplicial-lemma-hypercovering-X-descent-modules}
设 $\mathcal{C}$ 是具有纤维积的站点，$X \in \Ob(\mathcal{C})$，
$\mathcal{O}_\mathcal{C}$ 是环层，且 $K$ 是 $X$ 的超覆盖。
采用上述记号，
$$
a^* : \textit{Mod}(\mathcal{O}_X) \to \textit{Mod}(\mathcal{O})
$$
满忠实，其本质像为笛卡尔 $\mathcal{O}$-模。
函子 $a_*$ 给出拟逆。
\end{lemma}

\begin{proof}
由评注 \ref{spaces-simplicial-remark-semi-representable-ringed-over-object}、
\ref{spaces-simplicial-remark-augmentation-ringed-over-object} 以及本节引言中的讨论，
这由引理 \ref{spaces-simplicial-lemma-hypercovering-descent-modules} 得出。
\end{proof}

\begin{lemma}
\label{spaces-simplicial-lemma-hypercovering-X-descent-bounded-modules}
设 $\mathcal{C}$ 是具有纤维积的站点，$X \in \Ob(\mathcal{C})$，
$\mathcal{O}_\mathcal{C}$ 是环层，且 $K$ 是 $X$ 的超覆盖。
对 $E \in D(\mathcal{O}_X)$，映射
$$
E \longrightarrow Ra_*La^*E
$$
是同构。
\end{lemma}

\begin{proof}
由评注 \ref{spaces-simplicial-remark-semi-representable-ringed-over-object}、
\ref{spaces-simplicial-remark-augmentation-ringed-over-object} 以及本节引言中的讨论，
这由引理 \ref{spaces-simplicial-lemma-hypercovering-descent-bounded-modules} 得出。
\end{proof}

\begin{lemma}
\label{spaces-simplicial-lemma-compare-cohomology-hypercovering-X-modules}
设 $\mathcal{C}$ 是具有纤维积的站点，$X \in \Ob(\mathcal{C})$，
$\mathcal{O}_\mathcal{C}$ 是环层，且 $K$ 是 $X$ 的超覆盖。
则有典范同构
$$
R\Gamma(X, E) = R\Gamma((\mathcal{C}/K)_{total}, La^*E)
$$
% P10-E0211
其中 $E \in D(\mathcal{O}_\mathcal{C})$。
\end{lemma}

\begin{proof}
由评注 \ref{spaces-simplicial-remark-semi-representable-ringed-over-object}、
\ref{spaces-simplicial-remark-augmentation-ringed-over-object} 以及本节引言中的讨论，
这由引理 \ref{spaces-simplicial-lemma-compare-cohomology-hypercovering-modules} 得出。
\end{proof}

\begin{lemma}
\label{spaces-simplicial-lemma-hypercovering-X-equivalence-bounded-modules}
设 $\mathcal{C}$ 是具有纤维积的站点，$X \in \Ob(\mathcal{C})$，
$\mathcal{O}_\mathcal{C}$ 是环层，且 $K$ 是 $X$ 的超覆盖。
以 $\mathcal{A} \subset \textit{Mod}(\mathcal{O})$
表示笛卡尔 $\mathcal{O}$-模组成的弱 Serre 子范畴。
则函子 $La^*$ 给出等价
$$
D^+(\mathcal{O}_X) \longrightarrow D_\mathcal{A}^+(\mathcal{O})
$$
，其拟逆为 $Ra_*$。
\end{lemma}

\begin{proof}
由评注 \ref{spaces-simplicial-remark-semi-representable-ringed-over-object}、
\ref{spaces-simplicial-remark-augmentation-ringed-over-object} 以及本节引言中的讨论，
这由引理 \ref{spaces-simplicial-lemma-hypercovering-equivalence-bounded-modules} 得出。
\end{proof}










\section{站点的单纯对象给出的超覆盖}
\label{spaces-simplicial-section-hypercovering}

\noindent
设 $\mathcal{C}$ 是具有纤维积的站点，且 $X \in \Ob(\mathcal{C})$。
本节将阐明第 \ref{spaces-simplicial-section-cohomological-descent-hypercoverings-X}
节在超覆盖由站点的单纯对象给出时的结果。
设 $U$ 是 $\mathcal{C}$ 的单纯对象。照例，记
$U_n = U([n])$，并以 $f_\varphi : U_n \to U_m$ 表示对应于
$\varphi : [m] \to [n]$ 的态射 $f_\varphi = U(\varphi)$。
假设有增广
$$
a : U \to X
$$
由此得到单纯站点 $(\mathcal{C}/U)_{total}$ 和增广态射
$$
a : \Sh((\mathcal{C}/U)_{total}) \longrightarrow \Sh(\mathcal{C}/X)
$$
具体而言，由 $U$ 得到 $\text{SR}(\mathcal{C}, X)$ 的单纯对象
% P10-E0212
$K$，其次数 $n$ 部分为 $K_n = \{U_n \to X\}$；于是可以应用
评注 \ref{spaces-simplicial-remark-augmentation-over-object} 中的构造。
更明确地说，站点 $(\mathcal{C}/U)_{total}$ 的对象由
% P10-E0213
$V/U_n$ 给出；态射 $(\varphi, f) : V/U_n \to W/U_m$
由 $\Delta$ 中的态射 $\varphi : [m] \to [n]$ 和态射
$f : V \to W$ 给出，并要求下图交换：
$$
\xymatrix{
V \ar[r]_f \ar[d] & W \ar[d] \\
U_n \ar[r]^{f_\varphi} & U_m
}
$$
拓扑斯态射 $a$ 由余连续函子 $V/U_n \mapsto V/X$ 给出。
仅此而已！

\medskip\noindent
本节称增广 $a : U \to X$ 是
{\it $\mathcal{C}$ 中 $X$ 的超覆盖}，如果满足下列条件：
\begin{enumerate}
\item $\{U_0 \to X\}$ 是 $\mathcal{C}$ 的覆盖；
\item $\{U_1 \to U_0 \times_X U_0\}$ 是 $\mathcal{C}$ 的覆盖；
\item $\{U_{n + 1} \to (\text{cosk}_n\text{sk}_n U)_{n + 1}\}$
对 $n \geq 1$ 是 $\mathcal{C}$ 的覆盖。
\end{enumerate}
这等价于上述 $K$ 是 $X$ 的超覆盖这一条件；参见《超覆盖》例
\ref{hypercovering-example-hypercovering-in-C}。

\begin{lemma}
\label{spaces-simplicial-lemma-hypercovering-X-simple-descent-sheaves}
% P10-E0210
设 $\mathcal{C}$ 是具有纤维积的站点，$X \in \Ob(\mathcal{C})$，
且 $a : U \to X$ 是上述意义下 $\mathcal{C}$ 中 $X$ 的超覆盖。则
\begin{enumerate}
\item $a^{-1} : \Sh(\mathcal{C}/X) \to \Sh((\mathcal{C}/U)_{total})$
满忠实，其本质像为笛卡尔集合层；
\item $a^{-1} : \textit{Ab}(\mathcal{C}/X) \to
\textit{Ab}((\mathcal{C}/U)_{total})$
满忠实，其本质像为笛卡尔阿贝尔群层。
\end{enumerate}
在两种情形下，$a_*$ 都给出拟逆函子。
\end{lemma}

\begin{proof}
这是引理 \ref{spaces-simplicial-lemma-hypercovering-X-descent-sheaves} 的特殊情形。
\end{proof}

\begin{lemma}
\label{spaces-simplicial-lemma-hypercovering-X-simple-descent-bounded-abelian}
% P10-E0210
设 $\mathcal{C}$ 是具有纤维积的站点，$X \in \Ob(\mathcal{C})$，
且 $a : U \to X$ 是上述意义下 $\mathcal{C}$ 中 $X$ 的超覆盖。
对 $E \in D(\mathcal{C}/X)$，映射
$$
E \longrightarrow Ra_*a^{-1}E
$$
是同构。
\end{lemma}

\begin{proof}
这是引理 \ref{spaces-simplicial-lemma-hypercovering-X-descent-bounded-abelian}
的特殊情形。
\end{proof}

\begin{lemma}
\label{spaces-simplicial-lemma-compare-cohomology-hypercovering-X-simple}
设 $\mathcal{C}$ 是具有纤维积的站点，$X \in \Ob(\mathcal{C})$，
且 $a : U \to X$ 是上述意义下 $\mathcal{C}$ 中 $X$ 的超覆盖。
则有典范同构
$$
R\Gamma(X, E) = R\Gamma((\mathcal{C}/U)_{total}, a^{-1}E)
$$
其中 $E \in D(\mathcal{C}/X)$。
\end{lemma}

\begin{proof}
这是引理 \ref{spaces-simplicial-lemma-compare-cohomology-hypercovering-X} 的特殊情形。
\end{proof}

\begin{lemma}
\label{spaces-simplicial-lemma-hypercovering-X-simple-equivalence-bounded}
% P10-E0210
设 $\mathcal{C}$ 是具有纤维积的站点，$X \in \Ob(\mathcal{C})$，
且 $a : U \to X$ 是上述意义下 $\mathcal{C}$ 中 $X$ 的超覆盖。
以 $\mathcal{A} \subset \textit{Ab}((\mathcal{C}/U)_{total})$
表示笛卡尔阿贝尔层组成的弱 Serre 子范畴。
则函子 $a^{-1}$ 给出等价
$$
D^+(\mathcal{C}/X) \longrightarrow D_\mathcal{A}^+((\mathcal{C}/U)_{total})
$$
，其拟逆为 $Ra_*$。
\end{lemma}

\begin{proof}
% P10-E0214
这是引理 \ref{spaces-simplicial-lemma-hypercovering-X-equivalence-bounded} 的特殊情形。
\end{proof}

\begin{lemma}
\label{spaces-simplicial-lemma-sr-when-fibre-products}
设 $U$ 是具有纤维积的站点 $\mathcal{C}$ 的单纯对象。
\begin{enumerate}
\item $\mathcal{C}/U$ 在以站点为对象、以站点态射为态射的范畴中
具有单纯对象的结构；
\item 将引理 \ref{spaces-simplicial-lemma-simplicial-site-site} 的构造应用于 (1)
中的结构，重新得到上述站点 $(\mathcal{C}/U)_{total}$；
\item 若 $a : U \to X$ 是增广，则
$a_0 : \mathcal{C}/U_0 \to \mathcal{C}/X$ 是评注
\ref{spaces-simplicial-remark-augmentation-site} (A) 部分意义下的增广，并给出与上面相同的
拓扑斯态射
$a : \Sh((\mathcal{C}/U)_{total}) \to \Sh(\mathcal{C}/X)$
。
\end{enumerate}
\end{lemma}

\begin{proof}
给定 $\mathcal{C}$ 的对象之间的态射 $V \to W$，局部化态射
$j : \mathcal{C}/V \to \mathcal{C}/W$ 是基变换函子
$\mathcal{C}/W \to \mathcal{C}/V$ 的左伴随。基变换函子连续，
并诱导与 $j$ 相同的拓扑斯态射。参见《站点》引理
\ref{sites-lemma-relocalize-given-fibre-products}。这证明了 (1)。

\medskip\noindent
(2) 成立，因为在引理 \ref{spaces-simplicial-lemma-simplicial-site-site}
所构造的范畴中，态射 $V/U_n \to W/U_m$ 是 $U_n$ 上方的态射
$V \to W \times_{U_m, f_\varphi} U_n$；这等价于位于态射
$f_\varphi : U_n \to U_m$ 上方的态射 $f : V \to W$，
也就是上面定义的范畴 $(\mathcal{C}/U)_{total}$ 中的态射。
覆盖族的集合相等立即由定义得出。

\medskip\noindent
略去 (3) 的证明。
\end{proof}







\section{站点的单纯对象给出的超覆盖：模}
\label{spaces-simplicial-section-hypercovering-modules}

\noindent
设 $\mathcal{C}$ 是具有纤维积的站点，$X \in \Ob(\mathcal{C})$，
而 $\mathcal{O}_\mathcal{C}$ 是 $\mathcal{C}$ 上的环层。
设 $U \to X$ 是第 \ref{spaces-simplicial-section-hypercovering} 节所定义的
$\mathcal{C}$ 中 $X$ 的超覆盖。本节研究如下增广：
$$
a :
(\Sh((\mathcal{C}/U)_{total}), \mathcal{O})
\longrightarrow
(\Sh(\mathcal{C}/X), \mathcal{O}_X)
$$
把 $U$ 视为 $\mathcal{C}/X$ 的单纯半可表示对象，
% P10-E0212
其次数 $n$ 部分是单元素族 $\{U_n/X\}$，再应用评注
\ref{spaces-simplicial-remark-augmentation-ringed-over-object} 中的构造，便得到该增广。
因此，本节的所有结果立即由第
\ref{spaces-simplicial-section-cohomological-descent-hypercoverings-X-modules}
节中的相应结果得出。

\begin{lemma}
\label{spaces-simplicial-lemma-hypercovering-X-simple-descent-modules}
设 $\mathcal{C}$ 是具有纤维积的站点，$X \in \Ob(\mathcal{C})$，
$\mathcal{O}_\mathcal{C}$ 是环层，且 $U$ 是 $\mathcal{C}$ 中
$X$ 的超覆盖。采用上述记号，
$$
a^* : \textit{Mod}(\mathcal{O}_X) \to \textit{Mod}(\mathcal{O})
$$
满忠实，其本质像为笛卡尔 $\mathcal{O}$-模。
函子 $a_*$ 给出拟逆。
\end{lemma}

\begin{proof}
这是引理 \ref{spaces-simplicial-lemma-hypercovering-X-descent-modules} 的特殊情形。
\end{proof}

\begin{lemma}
\label{spaces-simplicial-lemma-hypercovering-X-simple-descent-bounded-modules}
设 $\mathcal{C}$ 是具有纤维积的站点，$X \in \Ob(\mathcal{C})$，
$\mathcal{O}_\mathcal{C}$ 是环层，且 $U$ 是 $\mathcal{C}$ 中
$X$ 的超覆盖。对 $E \in D(\mathcal{O}_X)$，映射
$$
E \longrightarrow Ra_*La^*E
$$
是同构。
\end{lemma}

\begin{proof}
这是引理 \ref{spaces-simplicial-lemma-hypercovering-X-descent-bounded-modules}
的特殊情形。
\end{proof}

\begin{lemma}
\label{spaces-simplicial-lemma-compare-cohomology-hypercovering-X-simple-modules}
设 $\mathcal{C}$ 是具有纤维积的站点，$X \in \Ob(\mathcal{C})$，
$\mathcal{O}_\mathcal{C}$ 是环层，且 $U$ 是 $\mathcal{C}$ 中
$X$ 的超覆盖。则有典范同构
$$
R\Gamma(X, E) = R\Gamma((\mathcal{C}/U)_{total}, La^*E)
$$
% P10-E0211
其中 $E \in D(\mathcal{O}_\mathcal{C})$。
\end{lemma}

\begin{proof}
这是引理 \ref{spaces-simplicial-lemma-compare-cohomology-hypercovering-X-modules}
的特殊情形。
\end{proof}

\begin{lemma}
\label{spaces-simplicial-lemma-hypercovering-X-simple-equivalence-bounded-modules}
设 $\mathcal{C}$ 是具有纤维积的站点，$X \in \Ob(\mathcal{C})$，
$\mathcal{O}_\mathcal{C}$ 是环层，且 $U$ 是 $\mathcal{C}$ 中
$X$ 的超覆盖。以 $\mathcal{A} \subset \textit{Mod}(\mathcal{O})$
表示笛卡尔 $\mathcal{O}$-模组成的弱 Serre 子范畴。
则函子 $La^*$ 给出等价
$$
D^+(\mathcal{O}_X) \longrightarrow D_\mathcal{A}^+(\mathcal{O})
$$
，其拟逆为 $Ra_*$。
\end{lemma}

\begin{proof}
这是引理 \ref{spaces-simplicial-lemma-hypercovering-X-equivalence-bounded-modules}
的特殊情形。
\end{proof}







\section{超覆盖的无界上同调下降}
\label{spaces-simplicial-section-unbounded-cohomological-descent}

\noindent
本节讨论无界上同调下降。结果本身将立即由前几节关于有界上同调下降
的结果，以及《站点上的上同调》引理
\ref{sites-cohomology-lemma-equivalence-unbounded-one} 和（或）
\ref{sites-cohomology-lemma-equivalence-unbounded-two} 得出；
真正的工作在于建立记号并选择适当的假设。
我们的讨论受 \cite{six-I} 中讨论的启发，不过细节颇为不同。

\medskip\noindent
设 $(\mathcal{C}, \mathcal{O}_\mathcal{C})$ 是赋环站点。
假设对 $\mathcal{C}$ 的每个对象 $U$，给定弱 Serre 子范畴
$\mathcal{A}_U \subset \textit{Mod}(\mathcal{O}_U)$，
并满足下列性质：
\begin{enumerate}
\item
\label{spaces-simplicial-item-restriction}
给定 $\mathcal{C}$ 的态射 $U \to V$，限制函子
$\textit{Mod}(\mathcal{O}_V) \to \textit{Mod}(\mathcal{O}_U)$
将 $\mathcal{A}_V$ 映入 $\mathcal{A}_U$；
\item
\label{spaces-simplicial-item-local}
给定 $\mathcal{C}$ 的覆盖 $\{U_i \to U\}_{i \in I}$，
$\textit{Mod}(\mathcal{O}_U)$ 的对象 $\mathcal{F}$ 属于
$\mathcal{A}_U$，当且仅当对所有 $i \in I$，
$\mathcal{F}$ 在 $\mathcal{C}/U_i$ 上的限制属于
$\mathcal{A}_{U_i}$；
\item
\label{spaces-simplicial-item-bounded-dimension}
存在子集 $\mathcal{B} \subset \Ob(\mathcal{C})$，使得
\begin{enumerate}
\item $\mathcal{C}$ 的每个对象都有一个覆盖，其各成员属于
$\mathcal{B}$；
\item 对每个 $V \in \mathcal{B}$，存在整数 $d_V$ 和 $V$
的覆盖组成的余终系 $\text{Cov}_V$，使得
$$
H^p(V_i, \mathcal{F}) = 0 \text{，当 }
\{V_i \to V\} \in \text{Cov}_V,\ p > d_V, \text{且 }
\mathcal{F} \in \Ob(\mathcal{A}_V)
$$
\end{enumerate}
\end{enumerate}
注意，我们要求这一点对 $\mathcal{A}_V$ 中的 $\mathcal{F}$ 成立，
而不仅仅对“整体”对象成立（因此它强于《站点上的上同调》情形
\ref{sites-cohomology-situation-olsson-laszlo} 中施加的条件）。
在这种情形下，存在弱 Serre 子范畴
$\mathcal{A} \subset \textit{Mod}(\mathcal{O}_\mathcal{C})$，
它由如下对象组成：对所有 $U \in \Ob(\mathcal{C})$，
其在 $\mathcal{C}/U$ 上的限制属于 $\mathcal{A}_U$。
此外，还有导出范畴 $D_\mathcal{A}(\mathcal{O}_\mathcal{C})$
和 $D_{\mathcal{A}_U}(\mathcal{O}_U)$，而限制函子在它们之间映射。

\begin{example}
\label{spaces-simplicial-example-quasi-coherent-spaces-etale}
设 $S$ 是概形，$X$ 是 $S$ 上的代数空间。令
$\mathcal{C} = X_{spaces, \etale}$ 为在 $X$ 上平展的代数空间范畴
上的平展站点；参见《空间的性质》定义
\ref{spaces-properties-definition-spaces-etale-site}。
以 $\mathcal{O}_\mathcal{C}$ 表示结构层，即由规则
$U \mapsto \Gamma(U, \mathcal{O}_U)$ 给出的层。
以 $\mathcal{A}_U$ 表示拟凝聚 $\mathcal{O}_U$-模的范畴。
令 $\mathcal{B} = \Ob(\mathcal{C})$；对 $V \in \mathcal{B}$，
置 $d_V = 0$，并以 $\text{Cov}_V$ 表示对所有 $i$，
$V_i$ 均为仿射的覆盖 $\{V_i \to V\}$。
此时假设 (1)、(2)、(3) 均成立。性质 (1) 和 (2) 见
《空间的性质》引理
\ref{spaces-properties-lemma-pullback-quasi-coherent} 和
\ref{spaces-properties-lemma-properties-quasi-coherent}；
(3) 中的消灭性由《概形的上同调》引理
\ref{coherent-lemma-quasi-coherent-affine-cohomology-zero}
以及《空间的上同调》第
\ref{spaces-cohomology-section-higher-direct-image} 节中的讨论得出。
\end{example}

\begin{example}
\label{spaces-simplicial-example-etale}
设 $S$ 是下列类型之一的概形：
\begin{enumerate}
\item 有限域的谱；
\item 可分闭域的谱；
\item 严格 Hensel 的 Noether 局部环的谱；
\item 剩余域有限的 Hensel Noether 局部环的谱；
\item % P10-E0215
在此添加更多情形。
\end{enumerate}
设 $\Lambda$ 是有限环，其阶在 $S$ 上可逆。令
$\mathcal{C} \subset (\Sch/S)_\etale$ 为由 $S$ 上局部有限型概形
组成的全子范畴，并赋予平展拓扑。令
$\mathcal{O}_\mathcal{C} = \underline{\Lambda}$ 为常值层。
置 $\mathcal{A}_U = \textit{Mod}(\mathcal{O}_U)$；换言之，
我们考虑所有 $\Lambda$-模的平展层。令
$\mathcal{B} \subset \Ob(\mathcal{C})$ 为拟紧对象的集合。
对 $V \in \mathcal{B}$，置
$$
d_V = 1 + 2\dim(S) +
\sup\nolimits_{v \in V}(\text{trdeg}_{\kappa(s)}(\kappa(v)) +
2 \dim \mathcal{O}_{V, v})
$$
并以 $\text{Cov}_V$ 表示对所有 $i$，$V_i$ 均拟紧的平展覆盖
$\{V_i \to V\}$。界 $d_V$ 的这一选择来自 Gabber
关于上同调维数的定理。要看出在此选择下条件 (3) 成立，
可使用 \cite[Expos\'e VIII-A, Corollary 1.2 and Lemma 2.2]{Traveaux}
以及关于域的上同调维数的初等论证。公式中加 $1$，
是因为我们的列表包含允许 $S$ 具有有限剩余域的情形。
稍后还会回到这个例子（% P10-E0216
插入将来的引用）。
\end{example}

\noindent
设 $(\mathcal{C}, \mathcal{O}_\mathcal{C})$ 是赋环站点。
假设给定满足条件 (\ref{spaces-simplicial-item-restriction}) 的弱 Serre 子范畴
$\mathcal{A}_U \subset \textit{Mod}(\mathcal{O}_U)$。则
\begin{enumerate}
\item 给定半可表示对象 $K = \{U_i\}_{i \in I}$，取
$\prod \mathcal{A}_{U_i} \subset
\prod \textit{Mod}(\mathcal{O}_{U_i}) = \textit{Mod}(\mathcal{O}_K)$，
便得到弱 Serre 子范畴
$\mathcal{A}_K \subset \textit{Mod}(\mathcal{O}_K)$；
\item 给定半可表示对象之间的态射 $f : K \to L$，拉回映射
% P10-E0217
$f^* : \textit{Mod}(\mathcal{O}_L) \to \textit{Mod}(\mathcal{O}_L)$
将 $\mathcal{A}_L$ 映入 $\mathcal{A}_K$。
\end{enumerate}
记号和说明见评注 \ref{spaces-simplicial-remark-semi-representable-ringed}。
特别地，给定单纯半可表示对象 $K$，评注
\ref{spaces-simplicial-remark-augmentation-ringed} 中 $\textit{Mod}(\mathcal{O})$
的对象 $\mathcal{F}$ 的限制 $\mathcal{F}_n$ 对所有 $n$
均属于 $\mathcal{A}_{K_n}$，其含义是无歧义的。

\begin{lemma}
\label{spaces-simplicial-lemma-hypercovering-equivalence-modules}
设 $(\mathcal{C}, \mathcal{O}_\mathcal{C})$ 是赋环站点。
假设给定满足上述条件 (\ref{spaces-simplicial-item-restriction})、
(\ref{spaces-simplicial-item-local}) 和 (\ref{spaces-simplicial-item-bounded-dimension}) 的弱 Serre 子范畴
$\mathcal{A}_U \subset \textit{Mod}(\mathcal{O}_U)$。
假设 $\mathcal{C}$ 具有等化子与纤维积，并设 $K$ 是超覆盖。
令 $((\mathcal{C}/K)_{total}, \mathcal{O})$ 如评注
\ref{spaces-simplicial-remark-augmentation-ringed} 所述。以
$\mathcal{A}_{total} \subset \textit{Mod}(\mathcal{O})$
表示由如下笛卡尔 $\mathcal{O}$-模 $\mathcal{F}$ 组成的弱 Serre
子范畴：对所有 $n$，其限制 $\mathcal{F}_n$ 属于
$\mathcal{A}_{K_n}$（含义如上所定义）。则函子 $La^*$ 给出等价
$$
D_\mathcal{A}(\mathcal{O}_\mathcal{C})
\longrightarrow
D_{\mathcal{A}_{total}}(\mathcal{O})
$$
，其拟逆为 $Ra_*$。
\end{lemma}

\begin{proof}
由引理 \ref{spaces-simplicial-lemma-Serre-subcat-cartesian-modules}，
笛卡尔 $\mathcal{O}$-模组成弱 Serre 子范畴
（评注 \ref{spaces-simplicial-remark-augmentation-ringed} 中的讨论表明所需假设成立）。
由于限制函子
% P10-E0219
$g_n^* : \textit{Mod}(\mathcal{O}) \to \textit{Mod}(\mathcal{O}_n)$
正合，$\mathcal{A}_{total}$ 也是弱 Serre 子范畴。

\medskip\noindent
我们来证明 $a^* : \mathcal{A} \to \mathcal{A}_{total}$ 是范畴的等价，
% P10-E0218
其逆由 $La_*$ 给出。由有界版本（引理
\ref{spaces-simplicial-lemma-hypercovering-equivalence-bounded-modules}），
我们已经知道 $La_*a^*\mathcal{F} = \mathcal{F}$。
对 $\mathcal{A}$ 中的 $\mathcal{F}$，显然
$a^*\mathcal{F}$ 属于 $\mathcal{A}_{total}$。反之，假设
$\mathcal{G} \in \mathcal{A}_{total}$。因为 $\mathcal{G}$ 是笛卡尔的，
由引理 \ref{spaces-simplicial-lemma-hypercovering-descent-modules}，存在某个
$\mathcal{O}_\mathcal{C}$-模 $\mathcal{F}$，使得
$\mathcal{G} = a^*\mathcal{F}$。我们要证明 $\mathcal{F}$ 属于
$\mathcal{A}$。取 $U \in \Ob(\mathcal{C})$。须证明
$\mathcal{F}$ 在 $\mathcal{C}/U$ 上的限制属于 $\mathcal{A}_U$。
照例，写 $K_0 = \{U_{0, i}\}_{i \in I_0}$。由于 $K$ 是超覆盖，
映射 $\coprod_{i \in I_0} h_{U_{0, i}} \to *$ 在层化后成为满射。
这意味着存在覆盖 $\{U_j \to U\}_{j \in J}$、映射
$\tau : J \to I_0$，并且对每个 $j \in J$，存在态射
$\varphi_j : U_j \to U_{0, \tau(j)}$。
由于 $\mathcal{G}_0 = a_0^*\mathcal{F}$，可知
$\mathcal{F}$ 在 $\mathcal{C}/U_j$ 上的限制，等于
$\mathcal{G}_0$ 的第 $\tau(j)$ 个分量经态射
% P10-E0220
$\varphi_j : U_j \to U_{0, \tau(i)}$ 在 $\mathcal{C}/U_j$
上的限制。因此，由 (\ref{spaces-simplicial-item-restriction})，
$\mathcal{F}|_{\mathcal{C}/U_j}$ 属于 $\mathcal{A}_{U_j}$；
进而由 (\ref{spaces-simplicial-item-local})，$\mathcal{F}|_{\mathcal{C}/U}$
属于 $\mathcal{A}_U$。

\medskip\noindent
特别地，引理的陈述是有意义的。现在，引理由《站点上的上同调》引理
\ref{sites-cohomology-lemma-equivalence-unbounded-one} 得出。
假设 (1) 显然成立（见评注 \ref{spaces-simplicial-remark-augmentation-ringed}）。
假设 (2) 和 (3) 已在上一段证明。假设 (4) 立即由
(\ref{spaces-simplicial-item-bounded-dimension}) 得出。对假设 (5)，令
$\mathcal{B}_{total}$ 为站点 $(\mathcal{C}/K)_{total}$ 中满足
$U \in \mathcal{B}$ 的对象 $U/U_{n, i}$ 的集合，其中
$\mathcal{B}$ 如 (\ref{spaces-simplicial-item-bounded-dimension}) 所述。
这里使用第 \ref{spaces-simplicial-section-simplicial-semi-representable} 节对站点
$(\mathcal{C}/K)_{total}$ 的描述。此外，置
$\text{Cov}_{U/U_{n, i}}$ 等于 $\text{Cov}_U$，并置
$d_{U/U_{n, i}}$ 等于 $d_U$，其中 $\text{Cov}_U$ 与 $d_U$
由 (\ref{spaces-simplicial-item-bounded-dimension}) 给出。我们断言，在这些选择下
条件 (5) 成立。这立即由引理
\ref{spaces-simplicial-lemma-sanity-check-simplicial-semi-representable}
和如下事实得出：$\mathcal{F} \in \mathcal{A}_{total}$ 蕴含
$\mathcal{F}_n \in \mathcal{A}_{K_n}$，从而
$\mathcal{F}_{n, i} \in \mathcal{A}_{U_{n, i}}$。
（担心阿贝尔层上同调与模层上同调之间差别的读者，可参阅
《站点上的上同调》引理
\ref{sites-cohomology-lemma-cohomology-modules-abelian-agree}。）
\end{proof}










\section{粘合复形}
\label{spaces-simplicial-section-glueing-complexes}

\noindent
本节是《上同调》第 \ref{cohomology-section-glueing-complexes} 节的续篇。
目标是证明 \cite[Theorem 3.2.4]{BBD} 的一个轻微推广。
我们的方法略有不同，因为会使用第 \ref{spaces-simplicial-section-glueing} 节和第
\ref{spaces-simplicial-section-glueing-modules} 节的材料。我们还将按 \cite{six-I}
中的方法重新证明无界版本。

\medskip\noindent
给读者的建议：不妨先看引理 \ref{spaces-simplicial-lemma-bbd-glueing-cech}
的陈述及其第二个证明。

\medskip\noindent
下面给出我们所关心的情形。

\begin{situation}
\label{spaces-simplicial-situation-locally-given}
设 $(\mathcal{C}, \mathcal{O}_\mathcal{C})$ 是赋环站点。给定
\begin{enumerate}
\item 范畴 $\mathcal{B}$ 和函子 $u : \mathcal{B} \to \mathcal{C}$；
\item 对 $U \in \Ob(\mathcal{B})$，对象
$E_U$ 属于 $D(\mathcal{O}_{u(U)})$；
\item 对 $\mathcal{B}$ 中的 $a : V \to U$，同构
$\rho_a : E_U|_{\mathcal{C}/u(V)} \to E_V$
属于 $D(\mathcal{O}_{u(V)})$。
\end{enumerate}
并且每当 $\mathcal{B}$ 中有可复合箭头
$b : W \to V$ 和 $a : V \to U$ 时，都有
$\rho_{a \circ b} = \rho_b \circ \rho_a|_{\mathcal{C}/u(W)}$。
\end{situation}

\noindent
若不再加假设，我们无法由此证明任何结论。一个重要情形是：
$\mathcal{B}$ 是全子范畴，并且 $\mathcal{C}$ 的每个对象都有
一个由 $\mathcal{B}$ 的对象组成的覆盖（\cite{BBD} 考虑的正是此情形）。
对我们而言，也必须允许并非如此的情形；主要的另一种情形是：
有站点态射 $f : \mathcal{C} \to \mathcal{D}$，而 $\mathcal{B}$
是 $\mathcal{D}$ 的全子范畴，使得 $\mathcal{D}$ 的每个对象
都有一个由 $\mathcal{B}$ 的对象组成的覆盖。

\medskip\noindent
在情形 \ref{spaces-simplicial-situation-locally-given} 中，{\it 解}是二元组
$(E, \rho_U)$，其中 $E$ 是 $D(\mathcal{O}_\mathcal{C})$ 的对象，
而对 $U \in \Ob(\mathcal{B})$，
$\rho_U : E|_{\mathcal{C}/u(U)} \to E_U$ 是同构；并且对
$\mathcal{B}$ 中的 $a : V \to U$，有
$\rho_a \circ \rho_U|_{\mathcal{C}/u(V)} = \rho_V$。

\begin{lemma}
\label{spaces-simplicial-lemma-prepare-bbd-glueing}
处于情形 \ref{spaces-simplicial-situation-locally-given}。假设 $E_U$ 在
$D(\mathcal{O}_{u(U)})$ 中所有负次数自 Ext 均为零。
设 $L$ 是 $\text{SR}(\mathcal{B})$ 的单纯对象。
考虑 $\text{SR}(\mathcal{C})$ 的单纯对象 $K = u(L)$，
并令 $((\mathcal{C}/K)_{total}, \mathcal{O})$ 如评注
\ref{spaces-simplicial-remark-augmentation-ringed} 所述。存在 $D(\mathcal{O})$
的笛卡尔对象 $E$，使得写 $L_n = \{U_{n, i}\}_{i \in I_n}$ 时，
$E$ 在 $D(\mathcal{O}_{\mathcal{C}/u(U_{n, i})})$ 上的限制
相容地等于 $E_{U_{n, i}}$（细节见证明）。此外，$E$
在唯一同构的意义下唯一。
\end{lemma}

\begin{proof}
回忆
$\Sh(\mathcal{C}/K_n) = \prod_{i \in I_n} \Sh(\mathcal{C}/u(U_{n, i}))$
，模范畴也有类似分解。模层的导出范畴也继承这一乘积分解；
而且该分解与单纯半可表示对象 $K$ 中的态射相容。
参见第 \ref{spaces-simplicial-section-semi-representable} 节。因此，可以在
$D(\mathcal{O}_n)$ 中置
$E_n = \prod_{i \in I_n} E_{U_{n, i}}$（“形式”乘积）。
取情形 \ref{spaces-simplicial-situation-locally-given} 中映射 $\rho_a$ 的形式乘积，
得到同构 $E_\varphi : f_\varphi^*E_n \to E_m$。
% P10-E0221
关于映射 $\rho_a$ 之复合的假设立即蕴含 $(E_n, E_\varphi)$
按定义 \ref{spaces-simplicial-definition-cartesian-derived-modules} 构成模的导出范畴的
单纯系统。引理中假设的负次数 Ext 消灭性蕴含：对
$n \geq 0$ 和 $t > 0$，有 $\Hom(E_n[t], E_n) = 0$。
因此，由引理 \ref{spaces-simplicial-lemma-cartesian-module-derived-from-simplicial}，
得到 $E$。在唯一同构的意义下的唯一性由引理
\ref{spaces-simplicial-lemma-nullity-cartesian-modules-derived} 和
\ref{spaces-simplicial-lemma-hom-cartesian-modules-derived} 得出。
\end{proof}

\begin{lemma}[BBD 粘合引理]
\label{spaces-simplicial-lemma-bbd-glueing}
处于情形 \ref{spaces-simplicial-situation-locally-given}。假设
\begin{enumerate}
\item $\mathcal{C}$ 具有等化子与纤维积；
\item 存在由连续函子 $u : \mathcal{D} \to \mathcal{C}$
给出的站点态射 $f : \mathcal{C} \to \mathcal{D}$，使得
\begin{enumerate}
\item $\mathcal{D}$ 具有等化子与纤维积，且 $u$ 与它们交换；
\item $\mathcal{B}$ 是 $\mathcal{D}$ 的全子范畴，而
$u : \mathcal{B} \to \mathcal{C}$ 是 $u$ 的限制；
\item $\mathcal{D}$ 的每个对象都有一个由 $\mathcal{B}$
的对象组成的覆盖；
\end{enumerate}
\item 对 $\mathcal{B}$ 中的所有 $U$，$E_U$ 在
$D(\mathcal{O}_{u(U)})$ 中所有负次数自 Ext 均为零；
\item 存在 $t \in \mathbf{Z}$，使得当 $i < t$ 且
$U \in \Ob(\mathcal{B})$ 时，$H^i(E_U) = 0$。
\end{enumerate}
则存在一个解，并且它在唯一同构的意义下唯一。
\end{lemma}

\begin{proof}
由《超覆盖》引理 \ref{hypercovering-lemma-hypercovering-site}，
存在站点 $\mathcal{D}$ 的超覆盖 $L$，使得
$L_n = \{U_{n, i}\}_{i \in I_n}$，其中
% P10-E0222
$U_{i, n} \in \Ob(\mathcal{B})$。置 $K = u(L)$。
应用引理 \ref{spaces-simplicial-lemma-prepare-bbd-glueing}，得到站点
$(\mathcal{C}/K)_{total}$ 上 $D(\mathcal{O})$ 的笛卡尔对象 $E$，
它在 $\mathcal{C}/u(U_{n, i})$ 上相容地限制为 $E_{U_{n, i}}$。
关于 $t$ 的假设蕴含 $E \in D^+(\mathcal{O})$。
由《超覆盖》引理
\ref{hypercovering-lemma-hypercovering-morphism-sites}，
$K$ 也是超覆盖。由引理
\ref{spaces-simplicial-lemma-hypercovering-equivalence-bounded-modules}，
存在 $D^+(\mathcal{O}_\mathcal{C})$ 中的 $F$，使得 $E = a^*F$。

\medskip\noindent
为了证明 $F$ 是解，我们将使用《超覆盖》引理
\ref{hypercovering-lemma-hypercovering-site} 的证明中给出的
$L_0$ 和 $L_1$ 的构造。（这略显不雅，但似乎没有完全直接的绕行办法。）

\medskip\noindent
具体地，$I_0 = \Ob(\mathcal{B})$，所以
$L_0 = \{U\}_{U \in \Ob(\mathcal{B})}$。因此，按我们对 $E$
的选择，将同构 $a^*F \to E$ 限制到 $\mathcal{C}/K_0$
的分量 $\mathcal{C}/u(U)$，便对 $U \in \Ob(\mathcal{B})$
定义了同构 $\rho_U : F|_{\mathcal{C}/u(U)} \to E_U$。

\medskip\noindent
% P10-E0224
为了证明 $\rho_U$ 满足与情形 \ref{spaces-simplicial-situation-locally-given}
中的映射 $\rho_a$ 相容这一要求，我们使用 $I_1$ 包含集合
$$
\Omega =
\{(U, V, W, a, b) \mid U, V, W \in \mathcal{B}, a : U \to V, b : U \to W\}
$$
并且对 $\Omega$ 中的 $i = (U, V, W, a, b)$，有
$U_{1, i} = U$。此外，两个态射 $K_1 \to K_0$ 的分量映射
$f_{\delta^1_0, i}$ 和 $f_{\delta^1_1, i}$ 是态射
% P10-E0223
$$
a : U \to V \quad\text{且}\quad b : U \to V
$$
因此，引理 \ref{spaces-simplicial-lemma-prepare-bbd-glueing} 中所述的相容性给出
$$
\rho_a \circ \rho_V|_{\mathcal{C}/u(U)} = \rho_U
\quad\text{且}\quad
\rho_b \circ \rho_W|_{\mathcal{C}/u(U)} = \rho_U
$$
例如取 $i = (U, V, U, a, \text{id}_U) \in \Omega$，
就得到所需的相容性。$F$ 的唯一性由前一引理中 $E$
的唯一性得出（略去一个小细节）。
\end{proof}

\begin{lemma}[无界 BBD 粘合引理]
\label{spaces-simplicial-lemma-bbd-unbounded-glueing}
处于情形 \ref{spaces-simplicial-situation-locally-given}。假设
\begin{enumerate}
\item $\mathcal{C}$ 具有等化子与纤维积；
\item 存在由连续函子 $u : \mathcal{D} \to \mathcal{C}$
给出的站点态射 $f : \mathcal{C} \to \mathcal{D}$，使得
\begin{enumerate}
\item $\mathcal{D}$ 具有等化子与纤维积，且 $u$ 与它们交换；
\item $\mathcal{B}$ 是 $\mathcal{D}$ 的全子范畴，而
$u : \mathcal{B} \to \mathcal{C}$ 是 $u$ 的限制；
\item $\mathcal{D}$ 的每个对象都有一个由 $\mathcal{B}$
的对象组成的覆盖；
\end{enumerate}
\item $E_U$ 在 $D(\mathcal{O}_{u(U)})$ 中所有负次数自 Ext 均为零；
\item 对所有 $U \in \Ob(\mathcal{C})$，存在满足条件
(\ref{spaces-simplicial-item-restriction})、(\ref{spaces-simplicial-item-local}) 和
(\ref{spaces-simplicial-item-bounded-dimension}) 的弱 Serre 子范畴
$\mathcal{A}_U \subset \textit{Mod}(\mathcal{O}_U)$；
\item % P10-E0225
$E_U \in D_{\mathcal{A}_U}(\mathcal{O}_U)$。
\end{enumerate}
则存在一个解，并且它在唯一同构的意义下唯一。
\end{lemma}

\begin{proof}
证明与引理 \ref{spaces-simplicial-lemma-bbd-glueing} 的证明{\bf 完全}相同。
唯一的改变是 $E$ 为 $D_{\mathcal{A}_{total}}(\mathcal{O})$ 的对象，
所以我们使用引理 \ref{spaces-simplicial-lemma-hypercovering-equivalence-modules}
得到满足 $E = a^*F$ 的 $F$，而不是使用引理
\ref{spaces-simplicial-lemma-hypercovering-equivalence-bounded-modules}。
\end{proof}

\noindent
下面给出上述一般理论的一个应用实例。

\begin{lemma}
\label{spaces-simplicial-lemma-bbd-glueing-cech}
\begin{reference}
Martin Olsson 于 2021 年 9 月 9 日发出的电子邮件。
\end{reference}
设 $(\mathcal{C}, \mathcal{O}_\mathcal{C})$ 是赋环站点，
并假设 $\mathcal{C}$ 具有纤维积。设
$\{U_i \to X\}_{i \in I}$ 是 $\mathcal{C}$ 中的覆盖。
对 $i \in I$，设 $E_i$ 是 $D(\mathcal{O}_{U_i})$ 的对象；
对 $i, j \in I$，设
$$
\rho_{ij} : E_i|_{\mathcal{C}/U_{ij}} \longrightarrow E_j|_{\mathcal{C}/U_{ij}}
$$
为 $D(\mathcal{O}_{U_{ij}})$ 中的同构，其中
$U_{ij} = U_i \times_X U_j$。假设
\begin{enumerate}
\item 对所有 $i, j, k \in I$，$\rho_{ij}$ 在
$U_i \times_X U_j \times_X U_k$ 上满足余圈条件；
\item $\SheafExt^p_{\mathcal{O}_{U_i}}(E_i, E_i) = 0$
对所有 $p < 0$ 和 $i \in I$ 成立；
\item 存在 $t \in \mathbf{Z}$，使得对 $p < t$ 及所有
$i \in I$，都有 $H^p(E_i) = 0$。
\end{enumerate}
则存在唯一的二元组 $(E, \rho_i)$，其中 $E$ 是
$D(\mathcal{O}_X)$ 的对象，而
$\rho_i : E|_{U_i} \to E_i$ 是 $D(\mathcal{O}_{U_i})$
中与 $\rho_{ij}$ 相容的同构。
\end{lemma}

\begin{proof}[第一个证明]
本证明由十分一般的引理 \ref{spaces-simplicial-lemma-bbd-glueing} 推出该引理。
% P10-E0226
我们建议读者改看第二个证明。

\medskip\noindent
可以用 $\mathcal{C}/X$ 替换 $\mathcal{C}$。
因此，我们可以并且确实假设 $X$ 是 $\mathcal{C}$ 的终对象，
且 $\mathcal{C}$ 具有所有有限极限。

\medskip\noindent
令 $\mathcal{B}$ 为 $\mathcal{C}$ 的全子范畴，其对象是满足下列条件的
$U \in \Ob(\mathcal{C})$：存在 $i(U) \in I$ 和态射
$a_U : U \to U_{i(U)}$。在 $D(\mathcal{O}_U)$ 中，以
$E_U = a_U^*E_{i(U)}$ 表示 $E_i$ 经 $a_U$ 的拉回（限制）。
% P10-E0227
给定 $\mathcal{B}$ 的态射 $a : U \to U'$，得到态射
$(a_{U'} \circ a, a_U) : U \to U_{i(U')} \times_X U_{i(U)} = U_{i(U')i(U)}$
，从而得到 $D(\mathcal{O}_U)$ 中的同构
$$
\rho_a :
a^*E_{U'} = a^*a_{U'}^*E_{i(U')}
\xrightarrow{(a_{U'} \circ a, a_U)^*\rho_{i(U')i(U)}}
a_{U}^*E_{i(U)} = E_{U}
$$
。数据 $\mathcal{B}, E_U, \rho_a$ 如情形
\ref{spaces-simplicial-situation-locally-given} 所述；同构 $\rho_a$
恰因引理陈述中的条件 (1) 而满足余圈条件（略去细节）。

\medskip\noindent
我们将把引理 \ref{spaces-simplicial-lemma-bbd-glueing} 应用于上述
$\mathcal{B}, E_U, \rho_a$，并取 $\mathcal{D} = \mathcal{C}$、
$f : \mathcal{C} \to \mathcal{D}$ 为恒同态射。
引理 \ref{spaces-simplicial-lemma-bbd-glueing} 的假设 (1) 和 (2)(a) 已在上面看到。
引理 \ref{spaces-simplicial-lemma-bbd-glueing} 的假设 (2)(b) 显然成立。
引理 \ref{spaces-simplicial-lemma-bbd-glueing} 的假设 (2)(c) 成立，因为
$\{U_i \to X\}$ 是覆盖\footnote{事实上，只需映射
$\coprod_{i \in I} h_{U_i} \to h_X$ 在层化后成为满射；
在这种情形下，该引理以同一证明成立。}。
引理 \ref{spaces-simplicial-lemma-bbd-glueing} 的假设 (3) 成立，因为我们已假设
$E_i$ 的所有负次数 Ext 层均消灭；这当然蕴含对位于 $U_i$
上方的任意对象 $U$，$E_i|_U$ 的负次数自 Ext 均为零。
引理 \ref{spaces-simplicial-lemma-bbd-glueing} 的假设 (4) 成立，因为我们已假设
每个 $E_i$ 的上同调层在 $t$ 的左侧为零。

\medskip\noindent
我们得到唯一的解 $(E, \rho_U)$。置 $\rho_i = \rho_{U_i}$，
便得到引理。
\end{proof}

\begin{proof}[第二个证明]
我们概述一个更直接的证明。以 $K$ 表示与覆盖
$\{U_i \to X\}_{i \in I}$ 关联的 $X$ 的 {\v C}ech 超覆盖；
参见《超覆盖》例 \ref{hypercovering-example-cech}。
例如，$K_0 = \{U_i \to X\}_{i \in I}$，而
$K_1 = \{U_i \times_X U_j \to X\}_{i, j \in I}$，依此类推。
令 $((\mathcal{C}/K)_{total}, \mathcal{O})$、$a$、$a_n$
如评注 \ref{spaces-simplicial-remark-augmentation-ringed-over-object} 所述。
对象 $E_i$ 在 $D(\mathcal{O}_0) = \prod D(\mathcal{O}_{U_i})$
中确定对象 $M_0$。类似地，同构 $\rho_{ij}$ 确定同构
$$
\alpha :
L(f_{\delta_1^1})^*M_0
\longrightarrow
L(f_{\delta_0^1})^*M_0
$$
，它属于 $D(\mathcal{O}_1)$ 并满足余圈条件。
由引理 \ref{spaces-simplicial-lemma-construct-cartesian-objects-derived-modules}，
得到导出范畴的笛卡尔单纯系统 $(M_n)$。
由所假设的负次数 Ext 层消灭性，可知对象 $M_n$
的负次数自 Ext 均消灭。因此，由引理
\ref{spaces-simplicial-lemma-cartesian-module-derived-from-simplicial}，得到
$D(\mathcal{O})$ 的笛卡尔对象 $M$，其关联单纯系统同构于 $(M_n)$。
由于 $M$ 的上同调层在次数 $< t$ 处为零，由引理
\ref{spaces-simplicial-lemma-hypercovering-X-equivalence-bounded-modules}，存在
$D(\mathcal{O}_X)$ 中的 $E$，使得 $M = La^*E$。
将同构 $La^*E \to M$ 限制到 $\mathcal{C}/U_i$，
便得到同构 $\rho_i$。略去它们与同构 $\rho_{ij}$ 相容性的验证。
\end{proof}












\section{拓扑中的真超覆盖}
\label{spaces-simplicial-section-proper-hypercovering}

\noindent
在由 Hausdorff 局部拟紧拓扑空间与连续映射组成的范畴
$\textit{LC}$ 中工作；参见《站点上的上同调》第
\ref{sites-cohomology-section-cohomology-LC} 节。
设 $X$ 是 $\textit{LC}$ 的对象，$U$ 是 $\textit{LC}$
的单纯对象。假设有增广
$$
a : U \to X
$$
如果满足下列条件，就称 $U$ 是 $X$ 的{\it 真超覆盖}：
\begin{enumerate}
\item $U_0 \to X$ 是真满射；
\item $U_1 \to U_0 \times_X U_0$ 是真满射；
\item $U_{n + 1} \to (\text{cosk}_n\text{sk}_n U)_{n + 1}$
对 $n \geq 1$ 是真满射。
\end{enumerate}
$\textit{LC}$ 具有所有有限极限，所以上述表述中使用的余骨架存在。
$$
\fbox{原理：可以用真超覆盖计算上同调。}
$$
证明该原理的一个关键思想，是在 $\textit{LC}$ 上找到一个比通常拓扑
更细的拓扑，使得：(a) 真满射定义覆盖；(b) 通常层相对于这个更细
拓扑的上同调与通常上同调一致。qc 拓扑满足性质 (a) 和 (b)；
参见《站点上的上同调》第
\ref{sites-cohomology-section-cohomology-LC} 节。
一旦有了 (a) 和 (b)，便可由本章前面的工作推出该原理。

\begin{lemma}
\label{spaces-simplicial-lemma-compare-simplicial-objects}
设 $U$ 是 $\textit{LC}$ 的单纯对象，$a : U \to X$ 是增广。
则有交换图
$$
\xymatrix{
\Sh((\textit{LC}_{qc}/U)_{total}) \ar[r]_-h \ar[d]_{a_{qc}} &
\Sh(U_{Zar}) \ar[d]^a \\
\Sh(\textit{LC}_{qc}/X) \ar[r]^-{h_{-1}} &
\Sh(X)
}
$$
其中左侧竖直箭头定义于第 \ref{spaces-simplicial-section-hypercovering} 节，
右侧竖直箭头定义于引理 \ref{spaces-simplicial-lemma-augmentation}。
\end{lemma}

\begin{proof}
写 $\Sh(X) = \Sh(X_{Zar})$。注意，
$(\textit{LC}_{qc}/U)_{total}$ 与 $U_{Zar}$ 都属于情形
\ref{spaces-simplicial-situation-simplicial-site} 的 A 类。对 $U_{Zar}$，
这立即由第 \ref{spaces-simplicial-section-simplicial-top} 节中的构造得出；
对 $(\textit{LC}_{qc}/U)_{total}$，则由引理
\ref{spaces-simplicial-lemma-sr-when-fibre-products} 得出。接着考虑函子
$U_{n, Zar} \to \textit{LC}_{qc}/U_n$, $U \mapsto U/U_n$
以及
$X_{Zar} \to \textit{LC}_{qc}/X$, $U \mapsto U/X$.
我们已在《站点上的上同调》第
\ref{sites-cohomology-section-cohomology-LC} 节看到，
它们定义站点态射。因此，得到评注
\ref{spaces-simplicial-remark-morphism-augmentation-simplicial-sites}
意义下与增广相容的单纯站点态射；应用引理
\ref{spaces-simplicial-lemma-morphism-augmentation-simplicial-sites} 即得结论。
\end{proof}

\begin{lemma}
\label{spaces-simplicial-lemma-descent-sheaves-for-proper-hypercovering}
设 $U$ 是 $\textit{LC}$ 的单纯对象，$a : U \to X$ 是增广。
若 $a : U \to X$ 给出 $X$ 的真超覆盖，则
$$
a^{-1} : \Sh(X) \to \Sh(U_{Zar})
\quad\text{且}\quad
a^{-1} : \textit{Ab}(X) \to \textit{Ab}(U_{Zar})
$$
满忠实，其本质像为笛卡尔层，拟逆由 $a_*$ 给出。
这里 $a : \Sh(U_{Zar}) \to \Sh(X)$ 如引理
\ref{spaces-simplicial-lemma-augmentation} 所述。
\end{lemma}

\begin{proof}
我们证明集合层的陈述。它几乎是已建立结果的形式推论。
考虑引理 \ref{spaces-simplicial-lemma-compare-simplicial-objects} 的图。
由《站点上的上同调》引理
\ref{sites-cohomology-lemma-describe-pullback-pi}，函子
$(h_{-1})^{-1}$ 满忠实，其拟逆为 $h_{-1, *}$。
对 $h$ 的分量 $h_n$ 也同样成立。由引理
\ref{spaces-simplicial-lemma-morphism-simplicial-sites} 对函子 $h^{-1}$ 和 $h_*$
的描述，可知 $h^{-1}$ 满忠实，其拟逆为 $h_*$。
$U$ 是 $\textit{LC}_{qc}$ 中 $X$ 的超覆盖
（定义见第 \ref{spaces-simplicial-section-hypercovering} 节）；这由《站点上的上同调》引理
\ref{sites-cohomology-lemma-proper-surjective-is-qc-covering} 得出。
由引理 \ref{spaces-simplicial-lemma-hypercovering-X-simple-descent-sheaves}，
$a_{qc}^{-1}$ 满忠实，其拟逆为 $a_{qc, *}$，本质像为
$(\textit{LC}_{qc}/U)_{total}$ 上的笛卡尔层。
现在，一个形式论证（沿图追逐）表明 $a^{-1}$ 满忠实。

\medskip\noindent
最后，假设 $\mathcal{G}$ 是 $U_{Zar}$ 上的笛卡尔层。
则 $h^{-1}\mathcal{G}$ 是 $\textit{LC}_{qc}/U$ 上的笛卡尔层。
因此，对 $\textit{LC}_{qc}/X$ 上的某个层 $\mathcal{H}$，有
$h^{-1}\mathcal{G} = a_{qc}^{-1}\mathcal{H}$。计算得
\begin{align*}
(h_{-1})^{-1}(a_*\mathcal{G})
& =
(h_{-1})^{-1}
\text{等化子}(
\xymatrix{
a_{0, *}\mathcal{G}_0
\ar@<1ex>[r] \ar@<-1ex>[r] &
a_{1, *}\mathcal{G}_1
}
) \\
& =
\text{等化子}(
\xymatrix{
(h_{-1})^{-1}a_{0, *}\mathcal{G}_0
\ar@<1ex>[r] \ar@<-1ex>[r] &
(h_{-1})^{-1}a_{1, *}\mathcal{G}_1
}
) \\
& =
\text{等化子}(
\xymatrix{
a_{qc, 0, *}h_0^{-1}\mathcal{G}_0
\ar@<1ex>[r] \ar@<-1ex>[r] &
a_{qc, 1, *}h_1^{-1}\mathcal{G}_1
}
) \\
& =
\text{等化子}(
\xymatrix{
a_{qc, 0, *}a_{qc, 0}^{-1}\mathcal{H}
\ar@<1ex>[r] \ar@<-1ex>[r] &
a_{qc, 1, *}a_{qc, 1}^{-1}\mathcal{H}
}
) \\
& =
a_{qc, *}a_{qc}^{-1}\mathcal{H} \\
& =
\mathcal{H}
\end{align*}
第一个等号由引理 \ref{spaces-simplicial-lemma-augmentation} 得出；第二个等号来自
$(h_{-1})^{-1}$ 是正合函子；第三个等号由《站点上的上同调》引理
\ref{sites-cohomology-lemma-push-pull-LC} 得出
（此处使用 $a_0 : U_0 \to X$ 和 $a_1: U_1 \to X$ 为真态射）；
第四个等号来自 $a_{qc}^{-1}\mathcal{H} = h^{-1}\mathcal{G}$；
第五个等号来自引理 \ref{spaces-simplicial-lemma-augmentation-site}；
第六个等号已在上面看到。由于
$a_{qc}^{-1}\mathcal{H} = h^{-1}\mathcal{G}$，可得
$h^{-1}\mathcal{G} \cong h^{-1}a^{-1}a_*\mathcal{G}$；
再利用 $h^{-1}$ 的满忠实性，证明即告完成。
\end{proof}

\begin{lemma}
\label{spaces-simplicial-lemma-cohomological-descent-for-proper-hypercovering}
设 $U$ 是 $\textit{LC}$ 的单纯对象，$a : U \to X$ 是增广。
若 $a : U \to X$ 给出 $X$ 的真超覆盖，则对 $K \in D^+(X)$，
$$
K \to Ra_*(a^{-1}K)
$$
是同构，其中 $a : \Sh(U_{Zar}) \to \Sh(X)$ 如引理
\ref{spaces-simplicial-lemma-augmentation} 所述。
\end{lemma}

\begin{proof}
考虑引理 \ref{spaces-simplicial-lemma-compare-simplicial-objects} 的图。
由《站点上的上同调》引理
\ref{sites-cohomology-lemma-cohomological-descent-LC}，
$Rh_{n, *}h_n^{-1}$ 是 $D^+(U_n)$ 上的恒同函子。
因此，由引理 \ref{spaces-simplicial-lemma-direct-image-morphism-simplicial-sites}，
$Rh_*h^{-1}$ 是 $D^+(U_{Zar})$ 上的恒同函子。我们有
\begin{align*}
Ra_*(a^{-1}K)
& =
Ra_*Rh_*h^{-1}a^{-1}K \\
& =
Rh_{-1, *}Ra_{qc, *}a_{qc}^{-1}(h_{-1})^{-1}K \\
& =
Rh_{-1, *}(h_{-1})^{-1}K \\
& =
K
\end{align*}
% P10-E0228
第一个等号来自上述讨论；第二个等号来自引理
\ref{spaces-simplicial-lemma-compare-simplicial-objects} 中图的交换性；
第三个等号来自引理
\ref{spaces-simplicial-lemma-hypercovering-X-simple-descent-bounded-abelian}
（由《站点上的上同调》引理
\ref{sites-cohomology-lemma-proper-surjective-is-qc-covering}，
$U$ 是 $\textit{LC}_{qc}$ 中 $X$ 的超覆盖）；最后一个等号
来自已经使用过的《站点上的上同调》引理
\ref{sites-cohomology-lemma-cohomological-descent-LC}。
\end{proof}

\begin{lemma}
\label{spaces-simplicial-lemma-compute-via-proper-hypercovering}
设 $U$ 是 $\textit{LC}$ 的单纯对象，$a : U \to X$ 是增广。
若 $U$ 是 $X$ 的真超覆盖，则
$$
R\Gamma(X, K) = R\Gamma(U_{Zar}, a^{-1}K)
$$
其中 $K \in D^+(X)$，而 $a : \Sh(U_{Zar}) \to \Sh(X)$
如引理 \ref{spaces-simplicial-lemma-augmentation} 所述。
\end{lemma}

\begin{proof}
这由引理 \ref{spaces-simplicial-lemma-cohomological-descent-for-proper-hypercovering} 得出，
因为按《站点上的上同调》评注
\ref{sites-cohomology-remark-before-Leray}，有
$R\Gamma(U_{Zar}, -) = R\Gamma(X, -) \circ Ra_*$。
\end{proof}

\begin{lemma}
\label{spaces-simplicial-lemma-proper-hypercovering-equivalence-bounded}
设 $U$ 是 $\textit{LC}$ 的单纯对象，$a : U \to X$ 是增广。
以 $\mathcal{A} \subset \textit{Ab}(U_{Zar})$
表示笛卡尔阿贝尔层组成的弱 Serre 子范畴。
若 $U$ 是 $X$ 的真超覆盖，则函子 $a^{-1}$ 给出等价
$$
D^+(X) \longrightarrow D_\mathcal{A}^+(U_{Zar})
$$
，其拟逆为 $Ra_*$；这里 $a : \Sh(U_{Zar}) \to \Sh(X)$
如引理 \ref{spaces-simplicial-lemma-augmentation} 所述。
\end{lemma}

\begin{proof}
由引理 \ref{spaces-simplicial-lemma-Serre-subcat-cartesian-modules}，
$\mathcal{A}$ 是弱 Serre 子范畴。该等价是迄今所得结果的形式推论。
应用引理 \ref{spaces-simplicial-lemma-descent-sheaves-for-proper-hypercovering}、
\ref{spaces-simplicial-lemma-cohomological-descent-for-proper-hypercovering} 以及
《站点上的上同调》引理 \ref{sites-cohomology-lemma-equivalence-bounded}。
\end{proof}

\begin{lemma}
\label{spaces-simplicial-lemma-spectral-sequence-proper-hypercovering}
设 $U$ 是 $\textit{LC}$ 的单纯对象，$a : U \to X$ 是增广。
设 $\mathcal{F}$ 是 $X$ 上的阿贝尔层，$\mathcal{F}_n$
是它到 $U_n$ 的拉回。若 $U$ 是 $X$ 的真超覆盖，
则存在典范谱序列
$$
E_1^{p, q} = H^q(U_p, \mathcal{F}_p)
$$
，它收敛到 $H^{p + q}(X, \mathcal{F})$。
\end{lemma}

\begin{proof}
这是引理 \ref{spaces-simplicial-lemma-compute-via-proper-hypercovering} 和
\ref{spaces-simplicial-lemma-simplicial-sheaf-cohomology} 的立即推论。
\end{proof}




\section{单纯概形}
\label{spaces-simplicial-section-simplicial}

\noindent
{\it 单纯概形}是概形范畴中的单纯对象；参见《单纯对象》定义
\ref{simplicial-definition-simplicial-object}。回忆，单纯概形形如
$$
\xymatrix{
X_2
\ar@<2ex>[r]
\ar@<0ex>[r]
\ar@<-2ex>[r]
&
X_1
\ar@<1ex>[r]
\ar@<-1ex>[r]
\ar@<1ex>[l]
\ar@<-1ex>[l]
&
X_0
\ar@<0ex>[l]
}
$$
这里有两个态射 $d^1_0, d^1_1 : X_1 \to X_0$ 和一个态射
$s^0_0 : X_0 \to X_1$，依此类推。这些态射满足一些必需关系，
例如 $d^1_0 \circ s^0_0 = \text{id}_{X_0} = d^1_1 \circ s^0_0$；
参见《单纯对象》引理
\ref{simplicial-lemma-characterize-simplicial-object}。
把 $d^n_i : X_n \to X_{n - 1}$ 看作“忘掉第 $i$ 个坐标的投影”，
把 $s^n_j : X_n \to X_{n + 1}$ 看作“重复第 $j$ 个坐标的对角映射”，
会很有帮助。

\medskip\noindent
{\it 单纯概形的态射} $h : X \to Y$，就是概形范畴中单纯对象的态射；
参见《单纯对象》定义 \ref{simplicial-definition-simplicial-object}。
因此，$h$ 由概形态射 $h_n : X_n \to Y_n$ 组成，并且只要等式有意义，
就有 $h_{n - 1} \circ d^n_j = d^n_j \circ h_n$ 和
$h_{n + 1} \circ s^n_j = s^n_j \circ h_n$。

\medskip\noindent
单纯概形 $X$ 的{\it 增广}是概形态射 $a_0 : X_0 \to S$，
满足 $a_0 \circ d^1_0 = a_0 \circ d^1_1$。
参见《单纯对象》第 \ref{simplicial-section-augmentation} 节。

\medskip\noindent
设 $X$ 是单纯概形。将第 \ref{spaces-simplicial-section-simplicial-top} 节的构造
应用于其底单纯拓扑空间，得到站点 $X_{Zar}$。另一方面，
对每个 $n$ 有小 Zariski 站点 $X_{n, Zar}$（《拓扑》定义
\ref{topologies-definition-big-small-Zariski}）；对每个态射
$\varphi : [m] \to [n]$，有与概形态射
$X(\varphi) : X_n \to X_m$ 关联的站点态射
$f_\varphi = X(\varphi)_{small} : X_{n, Zar} \to X_{m, Zar}$
（《拓扑》引理 \ref{topologies-lemma-morphism-big-small}）。
这给出站点范畴中的单纯对象 $\mathcal{C}$。在引理
\ref{spaces-simplicial-lemma-simplicial-site-site} 中，我们构造了关联站点
$\mathcal{C}_{total}$。把开浸入映到其像，定义了等价
$\mathcal{C}_{total} \to X_{Zar}$，并在此等价下识别层；即
$\Sh(\mathcal{C}_{total}) = \Sh(X_{Zar})$。
$\mathcal{C}_{total}$ 与 $X_{Zar}$ 之间的差别，类似于小 Zariski
站点 $S_{Zar}$ 与 $S$ 的底拓扑空间之间的差别。下文默认识别这些站点。

\medskip\noindent
设 $X_{Zar}$ 是与单纯概形 $X$ 关联的站点。
$X_{Zar}$ 上存在环层 $\mathcal{O}$，它在 $X_n$ 上的限制
是结构层 $\mathcal{O}_{X_n}$。这由引理
\ref{spaces-simplicial-lemma-describe-sheaves-simplicial-site} 或引理
\ref{spaces-simplicial-lemma-describe-sheaves-simplicial-site-site} 得出。
我们称{\it $\mathcal{O}$ 是单纯概形 $X$ 的结构层}。
至此，为单纯（赋环）站点发展的一切材料都适用；参见第
\ref{spaces-simplicial-section-simplicial-sites}、
\ref{spaces-simplicial-section-augmentation-simplicial-sites}、
\ref{spaces-simplicial-section-morphism-simplicial-sites}、
\ref{spaces-simplicial-section-simplicial-sites-modules}、
\ref{spaces-simplicial-section-cohomology-simplicial-sites}、
\ref{spaces-simplicial-section-cohomology-augmentation-simplicial-sites}、
\ref{spaces-simplicial-section-cohomology-simplicial-sites-modules}、
\ref{spaces-simplicial-section-cohomology-augmentation-ringed-simplicial-sites}、
\ref{spaces-simplicial-section-cartesian}、
\ref{spaces-simplicial-section-glueing} 和
\ref{spaces-simplicial-section-glueing-modules} 节。

\medskip\noindent
设 $X$ 是结构层为 $\mathcal{O}$ 的单纯概形。
与任何赋环拓扑斯上一样，可以定义
{\it $X_{Zar}$ 上的拟凝聚 $\mathcal{O}$-模}；参见
《站点上的模》定义 \ref{sites-modules-definition-site-local}。
不过，$X_{Zar}$ 上的拟凝聚 $\mathcal{O}$-模，恰是如下笛卡尔
$\mathcal{O}$-模 $\mathcal{F}$：其限制 $\mathcal{F}_n$
在 $X_n$ 上拟凝聚；参见引理 \ref{spaces-simplicial-lemma-quasi-coherent-sheaf}。

\medskip\noindent
设 $h : X \to Y$ 是单纯概形的态射。由引理
\ref{spaces-simplicial-lemma-simplicial-space-site-functorial}，或由引理
\ref{spaces-simplicial-lemma-morphism-simplicial-sites} 的证明，得到站点态射
$h_{Zar} : X_{Zar} \to Y_{Zar}$。回忆，$h_{Zar}^{-1}$ 和
$h_{Zar, *}$ 可用各分量简单描述；参见引理
\ref{spaces-simplicial-lemma-describe-functoriality} 或引理
\ref{spaces-simplicial-lemma-morphism-simplicial-sites}。
以 $\mathcal{O}_X$、$\mathcal{O}_Y$ 分别表示 $X$、$Y$ 的结构层。
定义
% P10-E0229
$h_{Zar}^\sharp : h_{Zar, *}\mathcal{O}_X \to \mathcal{O}_Y$
为 $Y_{Zar}$ 上的环层映射，它在 $Y_n$ 上由
$h_n^\sharp : h_{n, *}\mathcal{O}_{X_n} \to \mathcal{O}_{Y_n}$ 给出。
由此得到赋环站点的态射
$$
h_{Zar} : (X_{Zar}, \mathcal{O}_X) \longrightarrow (Y_{Zar}, \mathcal{O}_Y)
$$

\medskip\noindent
设 $X$ 是结构层为 $\mathcal{O}$ 的单纯概形。设 $S$ 是概形，
$a_0 : X_0 \to S$ 是 $X$ 的增广。由引理
\ref{spaces-simplicial-lemma-augmentation} 或引理 \ref{spaces-simplicial-lemma-augmentation-site}，
得到相应拓扑斯态射 $a : \Sh(X_{Zar}) \to \Sh(S)$。
注意，$a^{-1}\mathcal{G}$ 是 $X_{Zar}$ 上分量为
$a_n^{-1}\mathcal{G}$ 的层。因此，可以利用映射
$a_n^\sharp : a_n^{-1}\mathcal{O}_S \to \mathcal{O}_{X_n}$
定义映射 $a^\sharp : a^{-1}\mathcal{O}_S \to \mathcal{O}$；
等价地，由伴随定义映射
$a^\sharp : \mathcal{O}_S \to a_*\mathcal{O}$
（照例仍用同一名称）。这正处于第
\ref{spaces-simplicial-section-cohomology-augmentation-ringed-simplicial-sites}
节所讨论的情形。因此，得到赋环拓扑斯的态射
$$
a : (\Sh(X_{Zar}), \mathcal{O}) \longrightarrow (\Sh(S), \mathcal{O}_S)
$$

\medskip\noindent
最后作如下观察。设给定单纯概形 $X$ 与 $Y$ 的态射
$h : X \to Y$，其结构层分别为 $\mathcal{O}_X$、$\mathcal{O}_Y$；
并给定增广 $a_0 : X_0 \to X_{-1}$、$b_0 : Y_0 \to Y_{-1}$
和态射 $h_{-1} : X_{-1} \to Y_{-1}$，使得
$$
\xymatrix{
X_0 \ar[r]_{h_0} \ar[d]_{a_0} & Y_0 \ar[d]^{b_0} \\
X_{-1} \ar[r]^{h_{-1}} & Y_{-1}
}
$$
交换。由上面阐明的构造，得到如下赋环拓扑斯态射的交换图：
$$
\xymatrix{
(\Sh(X_{Zar}), \mathcal{O}_X) \ar[r]_{h_{Zar}} \ar[d]_a &
(\Sh(Y_{Zar}), \mathcal{O}_Y) \ar[d]^b \\
(\Sh(X_{-1}), \mathcal{O}_{X_{-1}}) \ar[r]^{h_{-1}} &
(\Sh(Y_{-1}), \mathcal{O}_{Y_{-1}})
}
$$








\section{用单纯概形表述下降}
\label{spaces-simplicial-section-simplicial-descent}

\noindent
笛卡尔态射定义如下。

\begin{definition}
\label{spaces-simplicial-definition-cartesian-morphism}
设 $a : Y \to X$ 是单纯概形的态射。如果对 $\Delta$
的每个态射 $\varphi : [n] \to [m]$，相应的图
$$
\xymatrix{
Y_m \ar[r]_a \ar[d]_{Y(\varphi)} & X_m \ar[d]^{X(\varphi)}\\
Y_n \ar[r]^{a} & X_n
}
$$
是概形范畴中的纤维方块，就称 $a$ 是{\it 笛卡尔的}，
或称{\it $Y$ 在 $X$ 上是笛卡尔的}。
\end{definition}

\noindent
笛卡尔态射与下降数据有关。首先证明一个一般引理，
描述固定单纯概形上的笛卡尔单纯概形之范畴。在该引理中，
以 $f^* : \Sch/X \to \Sch/Y$ 表示与概形态射
$f :Y \to X$ 关联的基变换函子。

\begin{lemma}
\label{spaces-simplicial-lemma-characterize-cartesian-schemes}
设 $X$ 是单纯概形。在 $X$ 上笛卡尔的单纯概形之范畴，
等价于二元组 $(V, \varphi)$ 的范畴，其中 $V$ 是 $X_0$ 上的概形，而
$$
\varphi :
V \times_{X_0, d^1_1} X_1
\longrightarrow
X_1 \times_{d^1_0, X_0} V
$$
是 $X_1$ 上的同构，满足 $(s_0^0)^*\varphi = \text{id}_V$，且
$$
(d^2_1)^*\varphi = (d^2_0)^*\varphi \circ (d^2_2)^*\varphi
$$
作为 $X_2$ 上的概形态射成立。
\end{lemma}

\begin{proof}
所显示等式的陈述有意义，因为
$d^1_1 \circ d^2_2 = d^1_1 \circ d^2_1$,
$d^1_1 \circ d^2_0 = d^1_0 \circ d^2_2$, and
$d^1_0 \circ d^2_0 = d^1_0 \circ d^2_1$ 作为态射
$X_2 \to X_0$ 成立；参见《单纯对象》评注
\ref{simplicial-remark-relations}。因此，可以把这些映射画成
$$
\xymatrix{
&
X_2 \times_{d^1_1 \circ d^2_0, X_0} V
\ar[r]_-{(d^2_0)^*\varphi} &
X_2 \times_{d^1_0 \circ d^2_0, X_0} V
\ar@{=}[rd] & \\
X_2 \times_{d^1_0 \circ d^2_2, X_0} V
\ar@{=}[ru] & & &
X_2 \times_{d^1_0 \circ d^2_1, X_0} V \\
&
X_2 \times_{d^1_1 \circ d^2_2, X_0} V
\ar[lu]^{(d^2_2)^*\varphi} \ar@{=}[r] &
X_2 \times_{d^1_1 \circ d^2_1, X_0} V
\ar[ru]_{(d^2_1)^*\varphi}
}
$$
，而该条件表示此图交换。显然，给定在 $X$ 上笛卡尔的单纯概形
$Y$，可以置 $V = Y_0$，并令 $\varphi$ 等于如下由笛卡尔结构
给出的恒等式之复合：
$$
V \times_{X_0, d^1_1} X_1 =
Y_0 \times_{X_0, d^1_1} X_1 = Y_1 =
X_1 \times_{X_0, d^1_0} Y_0 =
X_1 \times_{X_0, d^1_0} V
$$
为证明该函子是等价，构造一个拟逆。拟逆的构造类似于《下降》第
\ref{descent-section-descent-modules} 节所讨论的构造；我们从那里借用
记号 $\tau^n_i : [0] \to [n]$, $0 \mapsto i$ 和
$\tau^n_{ij} : [1] \to [n]$, $0 \mapsto i$, $1 \mapsto j$。
具体地，给定引理中的二元组 $(V, \varphi)$，置
$Y_n = X_n \times_{X(\tau^n_n), X_0} V$。给定
$\beta : [n] \to [m]$，定义
% P10-E0230
$V(\beta) : Y_m \to Y_n$ 为映射 $\varphi$ 与投影
$X_m \times_{X(\beta), X_n} Y_n \to Y_n$ 复合后，经
$X(\tau^m_{\beta(n)m})$ 的拉回。这是有意义的，因为
$$
X_m \times_{X(\tau^m_{\beta(n)m}), X_1} X_1 \times_{d^1_1, X_0} V
=
X_m \times_{X(\tau^m_m), X_0} V = Y_m
$$
以及
$$
X_m \times_{X(\tau^m_{\beta(n)m}), X_1} X_1 \times_{d^1_0, X_0} V =
X_m \times_{X(\tau^m_{\beta(n)}), X_0} V =
X_m \times_{X(\beta), X_n} Y_n.
$$
略去如下验证：上面所显示图的交换性蕴含这些映射正确复合。
也略去两个函子彼此为拟逆的验证。
\end{proof}

\begin{definition}
\label{spaces-simplicial-definition-fibre-products-simplicial-scheme}
设 $f : X \to S$ 是概形态射。与 $f$ {\it 关联的单纯概形}，
记为 $(X/S)_\bullet$，是函子
$\Delta^{opp} \to \Sch$, $[n] \mapsto X \times_S \ldots \times_S X$
，其描述见《单纯对象》例
\ref{simplicial-example-fibre-products-simplicial-object}。
\end{definition}

\noindent
因此，$(X/S)_n$ 是 $X$ 在 $S$ 上的 $(n + 1)$ 重纤维积。
态射 $d^1_0 : X \times_S X \to X$ 是映射
$(x_0, x_1) \mapsto x_1$，态射 $d^1_1$ 是另一个投影。
态射 $s^0_0$ 是对角态射 $X \to X \times_S X$。

\begin{lemma}
\label{spaces-simplicial-lemma-cartesian-over}
设 $f : X \to S$ 是概形态射，
$\pi : Y \to (X/S)_\bullet$ 是单纯概形的笛卡尔态射。
置 $V = Y_0$，并把它视为 $X$ 上的概形。态射
$d^1_0, d^1_1 : Y_1 \to Y_0$ 和态射
$\pi_1 : Y_1 \to X \times_S X$ 诱导同构
$$
\xymatrix{
V \times_S X & &
Y_1 \ar[ll]_-{(d^1_1, \text{pr}_1 \circ \pi_1)}
\ar[rr]^-{(\text{pr}_0 \circ \pi_1, d^1_0)} & &
X \times_S V.
}
$$
以 $\varphi : V \times_S X \to X \times_S V$ 表示由此得到的同构。
则二元组 $(V, \varphi)$ 是相对于 $X \to S$ 的下降数据。
\end{lemma}

\begin{proof}
这是引理 \ref{spaces-simplicial-lemma-characterize-cartesian-schemes}（一部分）
的特殊情形，因为该引理所显示的等式等价于《下降》定义
\ref{descent-definition-descent-datum} 的余圈条件。
\end{proof}

\begin{lemma}
\label{spaces-simplicial-lemma-cartesian-equivalent-descent-datum}
设 $f : X \to S$ 是概形态射。引理
\ref{spaces-simplicial-lemma-cartesian-over} 的构造
$$
\begin{matrix}
\text{笛卡尔概形的范畴} \\
\text{位于 } (X/S)_\bullet \text{ 上}
\end{matrix}
\longrightarrow
\begin{matrix}
\text{下降数据的范畴} \\
\text{相对于 } X/S
\end{matrix}
$$
是范畴的等价。
\end{lemma}

\begin{proof}
从左到右的函子由引理 \ref{spaces-simplicial-lemma-cartesian-over} 给出。
因此，这是引理 \ref{spaces-simplicial-lemma-characterize-cartesian-schemes} 的特殊情形。
\end{proof}

\noindent
可以如下重新解释《下降》引理 \ref{descent-lemma-pullback}
中的拉回。设给定单纯概形的态射 $f : X' \to X$ 和单纯概形的
笛卡尔态射 $Y \to X$。则单纯概形的纤维积（视为“拉回”）
$$
f^*Y = Y \times_X X'
$$
是在 $X'$ 上笛卡尔的单纯概形。设给定概形态射的交换图
$$
\xymatrix{
X' \ar[r]_f \ar[d] & X \ar[d] \\
S' \ar[r] & S.
}
$$
这给出单纯概形的态射
$$
f_\bullet : (X'/S')_\bullet \longrightarrow (X/S)_\bullet.
$$
我们断言，沿态射
$f_\bullet : (X'/S')_\bullet \to (X/S)_\bullet$ 的“拉回”
$f_\bullet^*$，经引理
\ref{spaces-simplicial-lemma-cartesian-equivalent-descent-datum} 的对应，
正是上述《下降》引理 \ref{descent-lemma-pullback}
中用下降数据定义的拉回。







\section{单纯概形上的拟凝聚模}
\label{spaces-simplicial-section-modules-simplicial}

\begin{lemma}
\label{spaces-simplicial-lemma-pullback-cartesian-module}
设 $f : V \to U$ 是单纯概形的态射。给定 $U_{Zar}$ 上的拟凝聚模
$\mathcal{F}$，其拉回 $f^*\mathcal{F}$ 是 $V_{Zar}$ 上的拟凝聚模。
\end{lemma}

\begin{proof}
回忆，$\mathcal{F}$ 是笛卡尔的，且 $\mathcal{F}_n$ 拟凝聚；
参见引理 \ref{spaces-simplicial-lemma-quasi-coherent-sheaf}。由引理
\ref{spaces-simplicial-lemma-describe-functoriality}，有
$(f^*\mathcal{F})_n = f_n^*\mathcal{F}_n$（略去一些细节）。
因此，$(f^*\mathcal{F})_n$ 拟凝聚。该事实与 $\mathcal{F}$
的笛卡尔性质共同蕴含 $f^*\mathcal{F}$ 的笛卡尔性质。
% P10-E0231
于是，再由引理 \ref{spaces-simplicial-lemma-quasi-coherent-sheaf}，
$\mathcal{F}$ 拟凝聚。
\end{proof}

\begin{lemma}
\label{spaces-simplicial-lemma-pushforward-cartesian-module}
设 $f : V \to U$ 是单纯概形的笛卡尔态射。
假设态射 $d^n_j : U_n \to U_{n - 1}$ 平坦，而态射
$V_n \to U_n$ 拟紧且拟分离。对 $V_{Zar}$ 上的拟凝聚模
$\mathcal{G}$，推前 $f_*\mathcal{G}$ 是 $U_{Zar}$ 上的拟凝聚模。
\end{lemma}

\begin{proof}
若 $\mathcal{F} = f_* \mathcal{G}$，则由引理
\ref{spaces-simplicial-lemma-describe-functoriality}，有
$\mathcal{F}_n = f_{n , *}\mathcal{G}_n$。
映射 $\mathcal{F}(\varphi)$ 用基变换映射定义；参见《上同调》第
\ref{cohomology-section-base-change-map} 节。由《概形》引理
\ref{schemes-lemma-push-forward-quasi-coherent}，以及由引理
\ref{spaces-simplicial-lemma-quasi-coherent-sheaf} 得知的 $\mathcal{G}_n$ 拟凝聚性，
层 $\mathcal{F}_n$ 拟凝聚。由《概形的上同调》引理
\ref{coherent-lemma-flat-base-change-cohomology}，以及由引理
\ref{spaces-simplicial-lemma-quasi-coherent-sheaf} 得知的 $\mathcal{G}$ 的笛卡尔性质，
% P10-E0232
沿退化态射 $d^n_j$ 的基变换映射是同构。
因此，由引理 \ref{spaces-simplicial-lemma-check-cartesian-module}，
$\mathcal{F}$ 是笛卡尔的；再由引理
\ref{spaces-simplicial-lemma-quasi-coherent-sheaf}，$\mathcal{F}$ 拟凝聚。
\end{proof}

\begin{lemma}
\label{spaces-simplicial-lemma-adjoint-functors-cartesian-modules}
设 $f : V \to U$ 是单纯概形的笛卡尔态射。
假设态射 $d^n_j : U_n \to U_{n - 1}$ 平坦，而态射
$V_n \to U_n$ 拟紧且拟分离。则 $f^*$ 与 $f_*$
构成 $U_{Zar}$ 和 $V_{Zar}$ 上拟凝聚模范畴之间的一对伴随函子。
\end{lemma}

\begin{proof}
引理 \ref{spaces-simplicial-lemma-pullback-cartesian-module} 和
\ref{spaces-simplicial-lemma-pushforward-cartesian-module} 已表明此陈述有意义。
伴随性立即由每个 $f_n^*$ 都左伴随于 $f_{n, *}$ 这一事实得出。
\end{proof}

\begin{lemma}
\label{spaces-simplicial-lemma-cartesian-modules-with-section}
设 $f : X \to S$ 是具有截面的概形态射\footnote{事实上，
只需假设 $f$ 在 $S$ 上 fpqc 局部地具有截面即可，因为由《下降》第
\ref{descent-section-fpqc-descent-quasi-coherent} 节，
拟凝聚模满足下降。}。令 $(X/S)_\bullet$ 为与 $X \to S$
关联的单纯概形；参见定义
\ref{spaces-simplicial-definition-fibre-products-simplicial-scheme}。
则拉回定义了拟凝聚 $\mathcal{O}_S$-模范畴与
$((X/S)_\bullet)_{Zar}$ 上拟凝聚模范畴之间的等价。
\end{lemma}

\begin{proof}
设 $\sigma : S \to X$ 是 $f$ 的截面。设 $(\mathcal{F}, \alpha)$
是引理 \ref{spaces-simplicial-lemma-characterize-cartesian-modules} 中的二元组。
置 $\mathcal{G} = \sigma^*\mathcal{F}$。考虑图
$$
\xymatrix{
X \ar[r]_-{(\sigma \circ f, 1)} \ar[d]_f &
X \times_S X \ar[d]^{\text{pr}_0} \ar[r]_-{\text{pr}_1} & X \\
S \ar[r]^\sigma & X
}
$$
注意，$\text{pr}_0 = d^1_1$ 且 $\text{pr}_1 = d^1_0$。
因此，$(\sigma \circ f, 1)^*\alpha$ 定义同构
$$
f^*\mathcal{G} = (\sigma \circ f, 1)^*\text{pr}_0^*\mathcal{F}
\longrightarrow
(\sigma \circ f, 1)^*\text{pr}_1^*\mathcal{F} = \mathcal{F}
$$
略去如下验证：该同构与 $\alpha$ 以及典范同构
$\text{pr}_0^*f^*\mathcal{G} \to \text{pr}_1^*f^*\mathcal{G}$ 相容。
\end{proof}






\section{群胚与单纯概形}
\label{spaces-simplicial-section-groupoids-simplicial}

\noindent
给定概形中的群胚，可以构造单纯概形。我们将看到，
群胚上的拟凝聚层范畴等价于关联单纯概形上的笛卡尔拟凝聚层范畴。

\begin{lemma}
\label{spaces-simplicial-lemma-groupoid-simplicial}
设 $(U, R, s, t, c, e, i)$ 是 $S$ 上的群胚概形。
则存在 $S$ 上的单纯概形 $X$，满足下列性质：
\begin{enumerate}
\item $X_0 = U$, $X_1 = R$, $X_2 = R \times_{s, U, t} R$；
\item $s_0^0 = e : X_0 \to X_1$；
\item $d^1_0 = s : X_1 \to X_0$, $d^1_1 = t : X_1 \to X_0$；
\item $s_0^1 = (e \circ t, 1) : X_1 \to X_2$,
$s_1^1 = (1, e \circ t) : X_1 \to X_2$；% P10-E0233
\item $d^2_0 = \text{pr}_1 : X_2 \to X_1$,
$d^2_1 = c : X_2 \to X_1$,
$d^2_2 = \text{pr}_0$；
\item $X = \text{cosk}_2 \text{sk}_2 X$.
\end{enumerate}
对所有 $n$，$X_n = R \times_{s, U, t} \ldots \times_{s, U, t} R$，
其中有 $n$ 个因子。映射 $d^n_j : X_n \to X_{n - 1}$
在点函子上由下式给出：
$$
(r_1, \ldots, r_n) \longmapsto (r_1, \ldots, c(r_j, r_{j + 1}), \ldots, r_n)
$$
这里 $1 \leq j \leq n - 1$；而
$d^n_0(r_1, \ldots, r_n) = (r_2, \ldots, r_n)$，且
$d^n_n(r_1, \ldots, r_n) = (r_1, \ldots, r_{n - 1})$。
\end{lemma}

\begin{proof}
只需验证 (1)、(2)、(3)、(4)、(5) 中规定的规则定义了
$S$ 上的 $2$-截断单纯概形 $U'$，因为此时 (6) 允许置
$X = \text{cosk}_2 U'$；参见《单纯对象》引理
\ref{simplicial-lemma-existence-cosk}。采用点函子方法，
只需验证：若 $(\text{对象}, \text{箭头}, s, t, c, e, i)$ 是群胚，则
$$
\xymatrix{
\text{箭头} \times_{s, \text{对象}, t} \text{箭头}
\ar@<8ex>[d]^{\text{pr}_0}
\ar@<0ex>[d]_c
\ar@<-8ex>[d]_{\text{pr}_1}
\\
\text{箭头}
\ar@<4ex>[d]^t
\ar@<-4ex>[d]_s
\ar@<4ex>[u]^{1, e}
\ar@<-4ex>[u]_{e, 1}
\\
\text{对象}
\ar@<0ex>[u]_e
}
$$
是 $2$-截断单纯集。略去细节。

\medskip\noindent
最后，对 $n > 2$ 的 $X_n$ 的描述，由 $X_0$、$X_1$、$X_2$
的描述以及《单纯对象》评注
\ref{simplicial-remark-inductive-coskeleton} 和引理
\ref{simplicial-lemma-work-out} 归纳得出。也可以证明，
将 $\text{cosk}_2$ 应用于上面所显示的 $2$-截断单纯集，
得到的单纯集之第 $n$ 项等于
$\text{箭头} \times_{s, \text{对象}, t} \ldots
\times_{s, \text{对象}, t} \text{箭头}$，有 $n$ 个因子，
而退化映射如引理所述。略去一些细节。
\end{proof}

\begin{lemma}
\label{spaces-simplicial-lemma-quasi-coherent-groupoid-simplicial}
设 $S$ 是概形，$(U, R, s, t, c)$ 是 $S$ 上的群胚概形，
$X$ 是引理 \ref{spaces-simplicial-lemma-groupoid-simplicial} 构造的 $S$ 上单纯概形。
则 $(U, R, s, t, c)$ 上的拟凝聚模范畴，
等价于 $X_{Zar}$ 上的拟凝聚模范畴。
\end{lemma}

\begin{proof}
这由引理 \ref{spaces-simplicial-lemma-quasi-coherent-sheaf}、
\ref{spaces-simplicial-lemma-characterize-cartesian-modules} 和《群胚》定义
\ref{groupoids-definition-groupoid-module} 显然成立。
\end{proof}

\noindent
在下一个引理中，将使用定义 \ref{spaces-simplicial-definition-cartesian-morphism}
所定义的单纯概形之笛卡尔态射 $V \to U$ 的概念。

\begin{lemma}
\label{spaces-simplicial-lemma-quasi-coherent-groupoid-R-cartesian}
设 $(U, R, s, t, c)$ 是概形 $S$ 上的群胚概形，
$X$ 是引理 \ref{spaces-simplicial-lemma-groupoid-simplicial} 构造的 $S$ 上单纯概形。
令 $(R/U)_\bullet$ 为与 $s : R \to U$ 关联的单纯概形；
参见定义 \ref{spaces-simplicial-definition-fibre-products-simplicial-scheme}。
则存在单纯概形的笛卡尔态射
$t_\bullet : (R/U)_\bullet \to X$，其低次数态射由下图给出：
$$
\xymatrix{
R \times_{s, U, s} R \times_{s, U, s} R
\ar@<3ex>[r]_-{\text{pr}_{12}}
\ar@<0ex>[r]_-{\text{pr}_{02}}
\ar@<-3ex>[r]_-{\text{pr}_{01}}
\ar[dd]_{(r_0, r_1, r_2) \mapsto (r_0 \circ r_1^{-1}, r_1 \circ r_2^{-1})} &
R \times_{s, U, s} R
\ar@<1ex>[r]_-{\text{pr}_1} \ar@<-2ex>[r]_-{\text{pr}_0}
\ar[dd]_{(r_0, r_1) \mapsto r_0 \circ r_1^{-1}} &
R \ar[dd]^t
\\
\\
R \times_{s, U, t} R
\ar@<3ex>[r]_{\text{pr}_1}
\ar@<0ex>[r]_c
\ar@<-3ex>[r]_{\text{pr}_0} &
R \ar@<1ex>[r]_s \ar@<-2ex>[r]_t &
U
}
$$
\end{lemma}

\begin{proof}
对任意 $n$，按规则
$$
(r_0, \ldots, r_n)
\longrightarrow
(r_0 \circ r_1^{-1}, \ldots, r_{n - 1} \circ r_n^{-1})
$$
定义 $(R/U)_\bullet \to X_n$。由引理
\ref{spaces-simplicial-lemma-groupoid-simplicial} 对退化映射的描述，
它与退化映射的相容性显然成立。
% P10-E0234
略去这些映射保持态射 $s^n_j$ 的验证。《群胚》引理
\ref{groupoids-lemma-diagram-pull}（交换 $s$ 与 $t$ 的角色）
表明右侧两个方块是笛卡尔的。完全类似地，可证所有其他方块
也都是笛卡尔的。因此，该态射是笛卡尔的。
\end{proof}




\section{下降数据给出等价关系}
\label{spaces-simplicial-section-equivalence-relation}

\noindent
在第 \ref{spaces-simplicial-section-simplicial-descent} 节中，我们看到如何用
$(X/S)_\bullet$ 上的笛卡尔单纯概形，表述相对于 $X \to S$
的下降数据。现在把它与等价关系联系起来。

\begin{lemma}
\label{spaces-simplicial-lemma-equivalence-relation}
设 $f : X \to S$ 是概形态射，
$\pi : Y \to (X/S)_\bullet$ 是单纯概形的笛卡尔态射；
参见定义 \ref{spaces-simplicial-definition-cartesian-morphism} 和
\ref{spaces-simplicial-definition-fibre-products-simplicial-scheme}。则态射
$$
j = (d^1_1, d^1_0) : Y_1 \to Y_0 \times_S Y_0
$$
定义 $S$ 上 $Y_0$ 的等价关系；参见《群胚》定义
\ref{groupoids-definition-equivalence-relation}。
\end{lemma}

\begin{proof}
注意，$j$ 是单态射。事实上，由于 $\pi$ 是笛卡尔的，
复合 $Y_1 \to Y_0 \times_S Y_0 \to Y_0 \times_S X$ 是同构。

\medskip\noindent
考虑态射
$$
(d^2_2, d^2_0) : Y_2 \to Y_1 \times_{d^1_0, Y_0, d^1_1} Y_1.
$$
这是可行的，因为 $d_0 \circ d_2 = d_1 \circ d_0$；
参见《单纯对象》评注 \ref{simplicial-remark-relations}。
此外，它是 $(X/S)_2$ 上的态射。由于
$Y \to (X/S)_\bullet$ 是笛卡尔的，它是同构。
例如，右边同构于
$Y_0 \times_{\pi_0, X, \text{pr}_1} (X \times_S X \times_S X) =
X \times_S Y_0 \times_S X$
，因为 $\pi$ 是笛卡尔的。略去细节。

\medskip\noindent
仿照《群胚》定义 \ref{groupoids-definition-equivalence-relation}，
记 $t = \text{pr}_0 \circ j = d^1_1$ 和
$s = \text{pr}_1 \circ j = d^1_0$。上述同构与态射
$d^2_1 : Y_2 \to Y_1$ 共同给出复合态射
$$
c : Y_1 \times_{s, Y_0, t} Y_1 \longrightarrow Y_1
$$
，它位于 $Y_0 \times_S Y_0$ 上方。这立即蕴含：对任意概形
$T/S$，关系 $Y_1(T) \subset Y_0(T) \times Y_0(T)$ 具有传递性。

\medskip\noindent
自反性来自如下事实：态射 $j$ 在对角
$\Delta : X \to X \times_S X$ 上的限制是同构
（再次使用 $\pi$ 的笛卡尔性质）。

\medskip\noindent
为证明对称性，考虑态射
$$
(d^2_2, d^2_1) : Y_2 \to Y_1 \times_{d^1_1, Y_0, d^1_1} Y_1.
$$
这是可行的，因为 $d_1 \circ d_2 = d_1 \circ d_1$；
参见《单纯对象》评注 \ref{simplicial-remark-relations}。
由于 $Y \to (X/S)_\bullet$ 是笛卡尔的，它是同构。
例如，右边同构于
$Y_0 \times_{\pi_0, X, \text{pr}_0} (X \times_S X \times_S X) =
Y_0 \times_S X \times_S X$
，因为 $\pi$ 是笛卡尔的。略去细节。

\medskip\noindent
设 $T/S$ 是概形。对 $a, b \in Y_0(T)$，规定 $a \sim b$
与 $(a, b) \in Y_1(T)$ 同义。上述同构 $(d^2_2, d^2_1)$
蕴含：若 $a \sim b$ 且 $a \sim c$，则 $b \sim c$。
结合自反性可知，$\sim$ 是等价关系。
\end{proof}







\section{一个示例情形}
\label{spaces-simplicial-section-example}

\noindent
本节证明，Artin 环的谱之不交并可以沿概形的拟紧平坦满射下降。

\begin{lemma}
\label{spaces-simplicial-lemma-equivalence-classes-points}
设 $X \to S$ 是概形态射，并假设
$Y \to (X/S)_\bullet$ 是单纯概形的笛卡尔态射。
对点 $y \in Y_0$，定义 $Y_0$ 的子集
$$
T_y = \{y' \in Y_0 \mid \exists\ y_1 \in Y_1:
d^1_1(y_1) = y, d^1_0(y_1) = y'\}
$$
。则 $y \in T_y$，且
$T_y \cap T_{y'} \not = \emptyset \Rightarrow T_y = T_{y'}$。
\end{lemma}

\begin{proof}
合用引理 \ref{spaces-simplicial-lemma-equivalence-relation} 和《群胚》引理
\ref{groupoids-lemma-pre-equivalence-equivalence-relation-points}。
\end{proof}

\begin{lemma}
\label{spaces-simplicial-lemma-quasi-compact}
设 $X \to S$ 是概形态射，并假设
$Y \to (X/S)_\bullet$ 是单纯概形的笛卡尔态射。
设 $y \in Y_0$ 是点。若 $X \to S$ 拟紧，则
$$
T_y = \{y' \in Y_0 \mid \exists\ y_1 \in Y_1:
d^1_1(y_1) = y, d^1_0(y_1) = y'\}
$$
是 $Y_0$ 的拟紧子集。
\end{lemma}

\begin{proof}
设 $F_y$ 是 $d^1_1 : Y_1 \to Y_0$ 在 $y$ 处的概形论纤维。
则 $T_y$ 是如下态射的像：
$$
\xymatrix{
F_y \ar[r] \ar[d] &
Y_1 \ar[r]^{d^1_0} \ar[d]^{d^1_1} &
Y_0 \\
y \ar[r] &
Y_0 &
}
$$
注意，$F_y$ 拟紧。这证明了引理。
\end{proof}

\begin{lemma}
\label{spaces-simplicial-lemma-descent-disjoint-union-Artinian-along-fields}
设 $X \to S$ 是拟紧平坦满射，$(V, \varphi)$ 是相对于
$X \to S$ 的下降数据。若 $V$ 是 Artin 环的谱之不交并，
则 $(V, \varphi)$ 有效。
\end{lemma}

\begin{proof}
由引理 \ref{spaces-simplicial-lemma-cartesian-equivalent-descent-datum}，
设 $Y \to (X/S)_\bullet$ 为对应于 $(V, \varphi)$ 的单纯概形之
笛卡尔态射。注意，$Y_0 = V$。写
$V = \coprod_{i \in I} \Spec(A_i)$，其中每个 $A_i$ 是局部 Artin 环。
此外，对所有 $i \in I$，设 $v_i \in V$ 为 $\Spec(A_i)$
的唯一闭点。采用上面引理
\ref{spaces-simplicial-lemma-equivalence-classes-points} 的记号，规定
$i \sim j$ 当且仅当 $v_i \in T_{v_j}$。由引理
\ref{spaces-simplicial-lemma-equivalence-classes-points} 和 \ref{spaces-simplicial-lemma-quasi-compact}，
这是等价类均有限的等价关系。令 $\overline{I} = I/\sim$。
于是可以写成
$V = \coprod_{\overline{i} \in \overline{I}} V_{\overline{i}}$
，其中
$V_{\overline{i}} = \coprod_{i \in \overline{i}} \Spec(A_i)$.
由构造，$\varphi : V \times_S X \to X \times_S V$
把开闭子空间 $V_{\overline{i}} \times_S X$ 映入开闭子空间
$X \times_S V_{\overline{i}}$。换言之，得到下降数据
$(V_{\overline{i}}, \varphi_{\overline{i}})$，而 $(V, \varphi)$
是它们在下降数据范畴中的余积。由于每个 $V_{\overline{i}}$
都是局部 Artin 环的谱的有限并，态射 $V_{\overline{i}} \to X$
是仿射的；参见《态射》引理
\ref{morphisms-lemma-Artinian-affine}。由于 $\{X \to S\}$
是 fpqc 覆盖，由《下降》引理 \ref{descent-lemma-affine}，
所有下降数据 $(V_{\overline{i}}, \varphi_{\overline{i}})$ 都有效。
\end{proof}

\noindent
诚然，上述引理的适用范围十分有限！









\section{单纯代数空间}
\label{spaces-simplicial-section-simplicial-algebraic-spaces}

\noindent
设 $S$ 是概形。{\it 单纯代数空间}是 $S$ 上代数空间范畴中的单纯对象；
参见《单纯对象》定义 \ref{simplicial-definition-simplicial-object}。
回忆，单纯代数空间形如
$$
\xymatrix{
X_2
\ar@<2ex>[r]
\ar@<0ex>[r]
\ar@<-2ex>[r]
&
X_1
\ar@<1ex>[r]
\ar@<-1ex>[r]
\ar@<1ex>[l]
\ar@<-1ex>[l]
&
X_0
\ar@<0ex>[l]
}
$$
这里有两个态射 $d^1_0, d^1_1 : X_1 \to X_0$ 和一个态射
$s^0_0 : X_0 \to X_1$，依此类推。这些态射满足一些必需关系，
例如 $d^1_0 \circ s^0_0 = \text{id}_{X_0} = d^1_1 \circ s^0_0$；
参见《单纯对象》引理
\ref{simplicial-lemma-characterize-simplicial-object}。
把 $d^n_i : X_n \to X_{n - 1}$ 看作“忘掉第 $i$ 个坐标的投影”，
把 $s^n_j : X_n \to X_{n + 1}$ 看作“重复第 $j$ 个坐标的对角映射”，
会很有帮助。

\medskip\noindent
{\it 单纯代数空间的态射} $h : X \to Y$，就是 $S$ 上代数空间范畴中
单纯对象的态射；参见《单纯对象》定义
\ref{simplicial-definition-simplicial-object}。因此，$h$
由代数空间态射 $h_n : X_n \to Y_n$ 组成，并且只要等式有意义，
就有 $h_{n - 1} \circ d^n_j = d^n_j \circ h_n$ 和
$h_{n + 1} \circ s^n_j = s^n_j \circ h_n$。

\medskip\noindent
单纯代数空间 $X$ 的{\it 增广} $a : X \to X_{-1}$，
由满足 $a_0 \circ d^1_0 = a_0 \circ d^1_1$ 的代数空间态射
$a_0 : X_0 \to X_{-1}$ 给出。参见《单纯对象》第
\ref{simplicial-section-augmentation} 节。在这种情形下，
对 $n \geq 0$，总以 $a_n : X_n \to X_{-1}$ 表示所诱导的态射。

\medskip\noindent
设 $X$ 是单纯代数空间。对每个 $n$，有站点
$X_{n, spaces, \etale}$（《空间的性质》定义
\ref{spaces-properties-definition-spaces-etale-site}）；
对每个态射 $\varphi : [m] \to [n]$，有站点态射
$$
f_\varphi = X(\varphi)_{spaces, \etale} :
X_{n, spaces, \etale} \to X_{m, spaces, \etale},
$$
，它与代数空间态射 $X(\varphi) : X_n \to X_m$ 关联
（《空间的性质》引理
\ref{spaces-properties-lemma-functoriality-etale-site}）。
这给出站点范畴中的单纯对象。在引理
\ref{spaces-simplicial-lemma-simplicial-site-site} 中，我们构造了关联站点，
记为 $X_{spaces, \etale}$。站点 $X_{spaces, \etale}$ 的对象，
% P10-E0235
是对某个 $n$ 在 $X_n$ 上平展的代数空间 $U$；态射
$(\varphi, f) : U/X_n \to V/X_m$ 由 $\Delta$ 中的态射
$\varphi : [m] \to [n]$ 和代数空间态射 $f : U \to V$ 给出，
并要求下图交换：
$$
\xymatrix{
U \ar[r]_f \ar[d] & V \ar[d] \\
X_n \ar[r]^{f_\varphi} & X_m
}
$$
考虑全子范畴
$$
X_{affine, \etale} \subset X_\etale \subset X_{spaces, \etale}
$$
，其对象是 $U/X_n$，其中 $U$ 分别为仿射的或为概形。
给这些范畴赋予其自然拓扑（参见《空间的性质》引理
\ref{spaces-properties-lemma-alternative}、定义
\ref{spaces-properties-definition-etale-site} 和引理
\ref{spaces-properties-lemma-compare-etale-sites}），
这些包含函子定义拓扑斯等价
$$
\Sh(X_{affine, \etale}) = \Sh(X_\etale) = \Sh(X_{spaces, \etale})
$$
下文默认识别这些拓扑斯。称 $X_\etale$ 为
{\it $X$ 的小平展站点}，称其拓扑斯为{\it $X$ 的小平展拓扑斯}。

\medskip\noindent
设 $X_\etale$ 是单纯代数空间 $X$ 的小平展站点。
$X_\etale$ 上存在环层 $\mathcal{O}$，它在 $X_n$ 上的限制
是结构层 $\mathcal{O}_{X_n}$。这由引理
\ref{spaces-simplicial-lemma-describe-sheaves-simplicial-site-site} 得出。
我们称{\it $\mathcal{O}$ 是单纯代数空间 $X$ 的结构层}。
至此，为单纯（赋环）站点发展的一切材料都适用；参见第
\ref{spaces-simplicial-section-simplicial-sites}、
\ref{spaces-simplicial-section-augmentation-simplicial-sites}、
\ref{spaces-simplicial-section-morphism-simplicial-sites}、
\ref{spaces-simplicial-section-simplicial-sites-modules}、
\ref{spaces-simplicial-section-cohomology-simplicial-sites}、
\ref{spaces-simplicial-section-cohomology-augmentation-simplicial-sites}、
\ref{spaces-simplicial-section-cohomology-simplicial-sites-modules}、
\ref{spaces-simplicial-section-cohomology-augmentation-ringed-simplicial-sites}、
\ref{spaces-simplicial-section-cartesian}、
\ref{spaces-simplicial-section-glueing} 和
\ref{spaces-simplicial-section-glueing-modules} 节。

\medskip\noindent
设 $X$ 是结构层为 $\mathcal{O}$ 的单纯代数空间。
与任何赋环拓扑斯上一样，可以定义
{\it $X_\etale$ 上的拟凝聚 $\mathcal{O}$-模}；参见
《站点上的模》定义 \ref{sites-modules-definition-site-local}。
不过，$X_\etale$ 上的拟凝聚 $\mathcal{O}$-模，恰是如下笛卡尔
$\mathcal{O}$-模 $\mathcal{F}$：其限制 $\mathcal{F}_n$
在 $X_n$ 上拟凝聚；参见引理 \ref{spaces-simplicial-lemma-quasi-coherent-sheaf}。

\medskip\noindent
设 $h : X \to Y$ 是 $S$ 上单纯代数空间的态射。将引理
\ref{spaces-simplicial-lemma-morphism-simplicial-sites} 应用于站点态射
% P10-E0237
$(h_n)_{spaces, \etale} : X_{spaces, \etale} \to Y_{spaces, \etale}$
（《空间的性质》引理
\ref{spaces-properties-lemma-functoriality-etale-site}），
得到小平展拓扑斯的态射
$h_\etale : \Sh(X_\etale) \to \Sh(Y_\etale)$。
回忆，$h_\etale^{-1}$ 和 $h_{\etale, *}$ 可用各分量简单描述；
参见引理 \ref{spaces-simplicial-lemma-morphism-simplicial-sites}。
以 $\mathcal{O}_X$、$\mathcal{O}_Y$ 分别表示 $X$、$Y$ 的结构层。
定义
% P10-E0229
$h_\etale^\sharp : h_{\etale, *}\mathcal{O}_X \to \mathcal{O}_Y$
为 $Y_\etale$ 上的环层映射，它在 $Y_n$ 上由
$h_n^\sharp : h_{n, *}\mathcal{O}_{X_n} \to \mathcal{O}_{Y_n}$ 给出。
由此得到赋环拓扑斯的态射
$$
h_\etale :
(\Sh(X_\etale), \mathcal{O}_X)
\longrightarrow
(\Sh(Y_\etale), \mathcal{O}_Y)
$$

\medskip\noindent
设 $X$ 是结构层为 $\mathcal{O}$ 的单纯代数空间。
设 $X_{-1}$ 是 $S$ 上的代数空间，$a_0 : X_0 \to X_{-1}$
是 $X$ 的增广。将引理 \ref{spaces-simplicial-lemma-augmentation-site}
应用于站点态射
$(a_0)_{spaces, \etale} :
X_{0, spaces, \etale} \to X_{-1, spaces, \etale}$
，得到相应拓扑斯态射
$a : \Sh(X_\etale) \to \Sh(X_{-1, \etale})$。
注意，$a^{-1}\mathcal{G}$ 是 $X_\etale$ 上分量为
$a_n^{-1}\mathcal{G}$ 的层。因此，可以利用映射
$a_n^\sharp : a_n^{-1}\mathcal{O}_{X_{-1}} \to \mathcal{O}_{X_n}$
定义映射 $a^\sharp : a^{-1}\mathcal{O}_{X_{-1}} \to \mathcal{O}$；
等价地，由伴随定义映射
$a^\sharp : \mathcal{O}_{X_{-1}} \to a_*\mathcal{O}$
（照例仍用同一名称）。这正处于第
\ref{spaces-simplicial-section-cohomology-augmentation-ringed-simplicial-sites}
节所讨论的情形。因此，得到赋环拓扑斯的态射
$$
a :
(\Sh(X_\etale), \mathcal{O})
\longrightarrow
(\Sh(X_{-1}), \mathcal{O}_{X_{-1}})
$$

\medskip\noindent
最后作如下观察。设给定单纯代数空间 $X$ 与 $Y$ 的态射
$h : X \to Y$，其结构层分别为 $\mathcal{O}_X$、$\mathcal{O}_Y$；
并给定增广 $a_0 : X_0 \to X_{-1}$、$b_0 : Y_0 \to Y_{-1}$
和态射 $h_{-1} : X_{-1} \to Y_{-1}$，使得
$$
\xymatrix{
X_0 \ar[r]_{h_0} \ar[d]_{a_0} & Y_0 \ar[d]^{b_0} \\
X_{-1} \ar[r]^{h_{-1}} & Y_{-1}
}
$$
交换。由上面阐明的构造，得到如下赋环拓扑斯态射的交换图：
% P10-E0236
$$
\xymatrix{
(\Sh(X_\etale), \mathcal{O}_X) \ar[r]_{h_\etale} \ar[d]_a &
(\Sh(Y_\etale), \mathcal{O}_Y) \ar[d]^b \\
(\Sh(X_{-1}), \mathcal{O}_{X_{-1}}) \ar[r]^{h_{-1}} &
(\Sh(Y_{-1}), \mathcal{O}_{Y_{-1}})
}
$$




\section{代数空间的 fppf 超覆盖}
\label{spaces-simplicial-section-fppf-hypercovering}

\noindent
本节是第 \ref{spaces-simplicial-section-proper-hypercovering} 节在代数空间和 fppf
超覆盖情形下的类似版本。有意如此做的读者，可以处处把“代数空间”
替换为“概形”，得到同样有效的结果。这样做的优点是，可以把对
《空间的上同调进阶》第
\ref{spaces-more-cohomology-section-fppf-etale} 节的引用，
替换为对《平展上同调》第
\ref{etale-cohomology-section-fppf-etale} 节的引用。

\medskip\noindent
固定基概形 $S$。设 $X$ 是 $S$ 上的代数空间，
$U$ 是 $S$ 上的单纯代数空间。假设有增广
$$
a : U \to X
$$
参见第 \ref{spaces-simplicial-section-simplicial-algebraic-spaces} 节。
如果满足下列条件，就称 $U$ 是 $X$ 的{\it fppf 超覆盖}：
\begin{enumerate}
\item $U_0 \to X$ 平坦、局部有限呈示且满射；
\item $U_1 \to U_0 \times_X U_0$ 平坦、局部有限呈示且满射；
\item $U_{n + 1} \to (\text{cosk}_n\text{sk}_n U)_{n + 1}$
对 $n \geq 1$ 平坦、局部有限呈示且满射。
\end{enumerate}
$S$ 上的代数空间范畴具有所有有限极限，所以上述表述中使用的余骨架存在。
$$
\fbox{原理：可以用 fppf 超覆盖计算平展上同调。}
$$
证明该原理的关键思想，是比较范畴 $\textit{Spaces}/S$ 上的 fppf
拓扑与平展拓扑。具体地，fppf 拓扑比平展拓扑更细，并且：
(a) 平坦、局部有限呈示的满射定义 fppf 覆盖；
(b) 从小平展站点拉回的层之 fppf 上同调与平展上同调一致，
正如《空间的上同调进阶》第
\ref{spaces-more-cohomology-section-fppf-etale} 节所述。

\begin{lemma}
\label{spaces-simplicial-lemma-compare-simplicial-objects-fppf-etale}
设 $S$ 是概形，$X$ 是 $S$ 上的代数空间，
$U$ 是 $S$ 上的单纯代数空间，$a : U \to X$ 是增广。
则有交换图
$$
\xymatrix{
\Sh((\textit{Spaces}/U)_{fppf, total}) \ar[r]_-h \ar[d]_{a_{fppf}} &
\Sh(U_\etale) \ar[d]^a \\
\Sh((\textit{Spaces}/X)_{fppf}) \ar[r]^-{h_{-1}} &
\Sh(X_\etale)
}
$$
其中左侧竖直箭头定义于第 \ref{spaces-simplicial-section-hypercovering} 节，
右侧竖直箭头定义于第
\ref{spaces-simplicial-section-simplicial-algebraic-spaces} 节。
\end{lemma}

\begin{proof}
记号 $(\textit{Spaces}/U)_{fppf, total}$ 表示，我们对站点
$(\textit{Spaces}/S)_{fppf}$ 及其单纯对象 $U$ 使用第
\ref{spaces-simplicial-section-hypercovering} 节的构造\footnote{也可以使用平展拓扑，
此时记为 $(\textit{Spaces}/U)_{\etale, total}$。}。
对右边的拓扑斯，我们将使用站点 $X_{spaces, \etale}$ 和
$U_{spaces, \etale}$；这是允许的，见第
\ref{spaces-simplicial-section-simplicial-algebraic-spaces} 节中的讨论。

\medskip\noindent
注意，$(\textit{Spaces}/U)_{fppf, total}$ 与
$U_{spaces, \etale}$ 都属于情形
\ref{spaces-simplicial-situation-simplicial-site} 的 A 类。对 $U_\etale$，
这立即由第 \ref{spaces-simplicial-section-simplicial-algebraic-spaces} 节中的构造得出；
对 $(\textit{Spaces}/U)_{fppf, total}$，则由引理
\ref{spaces-simplicial-lemma-sr-when-fibre-products} 得出。接着考虑函子
$U_{n, spaces, \etale} \to (\textit{Spaces}/U_n)_{fppf}$, $U \mapsto U/U_n$
以及
$X_{spaces, \etale} \to (\textit{Spaces}/X)_{fppf}$, $U \mapsto U/X$.
我们已在《空间的上同调进阶》第
\ref{spaces-more-cohomology-section-fppf-etale} 节看到，
它们定义站点态射；在那里分别记为
$a_{U_n} = \epsilon_{U_n} \circ \pi_{u_n}$ 和
$a_X = \epsilon_X \circ \pi_X$。因此，得到评注
\ref{spaces-simplicial-remark-morphism-augmentation-simplicial-sites}
意义下与增广相容的单纯站点态射；应用引理
\ref{spaces-simplicial-lemma-morphism-augmentation-simplicial-sites} 即得结论。
\end{proof}

\begin{lemma}
\label{spaces-simplicial-lemma-descent-sheaves-for-fppf-hypercovering}
设 $S$ 是概形，$X$ 是 $S$ 上的代数空间，
$U$ 是 $S$ 上的单纯代数空间，$a : U \to X$ 是增广。
若 $a : U \to X$ 是 $X$ 的 fppf 超覆盖，则
$$
a^{-1} : \Sh(X_\etale) \to \Sh(U_\etale)
\quad\text{且}\quad
a^{-1} : \textit{Ab}(X_\etale) \to \textit{Ab}(U_\etale)
$$
满忠实，其本质像为笛卡尔层，拟逆由 $a_*$ 给出。
这里 $a : \Sh(U_\etale) \to \Sh(X_\etale)$ 如第
\ref{spaces-simplicial-section-simplicial-algebraic-spaces} 节所述。
\end{lemma}

\begin{proof}
我们证明集合层的陈述。它几乎是已建立结果的形式推论。
考虑引理 \ref{spaces-simplicial-lemma-compare-simplicial-objects-fppf-etale} 的图。
该引理的证明表明，$h_{-1}$ 是《空间的上同调进阶》第
\ref{spaces-more-cohomology-section-fppf-etale} 节中的态射 $a_X$。
因此，由《空间的上同调进阶》引理
\ref{spaces-more-cohomology-lemma-comparison-fppf-etale}，
$(h_{-1})^{-1}$ 满忠实，其拟逆为 $h_{-1, *}$。
对 $h$ 的分量 $h_n$ 也同样成立。由引理
\ref{spaces-simplicial-lemma-morphism-simplicial-sites} 对函子 $h^{-1}$ 和 $h_*$
的描述，可知 $h^{-1}$ 满忠实，其拟逆为 $h_*$。
$U$ 是 $(\textit{Spaces}/S)_{fppf}$ 中 $X$ 的超覆盖，
其定义见第 \ref{spaces-simplicial-section-hypercovering} 节。
由引理 \ref{spaces-simplicial-lemma-hypercovering-X-simple-descent-sheaves}，
$a_{fppf}^{-1}$ 满忠实，其拟逆为 $a_{fppf, *}$，本质像为
$(\textit{Spaces}/U)_{fppf, total}$ 上的笛卡尔层。
现在，一个形式论证（沿图追逐）表明 $a^{-1}$ 满忠实。

\medskip\noindent
最后，假设 $\mathcal{G}$ 是 $U_\etale$ 上的笛卡尔层。
则 $h^{-1}\mathcal{G}$ 是 $(\textit{Spaces}/U)_{fppf, total}$
上的笛卡尔层。因此，对 $(\textit{Spaces}/X)_{fppf}$ 上的某个层
$\mathcal{H}$，有
$h^{-1}\mathcal{G} = a_{fppf}^{-1}\mathcal{H}$。
特别地，$h_0^{-1}\mathcal{G}_0 =
(a_{0, big, fppf})^{-1}\mathcal{H}$。回忆 $h_0 = a_{U_0}$，
且 $U_0 \to X$ 平坦、局部有限呈示并且满射；由
《空间的上同调进阶》引理
\ref{spaces-more-cohomology-lemma-descent-sheaf-fppf-etale}，
存在 $X_\etale$ 上的层 $\mathcal{F}$ 及同构
$\mathcal{H} = (h_{-1})^{-1}\mathcal{F}$。
由于 $a_{fppf}^{-1}\mathcal{H} = h^{-1}\mathcal{G}$，可得
$h^{-1}\mathcal{G} \cong h^{-1}a^{-1}\mathcal{F}$。
由 $h^{-1}$ 的满忠实性，得到
$a^{-1}\mathcal{F} \cong \mathcal{G}$。

\medskip\noindent
固定同构 $\theta : a^{-1}\mathcal{F} \to \mathcal{G}$。
为完成证明，须证明 $\mathcal{G} = a^{-1}a_*\mathcal{G}$
（目的是证明拟逆由 $a_*$ 给出；其余内容均已在上面证明）。
由于 $a^{-1}$ 满忠实，由《范畴》引理
\ref{categories-lemma-adjoint-fully-faithful}，有
$\text{id} \cong a_*a^{-1}$。因此，
$\mathcal{F} \cong a_*a^{-1}\mathcal{F}$ 与
$a_*\theta : a_*a^{-1}\mathcal{F} \to a_*\mathcal{G}$
合成同构 $\mathcal{F} \to a_*\mathcal{G}$。
将其沿 $a$ 拉回并与 $\theta^{-1}$ 预复合，便得到所需同构。
\end{proof}

\begin{lemma}
\label{spaces-simplicial-lemma-cohomological-descent-for-fppf-hypercovering}
设 $S$ 是概形，$X$ 是 $S$ 上的代数空间，
$U$ 是 $S$ 上的单纯代数空间，$a : U \to X$ 是增广。
若 $a : U \to X$ 是 $X$ 的 fppf 超覆盖，则对
$K \in D^+(X_\etale)$，
$$
K \to Ra_*(a^{-1}K)
$$
是同构。这里 $a : \Sh(U_\etale) \to \Sh(X_\etale)$ 如第
\ref{spaces-simplicial-section-simplicial-algebraic-spaces} 节所述。
\end{lemma}

\begin{proof}
考虑引理 \ref{spaces-simplicial-lemma-compare-simplicial-objects-fppf-etale} 的图。
由《空间的上同调进阶》引理
\ref{spaces-more-cohomology-lemma-cohomological-descent-etale-fppf}，
$Rh_{n, *}h_n^{-1}$ 是 $D^+(U_{n, \etale})$ 上的恒同函子。
因此，由引理 \ref{spaces-simplicial-lemma-direct-image-morphism-simplicial-sites}，
$Rh_*h^{-1}$ 是 $D^+(U_\etale)$ 上的恒同函子。我们有
\begin{align*}
Ra_*(a^{-1}K)
& =
Ra_*Rh_*h^{-1}a^{-1}K \\
& =
Rh_{-1, *}Ra_{fppf, *}a_{fppf}^{-1}(h_{-1})^{-1}K \\
& =
Rh_{-1, *}(h_{-1})^{-1}K \\
& =
K
\end{align*}
% P10-E0228
第一个等号来自上述讨论；第二个等号来自
% P10-E0238
引理 \ref{spaces-simplicial-lemma-compare-simplicial-objects} 中图的交换性；
第三个等号来自引理
\ref{spaces-simplicial-lemma-hypercovering-X-simple-descent-bounded-abelian}，
因为 $U$ 是 $(\textit{Spaces}/S)_{fppf}$ 中 $X$ 的超覆盖；
最后一个等号来自已经使用过的《空间的上同调进阶》引理
\ref{spaces-more-cohomology-lemma-cohomological-descent-etale-fppf}。
\end{proof}

\begin{lemma}
\label{spaces-simplicial-lemma-compute-via-fppf-hypercovering}
设 $S$ 是概形，$X$ 是 $S$ 上的代数空间，
$U$ 是 $S$ 上的单纯代数空间，$a : U \to X$ 是增广。
若 $a : U \to X$ 是 $X$ 的 fppf 超覆盖，则
$$
R\Gamma(X_\etale, K) = R\Gamma(U_\etale, a^{-1}K)
$$
其中 $K \in D^+(X_\etale)$。这里
$a : \Sh(U_\etale) \to \Sh(X_\etale)$ 如第
\ref{spaces-simplicial-section-simplicial-algebraic-spaces} 节所述。
\end{lemma}

\begin{proof}
这由引理 \ref{spaces-simplicial-lemma-cohomological-descent-for-fppf-hypercovering} 得出，
因为按《站点上的上同调》评注
\ref{sites-cohomology-remark-before-Leray}，有
$R\Gamma(U_\etale, -) = R\Gamma(X_\etale, -) \circ Ra_*$。
\end{proof}

\begin{lemma}
\label{spaces-simplicial-lemma-fppf-hypercovering-equivalence-bounded}
设 $S$ 是概形，$X$ 是 $S$ 上的代数空间，
$U$ 是 $S$ 上的单纯代数空间，$a : U \to X$ 是增广。
以 $\mathcal{A} \subset \textit{Ab}(U_\etale)$
表示笛卡尔阿贝尔层组成的弱 Serre 子范畴。
若 $U$ 是 $X$ 的 fppf 超覆盖，则函子 $a^{-1}$ 给出等价
$$
D^+(X_\etale) \longrightarrow D_\mathcal{A}^+(U_\etale)
$$
，其拟逆为 $Ra_*$。这里 $a : \Sh(U_\etale) \to \Sh(X_\etale)$
如第 \ref{spaces-simplicial-section-simplicial-algebraic-spaces} 节所述。
\end{lemma}

\begin{proof}
由引理 \ref{spaces-simplicial-lemma-Serre-subcat-cartesian-modules}，
$\mathcal{A}$ 是弱 Serre 子范畴。该等价是迄今所得结果的形式推论。
应用引理 \ref{spaces-simplicial-lemma-descent-sheaves-for-fppf-hypercovering}、
\ref{spaces-simplicial-lemma-cohomological-descent-for-fppf-hypercovering} 以及
《站点上的上同调》引理 \ref{sites-cohomology-lemma-equivalence-bounded}。
\end{proof}

\begin{lemma}
\label{spaces-simplicial-lemma-spectral-sequence-fppf-hypercovering}
设 $S$ 是概形，$X$ 是 $S$ 上的代数空间，
$U$ 是 $S$ 上的单纯代数空间，$a : U \to X$ 是增广。
设 $\mathcal{F}$ 是 $X_\etale$ 上的阿贝尔层，$\mathcal{F}_n$
是它到 $U_{n, \etale}$ 的拉回。若 $U$ 是 $X$ 的 fppf 超覆盖，
则存在典范谱序列
$$
E_1^{p, q} = H^q_\etale(U_p, \mathcal{F}_p)
$$
，它收敛到 $H^{p + q}_\etale(X, \mathcal{F})$。
\end{lemma}

\begin{proof}
这是引理 \ref{spaces-simplicial-lemma-compute-via-fppf-hypercovering} 和
\ref{spaces-simplicial-lemma-simplicial-sheaf-cohomology-site} 的立即推论。
\end{proof}








\section{代数空间的 fppf 超覆盖：模}
\label{spaces-simplicial-section-fppf-hypercovering-modules}

\noindent
我们继续第 \ref{spaces-simplicial-section-fppf-hypercovering} 节所开始的关于 fppf
超覆盖的（上同调）下降之讨论，但在本节中我们讨论模层的情形。
我们主要讨论拟凝聚模；事实证明，对它们可以进行无界上同调下降。

\begin{lemma}
\label{spaces-simplicial-lemma-compare-simplicial-objects-fppf-etale-modules}
设 $S$ 是概形。设 $X$ 是 $S$ 上的代数空间。
设 $U$ 是 $S$ 上的单纯代数空间。
设 $a : U \to X$ 是增广。存在如下赋环拓扑斯交换图
$$
\xymatrix{
(\Sh((\textit{Spaces}/U)_{fppf, total}), \mathcal{O}_{big, total})
\ar[r]_-h \ar[d]_{a_{fppf}} &
(\Sh(U_\etale), \mathcal{O}_U) \ar[d]^a \\
(\Sh((\textit{Spaces}/X)_{fppf}), \mathcal{O}_{big}) \ar[r]^-{h_{-1}} &
(\Sh(X_\etale), \mathcal{O}_X)
}
$$
其中左侧竖直箭头定义于第 \ref{spaces-simplicial-section-hypercovering-modules} 节，
右侧竖直箭头定义于第 \ref{spaces-simplicial-section-simplicial-algebraic-spaces} 节。
\end{lemma}

\begin{proof}
关于底层拓扑斯图，参见引理
\ref{spaces-simplicial-lemma-compare-simplicial-objects-fppf-etale} 的证明中的讨论。
层 $\mathcal{O}_U$ 是第 \ref{spaces-simplicial-section-simplicial-algebraic-spaces} 节
所定义的单纯代数空间 $U$ 的结构层。层 $\mathcal{O}_X$
是代数空间 $X$ 的通常结构层。环层 $\mathcal{O}_{big, total}$
和 $\mathcal{O}_{big}$ 以第 \ref{spaces-simplicial-section-hypercovering-modules} 节
所说明的方式来自 $(\textit{Spaces}/S)_{fppf}$ 上的结构层；
该节还将 $a_{fppf}$ 构造为赋环拓扑斯的态射。
由空间上上同调进阶的第
\ref{spaces-more-cohomology-section-fppf-etale-modules} 节，
分量态射 $h_n = a_{U_n}$ 和 $h_{-1} = a_X$
是赋环拓扑斯的态射。最后，用于定义 $h$ 的连续函子
$u : U_{spaces, \etale} \to (\textit{Spaces}/U)_{fppf, total}$
\footnote{这是在引理
\ref{spaces-simplicial-lemma-compare-simplicial-objects-fppf-etale} 的证明中，
通过应用引理 \ref{spaces-simplicial-lemma-morphism-augmentation-simplicial-sites}
完成的。}
由 $V/U_n \mapsto V/U_n$ 给出，故
$h_*\mathcal{O}_{big, total} = \mathcal{O}_U$。
我们据此按注 \ref{spaces-simplicial-remark-morphism-simplicial-sites-modules}
赋予 $h$ 赋环单纯站点之态射的结构。
于是由引理 \ref{spaces-simplicial-lemma-morphism-simplicial-sites-modules}，
得到作为赋环拓扑斯态射的 $h$。请注意，态射 $h_n$
确实与上面所述的态射 $a_{U_n}$ 一致。
我们略去该图作为赋环拓扑斯之图的交换性验证
（我们已经知道它作为拓扑斯之图是交换的）。
\end{proof}

\begin{lemma}
\label{spaces-simplicial-lemma-descent-qcoh-for-fppf-hypercovering}
设 $S$ 是概形，$X$ 是 $S$ 上的代数空间。
设 $U$ 是 $S$ 上的单纯代数空间，$a : U \to X$ 是增广。
若 $a : U \to X$ 是 $X$ 的 fppf 超覆盖，则
$$
a^* : \QCoh(\mathcal{O}_X) \to \QCoh(\mathcal{O}_U)
$$
% P10-E0239
是一个全忠实的等价，其拟逆由 $a_*$ 给出。
这里 $a : \Sh(U_\etale) \to \Sh(X_\etale)$
如第 \ref{spaces-simplicial-section-simplicial-algebraic-spaces} 节所述。
\end{lemma}

\begin{proof}
考虑引理
\ref{spaces-simplicial-lemma-compare-simplicial-objects-fppf-etale-modules} 中的图。
在该引理的证明中，我们已看到 $h_{-1}$ 是空间上上同调进阶的
第 \ref{spaces-more-cohomology-section-fppf-etale-modules} 节
中的态射 $a_X$。因此，由空间上上同调进阶的引理
\ref{spaces-more-cohomology-lemma-review-quasi-coherent}，
$$
(h_{-1})^* :
\QCoh(\mathcal{O}_X)
\longrightarrow
\QCoh(\mathcal{O}_{big})
$$
是以 $h_{-1, *}$ 为拟逆的等价。对 $h$ 的分量 $h_n$
同样成立。回忆 $\QCoh(\mathcal{O}_U)$ 和
$\QCoh(\mathcal{O}_{big, total})$ 由各分量拟凝聚的笛卡尔模组成，
见引理 \ref{spaces-simplicial-lemma-quasi-coherent-sheaf}。
由于引理 \ref{spaces-simplicial-lemma-morphism-simplicial-sites-modules} 中的函子
$h^*$ 和 $h_*$ 在各分量上与函子 $h_n^*$ 和 $h_{n, *}$ 一致，
我们得到
$$
h^* :
\QCoh(\mathcal{O}_U)
\longrightarrow 
\QCoh(\mathcal{O}_{big, total})
$$
是以 $h_*$ 为拟逆的等价。注意，按第
\ref{spaces-simplicial-section-hypercovering} 节的定义，$U$ 是
$(\textit{Spaces}/S)_{fppf}$ 中 $X$ 的超覆盖。
由引理 \ref{spaces-simplicial-lemma-hypercovering-X-simple-descent-modules}，
$a_{fppf}^*$ 是全忠实的，拟逆为 $a_{fppf, *}$，
且其本质像为 $\mathcal{O}_{fppf, total}$-模的笛卡尔层。
因此，由引理 \ref{spaces-simplicial-lemma-quasi-coherent-sheaf} 对
$\QCoh(\mathcal{O}_{big})$ 和
$\QCoh(\mathcal{O}_{big, total})$ 的描述，我们得到等价
$$
a_{fppf}^* :
\QCoh(\mathcal{O}_{big})
\longrightarrow
\QCoh(\mathcal{O}_{big, total})
$$
，其拟逆由 $a_{fppf, *}$ 给出。现在，一个形式论证
（沿图追逐）表明，$a^*$ 在 $\QCoh(\mathcal{O}_X)$ 上全忠实，
且其像包含在 $\QCoh(\mathcal{O}_U)$ 中。

\medskip\noindent
最后，设 $\mathcal{G}$ 属于 $\QCoh(\mathcal{O}_U)$。
则 $h^*\mathcal{G}$ 属于 $\QCoh(\mathcal{O}_{big, total})$。
因此 $h^*\mathcal{G} = a_{fppf}^*\mathcal{H}$，其中
$\mathcal{H} = a_{fppf, *}h^*\mathcal{G}$
属于 $\QCoh(\mathcal{O}_{big})$（见上文）。
进而，$\mathcal{H} = (h_{-1})^*\mathcal{F}$，其中
$\mathcal{F} = h_{-1, *}\mathcal{H}$ 属于
$\QCoh(\mathcal{O}_X)$。沿该图追逐可得
$h^*\mathcal{G} \cong h^*a^*\mathcal{F}$。
由 $h^*$ 的全忠实性，我们得到
$a^*\mathcal{F} \cong \mathcal{G}$。
又因为 $\mathcal{F} = h_{-1, *}a_{fppf, *}h^*\mathcal{G} =
a_*h_*h^*\mathcal{G} = a_*\mathcal{G}$，
所以也得到拟逆由 $a_*$ 给出的断言。
\end{proof}

\begin{lemma}
\label{spaces-simplicial-lemma-cohomological-descent-qcoh-for-fppf-hypercovering}
设 $S$ 是概形，$X$ 是 $S$ 上的代数空间。
设 $U$ 是 $S$ 上的单纯代数空间，$a : U \to X$ 是增广。
若 $a : U \to X$ 是 $X$ 的 fppf 超覆盖，则对任意拟凝聚
$\mathcal{O}_X$-模 $\mathcal{F}$，映射
$$
\mathcal{F} \to Ra_*(a^*\mathcal{F})
$$
是同构。这里 $a : \Sh(U_\etale) \to \Sh(X_\etale)$
如第 \ref{spaces-simplicial-section-simplicial-algebraic-spaces} 节所述。
\end{lemma}

\begin{proof}
% P10-E0238
考虑引理 \ref{spaces-simplicial-lemma-compare-simplicial-objects-fppf-etale} 中的图。
令 $\mathcal{F}_n = a_n^*\mathcal{F}$ 为
$a^*\mathcal{F}$ 的第 $n$ 个分量。它是拟凝聚
$\mathcal{O}_{U_n}$-模。由空间上上同调进阶的引理
\ref{spaces-more-cohomology-lemma-cohomological-descent-etale-fppf-modules}，
有 $\mathcal{F}_n = Rh_{n, *}h_n^*\mathcal{F}_n$。
因而由引理 \ref{spaces-simplicial-lemma-direct-image-morphism-simplicial-sites-modules}，
$a^*\mathcal{F} = Rh_*h^*a^*\mathcal{F}$。我们有
\begin{align*}
Ra_*(a^*\mathcal{F})
& =
Ra_*Rh_*h^*a^*\mathcal{F} \\
& =
Rh_{-1, *}Ra_{fppf, *}a_{fppf}^*(h_{-1})^*\mathcal{F} \\
& =
Rh_{-1, *}(h_{-1})^*\mathcal{F} \\
& =
\mathcal{F}
\end{align*}
% P10-E0228
第一个等号来自上面的讨论；第二个等号来自
引理 \ref{spaces-simplicial-lemma-compare-simplicial-objects} 中图的交换性；
第三个等号来自引理
\ref{spaces-simplicial-lemma-hypercovering-X-simple-descent-bounded-modules}，
因为 $U$ 是 $(\textit{Spaces}/S)_{fppf}$ 中 $X$ 的超覆盖，
且 $a_{fppf}$ 平坦，从而
$La_{fppf}^* = a_{fppf}^*$
（即 $a_{fppf}^{-1}\mathcal{O}_{big} =
\mathcal{O}_{big, total}$，见注 \ref{spaces-simplicial-remark-augmentation-ringed}）；
最后一个等号来自已经用过的空间上上同调进阶的引理
\ref{spaces-more-cohomology-lemma-cohomological-descent-etale-fppf-modules}。
\end{proof}

\begin{lemma}
\label{spaces-simplicial-lemma-coh-descent-qcoh-unbounded-for-fppf-hypercovering}
设 $S$ 是概形，$X$ 是 $S$ 上的代数空间。
设 $U$ 是 $S$ 上的单纯代数空间，$a : U \to X$ 是增广。
假设 $a : U \to X$ 是 $X$ 的 fppf 超覆盖。
则 $\QCoh(\mathcal{O}_U)$ 是
$\textit{Mod}(\mathcal{O}_U)$ 的弱 Serre 子范畴，且
$$
a^* : D_\QCoh(\mathcal{O}_X) \longrightarrow D_\QCoh(\mathcal{O}_U)
$$
是范畴等价，其拟逆由 $Ra_*$ 给出。
这里 $a : \Sh(U_\etale) \to \Sh(X_\etale)$
如第 \ref{spaces-simplicial-section-simplicial-algebraic-spaces} 节所述。
\end{lemma}

\begin{proof}
首先注意，由超覆盖的注 \ref{hypercovering-remark-P-covering}，
映射 $a_n : U_n \to X$ 和 $d^n_i : U_n \to U_{n - 1}$
是平坦、局部有限呈示且满的。

\medskip\noindent
回忆，$\mathcal{O}_U$-模 $\mathcal{F}$ 拟凝聚，当且仅当
它是笛卡尔的，并且对所有 $n$，$\mathcal{F}_n$ 都拟凝聚；
见引理 \ref{spaces-simplicial-lemma-quasi-coherent-sheaf}。
由引理 \ref{spaces-simplicial-lemma-Serre-subcat-cartesian-modules}
（以及上面所示映射 $d^n_i : U_n \to U_{n - 1}$ 的平坦性），
% P10-E0240
笛卡尔模构成 $\textit{Mod}(\mathcal{O}_U)$ 的弱 Serre 子范畴。
另一方面，对每个 $n$，
$\QCoh(\mathcal{O}_{U_n}) \subset \textit{Mod}(\mathcal{O}_{U_n})$
是弱 Serre 子范畴（空间的性质，引理
\ref{spaces-properties-lemma-properties-quasi-coherent}）。
合并二者可见
$\QCoh(\mathcal{O}_U) \subset \textit{Mod}(\mathcal{O}_U)$
是弱 Serre 子范畴。

\medskip\noindent
为完成证明，我们逐一验证站点上的上同调的引理
\ref{sites-cohomology-lemma-equivalence-unbounded-one}
中的条件 (1) -- (5)。

\medskip\noindent
% P10-E0241
关于 (1)。由引理 \ref{spaces-simplicial-lemma-flat-augmentation-modules}，
$a_n$ 平坦（见上文）蕴含 $a$ 平坦，故该条件成立。

\medskip\noindent
关于 (2)。这正是引理
\ref{spaces-simplicial-lemma-descent-qcoh-for-fppf-hypercovering} 的内容。

\medskip\noindent
关于 (3)。这正是引理
\ref{spaces-simplicial-lemma-cohomological-descent-qcoh-for-fppf-hypercovering} 的内容。

\medskip\noindent
关于 (4)。回忆，我们既可用站点 $U_\etale$，也可用
$U_{spaces, \etale}$ 来定义小平展拓扑斯
$\Sh(U_\etale)$；见第 \ref{spaces-simplicial-section-simplicial-algebraic-spaces} 节。
站点上的上同调的情形
\ref{sites-cohomology-situation-olsson-laszlo} 的假设对三元组
$(U_{spaces, \etale}, \mathcal{O}_U, \QCoh(\mathcal{O}_U))$
成立；同理也对三元组
$(U_\etale, \mathcal{O}_U, \QCoh(\mathcal{O}_U))$.
成立。事实上，取
$$
\mathcal{B} \subset \Ob(U_\etale) \subset \Ob(U_{spaces, \etale})
$$
为仿射对象之集。对 $V/U_n \in \mathcal{B}$，
取 $d_{V/U_n} = 0$，并令 $\text{Cov}_{V/U_n}$ 为
$V_i$ 仿射的平展覆盖 $\{V_i \to V\}$ 之集。
于是得到所需的消失，因为对
$\mathcal{F} \in \QCoh(\mathcal{O}_U)$
和任意 $V/U_n \in \mathcal{B}$，由引理
\ref{spaces-simplicial-lemma-sanity-check-modules} 有
$$
H^p(V/U_n, \mathcal{F}) = H^p(V, \mathcal{F}_n)
$$
。右端是 $U_n$ 上拟凝聚层 $\mathcal{F}_n$
在 $U_{n, \etale}$ 的仿射对象 $V$ 上的上同调。
由空间的上同调的第
\ref{spaces-cohomology-section-higher-direct-image} 节中的讨论
以及概形的上同调的引理
\ref{coherent-lemma-quasi-coherent-affine-cohomology-zero}，
当 $p > 0$ 时它消失。

\medskip\noindent
关于 (5)。取 $\mathcal{B} \subset \Ob(X_{spaces, \etale})$
为仿射对象之集，并应用上面给出的文献，即得该条件。
\end{proof}

\begin{lemma}
\label{spaces-simplicial-lemma-compute-via-fppf-hypercovering-modules}
设 $S$ 是概形，$X$ 是 $S$ 上的代数空间。
设 $U$ 是 $S$ 上的单纯代数空间，$a : U \to X$ 是增广。
若 $a : U \to X$ 是 $X$ 的 fppf 超覆盖，则
$$
R\Gamma(X_\etale, K) = R\Gamma(U_\etale, a^*K)
$$
对 $K \in D_\QCoh(\mathcal{O}_X)$ 成立。
这里 $a : \Sh(U_\etale) \to \Sh(X_\etale)$
如第 \ref{spaces-simplicial-section-simplicial-algebraic-spaces} 节所述。
\end{lemma}

\begin{proof}
这由引理
\ref{spaces-simplicial-lemma-coh-descent-qcoh-unbounded-for-fppf-hypercovering} 得出，
因为由站点上的上同调的注
\ref{sites-cohomology-remark-before-Leray}，
$R\Gamma(U_\etale, -) = R\Gamma(X_\etale, -) \circ Ra_*$。
\end{proof}

\begin{lemma}
\label{spaces-simplicial-lemma-spectral-sequence-fppf-hypercovering-modules}
设 $S$ 是概形，$X$ 是 $S$ 上的代数空间。
设 $U$ 是 $S$ 上的单纯代数空间，$a : U \to X$ 是增广。
% P10-E0242
设 $\mathcal{F}$ 是拟凝聚 $\mathcal{O}_X$-模，
并设 $\mathcal{F}_n$ 是它在 $U_{n, \etale}$ 上的拉回。
若 $U$ 是 $X$ 的 fppf 超覆盖，则存在典范谱序列
$$
E_1^{p, q} = H^q_\etale(U_p, \mathcal{F}_p)
$$
，它收敛于 $H^{p + q}_\etale(X, \mathcal{F})$。
\end{lemma}

\begin{proof}
这是引理 \ref{spaces-simplicial-lemma-compute-via-fppf-hypercovering-modules}
和 \ref{spaces-simplicial-lemma-simplicial-module-cohomology-site} 的直接推论。
\end{proof}







\section{复形的 fppf 下降}
\label{spaces-simplicial-section-fppf-descent-derived}

\noindent
在本节中，我们把前面证明的若干结果汇集起来，
用于代数空间的 fppf 覆盖和拟凝聚模的导出范畴。

\begin{lemma}
\label{spaces-simplicial-lemma-fppf-neg-ext-zero-hom}
设 $X$ 是概形 $S$ 上的代数空间。
设 $K, E \in D_\QCoh(\mathcal{O}_X)$。
设 $a : U \to X$ 是 fppf 超覆盖。
假设对所有 $n \geq 0$ 都有
$$
\Ext_{\mathcal{O}_{U_n}}^i(La_n^*K, La_n^*E) = 0
\text{，当 } i < 0
$$
。则
\begin{enumerate}
\item 当 $i < 0$ 时，$\Ext_{\mathcal{O}_X}^i(K, E) = 0$；并且
\item 存在正合序列
$$
0
\to
\Hom_{\mathcal{O}_X}(K, E)
\to
\Hom_{\mathcal{O}_{U_0}}(La_0^*K, La_0^*E)
\to
\Hom_{\mathcal{O}_{U_1}}(La_1^*K, La_1^*E)
$$
\end{enumerate}
\end{lemma}

\begin{proof}
记 $K_n = La_n^*K$ 且 $E_n = La_n^*E$。
那么，它们分别是与 $La^*K$ 和 $La^*E$ 关联的模之导出范畴的
单纯系统（定义 \ref{spaces-simplicial-definition-cartesian-derived-modules} 和
引理 \ref{spaces-simplicial-lemma-cartesian-objects-derived-modules}），
其中 $a : U_\etale \to X_\etale$
如第 \ref{spaces-simplicial-section-simplicial-algebraic-spaces} 节所述。
先证明 (2)。由引理
\ref{spaces-simplicial-lemma-coh-descent-qcoh-unbounded-for-fppf-hypercovering}，
$$
\Hom_{\mathcal{O}_X}(K, E) =
\Hom_{\mathcal{O}_U}(La^*K, La^*E)
$$
。因此，该序列具有如下形式：
$$
0
\to
\Hom_{\mathcal{O}_U}(La^*K, La^*E)
\to
\Hom_{\mathcal{O}_{U_0}}(K_0, E_0)
\to
\Hom_{\mathcal{O}_{U_1}}(K_1, E_1)
$$
由引理 \ref{spaces-simplicial-lemma-nullity-cartesian-modules-derived}，
第一个箭头是单射。由引理
\ref{spaces-simplicial-lemma-hom-cartesian-modules-derived}，
该箭头的像是第二个箭头的核。这就完成了 (2) 的证明。
对 $i > 0$ 将 (2) 应用于 $K[i]$ 和 $E$，即得 (1)。
\end{proof}

\begin{lemma}
\label{spaces-simplicial-lemma-fppf-glue-neg-ext-zero}
设 $X$ 是概形 $S$ 上的代数空间。
设 $a : U \to X$ 是 fppf 超覆盖。
假设给定 $K_0 \in D_\QCoh(U_0)$ 和同构
$$
\alpha :
L(f_{\delta_1^1})^*K_0
\longrightarrow
L(f_{\delta_0^1})^*K_0
$$
% P10-E0243
，它在 $U_1$ 上满足余圈条件。令
$\tau^n_i : [0] \to [n]$，$0 \mapsto i$，并令
$K_n = Lf_{\tau^n_n}^*K_0$。
假设当 $i < 0$ 时
$\Ext^i_{\mathcal{O}_{U_n}}(K_n, K_n) = 0$。
% P10-E0244
则存在对象 $K \in D_\QCoh(\mathcal{O}_X)$
以及与 $\alpha$ 相容的同构 $La_0^*K \to K$。
\end{lemma}

\begin{proof}
由引理 \ref{spaces-simplicial-lemma-construct-cartesian-objects-derived-modules}，
对象 $K_n$ 构成模之导出范畴的一个单纯系统的各成员。
于是得到对象
$K' \in D_\QCoh(\mathcal{O}_{U_\etale})$，
使得 $(K_n, K_\varphi)$ 是由 $K'$ 推出的系统；见引理
\ref{spaces-simplicial-lemma-cartesian-module-derived-from-simplicial}。
最后，应用引理
\ref{spaces-simplicial-lemma-coh-descent-qcoh-unbounded-for-fppf-hypercovering}，
可见对某个 $K \in D_\QCoh(\mathcal{O}_X)$，
$K' = La^*K$，这正是所需的。
\end{proof}











\section{代数空间的真超覆盖}
\label{spaces-simplicial-section-proper-hypercovering-spaces}

\noindent
本节给出代数空间情形下第
\ref{spaces-simplicial-section-proper-hypercovering} 节的对应版本。
愿意如此处理的读者可以在各处把“代数空间”替换为“概形”，
所得结果同样有效。这样做的好处是，可以把对空间上上同调进阶的
第 \ref{spaces-more-cohomology-section-ph-etale} 节的引用，
替换为对平展上同调的第
\ref{etale-cohomology-section-ph-etale} 节的引用。

\medskip\noindent
固定一个基概形 $S$。设 $X$ 是 $S$ 上的代数空间，
$U$ 是 $S$ 上的单纯代数空间。假设有增广
$$
a : U \to X
$$
；见第 \ref{spaces-simplicial-section-simplicial-algebraic-spaces} 节。
如果满足下列条件，就称 $U$ 是 $X$ 的{\it 真超覆盖}：
\begin{enumerate}
\item $U_0 \to X$ 是真且满的；
\item $U_1 \to U_0 \times_X U_0$ 是真且满的；
\item $U_{n + 1} \to (\text{cosk}_n\text{sk}_n U)_{n + 1}$
对 $n \geq 1$ 是真且满的。
\end{enumerate}
$S$ 上代数空间的范畴具有所有有限极限，因而上述表述中使用的
余骨架存在。
$$
\fbox{原理：真超覆盖可用于计算平展上同调。}
$$
证明该原理的关键思想，是比较范畴 $\textit{Spaces}/S$
上的 ph 拓扑与平展拓扑。具体地说，ph 拓扑强于平展拓扑，并且
(a) 真满映射给出 ph 覆盖；以及
(b) 从小平展站点拉回的层的 ph 上同调与平展上同调一致，
正如我们在空间上上同调进阶的第
\ref{spaces-more-cohomology-section-ph-etale} 节中所见。

\medskip\noindent
本节的所有结果都可推广到 $U \to X$ 仅为“ph 超覆盖”的情形，
即按第 \ref{spaces-simplicial-section-hypercovering} 节的定义，
它是站点 $(\textit{Spaces}/S)_{ph}$ 中 $X$ 的超覆盖。
若将来需要这一点，我们会在此处精确表述并证明它。

\begin{lemma}
\label{spaces-simplicial-lemma-compare-simplicial-objects-ph-etale}
设 $S$ 是概形，$X$ 是 $S$ 上的代数空间。
设 $U$ 是 $S$ 上的单纯代数空间。
设 $a : U \to X$ 是增广。存在交换图
$$
\xymatrix{
\Sh((\textit{Spaces}/U)_{ph, total}) \ar[r]_-h \ar[d]_{a_{ph}} &
\Sh(U_\etale) \ar[d]^a \\
\Sh((\textit{Spaces}/X)_{ph}) \ar[r]^-{h_{-1}} &
\Sh(X_\etale)
}
$$
，其中左侧竖直箭头定义于第 \ref{spaces-simplicial-section-hypercovering} 节，
右侧竖直箭头定义于第
\ref{spaces-simplicial-section-simplicial-algebraic-spaces} 节。
\end{lemma}

\begin{proof}
记号 $(\textit{Spaces}/U)_{ph, total}$ 表示我们把
第 \ref{spaces-simplicial-section-hypercovering} 节的构造用于站点
$(\textit{Spaces}/S)_{ph}$ 及其单纯对象 $U$
\footnote{这是为了与第 \ref{spaces-simplicial-section-fppf-hypercovering} 节中
用 fppf 拓扑定义的
$(\textit{Spaces}/U)_{fppf, total}$ 相区别。}。
对右侧的拓扑斯，我们将使用站点
$X_{spaces, \etale}$ 和 $U_{spaces, \etale}$；
这是容许的，见第 \ref{spaces-simplicial-section-simplicial-algebraic-spaces} 节中的讨论。

\medskip\noindent
注意，$(\textit{Spaces}/U)_{ph, total}$ 和
$U_{spaces, \etale}$ 都属于情形
\ref{spaces-simplicial-situation-simplicial-site} 的情形 A。
对 $U_\etale$，这直接来自第
\ref{spaces-simplicial-section-simplicial-algebraic-spaces} 节中的构造；
对 $(\textit{Spaces}/U)_{ph, total}$，则由引理
\ref{spaces-simplicial-lemma-sr-when-fibre-products} 得出。
接下来考虑函子
$U_{n, spaces, \etale} \to (\textit{Spaces}/U_n)_{ph}$, $U \mapsto U/U_n$
以及
$X_{spaces, \etale} \to (\textit{Spaces}/X)_{ph}$, $U \mapsto U/X$.
。我们已在空间上上同调进阶的第
\ref{spaces-more-cohomology-section-ph-etale} 节中看到，
它们定义站点的态射；在那里这些态射记为
% P10-E0168
% P10-E0245
$a_{U_n} = \epsilon_{U_n} \circ \pi_{u_n}$
和 $a_X = \epsilon_X \circ \pi_X$。
因此，如注 \ref{spaces-simplicial-remark-morphism-augmentation-simplicial-sites}，
我们得到与增广相容的单纯站点态射；
应用引理 \ref{spaces-simplicial-lemma-morphism-augmentation-simplicial-sites}
即可得出结论。
\end{proof}

\begin{lemma}
\label{spaces-simplicial-lemma-descent-sheaves-for-ph-hypercovering}
设 $S$ 是概形，$X$ 是 $S$ 上的代数空间。
设 $U$ 是 $S$ 上的单纯代数空间，$a : U \to X$ 是增广。
若 $a : U \to X$ 是 $X$ 的真超覆盖，则
$$
a^{-1} : \Sh(X_\etale) \to \Sh(U_\etale)
\quad\text{且}\quad
a^{-1} : \textit{Ab}(X_\etale) \to \textit{Ab}(U_\etale)
$$
是全忠实的，其本质像为笛卡尔层，拟逆由 $a_*$ 给出。
这里 $a : \Sh(U_\etale) \to \Sh(X_\etale)$
如第 \ref{spaces-simplicial-section-simplicial-algebraic-spaces} 节所述。
\end{lemma}

\begin{proof}
我们对集合层证明该断言。它几乎是已经建立的结果的形式推论。
考虑引理 \ref{spaces-simplicial-lemma-compare-simplicial-objects-ph-etale} 中的图。
在该引理的证明中，我们已看到 $h_{-1}$ 是空间上上同调进阶的
第 \ref{spaces-more-cohomology-section-ph-etale} 节中的态射 $a_X$。
因此，由空间上上同调进阶的引理
\ref{spaces-more-cohomology-lemma-comparison-ph-etale}，
$(h_{-1})^{-1}$ 是全忠实的，拟逆为 $h_{-1, *}$。
对 $h$ 的分量 $h_n$ 同样成立。由引理
\ref{spaces-simplicial-lemma-morphism-simplicial-sites} 对函子 $h^{-1}$ 和 $h_*$
的描述，我们得到 $h^{-1}$ 全忠实，拟逆为 $h_*$。
注意，按第 \ref{spaces-simplicial-section-hypercovering} 节的定义，
$U$ 是 $(\textit{Spaces}/S)_{ph}$ 中 $X$ 的超覆盖，
因为由空间上的拓扑的引理
\ref{spaces-topologies-lemma-surjective-proper-ph}，
真满态射给出 ph 覆盖。
由引理 \ref{spaces-simplicial-lemma-hypercovering-X-simple-descent-sheaves}，
$a_{ph}^{-1}$ 是全忠实的，拟逆为 $a_{ph, *}$，
且其本质像为 $(\textit{Spaces}/U)_{ph, total}$ 上的笛卡尔层。
现在，一个形式论证（沿图追逐）表明 $a^{-1}$ 全忠实。

\medskip\noindent
最后，设 $\mathcal{G}$ 是 $U_\etale$ 上的笛卡尔层。
则 $h^{-1}\mathcal{G}$ 是
$(\textit{Spaces}/U)_{ph, total}$ 上的笛卡尔层。
因而对 $(\textit{Spaces}/X)_{ph}$ 上的某个层
$\mathcal{H}$，有
$h^{-1}\mathcal{G} = a_{ph}^{-1}\mathcal{H}$。
用稍显繁琐的记号计算如下：
\begin{align*}
(h_{-1})^{-1}(a_*\mathcal{G})
& =
(h_{-1})^{-1}
\text{等化子}(
\xymatrix{
a_{0, small, *}\mathcal{G}_0
\ar@<1ex>[r] \ar@<-1ex>[r] &
a_{1, small, *}\mathcal{G}_1
}
) \\
& =
\text{等化子}(
\xymatrix{
(h_{-1})^{-1}a_{0, small, *}\mathcal{G}_0
\ar@<1ex>[r] \ar@<-1ex>[r] &
(h_{-1})^{-1}a_{1, small, *}\mathcal{G}_1
}
) \\
& =
\text{等化子}(
\xymatrix{
a_{0, big, ph, *}h_0^{-1}\mathcal{G}_0
\ar@<1ex>[r] \ar@<-1ex>[r] &
a_{1, big, ph, *}h_1^{-1}\mathcal{G}_1
}
) \\
& =
\text{等化子}(
\xymatrix{
a_{0, big, ph, *}(a_{0, big, ph})^{-1}\mathcal{H}
\ar@<1ex>[r] \ar@<-1ex>[r] &
a_{1, big, ph, *}(a_{1, big, ph})^{-1}\mathcal{H}
}
) \\
& =
a_{ph, *}a_{ph}^{-1}\mathcal{H} \\
& =
\mathcal{H}
\end{align*}
这里，第一个等号来自引理 \ref{spaces-simplicial-lemma-augmentation-site}；
第二个等号来自 $(h_{-1})^{-1}$ 是正合函子；
第三个等号来自空间上上同调进阶的引理
\ref{spaces-more-cohomology-lemma-proper-push-pull-ph-etale}
（这里用到 $a_0 : U_0 \to X$ 和 $a_1: U_1 \to X$ 为真）；
第四个等号来自
$a_{ph}^{-1}\mathcal{H} = h^{-1}\mathcal{G}$；
第五个等号来自引理 \ref{spaces-simplicial-lemma-augmentation-site}；
第六个等号已在上面见到。由于
$a_{ph}^{-1}\mathcal{H} = h^{-1}\mathcal{G}$，
我们推出
$h^{-1}\mathcal{G} \cong h^{-1}a^{-1}a_*\mathcal{G}$；
再由 $h^{-1}$ 的全忠实性，证明完毕。
\end{proof}

\begin{lemma}
\label{spaces-simplicial-lemma-cohomological-descent-for-ph-hypercovering}
设 $S$ 是概形，$X$ 是 $S$ 上的代数空间。
设 $U$ 是 $S$ 上的单纯代数空间，$a : U \to X$ 是增广。
若 $a : U \to X$ 是 $X$ 的真超覆盖，则对
$K \in D^+(X_\etale)$，
$$
K \to Ra_*(a^{-1}K)
$$
是同构。这里 $a : \Sh(U_\etale) \to \Sh(X_\etale)$
如第 \ref{spaces-simplicial-section-simplicial-algebraic-spaces} 节所述。
\end{lemma}

\begin{proof}
考虑引理 \ref{spaces-simplicial-lemma-compare-simplicial-objects-ph-etale} 中的图。
由空间上上同调进阶的引理
\ref{spaces-more-cohomology-lemma-cohomological-descent-etale-ph}，
$Rh_{n, *}h_n^{-1}$ 是 $D^+(U_{n, \etale})$ 上的恒等函子。
因此，由引理 \ref{spaces-simplicial-lemma-direct-image-morphism-simplicial-sites}，
$Rh_*h^{-1}$ 是 $D^+(U_\etale)$ 上的恒等函子。我们有
\begin{align*}
Ra_*(a^{-1}K)
& =
Ra_*Rh_*h^{-1}a^{-1}K \\
& =
Rh_{-1, *}Ra_{ph, *}a_{ph}^{-1}(h_{-1})^{-1}K \\
& =
Rh_{-1, *}(h_{-1})^{-1}K \\
& =
K
\end{align*}
% P10-E0228
第一个等号来自上面的讨论；第二个等号来自
% P10-E0238
引理 \ref{spaces-simplicial-lemma-compare-simplicial-objects} 中图的交换性；
第三个等号来自引理
\ref{spaces-simplicial-lemma-hypercovering-X-simple-descent-bounded-abelian}，
因为由空间上的拓扑的引理
\ref{spaces-topologies-lemma-surjective-proper-ph}，
$U$ 是 $(\textit{Spaces}/S)_{ph}$ 中 $X$ 的超覆盖；
最后一个等号来自已经用过的空间上上同调进阶的引理
\ref{spaces-more-cohomology-lemma-cohomological-descent-etale-ph}。
\end{proof}

\begin{lemma}
\label{spaces-simplicial-lemma-compute-via-ph-hypercovering}
设 $S$ 是概形，$X$ 是 $S$ 上的代数空间。
设 $U$ 是 $S$ 上的单纯代数空间，$a : U \to X$ 是增广。
若 $a : U \to X$ 是 $X$ 的真超覆盖，则
$$
R\Gamma(X_\etale, K) = R\Gamma(U_\etale, a^{-1}K)
$$
对 $K \in D^+(X_\etale)$ 成立。
这里 $a : \Sh(U_\etale) \to \Sh(X_\etale)$
如第 \ref{spaces-simplicial-section-simplicial-algebraic-spaces} 节所述。
\end{lemma}

\begin{proof}
这由引理 \ref{spaces-simplicial-lemma-cohomological-descent-for-ph-hypercovering} 得出，
因为由站点上的上同调的注
\ref{sites-cohomology-remark-before-Leray}，
$R\Gamma(U_\etale, -) = R\Gamma(X_\etale, -) \circ Ra_*$。
\end{proof}

\begin{lemma}
\label{spaces-simplicial-lemma-ph-hypercovering-equivalence-bounded}
设 $S$ 是概形，$X$ 是 $S$ 上的代数空间。
设 $U$ 是 $S$ 上的单纯代数空间，$a : U \to X$ 是增广。
以 $\mathcal{A} \subset \textit{Ab}(U_\etale)$
表示笛卡尔阿贝尔层所成的弱 Serre 子范畴。
若 $U$ 是 $X$ 的真超覆盖，则函子 $a^{-1}$ 定义等价
$$
D^+(X_\etale) \longrightarrow D_\mathcal{A}^+(U_\etale)
$$
，其拟逆为 $Ra_*$。这里
$a : \Sh(U_\etale) \to \Sh(X_\etale)$
如第 \ref{spaces-simplicial-section-simplicial-algebraic-spaces} 节所述。
\end{lemma}

\begin{proof}
由引理 \ref{spaces-simplicial-lemma-Serre-subcat-cartesian-modules}，
$\mathcal{A}$ 是弱 Serre 子范畴。
该等价是目前所得结果的形式推论。应用引理
\ref{spaces-simplicial-lemma-descent-sheaves-for-ph-hypercovering}、
\ref{spaces-simplicial-lemma-cohomological-descent-for-ph-hypercovering}
以及站点上的上同调的引理
\ref{sites-cohomology-lemma-equivalence-bounded}。
\end{proof}

\begin{lemma}
\label{spaces-simplicial-lemma-spectral-sequence-ph-hypercovering}
设 $S$ 是概形，$X$ 是 $S$ 上的代数空间。
设 $U$ 是 $S$ 上的单纯代数空间，$a : U \to X$ 是增广。
设 $\mathcal{F}$ 是 $X_\etale$ 上的阿贝尔层，
并设 $\mathcal{F}_n$ 是它在 $U_{n, \etale}$ 上的拉回。
若 $U$ 是 $X$ 的 ph 超覆盖，则存在典范谱序列
$$
E_1^{p, q} = H^q_\etale(U_p, \mathcal{F}_p)
$$
，它收敛于 $H^{p + q}_\etale(X, \mathcal{F})$。
\end{lemma}

\begin{proof}
这是引理 \ref{spaces-simplicial-lemma-compute-via-ph-hypercovering}
和 \ref{spaces-simplicial-lemma-simplicial-sheaf-cohomology-site} 的直接推论。
\end{proof}
